%=====================================================================
%  DATABASE SYSTEMS --- FALL 2026
%  Unified Master Lecture Handout
%
%  Taught as CS-103, BSIT 2nd semester, 4 credit hours (3+1) --- three
%  lecture credits plus a CREDITED laboratory.  The registrar's CBT
%  sheet (Summer_2026/question_papers/CBT- IT-Summer 2026 Final.xlsx,
%  row 54) is the authority for the code and the semester.  Unlike every
%  sibling course, there is no INFT- twin code: CS-103 is the only code
%  the department uses for it.
%
%  Written against the HEC outline published in BS & MS (CS, SE & IT)
%  Curriculum 2017, p. 88, and REPRINTED WORD FOR WORD in BS Curriculum
%  Computing Disciplines 2023, p. 58 -- the two texts differ by a single
%  added period across ninety words, six years apart.
%
%  The 2025 edition DELETES the syllabus outright.  It covers only BS
%  Computer Science and the Associate Degree (no BSIT at all), carries
%  no Course Content section for any major course, and supplies five
%  rewritten CLO bullets with no Bloom tags, delegating content,
%  prerequisites, pedagogy, assessment and textbooks to the
%  universities.  Those five CLOs are nevertheless the reason Chapter 23
%  exists, because they name security -- which 2017 and 2023 do not.
%
%  Six defects in the published documents are recorded in the Course
%  Specification rather than inherited, including a prerequisite that
%  contradicts itself inside one booklet and a CLO table that is
%  typographically mangled in print.
%
%  THE AUDIENCE IS SECOND SEMESTER -- three semesters earlier than any
%  other course in this series.  Data Structures has not been taken, so
%  Chapters 21 and 22 are self-contained.
%
%  Build:  latexmk -pdf DBS_Fall2026_Handout.tex
%     or:  .\build.ps1
%
%  Structure: 6 parts, 23 chapters, 7 appendices.
%  All style and pedagogy macros live in dbshandout.sty.
%=====================================================================
\documentclass[11pt,a4paper,oneside]{report}
\usepackage{dbshandout}

\begin{document}

%---------------------------------------------------------------------
%=====================================================================
% FRONT MATTER
%=====================================================================
\pagenumbering{roman}
\structuralfigures{F.}
\thispagestyle{empty}

\begin{center}
\vspace*{1.0cm}

{\color{secondary}\rule{0.85\textwidth}{2pt}}\\[6pt]
{\large\textsc{\color{secondary}Bachelor of Science in Information Technology}}\\[18pt]

{\Huge\bfseries\color{primary}DATABASE SYSTEMS}\\[10pt]
{\LARGE\bfseries\color{primary}Designing, Querying and Trusting\\[4pt]a Relational Database}\\[16pt]

{\color{secondary}\rule{0.85\textwidth}{2pt}}\\[16pt]

{\Large\textsc{Unified Lecture Handout}}\\[6pt]
{\large\color{gray!70!black}Complete Course in One Volume $\bullet$ Fall 2026}\\[24pt]

% --- the claim the whole book rests on ---
\begin{tikzpicture}
  \node[rounded corners=4pt, draw=primary, line width=1.1pt, fill=primary!6,
        text=primary, align=center, font=\small\bfseries,
        minimum width=11.6cm, minimum height=1.5cm]
       {\faIcon{database}\\[2pt]A database is not a place to put data.
        It is a set of promises about data.};
\end{tikzpicture}\\[10pt]
{\small\itshape\color{gray!70!black}This course is how each promise is made
--- and, every time, what the database does when you do not ask it to keep
one.}

\vspace{1.0cm}

\begin{tcolorbox}[enhanced, width=0.86\textwidth, colback=lightbg,
  colframe=primary, boxrule=0pt, leftrule=4pt, arc=2pt, top=6pt, bottom=6pt]
\small
\textbf{\color{primary}What this volume covers.} Every clause of the HEC
outline for \emph{Database Systems}, traced chapter by chapter in the Course
Specification overleaf.

\smallskip
That outline is published in the \textbf{2017} curriculum (p.~88) and
reprinted \emph{word for word} in the \textbf{2023} curriculum (p.~58) --- the
two differ by a single added full stop across ninety words, six years apart.
The \textbf{2025} curriculum deletes it: it covers only BS Computer Science,
publishes no course content for any major course at all, and supplies five
rewritten learning outcomes with no Bloom levels. All of that is recorded in
the Course Specification rather than smoothed over.

\smallskip
The course is taught on a \textbf{running database}. One schema is designed in
Chapter~6, created in Chapter~10, and every chapter afterwards adds to it or
asks something of it; by Chapter~23 it is the semester project. The tool is
\textbf{SQLite}, which you already have.
\end{tcolorbox}

\vspace{0.9cm}
{\small\color{gray!60!black}Six Parts $\bullet$ Twenty-Three Chapters $\bullet$ Seven Appendices}
\vspace{1cm}
\end{center}

\newpage

%=====================================================================
\chapter*{How to Use This Handout}
\addcontentsline{toc}{chapter}{How to Use This Handout}
\markboth{How to Use This Handout}{}

There are two ways to teach this subject badly.

The first is to teach it as theory. Relational algebra on a whiteboard,
normal forms as definitions to be recited, ACID as four words beginning with
the right letters. A student who has learned it this way can answer an
examination question about third normal form and cannot look at a table and
say whether it is in it.

The second is to teach it as a tool. Open the software, click here, type a
\sqlkw{SELECT}, admire the rows. A student who has learned it this way can
produce an answer and has no way of knowing whether it is the right one,
because nothing in the course ever told them what the database was promising.

\textbf{This handout does neither, and the reason is in the title.} A database
makes promises --- that this attribute is a number, that this row refers to a
row that exists, that either all of this happened or none of it did --- and
every one of those promises is \emph{checkable}. So the method of this book is
to state the promise, then ask the database to break it, and see what happens.
Eleven chapters end with a figure produced that way on a real machine. They
are marked, and the file that produces each one is named.

\section*{The furniture}

Every chapter has the same parts, in the same order, so that you can find
what you need without reading from the beginning.

\rowcolors{2}{white}{lightbg}
\begin{longtable}{@{}L{4.2cm}L{10.4cm}@{}}
\toprule
\rowcolor{primary!15}
\textbf{You will see} & \textbf{It means} \\
\midrule
\endhead
\faIcon{bullseye}~Learning Outcomes & What you should be able to \emph{do}
afterwards. Five per chapter, each with a verb you can be examined on \\
\faIcon{link}~Before You Start & What this chapter assumes, and where it was
taught \\
\faIcon{tags}~Key Terms & Every term defined here. If one is unfamiliar at the
end, the chapter did not work \\
\faIcon{book}~Definition & A term being given its precise meaning. These are
the sentences to learn exactly \\
\faIcon{lightbulb}~Concept & An idea, a rule of thumb, or a way of deciding \\
\faIcon{triangle-exclamation}~Watch Out & Something that is true, surprising,
and will cost you if you do not know it \\
\faIcon{code}~Example & A worked case, usually from the running database \\
\faIcon{pen-to-square}~Worked & Arithmetic, every step shown. Nothing is left
``by inspection'' \\
\faIcon{tachometer-alt}~Measured & A figure produced by running the named file
in \texttt{code/}, on the machine in Appendix~A. You can reproduce every one \\
\faIcon{circle-exclamation}~Common Pitfalls & What goes wrong, collected, so
that you meet it here rather than in the laboratory \\
\faIcon{flask}~Try It Yourself & The laboratory. Numbered steps, run in order \\
\faIcon{circle-question}~Checkpoint & Three questions. Answers in Appendix~C,
and they are full answers rather than a line each \\
\faIcon{list-check}~What to Take Away & The chapter in eight or nine lines \\
\faIcon{pen-fancy}~Review Questions & Seven: two multiple choice, three short,
two applied \\
\bottomrule
\end{longtable}

\section*{The laboratory, and what it needs}

\textbf{Nothing to install.} The whole course runs on \textbf{SQLite}, which
is already on every machine you are likely to use: it ships inside Python, it
ships with most operating systems, and a database is a single file you can
copy, email or delete. Appendix~A sets it up in one sitting, which for most
readers means checking that it is already there.

That choice is not a compromise with the curriculum. \textbf{HEC names no
database product in any edition} --- the 2017 and 2023 outlines say only
``in any DBMS'', and 2025 says ``popular database management systems''. What
SQLite can and cannot show you is stated plainly in \chref{whynotfiles} and
again in \chref{project}, because a tool's limits are part of what you are
learning.

\section*{The database we build}

One schema, carried through the whole course. You design it in
\chref{ermodel} and \chref{cardinality}, map it to tables in \chref{mapping},
create it in \chref{ddl}, fill it, query it for five chapters, find out in
\chref{normalforms} that it was wrong, fix it, index it, and defend it. It is
a university enrolment database --- students, courses, sections, instructors,
enrolments, grades --- chosen because you already know how it is supposed to
behave, so no lecture time goes on explaining the subject matter.

\section*{If you read nothing else}

\begin{tight}
  \item \chref{constraints}, on the difference between \emph{declaring} a key
        and \emph{having} one.
  \item \chref{normalforms}, which is the chapter this course exists for.
  \item \chref{indexes}, where the rule everybody repeats stops being true and
        you can see the exact row count at which it does.
\end{tight}

\newpage

%=====================================================================
\chapter*{Course Specification}
\addcontentsline{toc}{chapter}{Course Specification}
\markboth{Course Specification}{}

This chapter exists so that a student, a lecturer or a reviewer can see
exactly what this handout claims to cover and on what authority. Everything in
it is traceable to a published document, and where those documents are wrong,
unclear or contradict one another, that is recorded here rather than quietly
corrected.

\section*{Particulars}

\rowcolors{2}{white}{lightbg}
\begin{longtable}{@{}L{3.6cm}L{11.0cm}@{}}
\toprule
\rowcolor{primary!15}
\textbf{Item} & \textbf{Value} \\
\midrule
\endhead
Course title & \emph{Database Systems}. HEC 2023's description page calls it
\emph{Database System}, singular, while every table in the same booklet says
\emph{Systems}. The plural is used here \\
Course code & \textbf{CS-103}. Unlike every sibling course in this series,
there is no \texttt{INFT-} equivalent: the registrar uses CS-103 and nothing
else \\
Programme & BS Information Technology, \textbf{2nd semester}. This matches HEC
2023, which places the course in semester~2 of every computing degree. HEC
2017 placed it at BSIT semester~V; the discrepancy is recorded, not resolved \\
Credit hours & \textbf{4 (3+1)} --- three lecture credits and \textbf{one
credited laboratory}. The laboratory is paid for, which is why every chapter
has one and why they are long enough to be worth the hour \\
Prerequisite & \textbf{None.} HEC 2017's own booklet gives two different
answers --- see defect~4 \\
Duration & Sixteen teaching weeks, midterm in Week~8, project due Week~16
(Appendix~G) \\
Teaching methodology & As published (HEC 2017): lectures, written assignments,
\textbf{practical labs}, semester project, presentations \\
Assessment & As published (HEC 2017): sessional exam, home assignments,
quizzes, project, presentations, final exam \\
Set texts & As published: Connolly and Begg; Garcia-Molina, Ullman and Widom;
Silberschatz, Korth and Sudarshan; Ramakrishnan and Gehrke. No publication
years are given for any of the four, in any edition --- see Appendix~E \\
\bottomrule
\end{longtable}

\section*{Where the outline comes from}

\rowcolors{2}{white}{lightbg}
\begin{longtable}{@{}L{1.6cm}L{4.4cm}L{3.2cm}L{5.0cm}@{}}
\toprule
\rowcolor{primary!15}
\textbf{Edition} & \textbf{Status} & \textbf{Semester} &
\textbf{Course content published?} \\
\midrule
\endhead
\textbf{2017} & Computing Core --- compulsory in BSCS, BSSE \emph{and} BSIT &
IV, IV, \textbf{V} & \textbf{Yes} --- p.~88 \\
\textbf{2023} & Computing Core --- compulsory in all ten computing degrees and
the Associate Degree & \textbf{2} & \textbf{Yes} --- p.~58, \emph{the 2017 text
reprinted} \\
2025 & Major Compulsory --- but the edition covers only BS Computer Science
and the Associate Degree. \textbf{There is no BSIT in it at all} & III &
\textbf{No} --- deleted; five CLO bullets only \\
\bottomrule
\end{longtable}

\begin{alertbox}[Six defects in the published documents]
Recorded rather than inherited. Each is stated with the correction this
handout uses.

\begin{tightnum}
  \item \textbf{The syllabus is frozen.} The 2023 course outline is the 2017
        course outline word for word, with \textbf{one added full stop}
        (``vs'' becomes ``vs.'') across ninety words --- six years apart. Then
        2025 removes it entirely. A field that changed a great deal in that
        period is described by a text that did not change at all.
  \item \textbf{The outline normalises before it designs.} The clause list
        reaches \emph{normalization, functional dependencies, normal forms}
        \textbf{before} \emph{entity relationship model \dots\
        entity-relationship diagrams}. You cannot normalise a schema you have
        not designed. All four of HEC's own reference books teach the
        other order. \textbf{This handout reorders}: design in Part~II, map in
        Part~III, normalise in Part~V. The clause numbers in the table below
        show exactly where each one moved to.
  \item \textbf{One clause is printed twice.} ``attributes'' appears in the
        relational-model run and again in the entity-relationship run. Thirty
        three clauses, thirty two distinct. Both occurrences are traced below
        and the second is marked as a duplicate.
  \item \textbf{The prerequisite contradicts itself inside one booklet.} HEC
        2017's course description page says \texttt{Prerequisites: None}. All
        three of its scheme-of-studies tables --- BSCS p.~25, BSIT p.~31, BSSE
        p.~39 --- say \textbf{Data Structures \& Algorithms}. At 2nd semester
        the department cannot assume Data Structures, so \textbf{this handout
        assumes none}, and \chref{indexes} and \chref{transactions} are written
        to be self-contained because of it.
  \item \textbf{The 2023 learning-outcome table is broken in print.} CLO-3's
        text stops mid-sentence; its remainder (``anomalies by normalizing
        database tables.'') is printed as CLO-4's text; and the real CLO-4 is
        left in a trailing row with \emph{no number and no Bloom level}. This
        is a typesetting fault in the published national curriculum, not a
        misreading. The outcomes are reconstructed below from the 2017 text,
        which is intact.
  \item \textbf{A Bloom level is wrong.} CLO-3 reads ``Identify functional
        dependencies and \emph{resolve} database anomalies by normalizing
        database tables'' and is tagged \textbf{C2}. Identifying is C2;
        resolving and normalising are C3 at least. Every sibling course in
        this series complains that HEC left the Bloom column blank. Here it is
        filled in --- and filled in incorrectly, which is harder to notice.
\end{tightnum}
\end{alertbox}

\section*{Course content, traced clause by clause}

The outline is reproduced below exactly as published, in its published order,
split into its thirty-three clauses. The right-hand column is where each is
taught. Clauses that move are the subject of defect~2.

\rowcolors{2}{white}{lightbg}
\begin{longtable}{@{}L{0.9cm}L{9.3cm}L{4.4cm}@{}}
\toprule
\rowcolor{primary!15}
\textbf{\#} & \textbf{Clause, as published} & \textbf{Taught in} \\
\midrule
\endhead
1 & Basic database concepts & \chref{whynotfiles} \\
2 & Database approach vs file based system & \chref{whynotfiles} \\
3 & database architecture & \chref{whodoeswhat} \\
4 & three level schema architecture & \chref{architecture} \\
5 & data independence & \chref{architecture} \\
6 & relational data model & \chref{relational} \\
7 & attributes & \chref{relational} \\
8 & schemas & \chref{relational} \\
9 & tuples & \chref{relational} \\
10 & domains & \chref{relational} \\
11 & relation instances & \chref{relational} \\
12 & keys of relations & \chref{constraints} \\
13 & integrity constraints & \chref{constraints} \\
14 & relational algebra & \chref{algebra} \\
15 & selection & \chref{algebra} \\
16 & projection & \chref{algebra} \\
17 & Cartesian product & \chref{algebra} \\
18 & types of joins & \chref{joins} \\
19 & \textbf{normalization} & \chref{normalforms} \emph{(moved later)} \\
20 & \textbf{functional dependencies} & \chref{fds} \emph{(moved later)} \\
21 & \textbf{normal forms} & \chref{normalforms} \emph{(moved later)} \\
22 & \textbf{entity relationship model} & \chref{ermodel} \emph{(moved
earlier)} \\
23 & entity sets & \chref{ermodel} \\
24 & attributes & \textbf{duplicate of 7} \\
25 & relationship & \chref{ermodel}, \chref{cardinality} \\
26 & entity-relationship diagrams & \chref{cardinality}, \chref{mapping} \\
27 & Structured Query Language (SQL) & \chref{ddl}, \chref{select} \\
28 & Joins and sub-queries in SQL & \chref{joins}, \chref{subqueries} \\
29 & Grouping and aggregation in SQL & \chref{grouping} \\
30 & concurrency control & \chref{transactions} \\
31 & database backup and recovery & \chref{transactions} \\
32 & indexes & \chref{indexes} \\
33 & NoSQL systems & \chref{project} \\
\bottomrule
\end{longtable}

\begin{conceptbox}[Three topics taught that the outline does not name]
A reviewer will ask, so the justification is here rather than implied.

\begin{tight}
  \item \textbf{\chref{evolution}, changing a schema that is already full.}
        The outline never mentions \sqlkw{ALTER}. A student who cannot change
        a schema cannot carry one database through sixteen weeks, and the
        project requires exactly that.
  \item \textbf{\chref{views}, views and triggers.} Absent from the clause
        list. Required by CLO-2, which asks for \emph{logical and physical}
        schemas, and assumed by the department's own
        \emph{E-Commerce Application Development} (INFT-4041), which lists
        this course as its prerequisite --- see Appendix~F.
  \item \textbf{\chref{project}, security and authorisation.} Absent from 2017
        and 2023 entirely. \textbf{HEC 2025 requires it by name}: its fifth
        learning outcome is ``Analyze the role of databases in information
        systems and their security considerations''. The newest edition
        justifies the addition.
\end{tight}
\end{conceptbox}

\section*{Course Learning Outcomes}

\textbf{HEC 2017's four, verbatim}, with the Domain and Bloom levels exactly as
published --- including the one that is wrong.

\rowcolors{2}{white}{lightbg}
\begin{longtable}{@{}L{1.4cm}L{9.6cm}L{1.6cm}L{2.0cm}@{}}
\toprule
\rowcolor{primary!15}
\textbf{CLO} & \textbf{As published} & \textbf{Level} & \textbf{Taught in} \\
\midrule
\endhead
CLO-1 & ``Explain fundamental database concepts.'' & C2 & Part~I \\
CLO-2 & ``Design conceptual, logical and physical database schemas using
different data models.'' & C5 & Parts~II, III \\
CLO-3 & ``Identify functional dependencies and resolve database anomalies by
normalizing database tables.'' & \textbf{C2}\,$^\dagger$ & Part~V \\
CLO-4 & ``Use Structured Query Language (SQL) for database definition and
manipulation in any DBMS.'' & C4 & Parts~III, IV \\
\bottomrule
\end{longtable}

\noindent$^\dagger$\,\textbf{As published, and incorrect} --- see defect~6.
Assessed in this course at \textbf{C3}.

\medskip
\noindent\textbf{Four added by the department}, each discharging one of HEC
2025's five outcomes, which the 2017 clause list does not reach.

\rowcolors{2}{white}{lightbg}
\begin{longtable}{@{}L{1.4cm}L{1.3cm}L{8.3cm}L{3.6cm}@{}}
\toprule
\rowcolor{primary!15}
\textbf{CLO} & \textbf{Level} & \textbf{Outcome} & \textbf{Discharges} \\
\midrule
\endhead
CLO-5 & C3 & Write SQL for retrieval, insertion, update and deletion against a
schema you did not design & HEC 2025 outcome~2 \\
CLO-6 & C2 & Explain transactions, concurrency and recovery, and demonstrate
what a transaction prevents & HEC 2025 outcome~3 \\
CLO-7 & C3 & Implement and administer a working database, including backup and
restore & HEC 2025 outcome~4 \\
CLO-8 & C4 & Analyse a database's role in a wider system and its security
posture, and justify the analysis from evidence & HEC 2025 outcome~5 \\
\bottomrule
\end{longtable}

\newpage

%=====================================================================
\chapter*{Prerequisite Self-Check}
\addcontentsline{toc}{chapter}{Prerequisite Self-Check}
\markboth{Prerequisite Self-Check}{}

This is a \textbf{second-semester} course. It assumes one semester of
programming and nothing else --- in particular it does \emph{not} assume Data
Structures, whatever HEC 2017's scheme-of-studies tables say (defect~4).

Twelve questions. If you can answer ten, start at \chref{whynotfiles}. If you
cannot, the right-hand column says where to look --- and none of it will take
more than an evening.

\rowcolors{2}{white}{lightbg}
\begin{longtable}{@{}L{0.8cm}L{9.4cm}L{4.4cm}@{}}
\toprule
\rowcolor{primary!15}
\textbf{\#} & \textbf{Question} & \textbf{If not} \\
\midrule
\endhead
1 & What is the difference between a file and a folder? & --- \\
2 & A text file holds one record per line, fields separated by commas. How
would you find every line whose third field is ``2026''? & Programming
Fundamentals \\
3 & Write a loop that reads a file line by line. & Programming Fundamentals \\
4 & What does it mean for two values to be \emph{equal} when one is the number
7 and the other is the text ``7''? & \chref{ddl} returns to this \\
5 & What is an array index? What happens if you use one that is too large? &
Programming Fundamentals \\
6 & Given sets $A=\{1,2,3\}$ and $B=\{2,3,4\}$, write $A \cup B$, $A \cap B$
and $A - B$. & Appendix~B \\
7 & What does $A \times B$ mean for two sets, and how many elements does it
have? & Appendix~B \\
8 & Is the statement ``$x > 5$ and $x < 3$'' ever true? & Appendix~B \\
9 & What is the truth value of ``not (P and Q)'' when P is true and Q is
false? & Appendix~B \\
10 & Two people edit the same file at the same time and both save. What
happens to the first person's changes? & \chref{transactions} \\
11 & Roughly how many times do you have to halve 1{,}000{,}000 to get to 1? &
Appendix~B; \chref{indexes} needs it \\
12 & What is a function, in the mathematical sense --- how many outputs may one
input have? & Appendix~B; \chref{fds} depends on it \\
\bottomrule
\end{longtable}

\begin{conceptbox}[Question 12 is the one that matters]
Everything in Part~V rests on it. A \emph{functional dependency} is nothing
more than the statement that one set of attributes is a function of another:
give me a student's roll number and there is exactly one name it can be paired
with. If question~12 was uncomfortable, read Appendix~B before \chref{fds},
not during it.
\end{conceptbox}

%---------------------------------------------------------------------

\part{What a Database Is}
%=====================================================================
% PART I OPENER
%=====================================================================
\thispagestyle{plain}
\structuralfigures{P1.}

\begin{center}
\begin{tcolorbox}[enhanced, width=0.92\textwidth, colback=primary!5,
  colframe=primary, boxrule=0pt, leftrule=5pt, arc=3pt, top=8pt, bottom=8pt]
\large\itshape\color{primary!60!black}
Every program you have written so far has stored its data in a file, and that worked. Part~I is about the point at which it stops working --- not in theory, but on a particular afternoon when two people save at once, or a column turns out to mean two different things, or somebody deletes a row that something else was relying on.
\end{tcolorbox}
\end{center}

\vspace{6pt}

\dsfigh{%
\begin{tikzpicture}[
  cbox/.style={rounded corners=3pt, draw=primary, line width=1pt, fill=primary!7,
               text=primary!55!black, font=\small\bfseries, align=center,
               minimum width=3.5cm, minimum height=1.0cm},
  hbox/.style={rounded corners=3pt, draw=primary, line width=1.8pt, fill=primary!16,
               text=primary!55!black, font=\small\bfseries, align=center,
               minimum width=3.5cm, minimum height=1.0cm},
  arr/.style={-{Stealth[length=2.4mm]}, primary!60, line width=1pt},
  note/.style={font=\scriptsize\itshape, text=primary!55!black, align=center,
               fill=white, inner sep=2pt}
]
  \node[cbox] (c0) at (0.00,0) {Ch.\ 1\\Why Not Files};
  \node[cbox] (c1) at (4.30,0) {Ch.\ 2\\Three Levels};
  \node[cbox] (c2) at (8.60,0) {Ch.\ 3\\Who Does What};
  \node[hbox] (c3) at (12.90,0) {Ch.\ 4\\Your First Database};
  \draw[arr] (c0) -- (c1);
  \draw[arr] (c1) -- (c2);
  \draw[arr] (c2) -- (c3);
  \node[note] at (2.15,-1.0) {the problem};
  \node[note] at (6.45,-1.0) {the idea that solves it};
  \node[note] at (10.75,-1.0) {who makes it work};
  \node[font=\scriptsize\itshape, text=accent, align=center] at (12.90,1.15)
       {the database you keep for the rest of the course starts here};
\end{tikzpicture}}%
{Reading path through Part~I. Chapters 1 to 3 are the argument; Chapter 4 is
the first one you can run, and by the end of it there is a file on your disk
that every later chapter adds to.}{part1map}

\newpage

%=====================================================================
\chapter{Why Not Just Files?}\label{ch:whynotfiles}
%=====================================================================
\locator{1}{Part I --- What a Database Is}

\begin{outcomes}
By the end of this chapter you will be able to:
\begin{tight}
  \item \textbf{Describe} how a file-based application stores and finds data,
        and \textbf{state} what it costs to add one field.
  \item \textbf{Name} the six problems a file-based approach has, and
        \textbf{give} a concrete example of each.
  \item \textbf{Explain} what a database and a database management system are,
        and \textbf{distinguish} the two.
  \item \textbf{Judge} when a file is the right answer, which is more often
        than this chapter's title suggests.
  \item \textbf{Predict} which of the six problems a given scenario will hit
        first.
\end{tight}
\end{outcomes}

\begin{prereq}
One semester of programming: reading and writing a file, loops, and what a
record looks like when it is a line of comma-separated text. Nothing else.
Appendix~A if you have not checked that SQLite is on your machine.
\end{prereq}

\begin{keyterms}
data $\bullet$ information $\bullet$ file-based system $\bullet$ database
$\bullet$ database management system (DBMS) $\bullet$ data redundancy
$\bullet$ data inconsistency $\bullet$ data isolation $\bullet$ integrity
constraint $\bullet$ atomicity $\bullet$ concurrent access $\bullet$ lost
update $\bullet$ data security $\bullet$ application program
\end{keyterms}

\section{The Program You Have Already Written}

Somewhere in your first semester you wrote a program that kept a list of
things in a file. It probably looked like this, and there is nothing wrong
with it:

\begin{lstlisting}[language=python]
# students.csv
# 2021-IT-001,Ayesha Khan,BSIT,3
# 2021-IT-002,Bilal Ahmed,BSIT,3

def find(roll_no):
    for line in open("students.csv"):
        fields = line.strip().split(",")
        if fields[0] == roll_no:
            return fields
    return None
\end{lstlisting}

It works. It is twelve lines. For a list of your own books, or a configuration
file, or the output of a program that nobody else reads, it is the right
answer and this chapter is not an argument against it.

This chapter is about the afternoon it stops being the right answer --- and
about the fact that it does not stop gradually. A file-based approach works,
and works, and works, and then the department asks for one more thing and
suddenly nothing works.

\begin{definitionbox}[Data, Information, and the difference]
\term{Data} are raw recorded facts: the characters \texttt{2021-IT-001} in a
file.

\smallskip
\term{Information} is data in a context that makes it answer a question:
\emph{the roll number of the student who scored highest in Database Systems
last term}.

\smallskip
A \term{file-based system} stores data and leaves every question to the
application program. A \term{database} stores data \emph{together with a
description of what it means and what must be true of it}, so that some of the
questions --- and all of the guarantees --- can be answered by the store
itself.

\smallskip
That last sentence is the whole subject. Everything in the remaining
twenty-two chapters is a particular case of it.
\end{definitionbox}

\section{Six Things That Go Wrong}

These are not six opinions. Each is a specific failure with a name, and you
will meet all six in a working life.

\dsfig{%
\begin{tikzpicture}[
  app/.style={rounded corners=3pt, draw=#1, line width=1.0pt, fill=#1!8,
              text=#1!60!black, font=\bfseries\scriptsize, align=center,
              text width=2.5cm, minimum height=0.9cm},
  fil/.style={draw=neutral, line width=0.9pt, fill=neutral!8,
              text=neutral, font=\scriptsize, align=center,
              text width=2.3cm, minimum height=0.8cm},
  arr/.style={-{Stealth[length=2.0mm]}, neutral!70, line width=0.8pt},
  note/.style={font=\scriptsize\itshape, text=accent, align=center,
               text width=5.0cm}
]
  % --- file-based, on the left
  \node[font=\bfseries\small, text=accent] at (3.0,3.5) {File-based};
  \node[app=accent] (a1) at (0.6,2.5) {Enrolment\\program};
  \node[app=accent] (a2) at (3.0,2.5) {Results\\program};
  \node[app=accent] (a3) at (5.4,2.5) {Fees\\program};
  \node[fil] (f1) at (0.6,0.9) {students.csv};
  \node[fil] (f2) at (3.0,0.9) {marks.csv};
  \node[fil] (f3) at (5.4,0.9) {fees.csv};
  \draw[arr] (a1) -- (f1);
  \draw[arr] (a2) -- (f2);
  \draw[arr] (a3) -- (f3);
  \node[note] at (3.0,-0.3)
       {Each program owns its own file and its own idea of what a student is.
        The name is stored three times.};

  % --- database, on the right
  \node[font=\bfseries\small, text=greenhl] at (11.4,3.5) {Database};
  \node[app=greenhl] (b1) at (9.0,2.5) {Enrolment\\program};
  \node[app=greenhl] (b2) at (11.4,2.5) {Results\\program};
  \node[app=greenhl] (b3) at (13.8,2.5) {Fees\\program};
  \node[rounded corners=3pt, draw=greenhl, line width=1.2pt, fill=greenhl!14,
        text=greenhl!60!black, font=\bfseries\scriptsize, align=center,
        text width=5.6cm, minimum height=0.9cm] (dbms) at (11.4,1.3) {DBMS};
  \node[fil, text width=5.6cm, draw=greenhl] (db) at (11.4,0.1)
       {one database};
  \foreach \b in {b1,b2,b3}{\draw[arr] (\b) -- (dbms);}
  \draw[arr] (dbms) -- (db);
  \node[note, text=greenhl!60!black] at (11.4,-1.1)
       {One store, one description of what a student is, and one place the
        rules are kept.};
\end{tikzpicture}}%
{The same three programs, two ways. The difference is not that the right-hand
side has fewer files --- it is that the rules about what a student is have
moved out of the programs and into one place.}{filesvsdb}

\begin{alertbox}[The six, with a case for each]
\begin{tightnum}
  \item \textbf{Data redundancy.} The student's name is in
        \texttt{students.csv}, in \texttt{marks.csv} so the results sheet can
        print it, and in \texttt{fees.csv} so the receipt can. Three copies,
        and disk space is the least of it.
  \item \textbf{Data inconsistency.} A student marries and changes her name.
        Two of the three files are updated. The college now holds two different
        answers to one question, and \emph{nothing in the system is wrong} ---
        each file is internally fine. This is redundancy's actual cost.
  \item \textbf{Data isolation.} The question ``which students owe fees and
        are enrolled in a course that has already started?'' needs two files in
        different formats. There is no way to ask it; there is only a program
        somebody has to write.
  \item \textbf{Integrity problems.} ``A student's semester is between 1 and
        8.'' Where does that rule live? In the enrolment program. And in the
        results program, where somebody wrote $1$ to $10$. The rule is a
        comment in three places and enforced in none.
  \item \textbf{Atomicity problems.} Transferring a student between sections
        means removing a line from one file and adding it to another. The
        machine loses power in between. The student is now in no section at
        all, and nothing can tell you that happened.
  \item \textbf{Concurrent access and security.} Two clerks open the file, both
        edit, both save; the second overwrites the first and neither is told
        --- a \term{lost update}. And the fees clerk, who needs one column, has
        the whole file including every home address.
\end{tightnum}
\end{alertbox}

\begin{conceptbox}[Which one bites first, and why it is never the one you
expect]
Students asked to rank these almost always put redundancy first, because it is
the one that is visible --- you can \emph{see} the duplicated column.

\smallskip
In practice the order is almost exactly reversed. Redundancy costs disk, which
is cheap. \textbf{Inconsistency} is what gets noticed, because somebody
receives a letter addressed to a name they no longer use. \textbf{Isolation}
is what gets the department to pay for a new system, because the question they
want answered cannot be answered at all. And \textbf{atomicity} is the one
nobody finds until the day it matters, because a half-finished transfer looks
exactly like a student who was never enrolled.

\smallskip
Keep this ranking in mind in \chref{project}, when you have to justify your
design to somebody who will ask what problem it solves.
\end{conceptbox}

\section{What a Database Management System Is}

\begin{definitionbox}[Database, DBMS, Database System]
A \term{database} is a shared, organised collection of data, together with a
description of that data --- what it is called, what type it is, what must be
true of it, and how its parts relate.

\smallskip
A \term{database management system} (DBMS) is the \emph{software} that creates,
maintains and controls access to a database. SQLite is a DBMS; so are Oracle,
PostgreSQL, MySQL and SQL Server.

\smallskip
A \term{database system} is the whole arrangement: the data, the DBMS, the
application programs that use it, and the people.

\smallskip
These three are routinely used as if they meant the same thing, including by
people who should know better. In an examination they do not, and the
distinction is CLO-1's business.
\end{definitionbox}

The DBMS is what the right-hand side of Figure~\ref{fig:filesvsdb} adds. It is
a single piece of software that every program goes through, and because
everything goes through it, it is the one place where a rule can be stated once
and be true for everybody.

\begin{examplebox}[What moves into the DBMS, problem by problem]
\begin{tabular}{@{}L{4.6cm}L{9.6cm}@{}}
\textbf{Redundancy} & One copy of each fact, and other tables refer to it
(\chref{mapping}) \\
\textbf{Inconsistency} & Follows from the above --- there is nothing to
disagree with \\
\textbf{Isolation} & One query language over all the data (\chref{algebra}
onward) \\
\textbf{Integrity} & Constraints declared once, in the schema
(\chref{constraints}) \\
\textbf{Atomicity} & Transactions: all of it or none of it
(\chref{transactions}) \\
\textbf{Concurrency, security} & Locking, and authorisation
(\chref{transactions}, \chref{project}) \\
\end{tabular}

\smallskip
That table is this book's table of contents, rearranged. Every chapter from
here is one row of it, done properly.
\end{examplebox}

\section{When a File Is Still the Right Answer}

A course called Database Systems has an obvious bias, and you should know
about it.

\begin{conceptbox}[Use a file when every one of these is true]
\begin{tight}
  \item \textbf{One program reads it, and one writes it.} No concurrency, so
        no lost updates.
  \item \textbf{It is small enough to read entirely.} If you always want all of
        it, an index buys nothing.
  \item \textbf{The questions never change.} Isolation only hurts when somebody
        asks something new.
  \item \textbf{Losing it is survivable}, or it can be regenerated.
\end{tight}

\smallskip
Configuration files, log files, a cache, the output of a nightly job --- all
files, all correctly so. \textbf{The failure is not using a file; it is using a
file after the fourth condition stopped holding and not noticing.}
\end{conceptbox}

\begin{pitfall}
\begin{tight}
  \item \textbf{Treating redundancy as the main problem.} It is the cheapest of
        the six and the easiest to see, which is why it gets the attention.
  \item \textbf{Saying ``database'' when you mean ``DBMS''.} The data is not
        the software.
  \item \textbf{Believing a DBMS removes the need to think about integrity.}
        It gives you somewhere to put the rules. It does not write them, and
        \chref{constraints} measures what happens when you declare one and it
        is not enforced.
  \item \textbf{Assuming a database is always better.} Four conditions above
        say otherwise.
  \item \textbf{Calling a spreadsheet a database.} It has rows and columns and
        none of the six guarantees.
\end{tight}
\end{pitfall}

\begin{lab}[Break a file-based system on purpose]
No database yet. This laboratory is about earning the next twenty-two
chapters.

\begin{lstlisting}[language=bash]
# --- 1. Make the files -----------------------------------------------
# Three CSV files, by hand, in a text editor. Five students each.
#   students.csv   roll_no,name,programme,semester
#   marks.csv      roll_no,name,course,marks
#   fees.csv       roll_no,name,amount,paid

# --- 2. Count the redundancy -----------------------------------------
# How many times does one student's NAME appear across the three files?
# Multiply by 4,000 students. Write the number down.

# --- 3. Create an inconsistency --------------------------------------
# Change one student's name in students.csv and marks.csv but NOT in
# fees.csv. Now answer: what is that student's name? Justify it.

# --- 4. Hit data isolation -------------------------------------------
# Write the program that answers: "which students owe fees AND are
# enrolled in a course where they scored below 50?"
# Time yourself. Do not finish it if it takes more than twenty minutes
# -- stopping is the result.

# --- 5. Break atomicity ----------------------------------------------
# Write a script that moves a student from one file to another, with a
# deliberate crash between the two writes. Run it. Describe the state.

# --- 6. Simulate a lost update ---------------------------------------
# Open students.csv in two editors. Change a different row in each.
# Save both. Report what survived.

# --- 7. Write the memo ------------------------------------------------
# One page to a head of department who has asked why the existing
# system should be replaced. Use YOUR numbers from steps 2 and 4.
# No jargon: the reader has not taken this course.
\end{lstlisting}

\smallskip
\textbf{What to notice.} Step 4 is the one that convinces people. Redundancy is
an argument; a question that cannot be answered at all is a business problem,
and it is the one that gets a project funded.
\end{lab}

\begin{checkpoint}
\begin{tightnum}
  \item A college keeps student records in one file and library loans in
        another, each maintained by one clerk who never touches the other.
        Which of the six problems does this arrangement have, and which does it
        genuinely escape? Say what would have to change for your answer to
        change.
  \item Distinguish a database from a database management system, and explain
        why ``we bought a database'' is a sentence that needs clarifying before
        you can act on it.
  \item Your colleague argues that redundancy is harmless because disk is
        cheap. Give the argument you would make in reply, name the problem
        redundancy actually causes, and give an example in which no single file
        is wrong and the system still is.
\end{tightnum}
\end{checkpoint}

\begin{chaptersummary}
\begin{tight}
  \item Data are recorded facts; information is data in a context that answers
        a question. A database stores data \emph{with a description of what it
        means and what must be true of it}, which is the whole subject.
  \item A file-based system works until it does not, and the change is abrupt
        rather than gradual.
  \item The six problems are redundancy, inconsistency, isolation, integrity,
        atomicity, and concurrent access with security.
  \item Redundancy is the visible one and the cheapest. Inconsistency is what
        gets noticed, isolation is what gets a project funded, and atomicity is
        the one nobody finds until it matters.
  \item A \emph{database} is the data; a \emph{DBMS} is the software; a
        \emph{database system} is the whole arrangement including the people.
  \item Everything goes through the DBMS, which is why it is the one place a
        rule can be stated once and be true for everybody.
  \item A file is still right when there is one reader, one writer, a fixed set
        of questions and nothing irreplaceable. The failure is not noticing
        when that stops being true.
  \item The table in Section 1.3 is this book's contents rearranged: each of
        the six problems is solved by a named later chapter.
\end{tight}
\end{chaptersummary}

\begin{reviewq}
  \item \textbf{(MCQ)} Two files hold a student's name and only one is updated.
        This is an example of: (a)~redundancy; (b)~inconsistency; (c)~isolation;
        (d)~an integrity constraint.
  \item \textbf{(MCQ)} The software that creates and controls access to a
        database is called: (a)~the database; (b)~the database system; (c)~the
        DBMS; (d)~the schema.
  \item \textbf{(Short)} Give a question about a university that cannot be
        answered from two separate CSV files, and say which of the six problems
        that is.
  \item \textbf{(Short)} Explain atomicity using an example that is not a bank
        transfer.
  \item \textbf{(Short)} State two conditions under which a plain file is a
        better choice than a database, and give a real example of each.
  \item \textbf{(Applied)} A clinic keeps appointments in a spreadsheet shared
        on a network drive, edited by three receptionists. Identify every one of
        the six problems it has, rank them by which will cause a complaint
        first, and justify your ranking.
  \item \textbf{(Applied)} You are asked to justify replacing a working
        file-based system. Write the three strongest arguments in order, say
        what evidence you would gather for each before the meeting, and state
        which argument you would drop if you only had five minutes.
\end{reviewq}

%=====================================================================
\chapter{Three Levels, and Why They Are Separate}\label{ch:architecture}
%=====================================================================
\locator{1}{Part I --- What a Database Is}

\begin{outcomes}
By the end of this chapter you will be able to:
\begin{tight}
  \item \textbf{Name} the three levels of the ANSI/SPARC architecture and
        \textbf{state} what each one describes.
  \item \textbf{Distinguish} a schema from an instance, and \textbf{explain}
        why the difference matters when something changes.
  \item \textbf{Define} logical and physical data independence, and
        \textbf{give} an example of a change that each one protects you from.
  \item \textbf{Identify} which level a given change belongs to.
  \item \textbf{Explain} why a system with no separation of levels is the
        file-based system of \chref{whynotfiles} again.
\end{tight}
\end{outcomes}

\begin{prereq}
\chref{whynotfiles}: the six problems, and what a DBMS is. Nothing else. This
chapter is the one piece of pure architecture in the course and it is here
because every later chapter quietly assumes it.
\end{prereq}

\begin{keyterms}
ANSI/SPARC architecture $\bullet$ external level $\bullet$ external schema
$\bullet$ user view $\bullet$ conceptual level $\bullet$ conceptual schema
$\bullet$ internal level $\bullet$ internal schema $\bullet$ physical level
$\bullet$ schema $\bullet$ instance $\bullet$ mapping $\bullet$ logical data
independence $\bullet$ physical data independence
\end{keyterms}

\section{One Database, Three Descriptions}

In \chref{whynotfiles} the fees program and the results program each had their
own file and their own idea of what a student was. Putting them behind one
DBMS fixes the duplication --- but it creates a new problem immediately.

The fees clerk wants a student to be a roll number, a name and an amount
owing. The examinations officer wants a roll number, a name and a list of
marks. The person who has to make the database fast wants to know that the
student table is sorted by roll number and has an index on surname. All three
are talking about the same data and none of them should have to care about the
other two.

The answer, standardised in 1975 and unchanged since, is to describe the
database \emph{three times}.

\begin{definitionbox}[The three levels]
The \term{external level} is what a particular user or program sees: their own
tailored view of the part of the database they need, in the shape they need
it. There are many external schemas, one per group of users.

\smallskip
The \term{conceptual level} is the whole database, described once, logically:
every entity, every attribute, every relationship and every constraint ---
and nothing at all about how it is stored. There is exactly one conceptual
schema.

\smallskip
The \term{internal level} --- also called the \term{physical level} --- is how it is actually stored: files, pages, record
layouts, indexes, compression. There is exactly one internal schema.

\smallskip
A \term{mapping} connects adjacent levels, so that the DBMS can translate a
request expressed against one into an operation on the next.
\end{definitionbox}

\dsfig{%
\begin{tikzpicture}[
  lvl/.style={rounded corners=3pt, draw=#1, line width=1.1pt, fill=#1!8,
              text=#1!60!black, font=\bfseries\small, align=center,
              text width=3.3cm, minimum height=1.1cm},
  ext/.style={rounded corners=2pt, draw=secondary, line width=0.9pt,
              fill=secondary!8, text=secondary!60!black,
              font=\scriptsize\bfseries, align=center,
              text width=2.5cm, minimum height=0.85cm},
  note/.style={font=\scriptsize, text=neutral, align=left, text width=5.4cm,
               anchor=west},
  arr/.style={-{Stealth[length=2.2mm]}, primary!55, line width=0.9pt},
  map/.style={font=\scriptsize\itshape, text=accent, align=center}
]
  \node[ext] (e1) at (-2.9,3.6) {Fees clerk};
  \node[ext] (e2) at (0.0,3.6)  {Examinations};
  \node[ext] (e3) at (2.9,3.6)  {Student portal};
  \node[lvl=primary]  (c) at (0,1.7) {Conceptual schema\\{\scriptsize
       \mdseries the whole database, once}};
  \node[lvl=purplehl] (i) at (0,-0.2) {Internal schema\\{\scriptsize
       \mdseries pages, records, indexes}};
  \node[draw=neutral, line width=0.9pt, fill=neutral!8, text=neutral,
        font=\scriptsize, align=center, text width=3.3cm,
        minimum height=0.6cm] (d) at (0,-1.7) {the stored database};

  \foreach \e in {e1,e2,e3}{\draw[arr] (\e) -- (c);}
  \draw[arr] (c) -- (i);
  \draw[arr] (i) -- (d);

  \node[map] at (3.6,2.65) {external/conceptual\\mapping};
  \node[map] at (3.6,0.75) {conceptual/internal\\mapping};

  \node[note, text=secondary!60!black] at (4.9,3.6)
       {Many of these. Each sees only what it needs, shaped how it needs it.};
  \node[note, text=primary!60!black] at (4.9,1.7)
       {\textbf{Exactly one.} Every entity, attribute, relationship and
        constraint --- and nothing about storage.};
  \node[note, text=purplehl!60!black] at (4.9,-0.2)
       {Exactly one. How it is really kept. Change this and nothing above
        should notice.};
\end{tikzpicture}}%
{The ANSI/SPARC three-level architecture. The two mappings are the whole
value of the arrangement: because they exist, a change at one level can be
absorbed instead of propagated.}{threelevels}

\section{Schema and Instance}

Two words that students use interchangeably for about a fortnight, and then
cannot afford to.

\begin{definitionbox}[Schema, Instance]
A \term{schema} is the \emph{description}: the names of the tables, their
columns, the types, the constraints. It is set when the database is designed
and changes rarely.

\smallskip
An \term{instance} (or \emph{state}, or \emph{extension}) is the data
\emph{in} the database at one moment. It changes every time anybody inserts a
row.

\smallskip
The analogy worth keeping: a schema is a variable's declaration, an instance
is its value. \texttt{int count;} is the schema; \texttt{7} is the instance.
\end{definitionbox}

\begin{examplebox}[The same distinction, three ways]
\begin{tabular}{@{}L{4.0cm}L{5.0cm}L{5.0cm}@{}}
& \textbf{Schema} & \textbf{Instance} \\
Programming & \texttt{int count;} & \texttt{7} \\
A database & \relation{Student}{\relkey{roll\_no}, name, programme} &
the 4{,}000 rows in it today \\
A form & the blank form & one filled-in copy \\
\end{tabular}

\smallskip
When somebody says ``the database changed'', the first question is always
which of the two they mean, because the answers are days apart in effort.
\end{examplebox}

\section{Data Independence, Which Is the Point of All This}

Three descriptions and two mappings are a lot of machinery. It buys exactly
one thing, and the thing is worth it.

\begin{definitionbox}[Logical and physical data independence]
\term{Physical data independence} is the ability to change the \emph{internal}
schema --- how data is stored --- without changing the conceptual schema or
any application program. Adding an index, moving to a different disk,
reorganising the file: none of these should require a line of application code
to change.

\smallskip
\term{Logical data independence} is the ability to change the
\emph{conceptual} schema --- adding a table, adding a column, splitting one
table into two --- without changing the external schemas or the programs that
use them.

\smallskip
Physical independence is easy and essentially achieved. \textbf{Logical
independence is hard and only partly achieved}, which is why it is worth
understanding where it fails.
\end{definitionbox}

\begin{conceptbox}[Which level does this change belong to?]
A question worth being able to answer instantly, because the cost of a change
is decided by its level.

\begin{tight}
  \item \textbf{``Add an index on surname.''} Internal. Nothing else changes.
        \chref{indexes} does exactly this and measures it.
  \item \textbf{``Add a telephone number to students.''} Conceptual. Existing
        views that never mentioned it carry on working --- this is logical
        independence succeeding.
  \item \textbf{``Split the student table into student and student-contact.''}
        Conceptual, and this is where logical independence is \emph{tested}:
        \chref{views} shows how a view can hide the split from programs that
        still want one table. \chref{normalforms} does this split for real.
  \item \textbf{``The fees clerk should not see home addresses.''} External. A
        view, and in a server DBMS a privilege (\chref{project}).
  \item \textbf{``Store dates as text rather than numbers.''} Looks internal.
        It is \emph{conceptual}, because it changes what comparisons mean ---
        and \chref{ddl} measures what SQLite does about that.
\end{tight}

\smallskip
The last one is the one that catches people, and it catches them because
``how it is stored'' sounds internal in English and is not.
\end{conceptbox}

\section{Where SQLite's Three Levels Are}

An architecture standardised in 1975 is worth nothing to you unless you can
point at its three levels in the tool on your own machine, which at first
glance appears to have only one.

\begin{alertbox}[All three are there; only two can be enforced]
It is worth knowing how an architecture from 1975 maps onto the tool in front
of you, because at first glance SQLite seems to have only one level.

\begin{tight}
  \item \textbf{Conceptual:} your \sqlkw{CREATE TABLE} statements. Exactly one
        set of them, and \texttt{.schema} prints it.
  \item \textbf{Internal:} the page structure, the B-trees, the indexes.
        \chref{firstdatabase} measures it directly and you never write it ---
        SQLite decides, and \sqlkw{PRAGMA} is the only window onto it.
  \item \textbf{External:} views (\chref{views}). SQLite has them, so the
        level exists --- but it has no users, so it cannot restrict anybody
        \emph{to} a view. The level is present and the enforcement is not,
        which is \chref{project}'s subject.
\end{tight}

\smallskip
\textbf{This is a real limitation and worth stating now rather than in week
sixteen.} You will build all three levels in this course. You will only be
able to \emph{enforce} two of them.
\end{alertbox}

\begin{pitfall}
\begin{tight}
  \item \textbf{Confusing schema with instance.} A schema change is a project;
        an instance change is a Tuesday.
  \item \textbf{Thinking there are many conceptual schemas.} There is one. Many
        external, one conceptual, one internal.
  \item \textbf{Believing physical independence means performance does not
        matter.} It means you can change storage \emph{without rewriting
        programs}. It does not mean the choice is free --- \chref{indexes}
        measures exactly how unfree it is.
  \item \textbf{Calling a storage-format change internal when it changes
        meaning.} If a comparison starts giving different answers, it was
        conceptual.
  \item \textbf{Assuming logical independence is automatic.} It is the hard
        one, and \chref{views} is how it is bought.
\end{tight}
\end{pitfall}

\begin{lab}[Find the three levels in a database you did not design]
\begin{lstlisting}[language=bash]
# --- 1. Make a database with something in it -------------------------
sqlite3 arch.db
\end{lstlisting}
then, in the shell:
\begin{lstlisting}[language=sql]
-- --- 2. State a conceptual schema -------------------------------
CREATE TABLE student (
  roll_no   TEXT PRIMARY KEY,
  name      TEXT NOT NULL,
  programme TEXT NOT NULL,
  semester  INTEGER NOT NULL
);
INSERT INTO student VALUES
  ('2021-IT-001','Ayesha Khan','BSIT',2),
  ('2021-IT-002','Bilal Ahmed','BSIT',2),
  ('2021-CS-007','Hira Nadeem','BSCS',4);

-- --- 3. Read the conceptual schema back -------------------------
.schema student
-- That text IS the conceptual schema. Nothing in it mentions storage.

-- --- 4. Build an external schema --------------------------------
CREATE VIEW fees_view AS
  SELECT roll_no, name FROM student;
SELECT * FROM fees_view;
-- A user given only this cannot see "programme" or "semester".

-- --- 5. Look at the internal schema -----------------------------
PRAGMA page_size;
PRAGMA page_count;
SELECT name, type FROM sqlite_master;
-- You did not write any of this and cannot usefully change it.

-- --- 6. Change the internal level and prove nothing else moved --
CREATE INDEX idx_name ON student(name);
SELECT * FROM fees_view;
-- Same answer. That is physical data independence, demonstrated.

-- --- 7. Now change the conceptual level -------------------------
ALTER TABLE student ADD COLUMN email TEXT;
SELECT * FROM fees_view;
-- Still works, and still has two columns. Logical independence.

-- --- 8. Break it ------------------------------------------------
-- Try to drop the column the view depends on:
ALTER TABLE student DROP COLUMN name;
-- Write down exactly what happens, and which level failed.
\end{lstlisting}

\smallskip
\textbf{What to notice.} Steps 6 and 7 are the two kinds of independence, each
demonstrated in one line. Step 8 is the limit of the second one, and the error
message is worth reading twice.
\end{lab}

\begin{checkpoint}
\begin{tightnum}
  \item Classify each of these by level, and say which kind of data
        independence (if any) protects the programs from it: adding an index;
        adding a column; renaming a column; storing the database on a faster
        disk; giving the library staff a view with no home addresses.
  \item Explain the difference between a schema and an instance to somebody who
        has written programs but never used a database, using an example that
        is not a database.
  \item SQLite has views but no users. Say precisely which of the three levels
        is affected by that, what you can still do, what you cannot, and how
        you would achieve the same effect in an application built on SQLite.
\end{tightnum}
\end{checkpoint}

\begin{chaptersummary}
\begin{tight}
  \item The ANSI/SPARC architecture describes one database three times:
        \emph{many} external schemas, \emph{one} conceptual schema, \emph{one}
        internal schema, joined by two mappings.
  \item The external level is what a user sees; the conceptual level is the
        whole database logically and says nothing about storage; the internal
        level is how it is really kept.
  \item A \emph{schema} is the description and changes rarely; an
        \emph{instance} is the data and changes constantly. A schema is a
        declaration, an instance is a value.
  \item \textbf{Physical data independence} lets the internal schema change
        without touching programs. It is essentially achieved.
  \item \textbf{Logical data independence} lets the conceptual schema change
        without touching external schemas. It is hard, partial, and bought
        with views (\chref{views}).
  \item The level a change belongs to decides what it costs. ``Store dates as
        text'' sounds internal and is conceptual, because it changes what
        comparisons mean.
  \item In SQLite the conceptual level is your \sqlkw{CREATE TABLE}
        statements, the internal level is pages and B-trees you never write,
        and the external level is views --- which exist but cannot be
        enforced, because there are no users to restrict.
\end{tight}
\end{chaptersummary}

\begin{reviewq}
  \item \textbf{(MCQ)} How many conceptual schemas does a database have?
        (a)~one per user; (b)~one per table; (c)~exactly one; (d)~one per
        application.
  \item \textbf{(MCQ)} Adding an index without changing any program
        demonstrates: (a)~logical data independence; (b)~physical data
        independence; (c)~an external schema; (d)~an instance change.
  \item \textbf{(Short)} Give one change that logical data independence
        protects you from and one that it does not, and say what makes the
        difference.
  \item \textbf{(Short)} Why does the architecture need \emph{mappings} as well
        as schemas? What would be lost without them?
  \item \textbf{(Short)} A colleague says the internal level does not matter
        because of physical data independence. Reply in two sentences.
  \item \textbf{(Applied)} A university wants three groups to use one student
        database: admissions, examinations, and the students themselves. Design
        the external schemas --- what each sees and what each must not --- and
        say which parts of your design SQLite can enforce and which it cannot.
  \item \textbf{(Applied)} Your conceptual schema must change: the single
        \texttt{address} column becomes four columns. Describe how you would
        make that change without breaking existing programs, what the view
        would look like, what it would cost, and at what point you would decide
        the view is no longer worth keeping.
\end{reviewq}

%=====================================================================
\chapter{What Is Inside a DBMS, and Who Does What}\label{ch:whodoeswhat}
%=====================================================================
\locator{1}{Part I --- What a Database Is}

\begin{outcomes}
By the end of this chapter you will be able to:
\begin{tight}
  \item \textbf{Name} the major components of a DBMS and \textbf{state} what
        each one is responsible for.
  \item \textbf{Trace} a single \sqlkw{SELECT} statement through those
        components, in order, from text to answer.
  \item \textbf{Explain} what the data dictionary is, and \textbf{query} it.
  \item \textbf{Distinguish} the four database roles --- administrator,
        designer, application programmer, end user --- by what each one
        decides.
  \item \textbf{Identify} which components SQLite has, which it lacks, and
        which role has to cover for the missing ones.
\end{tight}
\end{outcomes}

\begin{prereq}
\chref{architecture}: the three levels and the two mappings. This chapter is
largely about \emph{who performs} those mappings.
\end{prereq}

\begin{keyterms}
query processor $\bullet$ parser $\bullet$ query optimiser $\bullet$ execution
engine $\bullet$ storage manager $\bullet$ buffer manager $\bullet$ file
manager $\bullet$ transaction manager $\bullet$ lock manager $\bullet$ recovery
manager $\bullet$ authorisation manager $\bullet$ data dictionary $\bullet$
system catalog $\bullet$ metadata $\bullet$ DBA $\bullet$ database designer
$\bullet$ application programmer $\bullet$ end user
\end{keyterms}

\section{A DBMS Is Not One Program}

You type eleven words:

\begin{lstlisting}[language=sql]
SELECT name FROM student WHERE programme = 'BSIT' ORDER BY name;
\end{lstlisting}

and rows come back. Between those two events, six or seven distinct pieces of
software run, each with one job. It is worth knowing them by name, because for
the rest of the course every topic is really a chapter about one of them:
\chref{indexes} is about the optimiser and the storage manager arguing,
\chref{transactions} is about the transaction and recovery managers, and
\chref{project} is about the one component SQLite does not have at all.

\begin{definitionbox}[The components of a DBMS]
The \term{query processor} turns your statement into an answer. It contains a
\term{parser} (is this valid SQL, and what does it refer to?), a
\term{query optimiser} (of the many ways to compute this, which is cheapest?)
and an \term{execution engine} (do it).

\smallskip
The \term{storage manager} stands between the query processor and the disk. It
contains a \term{file manager} (where on disk is this?) and a
\term{buffer manager} (which pages are in memory, and which do I evict?).

\smallskip
The \term{transaction manager} makes groups of operations behave as one, with
a \term{lock manager} to keep concurrent users from corrupting each other's
work and a \term{recovery manager} to put the database back together after a
crash.

\smallskip
The \term{authorisation manager} decides whether this user is allowed to do
this at all.

\smallskip
The \term{data dictionary} --- also called the \term{system catalog} --- is
where the DBMS keeps the schema itself. It is a database about the database:
\term{metadata}, data about data.
\end{definitionbox}

\dsfig{%
\begin{tikzpicture}[
  comp/.style={rounded corners=3pt, draw=#1, line width=1pt, fill=#1!7,
               text=#1!60!black, font=\scriptsize\bfseries, align=center,
               text width=2.7cm, minimum height=0.8cm},
  grp/.style={rounded corners=4pt, draw=#1!40, line width=0.8pt, dashed,
              inner sep=6pt},
  ttl/.style={font=\scriptsize\bfseries, text=#1, anchor=west},
  arr/.style={-{Stealth[length=2mm]}, primary!55, line width=0.9pt},
  thin/.style={-{Stealth[length=1.6mm]}, neutral!70, line width=0.6pt}
]
  % users
  \node[comp=secondary] (u) at (0,5.0) {SQL from a user\\or a program};

  % query processor
  \node[comp=primary] (p)  at (0,3.7) {Parser};
  \node[comp=primary] (o)  at (0,2.7) {Optimiser};
  \node[comp=primary] (x)  at (0,1.7) {Execution engine};
  \node[grp=primary, fit=(p)(o)(x)] (qp) {};
  \node[ttl=primary] at ([xshift=2pt]qp.north east) {Query processor};

  % storage manager
  \node[comp=purplehl] (b) at (0,0.3) {Buffer manager};
  \node[comp=purplehl] (f) at (0,-0.7) {File manager};
  \node[grp=purplehl, fit=(b)(f)] (sm) {};
  \node[ttl=purplehl] at ([xshift=2pt]sm.north east) {Storage manager};

  \node[comp=neutral] (disk) at (0,-2.1) {the database file};

  % side services
  \node[comp=tealhl]  (t) at (-4.3,2.2) {Transaction\\manager};
  \node[comp=tealhl]  (l) at (-4.3,0.9) {Lock manager};
  \node[comp=tealhl]  (r) at (-4.3,-0.4) {Recovery\\manager};
  \node[comp=accent]  (a) at (-4.3,3.9) {Authorisation\\manager};
  \node[comp=goldhl]  (d) at (4.3,1.7) {Data dictionary\\{\scriptsize\mdseries
        the schema, stored}};

  \draw[arr] (u) -- (p);
  \draw[arr] (p) -- (o);
  \draw[arr] (o) -- (x);
  \draw[arr] (x) -- (b);
  \draw[arr] (b) -- (f);
  \draw[arr] (f) -- (disk);

  \draw[thin] (a) -- (p);
  \draw[thin] (t) -- (x);
  \draw[thin] (l) -- (b);
  \draw[thin] (r) -- (f);
  \draw[thin, {Stealth[length=1.6mm]}-{Stealth[length=1.6mm]}] (d) -- (o);
  \draw[thin, {Stealth[length=1.6mm]}-{Stealth[length=1.6mm]}] (d) -- (p);
\end{tikzpicture}}%
{The components of a DBMS, and the path a statement takes down through them.
The optimiser and the parser both consult the data dictionary --- the parser
to check that \texttt{student} exists and has a \texttt{name}, the optimiser to
find out how big it is and what indexes it has.}{dbmscomponents}

\section{Following One Statement Down}

\begin{worked}[Eleven words, six components]
\texttt{SELECT name FROM student WHERE programme = 'BSIT' ORDER BY name;}

\begin{tightnum}
  \item \textbf{Authorisation manager.} May this user read
        \texttt{student}? In a server DBMS this is checked first and can stop
        everything here. \emph{In SQLite there is no such check} ---
        if you can open the file, you can read all of it.
  \item \textbf{Parser.} Is the syntax legal? Does \texttt{student} exist?
        Does it have columns \texttt{name} and \texttt{programme}? Is
        \texttt{'BSIT'} comparable with \texttt{programme}? The answers come
        from the \textbf{data dictionary}, not from the data.
  \item \textbf{Optimiser.} Two plans at least. Read all 4{,}000 rows and
        discard the non-BSIT ones, then sort; or, if an index on
        \texttt{programme} exists, use it to fetch only the BSIT rows. And if
        an index on \texttt{name} exists, the \sqlkw{ORDER BY} may be free ---
        indexes are already in order. The optimiser picks by \emph{estimated}
        cost, using statistics from the dictionary. \chref{indexes} measures
        how good those estimates are.
  \item \textbf{Execution engine.} Runs the chosen plan, asking the storage
        manager for pages as it goes.
  \item \textbf{Buffer manager.} Page 7 is already in memory --- no disk read.
        Page 112 is not, so fetch it, and evict something to make room. This
        is the single biggest determinant of how fast the query is, and you
        never see it.
  \item \textbf{File manager.} Page 112 is at byte offset $111 \times 4096$ of
        \texttt{university.db}. Read 4{,}096 bytes.
\end{tightnum}

\smallskip
\textbf{The point of the exercise.} You wrote \emph{what} you wanted. Nothing in
your statement said how. Steps 3 to 6 were chosen for you, and the only way to
influence them is to change the schema --- which is why design, and not
cleverness in query writing, is most of this course.
\end{worked}

\section{The Data Dictionary Is Just More Tables}

The claim that a DBMS stores the schema in the database is easier to believe
once you have queried it.

\begin{examplebox}[The dictionary, in SQLite]
Every SQLite database has a table called \texttt{sqlite\_master}, and you can
select from it like any other:

\begin{lstlisting}[language=sql]
SELECT type, name, tbl_name FROM sqlite_master;
\end{lstlisting}

\begin{lstlisting}
type   name          tbl_name
-----  ------------  ----------
table  student       student
index  idx_name      student
view   fees_view     student
\end{lstlisting}

The \texttt{sql} column holds the exact text of the \sqlkw{CREATE} statement
that made each object --- which is why \texttt{.schema} is instant: it is a
\sqlkw{SELECT}, not a reconstruction.

\smallskip
Other systems do the same thing under other names:
\texttt{information\_schema} in MySQL and PostgreSQL,
\texttt{ALL\_TABLES} in Oracle. The idea is universal; the spelling is not.
\end{examplebox}

\begin{conceptbox}[Why this is more than a curiosity]
Because the schema is data, every tool that works on data works on the schema.
Three consequences you will use:

\begin{tight}
  \item \textbf{A program can discover a database it has never seen.} Anything
        that shows you a list of tables --- a GUI browser, an admin page, a
        reporting tool --- is selecting from the dictionary.
  \item \textbf{The schema can be backed up with the data}, because it is in
        the same file.
  \item \textbf{A dictionary query can audit a schema.} In \chref{constraints}
        you will use one to find every table that declares a foreign key, and
        the point of that exercise is that the declarations are all there and
        none of them are being enforced.
\end{tight}
\end{conceptbox}

\section{The Four Roles}

The components are software. These are people, and the distinction that
matters is not job title --- in a small organisation one person is all four ---
but \emph{which decisions belong to which hat}.

\begin{definitionbox}[Who decides what]
The \term{database administrator} (DBA) owns the running system: installation,
users and privileges, backups, recovery, performance, storage. The DBA decides
\emph{how it runs}.

\smallskip
The \term{database designer} owns the schema: entities, relationships,
constraints, normalisation. The designer decides \emph{what is true}.

\smallskip
The \term{application programmer} owns the programs that use the database.
They decide \emph{how it is reached}.

\smallskip
The \term{end user} owns the questions. They decide \emph{what it is for}, and
they are the only role in the list who need not know that any of the others
exist.
\end{definitionbox}

\begin{examplebox}[One request, four reactions]
\emph{``Students should be able to see their own marks on the portal, but not
anybody else's.''}

\begin{tabular}{@{}L{3.1cm}L{10.9cm}@{}}
\textbf{End user} & States the requirement, including the half of it that is a
prohibition. Notices first when it is wrong. \\
\textbf{Designer} & Is there a reliable way to tell whose marks these are?
Does \relation{enrolment}{\relkey{roll\_no}, \relkey{section\_id}, grade} tie
every mark to exactly one student? If not, the requirement cannot be met at
any price. \\
\textbf{DBA} & Creates the portal's database user and grants it the minimum:
read on a view, and nothing on the base table. \\
\textbf{Programmer} & Writes the query with the logged-in roll number as a
\emph{parameter}, never as pasted-in text --- the difference between those
two is the whole of SQL injection. \\
\end{tabular}

\smallskip
\textbf{Where it breaks.} If the designer gets it wrong, nobody downstream can
fix it: a mark not tied to one student cannot be shown to one student. This is
why design comes before everything else in this course, and why
\chref{ermodel} through \chref{constraints} are four chapters in which you
write almost no SQL.
\end{examplebox}

\begin{alertbox}[In this course you are three of the four, because SQLite has no DBA]
SQLite has no users, no privileges, no server to administer and no process to
tune. Several things on the DBA's list simply do not exist, and the rest move
onto somebody else's:

\begin{tight}
  \item \textbf{Users and privileges:} nobody's. The capability is absent. File
        permissions are the only access control there is, and
        \chref{project} is about what you get back when you move to a server.
  \item \textbf{Backup:} yours, and it is \texttt{copy university.db} --- which
        is genuinely all it takes, and is \chref{transactions}'s subject once
        you ask what happens if you copy it mid-write.
  \item \textbf{Performance:} yours. \texttt{ANALYZE} and
        \texttt{EXPLAIN QUERY PLAN} are the whole toolkit, and
        \chref{indexes} uses both.
  \item \textbf{Schema design, and writing the programs:} yours, which is the
        point --- you will hold the designer's hat and the programmer's hat at
        the same time, and the labs are arranged so you can feel the conflict.
\end{tight}

\smallskip
So: when this handout says ``the DBA would'', read it as ``on a real system
somebody else would, and this term it is you, with one hand tied''.
\end{alertbox}

\begin{pitfall}
\begin{tight}
  \item \textbf{Thinking the optimiser can be instructed.} In SQL you say what,
        never how. If the plan is bad, change the schema or the statistics ---
        not your tone of voice.
  \item \textbf{Confusing the data dictionary with the data.} The dictionary
        holds the schema. \texttt{SELECT COUNT(*) FROM sqlite\_master} counts
        tables, not students.
  \item \textbf{Assuming the roles are four people.} They are four sets of
        decisions. In this course they are all you.
  \item \textbf{Believing ``the application will check it''} is equivalent to a
        constraint in the schema. It is not, and \chref{constraints} is eight
        pages on why.
  \item \textbf{Expecting an authorisation manager in SQLite.} There is none.
        Open the file and you have everything.
\end{tight}
\end{pitfall}

\begin{lab}[Query the dictionary, and watch the optimiser change its mind]
\begin{lstlisting}[language=sql]
-- --- 1. Build something for the optimiser to think about -------
sqlite3 inside.db
CREATE TABLE student (
  roll_no   TEXT PRIMARY KEY,
  name      TEXT NOT NULL,
  programme TEXT NOT NULL
);
-- 300 rows, generated by a recursive CTE (explained in Ch 15):
INSERT INTO student
WITH RECURSIVE n(i) AS (SELECT 1 UNION ALL SELECT i+1 FROM n WHERE i<300)
SELECT printf('2021-IT-%03d', i),
       printf('Student %03d', i),
       CASE i%3 WHEN 0 THEN 'BSIT' WHEN 1 THEN 'BSCS' ELSE 'BSSE' END
FROM n;

-- --- 2. Read the dictionary ------------------------------------
SELECT type, name FROM sqlite_master;
SELECT sql FROM sqlite_master WHERE name = 'student';

-- --- 3. Ask the optimiser what it intends to do ----------------
EXPLAIN QUERY PLAN
  SELECT name FROM student WHERE programme = 'BSIT' ORDER BY name;
-- Write the plan down exactly. Note the words SCAN and ORDER BY.

-- --- 4. Give it an index and ask again -------------------------
CREATE INDEX idx_prog ON student(programme);
EXPLAIN QUERY PLAN
  SELECT name FROM student WHERE programme = 'BSIT' ORDER BY name;
-- The plan changed and you did not touch the query.

-- --- 5. Make the sort free too ---------------------------------
CREATE INDEX idx_name ON student(name);
EXPLAIN QUERY PLAN SELECT name FROM student ORDER BY name;
-- Does "USE TEMP B-TREE FOR ORDER BY" still appear?

-- --- 6. Look at what you have added to the dictionary ----------
SELECT type, name, tbl_name FROM sqlite_master;
-- Three new rows you never inserted.
\end{lstlisting}

\smallskip
\textbf{What to notice.} In steps 3 to 5 the statement was identical and the
plan was not. Everything that changed, changed in the schema. Keep the three
plans --- \chref{indexes} asks you to compare them against measured times, and
one of them will turn out to be a worse idea than it looks.
\end{lab}

\begin{checkpoint}
\begin{tightnum}
  \item A query runs in 40 milliseconds on Monday and 4 seconds on Friday, on
        the same data and with no change to the statement or the schema. Name
        two components that could be responsible and say how you would tell
        them apart.
  \item The data dictionary is stored as tables in the database it describes.
        Give one real advantage of that and one problem it creates.
  \item For each of the four roles, state one decision that is theirs alone,
        and then say what happens to that decision in this course --- where
        SQLite provides no DBA at all.
\end{tightnum}
\end{checkpoint}

\begin{chaptersummary}
\begin{tight}
  \item A DBMS is several cooperating components: a \textbf{query processor}
        (parser, optimiser, execution engine), a \textbf{storage manager}
        (buffer manager, file manager), a \textbf{transaction manager} with
        lock and recovery managers, an \textbf{authorisation manager}, and the
        \textbf{data dictionary}.
  \item A statement travels down through them. You specify \emph{what}; steps
        3 onward --- plan, pages, bytes --- are chosen for you.
  \item The parser consults the dictionary to check names and types; the
        optimiser consults it for sizes and indexes.
  \item The \textbf{data dictionary} is ordinary tables. In SQLite it is
        \texttt{sqlite\_master}, and \texttt{.schema} is a \sqlkw{SELECT} over
        it.
  \item Because the schema is data, programs can discover a database they have
        never seen, and a query can audit a schema.
  \item Four roles, defined by their decisions: the \textbf{DBA} decides how it
        runs, the \textbf{designer} decides what is true, the
        \textbf{programmer} decides how it is reached, the \textbf{end user}
        decides what it is for.
  \item A designer's error cannot be repaired downstream, which is why design
        precedes SQL in this course.
  \item SQLite has no authorisation manager and no DBA role. Backup,
        performance and design all land on you, and \chref{project} is the
        chapter about what a server gives back.
\end{tight}
\end{chaptersummary}

\begin{reviewq}
  \item \textbf{(MCQ)} Which component decides whether to use an index?
        (a)~the parser; (b)~the buffer manager; (c)~the query optimiser;
        (d)~the application programmer.
  \item \textbf{(MCQ)} \texttt{sqlite\_master} contains: (a)~the rows of every
        table; (b)~the schema of the database; (c)~the query log;
        (d)~the list of users.
  \item \textbf{(Short)} Explain why the parser needs the data dictionary, with
        a specific example of a statement it would reject because of
        something it found there.
  \item \textbf{(Short)} What does the buffer manager do, and why can it make
        the same query take very different times on identical data?
  \item \textbf{(Short)} Which of the four roles is responsible when a database
        permits two students to have the same roll number? Justify your answer
        in two sentences.
  \item \textbf{(Applied)} Take the statement \texttt{SELECT s.name, e.grade
        FROM student s JOIN enrolment e ON s.roll\_no = e.roll\_no WHERE
        e.grade = 'A';} and walk it through all six components, saying at each
        step what information that component needs and where it gets it.
  \item \textbf{(Applied)} A college wants one database used by the
        examinations office, the fees office and 4{,}000 students. Divide the
        work among the four roles: list what each must do before the system
        opens, then mark each item that SQLite cannot support and say what you
        would do about it.
\end{reviewq}

%=====================================================================
\chapter{Your First Database, in One Command}\label{ch:firstdatabase}
%=====================================================================
\locator{1}{Part I --- What a Database Is}

\begin{outcomes}
By the end of this chapter you will be able to:
\begin{tight}
  \item \textbf{Create} a database, a table and a row, and \textbf{read} the
        row back, without installing anything.
  \item \textbf{Read} a relation schema written in the standard notation and
        \textbf{write} the \sqlkw{CREATE TABLE} that implements it.
  \item \textbf{State} what a database costs on disk, and \textbf{explain} why
        the cost goes up in jumps rather than smoothly.
  \item \textbf{Use} the eight dot-commands you will need every week.
  \item \textbf{Account for} the gap between an empty database at 0 bytes and
        one table at 8{,}192 bytes.
\end{tight}
\end{outcomes}

\begin{prereq}
\chref{whodoeswhat}: the components, and the data dictionary. Appendix~A for
the toolkit --- if \texttt{sqlite3} does not start, read that first and then
come back.
\end{prereq}

\begin{keyterms}
\sqlkw{CREATE TABLE} $\bullet$ \sqlkw{INSERT} $\bullet$ \sqlkw{SELECT}
$\bullet$ relation schema $\bullet$ column type $\bullet$ \sqlkw{PRIMARY KEY}
$\bullet$ \sqlkw{NOT NULL} $\bullet$ dot-command $\bullet$ page $\bullet$ page
size $\bullet$ \sqlkw{VACUUM} $\bullet$ database file
\end{keyterms}

\section{Three Statements and a File}

Everything in \chref{whynotfiles} through \chref{whodoeswhat} was description.
This is the chapter where you have a database.

\begin{lstlisting}[language=bash]
sqlite3 university.db
\end{lstlisting}

That is the installation, the server start-up, the account creation and the
connection, all of it. There was no download. Then:

\begin{lstlisting}[language=sql]
CREATE TABLE student (
  roll_no   TEXT    PRIMARY KEY,
  name      TEXT    NOT NULL,
  programme TEXT    NOT NULL,
  semester  INTEGER NOT NULL
);

INSERT INTO student VALUES ('2021-IT-001','Ayesha Khan','BSIT',2);

SELECT name FROM student WHERE roll_no = '2021-IT-001';
\end{lstlisting}

\begin{lstlisting}
Ayesha Khan
\end{lstlisting}

You have now used all three levels of \chref{architecture} --- you wrote the
conceptual schema, SQLite built the internal one, and the \sqlkw{SELECT} was an
external view of one column of one row --- and all six components of
\chref{whodoeswhat} ran. It took under a tenth of a second, and the rest of
this chapter is about what that tenth of a second and that file actually cost.

\section{Writing the Schema Down Before You Type It}

The \sqlkw{CREATE TABLE} above is SQL. The thing it implements has a shorter
notation, used in every textbook and in the rest of this handout:

\begin{definitionbox}[Relation schema notation]
A \term{relation schema} is written as the relation name followed by its
attributes in brackets, with the primary key underlined:

\medskip
\centerline{\relation{student}{\relkey{roll\_no}, name, programme, semester}}

\medskip
It records the \emph{names} and the \emph{key} and nothing else --- no types,
no lengths, no storage. It is the conceptual level in one line, and it is what
you write on paper while designing, before any SQL exists.
\end{definitionbox}

\begin{examplebox}[From notation to SQL, and the three things the notation loses]
\relation{student}{\relkey{roll\_no}, name, programme, semester} becomes the
\sqlkw{CREATE TABLE} on the previous page. Going in that direction you must
\emph{add} three decisions the notation does not carry:

\begin{tight}
  \item \textbf{A type for each attribute.} \texttt{roll\_no} is
        \texttt{TEXT}, not \texttt{INTEGER}, because \texttt{2021-IT-001} is
        not a number --- it has hyphens and leading zeros, and you will never
        do arithmetic on it. Choosing \texttt{INTEGER} here is the single most
        common first-week mistake.
  \item \textbf{Which attributes may be absent.} \sqlkw{NOT NULL} on three of
        them. \chref{select} is about how much trouble the fourth kind ---
        the ones that may be missing --- causes.
  \item \textbf{The key, enforced.} The underline is a claim; \sqlkw{PRIMARY
        KEY} is the enforcement. \chref{constraints} is about the difference
        between those two, and about how often the gap between them is where
        the bugs are.
\end{tight}
\end{examplebox}

\dsfig{%
\begin{tikzpicture}[
  tbl/.style={rounded corners=3pt, draw=primary, line width=1.1pt,
              fill=primary!6, align=left, inner sep=0pt, text width=5.0cm},
  hdr/.style={font=\scriptsize\bfseries, text=white},
  row/.style={font=\scriptsize\ttfamily, text=primary!60!black},
  note/.style={font=\scriptsize\itshape, text=neutral, align=left}
]
  % header bar
  \fill[primary, rounded corners=3pt] (0,0) rectangle (5.0,0.52);
  \node[hdr, anchor=west] at (0.14,0.26) {student};
  \draw[primary, line width=1.1pt, rounded corners=3pt]
       (0,-3.0) rectangle (5.0,0.52);

  \node[row, anchor=west, text=primary, font=\scriptsize\bfseries\ttfamily]
       at (0.14,-0.22) {roll\_no \ \ TEXT \ \ PK};
  \node[row, anchor=west] at (0.14,-0.66) {name \ \ \ \ \ TEXT \ \ NOT NULL};
  \node[row, anchor=west] at (0.14,-1.10) {programme \ TEXT \ \ NOT NULL};
  \node[row, anchor=west] at (0.14,-1.54) {semester \ \ INT \ \ NOT NULL};
  \draw[primary!25, line width=0.6pt] (0.1,-1.80) -- (4.9,-1.80);
  \node[row, anchor=west, text=neutral] at (0.14,-2.10)
       {2021-IT-001 Ayesha Khan};
  \node[row, anchor=west, text=neutral] at (0.14,-2.50)
       {\ \ \ \ \ \ ... one row ...};

  \node[note, text width=4.9cm, anchor=west] at (5.75,-0.75)
       {\textbf{The schema.} Written once, changes rarely. This is what
        \texttt{.schema} prints and what the data dictionary stores.};
  \node[note, text width=4.9cm, anchor=west] at (5.75,-2.30)
       {\textbf{The instance.} One row today, 4{,}000 next month. The file
        grows; none of the four lines above change.};

  \draw[-{Stealth[length=2mm]}, primary!50, line width=0.8pt]
       (5.65,-0.75) -- (5.15,-0.75);
  \draw[-{Stealth[length=2mm]}, neutral!60, line width=0.8pt]
       (5.65,-2.30) -- (5.15,-2.30);
\end{tikzpicture}}%
{One table, with \chref{architecture}'s schema/instance distinction drawn on
it. The box above the rule is the schema and is four lines long whether the
table holds one row or a million.}{firsttable}

\section{What It Actually Cost}

A claim that a tool is lightweight is worth nothing unless somebody weighs it.
The harness \texttt{code/m04\_firstdatabase.py} does the three statements above
and measures them.

\begin{measured}[From nothing to a row, and what the file weighs]
\textsf{\footnotesize SQLite 3.51.2, Python 3.13.13, Intel i5-6600 (4 cores,
3.30\,GHz), 16\,GB RAM, Windows 11, mechanical disk.}

\medskip
\textbf{1. Time from no database at all to a row read back}

\begin{tabular}{@{}L{4.6cm}R{2.4cm}L{6.5cm}@{}}
\texttt{connect}  & 0.52\,ms  & the file does not exist yet \\
\texttt{create}   & 25.27\,ms & one \sqlkw{CREATE TABLE} \\
\texttt{insert}   & 23.80\,ms & one row \\
\texttt{select}   & 0.16\,ms  & read it back \\
\textbf{TOTAL}    & \textbf{49.75\,ms} & $\rightarrow$ \texttt{'Ayesha Khan'} \\
\end{tabular}

\smallskip
No server was started, no port opened, no password set, nothing installed.

\medskip
\textbf{2. Bytes on disk}

\begin{tabular}{@{}L{8.2cm}R{2.6cm}@{}}
an empty database, just connected & 0 \\
after one \sqlkw{CREATE TABLE}, no rows & 8{,}192 \\
1 row & 8{,}192 \\
100 rows & 8{,}192 \\
1{,}000 rows & 32{,}768 \\
10{,}000 rows & 262{,}144 \\
the same 10{,}000 rows after \sqlkw{VACUUM} & 262{,}144 \\
\end{tabular}

\medskip
\textbf{3. Where the growth happens} (page size 4{,}096 bytes)

\begin{tabular}{@{}R{2.0cm}R{2.4cm}R{1.8cm}R{3.4cm}@{}}
\textbf{at row} & \textbf{bytes} & \textbf{pages} & \textbf{bytes/row so far} \\
1   & 8{,}192  & 2 & 8192 \\
181 & 16{,}384 & 4 & 91 \\
351 & 20{,}480 & 5 & 58 \\
\end{tabular}

\medskip
\textbf{Read the three results, in order.}

\begin{tight}
  \item \textbf{An empty database is 0 bytes.} Connecting does not write
        anything. The file appears when the first statement that changes
        something commits --- which is already \chref{transactions}'s subject
        showing through.
  \item \textbf{The first table costs 8{,}192 bytes, and so does the hundredth
        row.} Two pages: one for SQLite's own header and
        \texttt{sqlite\_master} --- the data dictionary from
        \chref{whodoeswhat}, which has to live somewhere --- and one for your
        table. That second page then holds about 180 of these rows for nothing.
  \item \textbf{The file grows in jumps of exactly 4{,}096 bytes}, at rows 1,
        181 and 351. It does not creep up per row. The cost of storage is
        quantised, and the unit is the page.
  \item \textbf{\sqlkw{VACUUM} recovered nothing}, because nothing had been
        deleted. It is not a compressor; it reclaims space that deletions
        freed. Students run it hopefully and it is worth knowing now that
        hope is not the mechanism.
\end{tight}

\smallskip
\textbf{On re-running this.} The total time varies: repeated runs on this
machine gave 49.75\,ms and 70.38\,ms. Nothing changed but the operating
system's disk cache --- which is the buffer manager's neighbour from
\chref{whodoeswhat} making itself felt in week four. \textbf{The bytes do not
vary at all.} When you report a measurement, say which kind you have: the
timings are a range, the page counts are a fact.
\end{measured}

\begin{conceptbox}[Why 8{,}192 bytes for an empty table is the right answer]
It looks like waste --- 8\,KB to store nothing. It is the price of being a
database rather than a file, and it buys exactly two things:

\begin{tight}
  \item \textbf{The dictionary.} Page 1 holds the schema. Because of it, any
        program can open \texttt{university.db} and discover what is in it.
        The three CSV files in \chref{whynotfiles} had no such page, which is
        precisely why a program could not tell what they meant.
  \item \textbf{Room to grow in place.} Page 2 is mostly empty on purpose.
        Rows 2 through 180 cost nothing because the space is already there,
        and this is why inserting is fast.
\end{tight}

\smallskip
At 10{,}000 rows the overhead has stopped mattering: 262{,}144 bytes is about
26 bytes a row, and the rows themselves are about that size. \textbf{The fixed
cost is paid once and the marginal cost is the data.} That is the shape you
want from a storage system, and \chref{indexes} is the chapter where you
deliberately pay the fixed cost a second time to make reads fast.
\end{conceptbox}

\begin{alertbox}[The comparison that makes ``lightweight'' mean something]
``Lightweight'' is only a claim about something else. Counting what each
arrangement requires before the first row exists:

\begin{tabular}{@{}L{4.9cm}R{2.0cm}R{2.0cm}L{4.3cm}@{}}
& \textbf{SQLite} & \textbf{a server} & \\
installation steps     & 0 & several & download, install, configure \\
background processes   & 0 & $\geq 1$ & a service that must be running \\
network ports open     & 0 & $\geq 1$ & 3306, 5432, 1433 \dots \\
accounts to create     & 0 & $\geq 1$ & starting with the superuser \\
config files to edit   & 0 & $\geq 1$ & before anything is tuned \\
things that can be \textbf{misconfigured} & \textbf{0} & \textbf{all of the above} & \\
\end{tabular}

\smallskip
Every zero in the left column is a thing that cannot go wrong in week one ---
and also \textbf{a thing you will not learn to do this term}. That is the
trade, stated plainly: you get sixteen weeks on data and none on
administration. \chref{project} is where the right-hand column is taught by
contrast, and where you write the \sqlkw{GRANT} statements that SQLite will
not run.
\end{alertbox}

\section{The Eight Dot-Commands}

Dot-commands are instructions to the \texttt{sqlite3} shell, not SQL. They take
no semicolon. There are about sixty; you will use these eight all term.

\begin{examplebox}[The working set]
\begin{tabular}{@{}L{4.3cm}L{9.7cm}@{}}
\texttt{.open university.db} & open a database from inside the shell \\
\texttt{.tables}             & list the tables \\
\texttt{.schema student}     & print the \sqlkw{CREATE TABLE} text, exactly as stored \\
\texttt{.headers on}         & put column names above the output --- \textbf{set this first, always} \\
\texttt{.mode box}           & draw results as a table instead of pipe-separated text \\
\texttt{.read v1\_create.sql}& run a file of SQL \\
\texttt{.once out.txt}       & send the next result to a file \\
\texttt{.quit}               & leave \\
\end{tabular}

\smallskip
\texttt{.headers on} and \texttt{.mode box} together are the difference
between output you can read and output you cannot. Appendix~A has the full
list and the two-line reset procedure for when a lab database gets into a
state you do not want to debug.
\end{examplebox}

\begin{pitfall}
\begin{tight}
  \item \textbf{A semicolon after a dot-command.} \texttt{.tables;} is an
        error. SQL takes a semicolon; dot-commands do not.
  \item \textbf{Forgetting the semicolon after SQL.} The shell prints
        \texttt{...>} and waits. It is not frozen --- it wants the rest of your
        statement.
  \item \textbf{Making \texttt{roll\_no} an \texttt{INTEGER}.}
        \texttt{2021-IT-001} is text. This one costs a schema rewrite in week
        six.
  \item \textbf{Expecting \sqlkw{VACUUM} to shrink a database.} It reclaims
        deleted space only. With no deletions it recovers nothing, as measured
        above.
  \item \textbf{Running \texttt{sqlite3} with no filename.} You get a database
        in memory, you can work in it happily for an hour, and it vanishes on
        \texttt{.quit} with no warning.
  \item \textbf{Judging a database's size by the file alone.} 8{,}192 bytes
        held 100 rows and also 1. The file size tells you how many pages were
        allocated, not how much data there is.
\end{tight}
\end{pitfall}

\begin{lab}[Build the first table of the database you will keep all term]
This is the beginning of the semester project. Keep
\texttt{university.db} --- every lab from here adds to this one file, and
\chref{project} submits it.

\begin{lstlisting}[language=bash]
# --- 1. Create it, in the folder you will keep all term --------------
sqlite3 university.db
\end{lstlisting}
\begin{lstlisting}[language=sql]
-- --- 2. Make the output readable before anything else -----------
.headers on
.mode box

-- --- 3. The first table -----------------------------------------
CREATE TABLE student (
  roll_no   TEXT    PRIMARY KEY,
  name      TEXT    NOT NULL,
  programme TEXT    NOT NULL,
  semester  INTEGER NOT NULL
);

-- --- 4. Four rows -----------------------------------------------
INSERT INTO student VALUES
  ('2021-IT-001','Ayesha Khan','BSIT',2),
  ('2021-IT-002','Bilal Ahmed','BSIT',2),
  ('2021-CS-007','Hira Nadeem','BSCS',4),
  ('2021-SE-013','Usman Tariq','BSSE',2);

-- --- 5. Read them back ------------------------------------------
SELECT * FROM student;
SELECT name, programme FROM student WHERE semester = 2;

-- --- 6. Prove the key is enforced -------------------------------
INSERT INTO student VALUES ('2021-IT-001','Someone Else','BSIT',2);
-- Write down the exact error message. This is Chapter 8 arriving early.

-- --- 7. Leave the shell ----------------------------------------
.quit
\end{lstlisting}
\begin{lstlisting}[language=bash]
# --- 8. Measure your own machine ------------------------------------
# the file size, and the pages it is made of
ls -l university.db
sqlite3 university.db "PRAGMA page_size; PRAGMA page_count;"
python code/m04_firstdatabase.py
\end{lstlisting}

\smallskip
\textbf{Record and keep.} Your own total milliseconds, your own page size, and
your own three growth points. They will not match the measured box exactly ---
different disk, different cache --- and \textbf{the ones that should match
exactly are the byte counts}. If your page size is not 4{,}096, say so when
you hand in, because every size figure in this handout scales by it.
\end{lab}

\begin{checkpoint}
\begin{tightnum}
  \item An empty database is 0 bytes and one empty table is 8{,}192. Account
        for every one of those bytes --- what is in them, and why could it not
        be fewer?
  \item The file was 8{,}192 bytes at 100 rows and 32{,}768 at 1{,}000. Predict
        the size at 2{,}000 rows, show the arithmetic, then run it and explain
        any difference between your prediction and the result.
  \item Write \relation{instructor}{\relkey{emp\_id}, name, department,
        designation} as a \sqlkw{CREATE TABLE}, choosing a type for each
        attribute, deciding which may be \sqlkw{NULL}, and justifying both
        choices for \texttt{emp\_id} in one sentence each.
\end{tightnum}
\end{checkpoint}

\begin{chaptersummary}
\begin{tight}
  \item \texttt{sqlite3 university.db} is the whole of installation,
        start-up and connection. Three statements --- \sqlkw{CREATE TABLE},
        \sqlkw{INSERT}, \sqlkw{SELECT} --- are a working database.
  \item A \textbf{relation schema} is written \relation{student}{\relkey{roll\_no},
        name, programme, semester}: name, attributes, key underlined. Turning
        it into SQL means adding types, nullability and real key enforcement.
  \item Measured: \textbf{49.75\,ms} from nothing to a row read back; an
        \textbf{empty database is 0 bytes}; one \sqlkw{CREATE TABLE} costs
        \textbf{8{,}192 bytes} and then holds about 180 rows free; 10{,}000
        rows is \textbf{262{,}144 bytes}, about 26 bytes a row.
  \item The file grows in \textbf{4{,}096-byte page jumps}, at rows 1, 181 and
        351 --- not smoothly. Storage is quantised and the unit is the page.
  \item The 8\,KB fixed cost buys the data dictionary and room to grow in
        place. It is paid once; after that the marginal cost is the data.
  \item \sqlkw{VACUUM} reclaims space freed by deletions. With no deletions it
        recovers nothing.
  \item Timings vary with the OS disk cache (49.75 vs 70.38\,ms on the same
        machine); byte counts do not vary. Report which kind of number you
        have.
  \item ``Lightweight'' means zero installation steps, processes, ports,
        accounts and config files --- and therefore zero of those to
        misconfigure, and zero of them learned this term.
  \item Eight dot-commands cover the course. Set \texttt{.headers on} and
        \texttt{.mode box} first, every session.
\end{tight}
\end{chaptersummary}

\begin{reviewq}
  \item \textbf{(MCQ)} A newly connected, empty SQLite database occupies:
        (a)~0 bytes; (b)~4{,}096 bytes; (c)~8{,}192 bytes; (d)~one page per
        table.
  \item \textbf{(MCQ)} \sqlkw{VACUUM} on a database from which nothing has been
        deleted: (a)~halves the file; (b)~compresses the rows;
        (c)~recovers nothing; (d)~rebuilds the indexes only.
  \item \textbf{(Short)} Why is \texttt{roll\_no} declared \texttt{TEXT} and
        not \texttt{INTEGER}? Give two reasons, one about correctness and one
        about what you would eventually want to do with the value.
  \item \textbf{(Short)} Explain why the file size jumped at row 181 and not at
        row 150 or row 200.
  \item \textbf{(Short)} The same harness gave 49.75\,ms and 70.38\,ms on the
        same machine with the same data. Is one of them wrong? Answer in two
        sentences, saying what you would report.
  \item \textbf{(Applied)} Write the \sqlkw{CREATE TABLE} for
        \relation{course}{\relkey{course\_code}, title, credit\_hours,
        department}, then predict its empty size in bytes and the row count at
        which it will first grow. Run it, record both, and explain any
        difference between prediction and measurement.
  \item \textbf{(Applied)} A colleague proposes storing the university's data
        in one table of 40 columns ``to avoid joins''. Using the page
        measurements in this chapter, estimate what that costs in bytes at
        10{,}000 students compared with the four-column table above, and then
        give the argument against it that is \emph{not} about bytes ---
        \chref{normalforms} will measure both halves of this answer.
\end{reviewq}


\part{Designing It}
%=====================================================================
% PART II OPENER
%=====================================================================
\thispagestyle{plain}
\structuralfigures{P2.}

\begin{center}
\begin{tcolorbox}[enhanced, width=0.92\textwidth, colback=secondary!5,
  colframe=secondary, boxrule=0pt, leftrule=5pt, arc=3pt, top=8pt, bottom=8pt]
\large\itshape\color{secondary!60!black}
Before a database can be built it has to be decided, and the deciding is done on paper. Part~II is the design half of the course: what things exist, how they relate, and which of those relationships the database will be made to guarantee. HEC's outline puts this after normalisation. It belongs here, because you cannot normalise a schema you have not designed.
\end{tcolorbox}
\end{center}

\vspace{6pt}

\dsfigh{%
\begin{tikzpicture}[
  cbox/.style={rounded corners=3pt, draw=secondary, line width=1pt, fill=secondary!7,
               text=secondary!55!black, font=\small\bfseries, align=center,
               minimum width=3.5cm, minimum height=1.0cm},
  hbox/.style={rounded corners=3pt, draw=secondary, line width=1.8pt, fill=secondary!16,
               text=secondary!55!black, font=\small\bfseries, align=center,
               minimum width=3.5cm, minimum height=1.0cm},
  arr/.style={-{Stealth[length=2.4mm]}, secondary!60, line width=1pt},
  note/.style={font=\scriptsize\itshape, text=secondary!55!black, align=center,
               fill=white, inner sep=2pt}
]
  \node[cbox] (c0) at (0.00,0) {Ch.\ 5\\The ER Model};
  \node[cbox] (c1) at (4.30,0) {Ch.\ 6\\Cardinality};
  \node[cbox] (c2) at (8.60,0) {Ch.\ 7\\The Relational Model};
  \node[hbox] (c3) at (12.90,0) {Ch.\ 8\\Keys and Constraints};
  \draw[arr] (c0) -- (c1);
  \draw[arr] (c1) -- (c2);
  \draw[arr] (c2) -- (c3);
  \node[note] at (2.15,-1.0) {what exists};
  \node[note] at (6.45,-1.0) {how many, and must it};
  \node[note] at (10.75,-1.0) {tables, precisely};
  \node[font=\scriptsize\itshape, text=accent, align=center] at (12.90,1.15)
       {where declaring a rule and enforcing one turn out to differ};
\end{tikzpicture}}%
{Reading path through Part~II. Chapter 6 draws the university database this
whole course uses; Chapter 8 is the first measured chapter in which the
database does something you did not ask for.}{part2map}

\newpage

%=====================================================================
\chapter{The Entity-Relationship Model}\label{ch:ermodel}
%=====================================================================
\locator{2}{Part II --- Designing It}

\begin{outcomes}
By the end of this chapter you will be able to:
\begin{tight}
  \item \textbf{Distinguish} an entity, an entity type and an entity set, and
        \textbf{say} which of the three a table corresponds to.
  \item \textbf{Classify} an attribute as simple or composite, single- or
        multi-valued, stored or derived, and \textbf{explain} why each
        distinction changes what you build.
  \item \textbf{Identify} superkeys, candidate keys and the primary key of an
        entity type, and \textbf{justify} the choice between candidates.
  \item \textbf{Read and draw} a Chen ER diagram, and \textbf{state} the degree
        of a relationship.
  \item \textbf{Decide}, for a given noun in a requirements statement, whether
        it is an entity, an attribute or a relationship.
\end{tight}
\end{outcomes}

\begin{prereq}
\chref{whodoeswhat}: the designer's job --- deciding \emph{what is true}.
\chref{firstdatabase} for the \relation{}{} notation. No SQL is written in
this chapter; you will not touch \texttt{university.db} until
\chref{mapping}.
\end{prereq}

\begin{keyterms}
entity $\bullet$ entity type $\bullet$ entity set $\bullet$ attribute
$\bullet$ simple attribute $\bullet$ composite attribute $\bullet$
single-valued $\bullet$ multi-valued $\bullet$ stored attribute $\bullet$
derived attribute $\bullet$ null value $\bullet$ key attribute $\bullet$
superkey $\bullet$ candidate key $\bullet$ primary key $\bullet$ composite key
$\bullet$ surrogate key $\bullet$ relationship $\bullet$ relationship type
$\bullet$ degree $\bullet$ binary $\bullet$ ternary $\bullet$ recursive
relationship $\bullet$ role name $\bullet$ ER diagram $\bullet$ Chen notation
\end{keyterms}

\section{Design Happens on Paper}

You can create a table in eleven milliseconds (\chref{firstdatabase}). That is
exactly the problem: it is so cheap to make a table that people make the wrong
ones, and a wrong table is discovered in week fourteen when four programs
already read it.

The entity-relationship model is a way of arguing about a design before any of
it is expensive. It was proposed by Peter Chen in 1976 and it has outlived
every database product of that decade, for one reason: it is about the
\emph{university}, not about the database. Everything in this chapter could be
checked by the examinations officer, who has never heard of SQL.

\begin{alertbox}[HEC's outline puts this chapter after normalisation. This handout does not.]
The published course content lists \emph{normalization, functional
dependencies, normal forms} before \emph{entity relationship model, entity
sets, \dots, entity-relationship diagrams}. Taken literally, that normalises a
schema before designing one, which cannot be done --- normalisation
(\chref{normalforms}) takes an existing set of tables as its input.

\smallskip
All four of HEC's own recommended textbooks order it the other way. So does
this handout: design first (Part~II), map it to tables (Part~III), and only
then check the result for faults (Part~V). Every clause is still covered; the
Course Specification records the reordering and this is the one place it is
visible.
\end{alertbox}

\section{Entities: the Things}

\begin{definitionbox}[Entity, entity type, entity set]
An \term{entity} is one thing in the world that can be told apart from other
things: the student Ayesha Khan, the course CS-103, Tuesday's section of it.

\smallskip
An \term{entity type} is the \emph{description} of a kind of thing --- the
name \ent{student} together with the attributes every student has.

\smallskip
An \term{entity set} is all the entities of that type at one moment --- the
4{,}000 students enrolled today.

\smallskip
This is \chref{architecture}'s schema/instance distinction wearing different
words. The entity type is schema; the entity set is instance. A table
corresponds to the \textbf{entity type}; its rows are the entity set.
\end{definitionbox}

\begin{examplebox}[Three words people use for one thing, and when it matters]
\begin{tabular}{@{}L{3.6cm}L{5.0cm}L{5.4cm}@{}}
\textbf{Term} & \textbf{Example} & \textbf{Changes how often} \\
Entity        & Ayesha Khan, roll 2021-IT-001 & constantly \\
Entity set    & today's 4{,}000 students      & every admission and graduation \\
Entity type   & \ent{student}(roll\_no, name, programme, semester)
                                              & once a year, by argument \\
\end{tabular}

\smallskip
In conversation ``student'' means all three and no harm is done. In an
examination answer it must be the right one, and in a design meeting the
confusion is expensive: ``students change all the time'' is true of the set
and false of the type, and the two sentences lead to different databases.
\end{examplebox}

\section{Attributes, and the Four Distinctions That Matter}

An \term{attribute} is a property of an entity type. Students have a name and
a roll number. That much is obvious; the four ways attributes differ are not,
and each one changes what you have to build.

\begin{definitionbox}[Four distinctions]
\textbf{Simple or composite.} A \term{simple} (atomic) attribute cannot
usefully be divided: \texttt{semester}. A \term{composite} attribute has
meaningful parts: \texttt{address} = house, street, city, postcode.

\smallskip
\textbf{Single-valued or multi-valued.} \texttt{date\_of\_birth} has exactly
one value per student. \texttt{phone} may have several. A \term{multi-valued}
attribute is drawn with a double ellipse --- and it is the one kind that
\emph{cannot become a column}, which is \chref{mapping}'s problem.

\smallskip
\textbf{Stored or derived.} \texttt{date\_of\_birth} is \term{stored};
\texttt{age} is \term{derived} from it and today's date. A derived attribute
is drawn with a dashed ellipse and is normally \emph{not} stored at all.

\smallskip
\textbf{Key or non-key.} A \term{key attribute} has a different value for
every entity in the set. It is underlined.

\smallskip
A fifth property cuts across all four: an attribute may have no value for some
entity, which is the \term{null} of \chref{select} and a much deeper problem
than it looks.
\end{definitionbox}

\dsfig{%
\begin{tikzpicture}[
  ent/.style={draw=primary, line width=1.1pt, fill=primary!8,
              text=primary!60!black, font=\small\bfseries,
              minimum width=2.4cm, minimum height=0.85cm},
  att/.style={draw=secondary, line width=0.8pt, fill=white,
              text=secondary!60!black, font=\scriptsize, ellipse,
              inner sep=1.6pt, minimum width=1.5cm},
  keyatt/.style={att, text=primary, font=\scriptsize\bfseries},
  multi/.style={att, double, double distance=1.3pt, draw=purplehl,
                text=purplehl!70!black},
  deriv/.style={att, dashed, draw=neutral, text=neutral},
  lnk/.style={draw=secondary!50, line width=0.7pt},
  note/.style={font=\scriptsize\itshape, text=neutral}
]
  \node[ent] (s) at (0,0) {STUDENT};

  \node[keyatt] (k)  at (-3.9, 1.5) {\relkey{roll\_no}};
  \node[att]    (n)  at (-1.6, 2.1) {name};
  \node[att]    (pr) at ( 0.9, 2.1) {programme};
  \node[att]    (dob) at ( 3.3, 1.5) {date\_of\_birth};
  \node[deriv]  (ag) at ( 4.6, 0.1) {age};
  \node[multi]  (ph) at ( 3.1,-1.5) {phone};
  \node[att]    (ad) at (-1.4,-1.7) {address};
  \node[att]    (h)  at (-4.2,-0.8) {house};
  \node[att]    (st) at (-4.3,-1.8) {street};
  \node[att]    (ci) at (-3.4,-2.7) {city};

  \foreach \a in {k,n,pr,dob,ag,ph,ad}{\draw[lnk] (s) -- (\a);}
  \foreach \a in {h,st,ci}{\draw[lnk] (ad) -- (\a);}

  \node[note, anchor=west, text=purplehl] at (4.1,-2.2)
       {double = multi-valued};
  \node[note, anchor=west, text=neutral]  at (5.0,-0.5) {dashed = derived};
  \node[note, anchor=east, text=secondary!70!black] at (-4.6,-3.3)
       {composite: parts hang off the whole};
  \node[note, anchor=east, text=primary] at (-4.6,2.1) {underlined = key};
\end{tikzpicture}}%
{\ent{student} in Chen notation. Four of the five attribute kinds are
present. \texttt{phone} and \texttt{age} are the two that will not survive
contact with a table, and \chref{mapping} says what happens to each.}{erstudent}

\begin{conceptbox}[Why the distinctions are not decoration]
Each one has a consequence you will meet within three chapters.

\begin{tight}
  \item \textbf{Composite.} Store the parts, not the whole. A system that
        stores \texttt{address} as one text field can never answer ``how many
        students live in Bahawalpur'', and finding out costs a schema change
        after the data exists --- which is \chref{evolution}.
  \item \textbf{Multi-valued.} A column cannot hold three phone numbers. Not
        ``should not'': \emph{cannot}, if you want to search them. It becomes
        its own table, and that is a rule in \chref{mapping}.
  \item \textbf{Derived.} Storing \texttt{age} means storing a fact that
        becomes wrong overnight, silently, for about a quarter of the rows
        each year. Store the birth date and compute.
  \item \textbf{Key.} Everything else in the design hangs off this choice, and
        the next section is about how to make it.
\end{tight}
\end{conceptbox}

\section{Keys: Three Words, One Choice}

\begin{definitionbox}[Superkey, candidate key, primary key]
A \term{superkey} is any set of attributes whose values are different for
every entity. $\{$\texttt{roll\_no}$\}$ is one; so is
$\{$\texttt{roll\_no}, \texttt{name}$\}$, which is a superkey carrying a
passenger.

\smallskip
A \term{candidate key} is a superkey with no attribute to spare: remove any one
attribute and it stops being unique. $\{$\texttt{roll\_no}$\}$ is a candidate
key; $\{$\texttt{roll\_no}, \texttt{name}$\}$ is not.

\smallskip
The \term{primary key} is the one candidate key the designer picks to identify
entities throughout the database. The others are \term{alternate keys}.

\smallskip
A key made of more than one attribute is \term{composite}. A key invented by
the designer with no meaning in the world --- an auto-numbered id --- is a
\term{surrogate key}.
\end{definitionbox}

\begin{worked}[Finding every key of \ent{student}, and choosing one]
Attributes: \texttt{roll\_no}, \texttt{cnic}, \texttt{name},
\texttt{programme}, \texttt{semester}, \texttt{email}.

\medskip
\textbf{Step 1 --- what the university guarantees.} Not what today's data
happens to show. Ask, and write the answers down:
\begin{tight}
  \item Roll numbers are unique and never reissued. \emph{Guaranteed.}
  \item CNICs are unique. \emph{Guaranteed} --- but a first-semester student
        may not have one yet.
  \item Email addresses are issued one per student. \emph{Guaranteed} while
        enrolled; reused after five years.
  \item Names are not unique. Two students called Muhammad Ali is normal.
\end{tight}

\textbf{Step 2 --- candidate keys.} Three:
$\{$\texttt{roll\_no}$\}$, $\{$\texttt{cnic}$\}$, $\{$\texttt{email}$\}$.
$\{$\texttt{name}, \texttt{programme}, \texttt{semester}$\}$ is \emph{not} one
--- two BSIT second-semester students can share a name, and ``it has not
happened yet'' is not a guarantee.

\medskip
\textbf{Step 3 --- choose, on four grounds.}

\smallskip
\begin{tabular}{@{}L{2.2cm}L{2.4cm}L{2.4cm}L{2.2cm}L{3.0cm}@{}}
& \textbf{always present?} & \textbf{never changes?} & \textbf{short?}
& \textbf{verdict} \\
\texttt{roll\_no} & yes & yes & yes & \textbf{primary key} \\
\texttt{cnic}     & \textbf{no} & yes & yes & alternate, and nullable \\
\texttt{email}    & yes & \textbf{no} & no & alternate at best \\
\end{tabular}

\medskip
\textbf{The answer.} \relation{student}{\relkey{roll\_no}, cnic, name,
programme, semester, email}, with \texttt{cnic} and \texttt{email} declared
unique but not primary.

\smallskip
\textbf{What decided it.} Not uniqueness --- all three are unique. The primary
key must also be \emph{always present} and \emph{never changing}, because
every other table will refer to it (\chref{constraints}). \texttt{cnic} fails
the first test and \texttt{email} fails the second, and a primary key that can
change is a repair job spread across every table that referenced it.
\end{worked}

\begin{alertbox}[Three things students believe about keys, and what is wrong with each]
\begin{tight}
  \item \textbf{``The data is unique, so it is a key.''} A key is a
        \emph{constraint on all future data}, not an observation about today's.
        Check the current data by all means --- a duplicate proves it is not a
        key --- but no duplicate proves nothing.
  \item \textbf{``Add an auto-number id to every table and the problem goes
        away.''} The surrogate makes rows addressable; it does not make them
        \emph{distinguishable}. With an id and no other constraint you can
        insert the same student twice and the database is content. Surrogates
        are useful --- and whenever you use one, the real candidate key still
        has to be declared unique beside it.
  \item \textbf{``A composite key means bad design.''} It usually means the
        entity genuinely is identified by a combination. An enrolment is
        identified by \emph{which student} and \emph{which section}, and the
        pair is the honest key --- \chref{cardinality} and \chref{mapping}
        both rely on it.
\end{tight}
\end{alertbox}

\section{Relationships, and Degree}

\begin{definitionbox}[Relationship, relationship type, degree]
A \term{relationship} is an association among entities: Ayesha is enrolled in
Tuesday's section of CS-103.

\smallskip
A \term{relationship type} is the kind of association --- \ent{enrols},
between \ent{student} and \ent{section}. It is drawn as a diamond.

\smallskip
The \term{degree} of a relationship type is how many entity types take part:
\term{unary} (1), \term{binary} (2), \term{ternary} (3). Almost everything is
binary.

\smallskip
A \term{recursive} relationship is a unary one --- an entity type related to
itself, like a course being a prerequisite for another course. The two ends
need \term{role names} to tell them apart, because ``course'' appears twice.

\smallskip
A relationship type may have attributes of its own. \texttt{grade} belongs to
the \emph{enrolment}, not to the student and not to the section.
\end{definitionbox}

\dsfig{%
\begin{tikzpicture}[
  ent/.style={draw=primary, line width=1.1pt, fill=primary!8,
              text=primary!60!black, font=\scriptsize\bfseries,
              minimum width=2.0cm, minimum height=0.72cm},
  rel/.style={draw=greenhl, line width=1pt, fill=greenhl!8,
              text=greenhl!55!black, font=\scriptsize\bfseries,
              diamond, aspect=2.1, inner sep=1pt},
  att/.style={draw=secondary, line width=0.7pt, fill=white,
              text=secondary!60!black, font=\scriptsize, ellipse,
              inner sep=1.4pt},
  lnk/.style={draw=neutral!55, line width=0.7pt},
  rl/.style={font=\scriptsize\itshape, text=neutral},
  hdr/.style={font=\scriptsize\bfseries, text=primary!70!black}
]
  % --- binary
  \node[hdr] at (-4.6,1.9) {Binary};
  \node[ent] (st) at (-6.3,0.6) {STUDENT};
  \node[rel] (en) at (-4.0,0.6) {ENROLS};
  \node[ent] (se) at (-1.7,0.6) {SECTION};
  \node[att] (gr) at (-4.0,-0.9) {grade};
  \draw[lnk] (st) -- (en);  \draw[lnk] (en) -- (se);
  \draw[lnk] (en) -- (gr);
  \node[rl, text width=3.4cm, align=center] at (-4.0,-1.9)
       {the grade belongs to the \emph{relationship}, not to either end};

  % --- recursive
  \node[hdr] at (2.6,1.9) {Unary (recursive)};
  \node[ent] (co) at (2.6,0.6) {COURSE};
  \node[rel] (pr) at (2.6,-1.2) {REQUIRES};
  \draw[lnk] (co.west) .. controls (0.7,0.4) and (0.7,-1.1) .. (pr.west);
  \draw[lnk] (co.east) .. controls (4.5,0.4) and (4.5,-1.1) .. (pr.east);
  \node[rl, anchor=east] at (1.0,-0.3) {\emph{is prerequisite for}};
  \node[rl, anchor=west] at (4.2,-0.3) {\emph{depends on}};
  \node[rl, text width=3.8cm, align=center] at (2.6,-2.3)
       {one entity type, two ends: the role names are what tell them apart};
\end{tikzpicture}}%
{Two of the three degrees. \ent{enrols} carries an attribute of its own;
\ent{requires} joins \ent{course} to itself and is unreadable without
role names. The university database you are building contains both.}{erdegree}

\begin{examplebox}[Entity, attribute or relationship? The question that decides a design]
Every noun in a requirements document is one of the three, and getting it wrong
is the most common design error there is. Three tests, applied in order:

\begin{tight}
  \item \textbf{Does it have properties of its own?} If you can list two or
        three facts \emph{about} it, it is an entity.
  \item \textbf{Is there more than one of it per owner?} If yes, it is not an
        attribute.
  \item \textbf{Does it only exist because two other things met?} Then it is a
        relationship.
\end{tight}

\medskip
\begin{tabular}{@{}L{3.3cm}L{2.5cm}L{8.1cm}@{}}
\textbf{Noun} & \textbf{Verdict} & \textbf{Why} \\
roll number   & attribute    & no properties of its own \\
department    & \textbf{entity} & it has a name, a head, a building --- and
                                 fifty courses point at it \\
grade         & \textbf{relationship attribute} & exists only because a
                                 student met a section \\
prerequisite  & \textbf{relationship} & joins course to course \\
semester      & it depends  & an attribute of \ent{student} if it means
                              ``which year they are in''; an \emph{entity} if
                              the university runs Fall 2026 and Spring 2027 as
                              distinct things with dates \\
\end{tabular}

\smallskip
The last row is the honest one. The same word is an attribute in one
university and an entity in another, and no rule decides it --- the
\emph{requirement} does. That is why design is a conversation with the
examinations officer and not a procedure.
\end{examplebox}

\begin{pitfall}
\begin{tight}
  \item \textbf{Making an entity out of every noun.} ``Name'' is not an
        entity. Apply the three tests.
  \item \textbf{Hanging a relationship's attribute on an entity.} Putting
        \texttt{grade} on \ent{student} gives each student one grade for
        the whole degree. It belongs on \ent{enrols}.
  \item \textbf{Choosing a primary key that can change.} Email, phone, name.
        Every table that referenced it has to be repaired with it.
  \item \textbf{Treating today's data as a guarantee.} No duplicate names today
        does not make \texttt{name} a key.
  \item \textbf{Drawing a recursive relationship without role names.} The
        diagram is then genuinely ambiguous --- nobody can tell which end is
        the prerequisite.
  \item \textbf{Storing a derived attribute because it is convenient.}
        \texttt{age} is wrong for a quarter of your rows within a year, and
        nothing reports it.
\end{tight}
\end{pitfall}

\begin{lab}[Turn a page of English into an ER diagram]
No database this week. Paper, and one argument.

\begin{lstlisting}
# --- 1. The requirement, as the examinations officer states it -------

"Every student has a roll number, a name, a CNIC, one or more phone
 numbers, a postal address, and a date of birth. Students belong to
 one programme. A course has a code, a title and a number of credit
 hours, and belongs to a department. A course may require other
 courses as prerequisites. Each semester a course runs as one or
 more sections; a section has a number, a room and a timetable slot,
 and is taught by one instructor. Students enrol in sections and
 receive a grade. We need to know each student's age, and how many
 students are in each section."
\end{lstlisting}

\begin{lstlisting}
# --- 2. List every noun --------------------------------------------
    Write them in one column. There are about twenty.

# --- 3. Classify each one ------------------------------------------
    entity / attribute / relationship, using the three tests.
    Write ONE SENTENCE of justification for each. The sentence is
    the exercise; the label is not.

# --- 4. Mark the attribute kinds -----------------------------------
    composite, multi-valued, derived, key. There is at least one
    of each in the paragraph above.

# --- 5. Find the candidate keys ------------------------------------
    For STUDENT, COURSE and SECTION. For each, say what the
    university guarantees -- not what sample data shows.

# --- 6. Draw it, Chen notation -------------------------------------
    Rectangles, ellipses, diamonds. Role names on the recursive
    relationship. Do NOT draw cardinalities yet; that is Chapter 6.

# --- 7. Find the two traps -----------------------------------------
    Two things in the paragraph must NOT become columns, and one
    must not be stored at all. Name all three and say why.

# --- 8. Hand it to somebody else -----------------------------------
    Give your diagram, without the paragraph, to another student.
    Ask them to read it back as English. Every sentence they get
    wrong is a defect in your diagram, not in their reading.
\end{lstlisting}

\smallskip
\textbf{What to notice.} Step 8 is the one that matters and the one people
skip. An ER diagram's whole purpose is to be checked by somebody who cannot
read SQL. If it cannot be read back, it has failed at the only job it has ---
and it is still cheap to fix, which it will not be in \chref{mapping}.

\smallskip
\textbf{Keep this diagram.} Chapter~6 adds cardinalities to it, \chref{mapping}
turns it into tables, and it is part of your project submission.
\end{lab}

\begin{checkpoint}
\begin{tightnum}
  \item A colleague proposes \texttt{phone} as a column of \ent{student},
        because ``almost every student has one number''. Give the argument
        against in three sentences, and say what you would do instead and what
        it costs.
  \item \ent{instructor} has \texttt{emp\_id}, \texttt{cnic}, \texttt{name},
        \texttt{email}, \texttt{department}, \texttt{office\_no}. List every
        candidate key you can justify, state the guarantee each one rests on,
        choose a primary key and defend the choice on all four grounds in the
        worked example.
  \item Explain, to somebody who has written programs but never designed a
        database, why \texttt{grade} belongs to the \emph{relationship} between
        a student and a section rather than to either of them. Then say what
        goes wrong if it is put on \ent{student}, and what goes wrong if it
        is put on \ent{section}.
\end{tightnum}
\end{checkpoint}

\begin{chaptersummary}
\begin{tight}
  \item An \textbf{entity} is one thing; an \textbf{entity type} describes a
        kind of thing; an \textbf{entity set} is all of them today. Type is
        schema, set is instance, and a table corresponds to the type.
  \item Attributes differ in four ways: simple/\textbf{composite},
        single-/\textbf{multi-valued}, stored/\textbf{derived}, key/non-key ---
        and each distinction has a consequence in Part~III.
  \item A \textbf{multi-valued} attribute cannot become a column; a
        \textbf{derived} one should not be stored; a \textbf{composite} one is
        stored as its parts.
  \item A \textbf{superkey} is unique; a \textbf{candidate key} is unique with
        nothing to spare; the \textbf{primary key} is the candidate the
        designer picks, and the others are alternate keys.
  \item Choose a primary key on four grounds: unique, \textbf{always present},
        \textbf{never changing}, short. Uniqueness alone decides nothing when
        there are three candidates.
  \item A key is a constraint on \emph{all future data}. A duplicate disproves
        a key; the absence of one proves nothing.
  \item A \textbf{relationship} associates entities; its \textbf{degree} is how
        many entity types take part. A \textbf{recursive} relationship joins a
        type to itself and needs \textbf{role names} to be readable.
  \item A relationship can carry attributes of its own. \texttt{grade} belongs
        to the enrolment, not to the student and not to the section.
  \item Every noun is an entity, an attribute or a relationship. The same word
        can be different things in different universities; the requirement
        decides, not a rule.
  \item HEC's clause list puts normalisation before the ER model. This handout
        reverses the two, because normalisation takes tables as its input and
        there are none until Part~III.
\end{tight}
\end{chaptersummary}

\begin{reviewq}
  \item \textbf{(MCQ)} A double ellipse in Chen notation marks an attribute
        that is: (a)~derived; (b)~composite; (c)~multi-valued; (d)~a key.
  \item \textbf{(MCQ)} $\{$\texttt{roll\_no}, \texttt{name}$\}$, where
        \texttt{roll\_no} is unique, is: (a)~a candidate key; (b)~a superkey
        but not a candidate key; (c)~neither; (d)~the primary key.
  \item \textbf{(Short)} Give an example of a derived attribute other than age,
        and say what goes wrong if it is stored.
  \item \textbf{(Short)} What is the degree of the \ent{requires}
        relationship between courses, and why does it need role names when a
        binary relationship usually does not?
  \item \textbf{(Short)} A surrogate key is added to a table that already has a
        natural candidate key. What must still be declared, and what breaks if
        it is not?
  \item \textbf{(Applied)} A library wants to record books, copies of books,
        members and loans. Decide which of the four are entity types, identify
        the relationships and their degrees, choose a primary key for each
        entity type with its justification, and say which attributes of a loan
        belong to the relationship rather than to either end.
  \item \textbf{(Applied)} Take the \ent{student} diagram in
        Figure~\ref{fig:erstudent} and the requirement ``we must be able to
        list every student in a given city, and ring any of a student's
        numbers''. State precisely which parts of the diagram now cannot be
        implemented as columns, what you would build instead, and what the
        design would have cost if this requirement had arrived after the
        database was full.
\end{reviewq}

%=====================================================================
\chapter{Cardinality, Participation and Weak Entities}\label{ch:cardinality}
%=====================================================================
\locator{2}{Part II --- Designing It}

\begin{outcomes}
By the end of this chapter you will be able to:
\begin{tight}
  \item \textbf{Determine} the cardinality ratio of a relationship --- 1:1,
        1:N or M:N --- by asking two questions in a fixed order.
  \item \textbf{Distinguish} total from partial participation, and
        \textbf{state} what each one forbids.
  \item \textbf{Write} a relationship's constraints in $(\min,\max)$ notation
        and \textbf{convert} between that and the ratio-plus-participation
        form.
  \item \textbf{Recognise} a weak entity type, and \textbf{identify} its owner,
        its identifying relationship and its partial key.
  \item \textbf{Read} the same diagram in Chen and crow's-foot notation.
\end{tight}
\end{outcomes}

\begin{prereq}
\chref{ermodel}: entities, attributes, keys, relationships and degree. You
will extend the diagram you drew in that chapter's laboratory, so have it with
you.
\end{prereq}

\begin{keyterms}
cardinality ratio $\bullet$ one-to-one $\bullet$ one-to-many $\bullet$
many-to-many $\bullet$ participation constraint $\bullet$ total participation
$\bullet$ partial participation $\bullet$ existence dependency $\bullet$
structural constraint $\bullet$ $(\min,\max)$ notation $\bullet$ weak entity
type $\bullet$ strong entity type $\bullet$ owner entity $\bullet$ identifying
relationship $\bullet$ partial key $\bullet$ discriminator $\bullet$
crow's-foot notation
\end{keyterms}

\section{A Diagram Without Cardinalities Says Almost Nothing}

The diagram from \chref{ermodel} says students enrol in sections. It does not
say whether a student may take two sections of the same course, whether a
section can run with nobody in it, or whether an instructor must be teaching
something. Those are the facts an examinations officer actually argues about,
and all three are cardinality or participation.

The two constraints answer different questions and are routinely confused, so
they are worth separating before either is defined.

\begin{definitionbox}[Two different questions]
The \term{cardinality ratio} answers \emph{how many}: given one entity on this
side, how many on the other side may it be related to --- one, or more than
one?

\smallskip
The \term{participation constraint} answers \emph{must it at all}: can an
entity on this side exist without taking part in the relationship?

\smallskip
They are independent. Every combination occurs, and a relationship is not
specified until both have been answered for \emph{both} ends --- four answers
in all.
\end{definitionbox}

\section{Cardinality, by Asking Two Questions}

\begin{definitionbox}[The three ratios]
\term{One-to-one (1:1)}: each entity on either side relates to at most one on
the other. A department has one head; a person heads at most one department.

\smallskip
\term{One-to-many (1:N)}: one entity on the ``one'' side relates to many on
the ``many'' side, but each of those relates back to at most one. A department
offers many courses; each course belongs to one department.

\smallskip
\term{Many-to-many (M:N)}: entities on both sides may relate to many. A student
enrols in many sections; a section holds many students.
\end{definitionbox}

\begin{conceptbox}[The procedure, and it is only two questions]
Never guess the ratio. Ask these, in this order, out loud, with a concrete
entity in each:

\begin{tight}
  \item \textbf{``One \emph{student} --- how many \emph{sections}?''}
        \quad One, or many?
  \item \textbf{``One \emph{section} --- how many \emph{students}?''}
        \quad One, or many?
\end{tight}

\smallskip
Two answers, four combinations, three ratios:

\medskip
\begin{tabular}{@{}L{3.3cm}L{3.3cm}L{3.0cm}L{4.1cm}@{}}
\textbf{Q1} & \textbf{Q2} & \textbf{Ratio} & \textbf{Example} \\
one  & one  & 1:1 & department --- head \\
many & one  & 1:N & department --- course \\
one  & many & 1:N & \emph{the same, read backwards} \\
many & many & M:N & student --- section \\
\end{tabular}

\smallskip
\textbf{Rows 2 and 3 are the same relationship.} ``1:N'' is meaningless until
you say which end is the one. Always write it as \ent{department}
1---N \ent{course}, never as a bare ``1:N''.
\end{conceptbox}

\begin{examplebox}[The question that catches everybody]
\emph{``A student enrols in many sections. A section has many students.
Therefore M:N.''} --- correct. Now:

\medskip
\emph{``A student enrols in many \textbf{courses}. A course has many students.
Therefore M:N.''} --- also correct, and \textbf{it is the wrong relationship}.

\smallskip
A student does not enrol in a course. They enrol in a \emph{section} of a
course, in a particular semester, taught by a particular instructor, in a
particular room. Model the relationship to \ent{course} and you can never
record which section somebody attended, which instructor taught them, or that
they failed in Fall and passed in Spring.

\smallskip
\textbf{The lesson is not about cardinality.} It is that a cardinality
question asked about the wrong pair of entity types gives a confident,
well-formed, useless answer. Settle what the entities are
(\chref{ermodel}) before asking how many.
\end{examplebox}

\section{Participation: Must It Take Part at All?}

\begin{definitionbox}[Total and partial participation]
\term{Total participation} of an entity type in a relationship means
\emph{every} entity of that type must take part. It is drawn as a double line
and is also called an \term{existence dependency}: a section cannot exist
without a course.

\smallskip
\term{Partial participation} means some entities may take part and others need
not. A single line. A student may be enrolled in nothing --- newly admitted,
on leave --- and is still a student.

\smallskip
Total participation is the stronger claim, and therefore the one to be
sceptical about: it forbids something, and somebody will eventually need the
thing it forbids.
\end{definitionbox}

\begin{alertbox}[Total participation is a promise about Tuesday afternoon]
``Every section must be taught by an instructor'' sounds obviously true, and it
is how timetables are published. Then:

\begin{tight}
  \item A section is created in June and the instructor is appointed in August.
        For two months it has none.
  \item An instructor resigns mid-semester. For a week the section has none.
  \item A section is cancelled but its enrolment records must be kept.
\end{tight}

\smallskip
If participation is total, the database \emph{forbids} all three, and the clerk
who meets them invents a fake instructor called ``TBA'' --- which now appears
in the staff list, the workload report and the payroll query. A constraint
that gets worked around is worse than one that was never there, because the
data now lies and nothing says so.

\smallskip
\textbf{The decision rule.} Make participation total when the entity is
genuinely meaningless without the relationship --- a section without a course
is not a thing at all. Make it partial when the relationship is merely
\emph{usual}. ``Always'' and ``usually'' are different words and the database
can tell.
\end{alertbox}

\begin{definitionbox}[$(\min,\max)$ notation]
Chen's two constraints can be written as one pair on each side. Against each
entity type's line, $(\min,\max)$ --- together, the relationship's \term{structural
constraints} --- gives the smallest and largest number of
relationship instances \emph{one entity of that type} may take part in.

\smallskip
$\min = 0$ is partial participation; $\min \ge 1$ is total. $\max = 1$ is the
``one'' side; $\max = N$ is the ``many'' side.

\smallskip
\textbf{The notation is read against the entity it is nearest}, which is the
opposite of where the 1 and N labels go in classic Chen notation. This single
difference is responsible for most misread diagrams, and is why many
examiners insist on one or the other but not both.
\end{definitionbox}

\begin{worked}[The university's four relationships, specified completely]
Four answers each: ratio from both sides, participation from both sides.

\medskip
\textbf{1. \ent{department} offers \ent{course}}
\begin{tight}
  \item One department --- many courses. One course --- one department.
        \textbf{1:N.}
  \item Must every department offer a course? A new department may not yet.
        \textbf{Partial} (0,N).
  \item Must every course belong to a department? Yes; it is how it is
        numbered. \textbf{Total} (1,1).
\end{tight}
$\Rightarrow$ \ent{department}~(0,N)~---~\ent{offers}~---~(1,1)~\ent{course}

\medskip
\textbf{2. \ent{instructor} teaches \ent{section}}
\begin{tight}
  \item One instructor --- many sections. One section --- one instructor
        (this university's rule; some allow co-teaching, and then it is M:N).
        \textbf{1:N.}
  \item Must every instructor teach? On sabbatical, no. \textbf{Partial} (0,N).
  \item Must every section have an instructor? See the alert box above:
        \textbf{partial} (0,1), deliberately.
\end{tight}

\medskip
\textbf{3. \ent{student} enrols in \ent{section}}
\begin{tight}
  \item Many each way. \textbf{M:N}, and it carries \texttt{grade}.
  \item Must every student be enrolled in something? No. \textbf{Partial} (0,N).
  \item Must every section have students? A section with nobody in it is
        cancelled, not impossible. \textbf{Partial} (0,N).
\end{tight}

\medskip
\textbf{4. \ent{course} requires \ent{course}} (recursive)
\begin{tight}
  \item A course may require several; a course may be required by several.
        \textbf{M:N}, both ends partial (0,N) --- most courses have no
        prerequisites and are a prerequisite for nothing.
\end{tight}

\medskip
\textbf{What the four tell you.} Three of the four are partial on both ends.
Beginners' diagrams are covered in double lines, because every relationship
\emph{sounds} compulsory when you describe the normal case. The normal case is
not the constraint.
\end{worked}

\section{Weak Entities}

\begin{definitionbox}[Weak entity type, owner, identifying relationship, partial key]
A \term{weak entity type} has no key of its own. Its entities are
distinguishable only in combination with the entity that owns them.

\smallskip
The \term{owner} (or identifying) entity type is the strong one it depends on.
The \term{identifying relationship} joins them, drawn as a double diamond, and
the weak entity's participation in it is always \emph{total}.

\smallskip
The \term{partial key} (or \term{discriminator}) is the attribute that tells
one weak entity from another \emph{belonging to the same owner}. It is
underlined with a dashed line.

\smallskip
The full key of a weak entity is \textbf{the owner's primary key plus the
partial key}.
\end{definitionbox}

\begin{examplebox}[\ent{section} is weak, and this is why]
Sections are numbered 1, 2, 3 within a course. Section 1 of CS-103 and section
1 of CS-201 are different sections with the same number.

\begin{tight}
  \item \texttt{section\_no} alone is not a key --- dozens of sections are
        numbered 1.
  \item \texttt{section\_no} \emph{within one course} is unique. That makes it
        a \textbf{partial key}.
  \item A section cannot exist without a course. \textbf{Total participation}
        in the identifying relationship.
  \item Its real key is $\{$\texttt{course\_code}, \texttt{semester},
        \texttt{section\_no}$\}$ --- the owner's key, plus the semester,
        plus the discriminator.
\end{tight}

\smallskip
\relation{section}{\relkey{course\_code}, \relkey{semester},
\relkey{section\_no}, room, slot, emp\_id}

\smallskip
\textbf{The alternative, and why it is not better.} You could give
\ent{section} a surrogate \texttt{section\_id} and make it strong. Many real
systems do, and \chref{mapping} shows the mechanics. It does not make the
entity any less dependent --- it only hides the dependency from the key, and
the uniqueness of (course, semester, number) must still be declared or the
same section can be created twice.
\end{examplebox}

\dsfig{%
\begin{tikzpicture}[
  ent/.style={draw=primary, line width=1.1pt, fill=primary!8,
              text=primary!60!black, font=\scriptsize\bfseries,
              minimum width=1.9cm, minimum height=0.68cm},
  weak/.style={ent, draw=purplehl, fill=purplehl!8, text=purplehl!65!black,
               double, double distance=1.2pt},
  rel/.style={draw=greenhl, line width=0.9pt, fill=greenhl!8,
              text=greenhl!55!black, font=\scriptsize\bfseries,
              diamond, aspect=2.2, inner sep=0.5pt},
  irel/.style={rel, draw=purplehl, fill=purplehl!8, text=purplehl!65!black,
               double, double distance=1.2pt},
  one/.style={draw=neutral!60, line width=0.7pt},
  two/.style={draw=neutral!60, line width=0.7pt, double,
              double distance=1.1pt},
  mm/.style={font=\scriptsize\bfseries, text=accent, inner sep=1.2pt},
  lbl/.style={font=\scriptsize, text=neutral}
]
  % Centres 3.2 apart on both rows. The first draft used 2.3 and the TEACHES
  % diamond printed on top of the INSTRUCTOR box -- a diamond of aspect 2.2 is
  % wider than it looks, and three of the (min,max) labels landed on boxes for
  % the same reason. The labels now sit at pos=0.22, which also puts each one
  % nearest the entity it is read against, as the definition requires.
  \node[ent]  (dep) at (-7.6, 1.7) {DEPARTMENT};
  \node[rel]  (off) at (-3.8, 1.7) {OFFERS};
  \node[ent]  (crs) at ( 0.0, 1.7) {COURSE};
  \node[rel]  (req) at ( 0.0, 3.9) {REQUIRES};
  \node[irel] (has) at ( 0.0, 0.0) {HAS};
  \node[weak] (sec) at ( 0.0,-1.9) {SECTION};
  \node[ent]  (stu) at (-7.6,-1.9) {STUDENT};
  \node[rel]  (enr) at (-3.8,-1.9) {ENROLS};
  \node[rel]  (tch) at ( 3.8,-1.9) {TEACHES};
  \node[ent]  (ins) at ( 7.6,-1.9) {INSTRUCTOR};

  \draw[one] (dep) -- node[mm, above=1pt, pos=0.35] {(0,N)} (off);
  \draw[two] (off) -- node[mm, above=1pt, pos=0.65] {(1,1)} (crs);
  \draw[one] (crs) -- node[mm, right=1pt, pos=0.3] {(0,N)} (has);
  \draw[two] (has) -- node[mm, right=1pt, pos=0.7] {(1,1)} (sec);
  \draw[one] (stu) -- node[mm, above=1pt, pos=0.35] {(0,N)} (enr);
  \draw[one] (enr) -- node[mm, above=1pt, pos=0.65] {(0,N)} (sec);
  \draw[one] (sec) -- node[mm, above=1pt, pos=0.35] {(0,1)} (tch);
  \draw[one] (tch) -- node[mm, above=1pt, pos=0.65] {(0,N)} (ins);

  % The recursive relationship, as a loop above COURSE so that both of its
  % ends are visible and can carry role names. Drawn flat on the right of the
  % course box it read as a line passing over the top to nothing.
  \draw[one] (crs.north west) .. controls (-1.9,2.6) and (-1.9,3.9)
             .. (req.west);
  \draw[one] (crs.north east) .. controls ( 1.9,2.6) and ( 1.9,3.9)
             .. (req.east);
  \node[lbl, anchor=east, font=\scriptsize\itshape] at (-1.25,3.15)
       {is prerequisite for};
  \node[lbl, anchor=west, font=\scriptsize\itshape] at ( 1.25,3.15)
       {depends on};
  \node[mm, anchor=south east] at (-0.35,2.15) {(0,N)};
  \node[mm, anchor=south west] at ( 0.35,2.15) {(0,N)};

  \node[lbl, text=purplehl!70!black, text width=4.2cm, align=left,
        anchor=west] at (0.95,0.0)
       {double diamond and double line: identifying, and always total};
  \node[lbl, text=purplehl!70!black, anchor=north, text width=3.4cm,
        align=center] at (0.0,-2.45)
       {weak: no key of its own};
  % No explicit \\ here. An align=right node whose forced line does not fill
  % its text width is reported as an underfull hbox, twice, once per line.
  % Letting it wrap on its own costs nothing and the gate stays clean.
  \node[lbl, text=accent, anchor=east, align=right, text width=3.8cm]
       at (-4.6,0.3) {$(\min,\max)$ is read against the entity it touches};
\end{tikzpicture}}%
{The university database, fully constrained. Only two lines in the whole
diagram are double --- a course must belong to a department and a section must
belong to a course. Everything else is partial, which is the usual
proportion.}{ercomplete}

\begin{conceptbox}[Chen or crow's foot? Both, and here is the translation]
Textbooks and examination papers use Chen. Every tool you will meet at work
uses crow's foot, because it fits on a screen. They say the same things.

\medskip
\begin{tabular}{@{}L{4.4cm}L{4.6cm}L{4.7cm}@{}}
\textbf{Meaning} & \textbf{Chen} & \textbf{Crow's foot} \\
relationship & a diamond between boxes & a line between boxes \\
``many'' end & N on the line & the three-pronged foot \\
``one'' end & 1 on the line & a single cross-stroke \\
total participation & double line & a cross-stroke nearest the box \\
partial participation & single line & a circle nearest the box \\
attributes & ellipses around the box & listed inside the box \\
weak entity & double rectangle & usually not distinguished \\
M:N relationship & one diamond & \textbf{cannot be drawn} --- it is
resolved into two 1:N at once \\
\end{tabular}

\smallskip
\textbf{The last row is the real difference.} Crow's-foot notation has no way
to draw a many-to-many relationship, so a tool forces you to resolve it into a
third box before you have decided to. Chen lets you keep the design honest
until \chref{mapping}, where the resolution is a rule with a reason rather
than something the drawing program did to you.
\end{conceptbox}

\begin{pitfall}
\begin{tight}
  \item \textbf{Writing ``1:N'' without saying which end is the one.} The
        label alone is ambiguous; name the entity types.
  \item \textbf{Confusing cardinality with participation.} ``Every section has
        one instructor'' states both --- one is the ratio, every is the
        participation --- and people hear only the first.
  \item \textbf{Double lines everywhere.} Every relationship sounds compulsory
        when the normal case is described. Three of this database's four are
        partial on both ends.
  \item \textbf{Reading $(\min,\max)$ against the far entity.} It is read
        against the entity it touches, which is the opposite of the 1/N labels.
  \item \textbf{Asking the cardinality of the wrong pair.} Student-to-course is
        M:N and is the wrong relationship; student-to-section is the right one.
  \item \textbf{Giving a weak entity a surrogate key and calling it strong.}
        The dependency does not go away; it just stops being visible, and the
        uniqueness of the real key must still be declared.
\end{tight}
\end{pitfall}

\begin{lab}[Constrain your diagram, then try to break it]
Take the diagram from \chref{ermodel}'s laboratory. Still no database.

\begin{lstlisting}
# --- 1. Four answers per relationship -------------------------------
    For EVERY relationship, write all four:
        one <A> -- how many <B>?        one <B> -- how many <A>?
        must every <A> take part?       must every <B> take part?
    Four answers. Not three.

# --- 2. Mark the ratio ----------------------------------------------
    1:1, 1:N or M:N, naming the entity types -- not a bare "1:N".

# --- 3. Mark participation ------------------------------------------
    Double line for total, single for partial.

# --- 4. Rewrite it in (min,max) -------------------------------------
    Every line gets a pair at each end. Check each against step 1;
    they must agree, and if they do not, step 1 was guessed.

# --- 5. Find the weak entity ----------------------------------------
    There is at least one. Name its owner, its identifying
    relationship, its partial key and its full key.

# --- 6. Attack every double line ------------------------------------
    For EACH total participation, invent a real Tuesday on which the
    university needs to break it. If you can, make it partial.
    Write the Tuesday down next to the line; it is the evidence.

# --- 7. Attack every "one" ------------------------------------------
    For each 1 end, ask what happens if the university someday allows
    two. Co-teaching. Joint degrees. Two heads of department.
    Which of these would cost a schema change, and which would not?

# --- 8. Read it back ------------------------------------------------
    Turn the finished diagram into ten English sentences, one per
    constraint. Give them to the person who has the paragraph from
    Chapter 5's lab. Every disagreement is a decision you had not
    actually made.
\end{lstlisting}

\smallskip
\textbf{What to notice.} Steps 6 and 7 usually remove more constraints than
they confirm, and that is the expected result rather than a sign you did it
wrong. A design is not improved by forbidding more; it is improved by
forbidding exactly what the university actually forbids.
\end{lab}

\begin{checkpoint}
\begin{tightnum}
  \item A relationship is drawn with a double line at the \ent{section} end
        of \ent{teaches}. State in one sentence what the database will
        refuse to do, and give a realistic situation in which somebody needs it
        to do exactly that.
  \item \ent{enrolment} has $(0,N)$ at the \ent{student} end and $(0,N)$
        at the \ent{section} end. Rewrite this as a cardinality ratio plus
        two participation constraints, and say what would change if the
        \ent{section} end became $(1,N)$.
  \item Explain why \ent{section} is a weak entity type, name its partial
        key and its full key, and then say what you would have to declare
        instead if you chose to give it a surrogate \texttt{section\_id} ---
        and what breaks if you forget.
\end{tightnum}
\end{checkpoint}

\begin{chaptersummary}
\begin{tight}
  \item A relationship needs \textbf{four} answers: the ratio from each side
        and the participation from each side. Fewer than four is an
        unspecified relationship.
  \item The \textbf{cardinality ratio} is found by two questions --- ``one A,
        how many B?'' and ``one B, how many A?'' --- giving 1:1, 1:N or M:N.
        Always name which end is the one.
  \item \textbf{Total participation} (double line) means every entity must take
        part; \textbf{partial} means it need not. Total is the stronger claim
        and deserves the scepticism.
  \item A constraint that gets worked around is worse than no constraint: the
        data then contains a fake row and nothing marks it as fake.
  \item $(\min,\max)$ carries both facts at once: $\min=0$ is partial,
        $\min\ge 1$ is total, $\max=1$ is the one side. \textbf{It is read
        against the entity it touches}, unlike the 1/N labels.
  \item A \textbf{weak entity type} has no key of its own. It has an
        \textbf{owner}, an \textbf{identifying relationship} (double diamond,
        always total) and a \textbf{partial key}; its full key is the owner's
        key plus the partial key.
  \item \ent{section} is weak: section 1 exists in every course, and only
        (course, semester, number) identifies one.
  \item Crow's-foot notation says the same things more compactly but
        \textbf{cannot draw an M:N relationship}, so a tool resolves it for you
        before you have decided to.
  \item Asking a cardinality question about the wrong pair of entity types
        gives a confident and useless answer. Student-to-\emph{course} is the
        classic case.
\end{tight}
\end{chaptersummary}

\begin{reviewq}
  \item \textbf{(MCQ)} A double line between an entity type and a relationship
        means: (a)~many-to-many; (b)~total participation; (c)~a weak entity;
        (d)~an identifying relationship.
  \item \textbf{(MCQ)} In $(\min,\max)$ notation, $(1,1)$ against
        \ent{course} on the \ent{offers} line means: (a)~a department
        offers exactly one course; (b)~a course belongs to exactly one
        department; (c)~there is one department; (d)~the relationship is 1:1.
  \item \textbf{(Short)} Give an example of a 1:1 relationship in a university
        and say which end you would make total, with a reason.
  \item \textbf{(Short)} Why is the participation of a weak entity type in its
        identifying relationship always total? Answer from the definition of a
        weak entity.
  \item \textbf{(Short)} A tool that draws crow's-foot diagrams will not let
        you record a many-to-many relationship. What does it make you do
        instead, and what is lost by doing it at that moment rather than in
        \chref{mapping}?
  \item \textbf{(Applied)} A hospital records patients, doctors, appointments
        and prescriptions. Give all four answers for every relationship you
        identify, state which entity types are weak and why, and defend each
        total participation with a specific situation that would break it if
        you were wrong.
  \item \textbf{(Applied)} Your university decides to allow co-teaching: two
        instructors may share a section. State exactly which constraint in
        Figure~\ref{fig:ercomplete} changes, what the relationship becomes,
        what new attributes it may now need, and --- looking ahead to
        \chref{mapping} --- why this change is cheap on paper and expensive
        once the database is full.
\end{reviewq}

%=====================================================================
\chapter{The Relational Model}\label{ch:relational}
%=====================================================================
\locator{2}{Part II --- Designing It}

\begin{outcomes}
By the end of this chapter you will be able to:
\begin{tight}
  \item \textbf{Define} relation, tuple, attribute, domain, degree and
        cardinality, and \textbf{use} each word correctly.
  \item \textbf{State} the four properties of a relation and \textbf{show}
        which of them a SQL table does not have.
  \item \textbf{Distinguish} a relation schema from a relation instance, and
        \textbf{write} both in the standard notation.
  \item \textbf{Identify} superkeys, candidate keys and the primary key of a
        relation from its data \emph{and} its guarantees.
  \item \textbf{Explain} what a null value is, what it is not, and why it is
        the model's weakest point.
\end{tight}
\end{outcomes}

\begin{prereq}
\chref{ermodel} for keys --- the definitions are the same ones, restated for
relations. \chref{firstdatabase} for the \relation{}{} notation and for
\texttt{sqlite3}, which you will use at the end to break one of this
chapter's rules on purpose.
\end{prereq}

\begin{keyterms}
relational data model $\bullet$ relation $\bullet$ tuple $\bullet$ attribute
$\bullet$ domain $\bullet$ atomic value $\bullet$ relation schema $\bullet$
relation instance $\bullet$ degree $\bullet$ arity $\bullet$ cardinality
$\bullet$ relational database schema $\bullet$ superkey $\bullet$ candidate
key $\bullet$ primary key $\bullet$ foreign key $\bullet$ null value
$\bullet$ unknown $\bullet$ inapplicable
\end{keyterms}

\section{One Idea, and It Is Mathematics}

Codd's 1970 paper proposed that all data in a database be represented one way:
as relations. Not as files, not as linked records, not as trees. One structure,
with a mathematical definition, so that questions about data could be answered
by reasoning rather than by knowing how the file was laid out.

That is the whole of the relational model, and the reason every system in
\chref{whynotfiles}'s list --- Oracle, PostgreSQL, MySQL, SQL Server, SQLite
--- is a different implementation of the same five pages of theory.

\begin{definitionbox}[Domain, attribute, tuple, relation]
A \term{domain} is a named set of \term{atomic} values --- values the model
treats as indivisible. \texttt{semester} might have domain
$\{1,2,\dots,8\}$; \texttt{roll\_no} the set of strings matching the
university's pattern.

\smallskip
An \term{attribute} is a name for the role a domain plays in a relation. Two
attributes may share a domain: \texttt{start\_year} and \texttt{end\_year} are
both years and are not the same attribute.

\smallskip
A \term{tuple} is an ordered list of values, one from each attribute's domain.
Informally, a row.

\smallskip
A \term{relation} is a \emph{set} of tuples over the same attributes.
Informally, a table --- and the gap between ``informally'' and ``exactly'' is
section~7.3.
\end{definitionbox}

\begin{examplebox}[The vocabulary on one small relation]
\relation{student}{\relkey{roll\_no}, name, programme, semester}

\medskip
\begin{tabular}{@{}L{3.2cm}L{3.0cm}L{2.4cm}L{2.2cm}@{}}
\textbf{roll\_no} & \textbf{name} & \textbf{programme} & \textbf{semester} \\
2021-IT-001 & Ayesha Khan & BSIT & 2 \\
2021-IT-002 & Bilal Ahmed & BSIT & 2 \\
2021-CS-007 & Hira Nadeem & BSCS & 4 \\
\end{tabular}

\medskip
\begin{tight}
  \item The \textbf{relation schema} is the heading: the name plus the four
        attributes. It is written \relation{student}{\relkey{roll\_no}, name,
        programme, semester}.
  \item The \textbf{relation instance} is the three tuples. Tomorrow it is
        different; the schema is not.
  \item The \textbf{degree} (or \term{arity}) is \textbf{4} --- the number of
        attributes. It changes only when the schema changes.
  \item The \textbf{cardinality} is \textbf{3} --- the number of tuples. It
        changes every time somebody enrols.
  \item The \textbf{domain} of \texttt{semester} is the integers 1 to 8. The
        value 2 appears twice and that is fine; domains are not keys.
\end{tight}

\smallskip
A \term{relational database schema} is the set of all the relation
schemas in one database --- \textit{student}, \textit{course},
\textit{section} and the rest, taken together.

\smallskip
\textbf{Degree and cardinality are easy to swap by accident.} Degree counts
columns and is a property of the schema; cardinality counts rows and is a
property of the instance. ``A relation of degree 4 and cardinality 3'' names
the table above exactly.
\end{examplebox}

\dsfig{%
\begin{tikzpicture}[
  hd/.style={font=\scriptsize\bfseries, text=white},
  cell/.style={font=\scriptsize\ttfamily, text=primary!65!black},
  note/.style={font=\scriptsize\itshape, text=neutral, align=left},
  tag/.style={font=\scriptsize\bfseries, text=accent},
  br/.style={draw=accent!70, line width=0.8pt,
             decorate, decoration={brace, amplitude=4pt}}
]
  % heading band
  \fill[primary, rounded corners=2pt] (0,0) rectangle (9.2,0.55);
  \node[hd, anchor=west] at (0.15,0.28)  {roll\_no};
  \node[hd, anchor=west] at (2.55,0.28)  {name};
  \node[hd, anchor=west] at (5.15,0.28)  {programme};
  \node[hd, anchor=west] at (7.65,0.28)  {semester};
  % body
  \foreach \i/\a/\b/\c/\d in {%
      1/2021-IT-001/Ayesha Khan/BSIT/2,
      2/2021-IT-002/Bilal Ahmed/BSIT/2,
      3/2021-CS-007/Hira Nadeem/BSCS/4}{
    \pgfmathsetmacro{\yy}{-0.42*\i}
    \ifodd\i\fill[primary!5] (0,\yy-0.17) rectangle (9.2,\yy+0.23);\fi
    \node[cell, anchor=west] at (0.15,\yy) {\a};
    \node[cell, anchor=west] at (2.55,\yy) {\b};
    \node[cell, anchor=west] at (5.15,\yy) {\c};
    \node[cell, anchor=west] at (7.65,\yy) {\d};}
  \draw[primary!35, line width=0.8pt] (0,-1.43) rectangle (9.2,0.55);

  % annotations
  \draw[br] (9.45,0.5) -- node[tag, right=5pt] {} (9.45,-1.4);
  \node[tag, anchor=west, align=left] at (9.9,-0.45)
       {cardinality\\= 3 tuples};
  \draw[br] (0,0.85) -- (9.2,0.85);
  \node[tag, anchor=south] at (4.6,1.0) {degree = 4 attributes};

  \node[note, anchor=east, text width=3.3cm] at (-0.35,0.28)
       {\textbf{\color{primary}heading}\\the relation \emph{schema}};
  \node[note, anchor=east, text width=3.3cm] at (-0.35,-0.85)
       {\textbf{\color{primary}body}\\the relation \emph{instance}};
  \node[note, anchor=north west, text width=5.4cm] at (0,-1.75)
       {\textbf{one tuple} --- a value from each attribute's domain, every one
        of them \emph{atomic}};
  \node[note, anchor=north west, text width=5.2cm] at (5.4,-1.75)
       {\textbf{a set} --- no tuple repeats, and no tuple is first};
  \draw[-{Stealth[length=1.8mm]}, neutral!60] (1.6,-1.72) -- (1.3,-1.3);
  \draw[-{Stealth[length=1.8mm]}, neutral!60] (7.4,-1.72) -- (7.4,-1.45);
\end{tikzpicture}}%
{One relation, with every word of the vocabulary attached to the thing it
names. Degree belongs to the heading and changes only when the schema does;
cardinality belongs to the body and changes all day.}{relanatomy}

\section{The Four Properties of a Relation}

\begin{definitionbox}[What makes a relation a relation]
\begin{tightnum}
  \item \textbf{No duplicate tuples.} A relation is a \emph{set}. The same
        tuple cannot appear twice --- which is another way of saying every
        relation has a key, if only the whole tuple.
  \item \textbf{Tuples are unordered.} There is no first row. Two relations
        with the same tuples in a different sequence are the same relation.
  \item \textbf{Attributes are unordered.} Attributes are referred to by name,
        not position. \texttt{(name, roll\_no)} and \texttt{(roll\_no, name)}
        are the same heading.
  \item \textbf{Every value is atomic.} One value per attribute per tuple: no
        lists, no nested tables, no ``three phone numbers in one field''.
\end{tightnum}
\end{definitionbox}

\begin{conceptbox}[Why each property is worth having]
These are not arbitrary tidiness. Each one buys something specific.

\begin{tight}
  \item \textbf{No duplicates} means every fact is stored once and can be
        addressed. Two identical rows cannot be told apart, so neither can be
        updated or deleted individually --- which is the file-based system's
        redundancy problem (\chref{whynotfiles}) reappearing inside the
        database.
  \item \textbf{Unordered tuples} is what lets the optimiser
        (\chref{whodoeswhat}) choose any plan it likes. If row order were part
        of the meaning, no index could ever be used and \chref{indexes} would
        not exist.
  \item \textbf{Unordered attributes} is what lets you add a column without
        breaking programs --- the logical data independence of
        \chref{architecture}, and the reason \sqlkw{SELECT *} is a bad habit:
        it depends on an order the model says is not there.
  \item \textbf{Atomic values} is what makes the query language finite. If a
        field could hold a list, every query would need loops, and relational
        algebra (\chref{algebra}) could not be written down. This property is
        also \textbf{first normal form}, arrived at from the other direction
        in \chref{normalforms}.
\end{tight}
\end{conceptbox}

\begin{alertbox}[A SQL table is not a relation, and the difference is measurable]
SQL was designed around the relational model and departs from it in three
places. You should know all three, because each one is a bug waiting to
happen.

\medskip
\begin{tabular}{@{}L{4.0cm}L{4.6cm}L{5.1cm}@{}}
\textbf{Relation} & \textbf{SQL table} & \textbf{What follows} \\
no duplicate tuples & \textbf{duplicates allowed} unless a key is declared &
  a table without a key is not a relation at all \\
tuples unordered & \sqlkw{ORDER BY} imposes one, and without it the order is
  arbitrary & never rely on the order you happen to see \\
values atomic & mostly enforced, but a text column will hold
  \texttt{"0300-1234567, 0321-9876543"} quite happily &
  the model cannot stop you; only you can \\
\end{tabular}

\smallskip
The first is the serious one and you can demonstrate it in four lines:

\begin{lstlisting}[language=sql]
CREATE TABLE t (a TEXT, b TEXT);        -- no key declared
INSERT INTO t VALUES ('x','y'), ('x','y');
SELECT COUNT(*) FROM t;                 -- 2
DELETE FROM t WHERE a='x' AND b='y';    -- deletes BOTH
\end{lstlisting}

Two identical rows, and no \sqlkw{WHERE} clause in the language can address one
of them, because there is nothing that distinguishes them. \textbf{This is why
every table in this course declares a primary key}, and \chref{constraints} is
about making the declaration real.
\end{alertbox}

\section{Keys, Restated for Relations}

The definitions are \chref{ermodel}'s, which is the point: the ER model and the
relational model agree about keys, and that agreement is what makes
\chref{mapping} possible.

\begin{definitionbox}[Keys of a relation]
A \term{superkey} is a set of attributes $K$ such that no two tuples of any
legal instance agree on all of $K$.

\smallskip
A \term{candidate key} is a minimal superkey: no proper subset of it is a
superkey.

\smallskip
The \term{primary key} is the chosen candidate key. Its attributes may not be
null --- an identifier that is missing identifies nothing.

\smallskip
A \term{foreign key} is a set of attributes in one relation whose values must
appear as the primary key of another (or be null). It is how relations are
joined to each other, and it is \chref{constraints}'s subject.
\end{definitionbox}

\begin{worked}[Reading keys off an instance, and why you cannot]
\begin{tabular}{@{}L{2.8cm}L{2.6cm}L{2.2cm}L{2.0cm}L{2.4cm}@{}}
\textbf{roll\_no} & \textbf{name} & \textbf{prog} & \textbf{sem}
& \textbf{email} \\
2021-IT-001 & Ayesha Khan & BSIT & 2 & ayesha@iub.pk \\
2021-IT-002 & Bilal Ahmed & BSIT & 2 & bilal@iub.pk \\
2021-CS-007 & Hira Nadeem & BSCS & 4 & hira@iub.pk \\
\end{tabular}

\medskip
\textbf{Step 1 --- what is unique in this instance?} All of
\texttt{roll\_no}, \texttt{name}, \texttt{email}; also
$\{$\texttt{prog}, \texttt{sem}$\}$ is \emph{not} --- rows 1 and 2 agree.

\medskip
\textbf{Step 2 --- which of those are superkeys?} This is where the exercise
stops being mechanical. A superkey is a constraint on \emph{every legal
instance}, and three rows cannot establish one.
\begin{tight}
  \item \texttt{roll\_no}: the university guarantees it. \textbf{Superkey.}
  \item \texttt{email}: issued one per student. \textbf{Superkey} --- on that
        guarantee, not on these rows.
  \item \texttt{name}: unique \emph{here}. The university guarantees nothing.
        \textbf{Not a superkey}, and the data cannot tell you so.
\end{tight}

\medskip
\textbf{Step 3 --- minimality.} $\{$\texttt{roll\_no}$\}$ and
$\{$\texttt{email}$\}$ are both minimal, so both are candidate keys.
$\{$\texttt{roll\_no}, \texttt{sem}$\}$ is a superkey and not minimal.

\medskip
\textbf{Step 4 --- choose.} \texttt{roll\_no}, on \chref{ermodel}'s four
grounds: unique, always present, never changing, short. \texttt{email} fails
``never changing''.

\medskip
\textbf{The asymmetry worth remembering.} An instance can \emph{disprove} a key
instantly --- one duplicate and it is settled. It can never prove one. So:
check the data to find candidates, then ask a human to confirm each, and
record the answer. The examination version of this question gives you the
guarantees; the real version does not, and you have to go and ask.
\end{worked}

\section{Null, and Why It Is the Awkward Part}

\begin{definitionbox}[Null]
A \term{null} is a marker in a tuple meaning that no value from the domain is
recorded for that attribute. It is \textbf{not} zero, \textbf{not} the empty
string, and \textbf{not} a value of any kind --- it is the absence of one.

\smallskip
Codd distinguished two reasons for a null and most systems, including SQLite,
distinguish neither:
\begin{tight}
  \item \term{unknown} --- the student has a CNIC, we do not know it;
  \item \term{inapplicable} --- the student has no CNIC, because they are
        fifteen.
\end{tight}
One marker, two meanings, and no way to tell them apart in a query.
\end{definitionbox}

\begin{examplebox}[Three nulls in one row, meaning three different things]
\begin{tabular}{@{}L{2.9cm}L{3.4cm}L{7.6cm}@{}}
\textbf{Attribute} & \textbf{Value} & \textbf{What the null means} \\
\texttt{cnic}        & \textsc{null} & unknown --- they have one, it is not on file \\
\texttt{thesis\_title} & \textsc{null} & inapplicable --- second-semester
                                        students do not write one \\
\texttt{grade}       & \textsc{null} & not yet --- the semester is not over \\
\end{tabular}

\smallskip
All three print as blank and all three behave identically in every comparison.
A query that counts ``students with no CNIC on file'' will silently include
anybody whose CNIC was inapplicable, and nothing in the database records that
the two cases were different.

\smallskip
\textbf{What a designer does about it.} Where the distinction matters, make it
explicit: a separate \texttt{cnic\_status} attribute, or a different relation
for thesis students. Where it does not, allow the null and write the ambiguity
down. The one thing not to do is invent a magic value --- \texttt{'N/A'},
\texttt{0000-0000000}, \texttt{9999-12-31} --- which is a null that the
database does not know is a null, and which will be counted, averaged and
sorted as though it were real.
\end{examplebox}

\begin{alertbox}[Null is ahead of you in three later chapters]
This chapter only names the problem. It returns, each time worse:

\begin{tight}
  \item \chref{select}: \texttt{WHERE grade = NULL} matches nothing, ever ---
        not even the rows where grade is null. Comparison with null is neither
        true nor false but \emph{unknown}, and SQL's logic has three values.
  \item \chref{grouping}: \sqlkw{COUNT(*)} and \sqlkw{COUNT(grade)} give
        different answers on the same rows, and \sqlkw{AVG} silently ignores
        nulls rather than treating them as zero --- which is usually right and
        occasionally catastrophic.
  \item \chref{joins}: a row whose foreign key is null disappears from an inner
        join, so the section with no instructor vanishes from the timetable
        report and nobody notices until a Tuesday.
\end{tight}

\smallskip
Nulls are the single largest source of quietly wrong answers in SQL. Every one
you allow in a schema is a decision; allow them deliberately.
\end{alertbox}

\begin{pitfall}
\begin{tight}
  \item \textbf{Saying ``table'' when the question asks about a relation.} In
        an examination the distinction is the answer: a table may have
        duplicates and an order.
  \item \textbf{Swapping degree and cardinality.} Degree counts attributes,
        cardinality counts tuples.
  \item \textbf{Relying on row order.} Without \sqlkw{ORDER BY} there is none,
        and the order you see today comes from the plan, which
        \chref{indexes} will change.
  \item \textbf{Reading keys off sample data.} Data disproves keys; it never
        proves them.
  \item \textbf{Treating null as zero or as the empty string.} It is neither,
        and the arithmetic will show it.
  \item \textbf{Inventing a magic value instead of a null.} It gets averaged.
  \item \textbf{Creating a table with no primary key.} It is not a relation,
        and you have given up the ability to address a row.
\end{tight}
\end{pitfall}

\begin{lab}[Break three properties of a relation, deliberately]
\begin{lstlisting}[language=bash]
# --- 1. A scratch database; keep university.db clean ----------------
sqlite3 relmodel.db
\end{lstlisting}
\begin{lstlisting}[language=sql]
-- --- 2. A table that is not a relation --------------------------
.headers on
.mode box
CREATE TABLE nokey (roll_no TEXT, name TEXT);
INSERT INTO nokey VALUES ('2021-IT-001','Ayesha Khan');
INSERT INTO nokey VALUES ('2021-IT-001','Ayesha Khan');
SELECT COUNT(*) FROM nokey;
-- Property 1 broken. Now try to delete exactly one of them.
DELETE FROM nokey WHERE roll_no = '2021-IT-001' AND name = 'Ayesha Khan';
SELECT COUNT(*) FROM nokey;
-- Write down how many rows went, and say why you could not keep one.

-- --- 3. The same table, as a relation ---------------------------
CREATE TABLE withkey (roll_no TEXT PRIMARY KEY, name TEXT NOT NULL);
INSERT INTO withkey VALUES ('2021-IT-001','Ayesha Khan');
INSERT INTO withkey VALUES ('2021-IT-001','Ayesha Khan');
-- Record the error. One keyword is the whole difference.

-- --- 4. Prove tuples are unordered ------------------------------
INSERT INTO withkey VALUES ('2021-CS-007','Hira Nadeem');
INSERT INTO withkey VALUES ('2021-IT-002','Bilal Ahmed');
SELECT * FROM withkey;
CREATE INDEX idx_nm ON withkey(name);
SELECT * FROM withkey;
-- Same relation, same query. Did the order change? Either way,
-- say what you are entitled to assume about it.

-- --- 5. Break atomicity, which nothing will stop you doing ------
CREATE TABLE badphone (roll_no TEXT PRIMARY KEY, phones TEXT);
INSERT INTO badphone VALUES
  ('2021-IT-001','0300-1234567, 0321-9876543');
-- Now answer: which students have a number starting 0321?
SELECT roll_no FROM badphone WHERE phones LIKE '%0321%';
-- It appears to work. Now explain, in one sentence, why
-- '0321' in the MIDDLE of a number would also match, and what
-- you would have to write to do this correctly.

-- --- 6. Three nulls, one appearance -----------------------------
CREATE TABLE nulls (roll_no TEXT PRIMARY KEY, cnic TEXT, grade TEXT);
INSERT INTO nulls VALUES ('2021-IT-001', NULL, NULL);
SELECT * FROM nulls WHERE cnic IS NULL;
SELECT * FROM nulls WHERE cnic = NULL;
-- Two queries, one of them returns nothing. Which, and WHY?
-- This is the whole of Chapter 14 in two lines; write your
-- explanation down now and check it when you get there.
\end{lstlisting}

\smallskip
\textbf{What to notice.} Step 2 is the one to keep. Two identical rows, and the
language has no way to refer to one of them --- not because SQLite is
deficient, but because there is genuinely nothing that tells them apart. Every
primary key you declare for the rest of this course is buying your way out of
that position.
\end{lab}

\begin{checkpoint}
\begin{tightnum}
  \item A relation has degree 5 and cardinality 200. Say precisely what each
        number counts, which of the two can change without a schema change,
        and what operation changes the other.
  \item Give a specific query that returns a different answer depending on
        whether duplicate tuples are permitted, and say what you would declare
        to make the answer reliable.
  \item \texttt{grade} is null for a student who has not sat the examination
        and also for one whose paper was lost. Explain why the database cannot
        distinguish them, give two different designs that could, and say what
        each design costs.
\end{tightnum}
\end{checkpoint}

\begin{chaptersummary}
\begin{tight}
  \item A \textbf{domain} is a set of atomic values; an \textbf{attribute}
        names a domain's role; a \textbf{tuple} picks one value per attribute;
        a \textbf{relation} is a \emph{set} of tuples.
  \item \textbf{Degree} counts attributes and belongs to the schema;
        \textbf{cardinality} counts tuples and belongs to the instance.
  \item Four properties: \textbf{no duplicate tuples}, \textbf{tuples
        unordered}, \textbf{attributes unordered}, \textbf{values atomic}.
  \item Each property buys something: addressability, the optimiser's freedom,
        logical data independence, and a finite query language.
  \item \textbf{A SQL table is not a relation.} It permits duplicates unless a
        key is declared, exposes an order, and cannot stop you putting a list
        in a text column.
  \item A table with no primary key gives up the ability to address a row: two
        identical rows cannot be told apart by any \sqlkw{WHERE} clause.
  \item \textbf{Superkey}, \textbf{candidate key} and \textbf{primary key} mean
        in the relational model exactly what they meant in the ER model, which
        is what makes \chref{mapping} mechanical.
  \item An instance can \textbf{disprove} a key but never prove one. Use the
        data to find candidates and a human to confirm them.
  \item A \textbf{null} marks the absence of a value. It is not zero and not
        the empty string, it conflates \emph{unknown} with
        \emph{inapplicable}, and it is the source of quietly wrong answers in
        Chapters~\ref{ch:joins}, \ref{ch:select} and \ref{ch:grouping}.
  \item Never substitute a magic value for a null. It will be counted, averaged
        and sorted as though it were real.
\end{tight}
\end{chaptersummary}

\begin{reviewq}
  \item \textbf{(MCQ)} The number of attributes in a relation is its:
        (a)~cardinality; (b)~degree; (c)~domain; (d)~arity of its key.
  \item \textbf{(MCQ)} Which property of a relation does a SQL table most
        clearly fail to have? (a)~attributes unordered; (b)~values atomic;
        (c)~no duplicate tuples; (d)~a named schema.
  \item \textbf{(Short)} Two attributes share the domain ``year''. Explain why
        they are still two different attributes, with an example.
  \item \textbf{(Short)} Why can an instance disprove a candidate key but never
        establish one? Illustrate with \texttt{name}.
  \item \textbf{(Short)} State the difference between \texttt{WHERE cnic IS
        NULL} and \texttt{WHERE cnic = NULL}, and say which returns rows.
  \item \textbf{(Applied)} You are given a table of 50{,}000 rows and told only
        that ``it is the student list''. Describe the queries you would run to
        propose a primary key, say exactly what each result would and would not
        establish, and state what you would have to ask a person before
        declaring one.
  \item \textbf{(Applied)} A system stores a student's phone numbers in one
        comma-separated text column. List four distinct questions the
        university cannot reliably answer in that design, show the relation you
        would replace it with, and explain which of the four properties of a
        relation the original violated.
\end{reviewq}

%=====================================================================
\chapter{Keys and Integrity Constraints}\label{ch:constraints}
%=====================================================================
\locator{2}{Part II --- Designing It}

\begin{outcomes}
By the end of this chapter you will be able to:
\begin{tight}
  \item \textbf{State} the three model-based integrity constraints --- domain,
        entity and referential --- and \textbf{give} an example of each being
        violated.
  \item \textbf{Define} a foreign key and \textbf{explain} what referential
        integrity forbids.
  \item \textbf{Choose} an appropriate referential action (\sqlkw{CASCADE},
        \sqlkw{RESTRICT}, \sqlkw{SET NULL}) and \textbf{justify} the choice.
  \item \textbf{Demonstrate} that a declared constraint and an enforced
        constraint are not the same thing.
  \item \textbf{Verify} on any SQLite database whether its foreign keys are
        actually being enforced.
\end{tight}
\end{outcomes}

\begin{prereq}
\chref{relational}: relations, keys, nulls. \chref{cardinality} for
participation, which is the constraint this chapter implements. You will need
\texttt{sqlite3} and \texttt{code/m08\_constraints.py}.
\end{prereq}

\begin{keyterms}
integrity constraint $\bullet$ domain constraint $\bullet$ entity integrity
$\bullet$ referential integrity $\bullet$ foreign key $\bullet$ referencing
relation $\bullet$ referenced relation $\bullet$ orphan row $\bullet$
referential action $\bullet$ \sqlkw{CASCADE} $\bullet$ \sqlkw{RESTRICT}
$\bullet$ \sqlkw{SET NULL} $\bullet$ semantic constraint $\bullet$
\sqlkw{CHECK} $\bullet$ \sqlkw{PRAGMA} $\bullet$ declared versus enforced
\end{keyterms}

\section{Three Constraints the Model Gives You Free}

A schema is not only a list of attributes. It is a list of \emph{promises}, and
the front matter of this handout claims that a database is nothing but a set
of promises about data. This is the chapter where that is made literal.

\begin{definitionbox}[The three model-based constraints]
A \term{domain constraint} says every value of an attribute comes from its
domain. \texttt{semester} is an integer between 1 and 8; \texttt{'spring'} is
not a semester and neither is $-3$.

\smallskip
\term{Entity integrity} says no attribute of a primary key may be null. An
identifier that is missing identifies nothing, so the tuple could not be
addressed.

\smallskip
\term{Referential integrity} says that if a tuple refers to another relation,
the tuple it refers to must exist. A section claiming to belong to course
CS-999 when there is no CS-999 is an \term{orphan}.

\smallskip
These three come from the model itself. A fourth kind, \term{semantic}
constraints, come from the university --- ``a student may not enrol in more
than six sections'' --- and have to be stated case by case.
\end{definitionbox}

\begin{examplebox}[One row, three ways to be wrong]
Into \relation{section}{\relkey{course\_code}, \relkey{section\_no}, room,
emp\_id}:

\medskip
\begin{tabular}{@{}L{4.4cm}L{3.2cm}L{5.8cm}@{}}
\textbf{The row} & \textbf{Violates} & \textbf{Why it matters} \\
\texttt{('CS-103','two','R7')} & domain & \texttt{section\_no} is an
  integer; sorting and counting now fail \\
\texttt{(NULL,1,'R7')} & entity integrity & part of the primary key is
  missing, so the row cannot be addressed \\
\texttt{('CS-999',1,'R7')} & referential integrity & there is no course
  CS-999; the timetable now lists a section of nothing \\
\end{tabular}

\smallskip
All three are refusals a database can make \emph{for you}, on every insert,
for ever --- which is exactly what the application code in
\chref{whynotfiles} could not do, because each program had to remember to
check and one of them always forgot.
\end{examplebox}

\section{Foreign Keys}

\begin{definitionbox}[Foreign key]
A \term{foreign key} is a set of attributes $FK$ in a \term{referencing}
relation $R$ whose values must either
\begin{tight}
  \item appear as the primary key of some tuple in the \term{referenced}
        relation $S$, or
  \item be entirely null.
\end{tight}

\smallskip
The null case is not a loophole --- it is \chref{cardinality}'s partial
participation. A section with no instructor yet has a null \texttt{emp\_id},
and that is a design decision you already took.
\end{definitionbox}

\dsfig{%
\begin{tikzpicture}[
  hd/.style={font=\scriptsize\bfseries, text=white},
  cell/.style={font=\scriptsize\ttfamily, text=primary!65!black},
  bad/.style={font=\scriptsize\ttfamily, text=accent},
  ttl/.style={font=\scriptsize\bfseries, text=primary},
  note/.style={font=\scriptsize\itshape, text=neutral, align=left},
  ok/.style={-{Stealth[length=2mm]}, greenhl, line width=0.9pt},
  no/.style={-{Stealth[length=2mm]}, accent, line width=0.9pt, dashed}
]
  % --- parent: course, top right so the arrows have clear air on the left
  \node[ttl, anchor=west] at (2.2,1.42) {course \ \ (referenced)};
  \fill[primary, rounded corners=2pt] (2.2,0.55) rectangle (6.6,1.1);
  \node[hd, anchor=west] at (2.32,0.83) {course\_code};
  \node[hd, anchor=west] at (4.75,0.83) {title};
  \node[cell, anchor=west] at (2.32,0.3)  {CS-103};
  \node[cell, anchor=west] at (4.75,0.3)  {Databases};
  \node[cell, anchor=west] at (2.32,-0.1) {CS-201};
  \node[cell, anchor=west] at (4.75,-0.1) {Data Struct};
  \draw[primary!35, line width=0.8pt] (2.2,-0.33) rectangle (6.6,1.1);
  \node[note, anchor=north west, text=primary] at (2.2,-0.40)
       {course\_code is the primary key};

  % --- child: section
  \node[ttl, anchor=west] at (0,-1.18) {section \ \ (referencing)};
  \fill[accent!10] (0,-3.3) rectangle (6.6,-2.92);
  \fill[primary, rounded corners=2pt] (0,-2.05) rectangle (6.6,-1.5);
  \node[hd, anchor=west] at (0.12,-1.78) {course\_code};
  \node[hd, anchor=west] at (2.6,-1.78)  {section\_no};
  \node[hd, anchor=west] at (4.8,-1.78)  {emp\_id};
  \node[cell, anchor=west] at (0.12,-2.3) {CS-103};
  \node[cell, anchor=west] at (2.6,-2.3)  {1};
  \node[cell, anchor=west] at (4.8,-2.3)  {E-114};
  \node[cell, anchor=west] at (0.12,-2.7) {CS-103};
  \node[cell, anchor=west] at (2.6,-2.7)  {2};
  \node[cell, anchor=west, text=neutral] at (4.8,-2.7) {NULL};
  \node[bad, anchor=west] at (0.12,-3.11) {CS-999};
  \node[bad, anchor=west] at (2.6,-3.11)  {1};
  \node[bad, anchor=west] at (4.8,-3.11)  {E-114};
  \draw[primary!35, line width=0.8pt] (0,-3.33) rectangle (6.6,-1.5);

  % the two valid references, up the left-hand side
  \draw[ok] (0.08,-2.3) .. controls (-1.5,-1.9) and (-1.5,0.5) .. (2.12,0.3);
  \draw[ok] (0.08,-2.7) .. controls (-1.9,-2.1) and (-1.9,0.7) .. (2.12,0.24);
  % The orphan gets a mark beside its row, not an arrow. A dashed arrow
  % "pointing at nothing" had to curve round the outside of the figure and
  % crossed both annotations on the way.
  \node[bad, anchor=west] at (6.75,-3.11) {$\times$ no such course};

  % Both notes below the tables, side by side: nothing can overlap a row.
  \node[note, anchor=north west, text width=5.0cm] at (0,-3.62)
       {\texttt{NULL} is allowed by the definition --- it is Chapter~6's
        partial participation, not a hole in the rule};
  \node[note, anchor=north west, text=accent, text width=5.4cm] at (5.4,-3.62)
       {\textbf{the orphan.} CS-999 does not exist.
        \textbf{Measured: SQLite accepts this row by default.}};
  \draw[-{Stealth[length=1.8mm]}, neutral!60] (1.9,-3.58) -- (4.75,-2.7);
\end{tikzpicture}}%
{Referential integrity. Two of the three child rows point at a course that
exists; the null is permitted by the definition; the third is an orphan. The
whole of this chapter's measurement is that the third row goes in.}{fkorphan}

\begin{conceptbox}[Where foreign keys come from]
You do not invent them. Each one is a \chref{cardinality} constraint written
down in a different notation, and \chref{mapping} turns one into the other
mechanically.

\begin{tight}
  \item \textbf{A 1:N relationship} puts a foreign key on the \emph{many} side,
        pointing at the one. \ent{course} $\rightarrow$
        \ent{department}.
  \item \textbf{Total participation} makes that foreign key \sqlkw{NOT NULL}.
        Partial participation leaves it nullable.
  \item \textbf{A weak entity} gets a foreign key to its owner, and that
        foreign key is part of its own primary key.
  \item \textbf{An M:N relationship} has no foreign key, because it has no
        table yet. It becomes one in \chref{mapping}, with two.
\end{tight}

\smallskip
So the four answers you argued about in \chref{cardinality} are not an
academic exercise: each one lands as a specific line of \sqlkw{CREATE TABLE},
and you can check your diagram against your schema by reading them back.
\end{conceptbox}

\begin{definitionbox}[Referential actions]
Referential integrity can be broken from the \emph{other} end: delete the
course and its sections become orphans. A \term{referential action} says what
the database should do about it.

\medskip
\begin{tabular}{@{}L{3.1cm}L{5.0cm}L{5.9cm}@{}}
\textbf{Action} & \textbf{On deleting the parent} & \textbf{Use when} \\
\sqlkw{RESTRICT} / \sqlkw{NO ACTION} & refuse the delete & the default, and
  the right answer when you are unsure \\
\sqlkw{CASCADE} & delete the children too & the child is meaningless alone ---
  a weak entity \\
\sqlkw{SET NULL} & null the children's foreign key & the child survives
  without the parent, and the column is nullable \\
\end{tabular}

\smallskip
The same three apply \sqlkw{ON UPDATE}, for when a primary key value changes
--- which is one more reason \chref{ermodel} said to choose a key that never
changes.
\end{definitionbox}

\begin{alertbox}[\sqlkw{CASCADE} is a loaded weapon]
\sqlkw{ON DELETE CASCADE} is attractive because it makes a tidy-up problem
disappear. It also means that one \sqlkw{DELETE} statement can remove rows
from tables you did not name, in tables you did not think about, through
chains you did not draw.

\smallskip
Delete a department with cascade all the way down and you remove its courses,
their sections, every enrolment in them, and therefore the grade history of
students who graduated years ago. One statement, no warning, and the
transaction commits.

\smallskip
\textbf{The rule used in this course:} cascade only to a \emph{weak} entity ---
one that genuinely cannot exist without its owner --- and use \sqlkw{RESTRICT}
everywhere else. A refused delete is an inconvenience; a successful one is
sometimes a disaster, and only one of the two is reversible.
\end{alertbox}

\section{Declared Is Not Enforced}

Everything above is standard and true of every relational system. What follows
is specific to the tool on your machine, and it is the most important single
fact in this chapter.

\begin{measured}[A correct \sqlkw{REFERENCES} clause that does nothing]
\textsf{\footnotesize \texttt{code/m08\_constraints.py}, SQLite 3.51.2,
Intel i5-6600, 16\,GB, Windows 11.}

\medskip
\textbf{1. The same schema, the same \sqlkw{INSERT}, two connections}

\begin{tabular}{@{}L{5.6cm}R{2.4cm}L{5.6cm}@{}}
\textbf{\sqlkw{PRAGMA} foreign\_keys} & \textbf{orphans in} & \textbf{outcome} \\
default (untouched) & 1 & \textbf{accepted} \\
\sqlkw{ON}          & 0 & refused: \texttt{FOREIGN KEY constraint failed} \\
\end{tabular}

\smallskip
The \sqlkw{REFERENCES} clause is identical in both. Only the connection
differs.

\medskip
\textbf{2. The setting is per connection, not per database}

\begin{tabular}{@{}L{8.6cm}L{5.0cm}@{}}
connection A, after \sqlkw{PRAGMA} foreign\_keys = \sqlkw{ON} & reports 1 \\
connection A, orphan insert & refused \\
connection B, \emph{same file}, nothing set & reports 0 \\
connection B, orphan insert & \textbf{accepted} \\
\end{tabular}

\smallskip
Nothing in the file remembers that enforcement was ever wanted.

\medskip
\textbf{3. A \sqlkw{PRIMARY KEY} column that stores a null}

\begin{tabular}{@{}L{5.6cm}R{2.6cm}R{2.4cm}R{2.0cm}@{}}
\textbf{declared as} & \textbf{\sqlkw{INSERT} NULL} & \textbf{stored} & \textbf{rows} \\
\texttt{TEXT PRIMARY KEY}            & accepted & \textbf{null} & 1 \\
\texttt{INTEGER PRIMARY KEY}         & accepted & integer & 1 \\
\texttt{TEXT PRIMARY KEY}, \sqlkw{STRICT} & refused & --- & 0 \\
\texttt{TEXT PRIMARY KEY NOT NULL}   & refused & --- & 0 \\
\end{tabular}

\smallskip
Read the \emph{stored} column, not the first one. \texttt{INTEGER PRIMARY KEY}
accepted the null and stored an integer --- it is an alias for the rowid, so
the null was replaced, and no constraint was broken. \texttt{TEXT PRIMARY KEY}
accepted it and \textbf{stored null}, which breaks entity integrity --- and
that is the declaration every table in this course uses.

\medskip
\textbf{4. What enforcement costs} (20{,}000 valid child rows, 5 runs)

\begin{tabular}{@{}L{6.6cm}R{2.4cm}R{2.2cm}R{2.2cm}@{}}
& \textbf{median} & \textbf{min} & \textbf{max} \\
foreign\_keys \sqlkw{OFF} & 58.7\,ms & 55.1 & 63.8 \\
foreign\_keys \sqlkw{ON}  & 73.0\,ms & 67.3 & 79.6 \\
\textbf{cost of enforcement} & \textbf{14.3\,ms} & \multicolumn{2}{l}{\textbf{(+24\%)}} \\
\end{tabular}

\smallskip
\textbf{Read the four results together.} The declaration is accepted, stored,
and printed back by \texttt{.schema} --- and by default enforces nothing. It
costs 24\% on bulk inserts to turn on, which is a real cost and a small one:
note that the slowest run with enforcement on (79.6\,ms) is only 25\% worse
than the slowest without it, and that the two ranges do not overlap, so the
24\% is a result rather than noise. \textbf{A single run of this measurement
put the two settings 0.6\,ms apart}, which would have been reported as 3\%.
Repeating it is not a nicety.
\end{measured}

\begin{alertbox}[What this means for everything you write from here on]
\sqlkw{PRAGMA} \texttt{foreign\_keys} has been off by default since foreign
keys were added to SQLite in 2009, and it will stay off: turning it on would
break existing databases that have been quietly accumulating orphans.

\smallskip
The consequences you have to live with:

\begin{tight}
  \item \textbf{Set it on every connection}, as the first statement, before
        any transaction. It is a no-op inside one.
  \item \textbf{Your lab partner's connection is not yours.} A database can be
        enforced in your program and unenforced in theirs, at the same moment,
        on the same file.
  \item \textbf{Orphans already in the file are not detected} when you turn it
        on. Enforcement applies to \emph{new} statements only.
        \texttt{PRAGMA foreign\_key\_check} is how you find the old ones, and
        step 7 of the laboratory runs it.
  \item \textbf{A schema that looks correct may never have been enforced.}
        When you are handed an unfamiliar SQLite database, this is the first
        thing to check and the last thing anybody remembers.
\end{tight}

\smallskip
This is not a criticism of SQLite --- the behaviour is documented and the
reason is sound. It is a demonstration of something general: \textbf{a
constraint has two separate lives, one in the schema and one at run time, and
only the second one protects anything.}
\end{alertbox}

\begin{examplebox}[The university schema, with its promises declared]
\begin{lstlisting}[language=sql]
PRAGMA foreign_keys = ON;        -- first. every time.

CREATE TABLE department (
  dept_code TEXT PRIMARY KEY NOT NULL,
  name      TEXT NOT NULL
) STRICT;
\end{lstlisting}

A course belongs to exactly one department, and must:

\begin{lstlisting}[language=sql]
CREATE TABLE course (
  course_code  TEXT PRIMARY KEY NOT NULL,
  title        TEXT NOT NULL,
  credit_hours INTEGER NOT NULL CHECK (credit_hours BETWEEN 1 AND 4),
  dept_code    TEXT NOT NULL              -- total participation
               REFERENCES department(dept_code)
               ON DELETE RESTRICT         -- a department with courses stays
) STRICT;
\end{lstlisting}

A section is the weak entity of \chref{cardinality}, so its owner's key is part
of its own and the cascade is correct here and nowhere else:

\begin{lstlisting}[language=sql]
CREATE TABLE section (
  course_code TEXT    NOT NULL
              REFERENCES course(course_code)
              ON DELETE CASCADE,          -- weak entity: cascade is correct
  semester    TEXT    NOT NULL,
  section_no  INTEGER NOT NULL CHECK (section_no > 0),
  room        TEXT,
  emp_id      TEXT                        -- nullable: partial participation
              REFERENCES instructor(emp_id)
              ON DELETE SET NULL,         -- the section survives; the
  PRIMARY KEY (course_code, semester, section_no)   -- teacher changed
) STRICT;
\end{lstlisting}

\smallskip
\textbf{Every line is a decision you have already made.} \texttt{NOT NULL} on
\texttt{dept\_code} is the total participation from \chref{cardinality}'s
worked example; the nullable \texttt{emp\_id} is the partial participation
argued for in the same place; the three-part primary key is the weak entity's
owner key plus discriminator; and each of the three referential actions
follows the rule in this chapter. Nothing here is a matter of taste.
\end{examplebox}

\begin{pitfall}
\begin{tight}
  \item \textbf{Assuming \sqlkw{REFERENCES} does something.} By default, in
        SQLite, it does not. Check, every time.
  \item \textbf{Setting the pragma once and assuming it stuck.} It is per
        connection and lasts until that connection closes.
  \item \textbf{Setting it inside a transaction.} It is silently ignored there.
  \item \textbf{Declaring \texttt{TEXT PRIMARY KEY} without \sqlkw{NOT NULL}.}
        It will store a null, as measured. Add \sqlkw{NOT NULL} or
        \sqlkw{STRICT}, and this course does both.
  \item \textbf{Cascading beyond a weak entity.} One delete, a chain you did
        not draw, and a committed transaction.
  \item \textbf{Believing the application will check it.} Three programs, one
        of them written next year by somebody else. This is
        \chref{whynotfiles}'s whole argument.
  \item \textbf{Turning enforcement on and assuming the old data is clean.} It
        applies to new statements only; run
        \texttt{PRAGMA foreign\_key\_check}.
\end{tight}
\end{pitfall}

\begin{lab}[Insert an orphan into a correct schema, then make it impossible]
\begin{lstlisting}[language=bash]
# --- 1. A scratch database ------------------------------------------
sqlite3 fk.db
\end{lstlisting}
\begin{lstlisting}[language=sql]
-- --- 2. A schema that declares referential integrity ------------
.headers on
.mode box
CREATE TABLE course (
  course_code TEXT PRIMARY KEY NOT NULL,
  title       TEXT NOT NULL
) STRICT;
CREATE TABLE section (
  course_code TEXT    NOT NULL REFERENCES course(course_code),
  section_no  INTEGER NOT NULL,
  room        TEXT,
  PRIMARY KEY (course_code, section_no)
) STRICT;
INSERT INTO course VALUES ('CS-103','Database Systems');

-- --- 3. Confirm the declaration is really there -----------------
.schema section
PRAGMA foreign_key_list(section);
-- It is declared. Written down. Reported back to you.

-- --- 4. Now break it --------------------------------------------
PRAGMA foreign_keys;              -- what does it say?
INSERT INTO section VALUES ('CS-999', 1, 'Room 7');
SELECT * FROM section;
-- A section of a course that does not exist. No error.

-- --- 5. Turn enforcement on and try again -----------------------
PRAGMA foreign_keys = ON;
PRAGMA foreign_keys;
INSERT INTO section VALUES ('CS-888', 1, 'Room 8');
-- Record the exact error message.

-- --- 6. The orphan from step 4 is still there -------------------
SELECT * FROM section;
-- Enforcement did not clean up. It never looks backwards.

-- --- 7. Find every orphan already in the file -------------------
PRAGMA foreign_key_check;
-- Learn this one. It is how you audit a database you were handed.

-- --- 8. Prove it is per connection ------------------------------
.quit
\end{lstlisting}
\begin{lstlisting}[language=bash]
# --- 9. A brand new connection to the SAME file ---------------------
sqlite3 fk.db "PRAGMA foreign_keys;"
sqlite3 fk.db "INSERT INTO section VALUES ('CS-777',1,'R9');"
sqlite3 fk.db "SELECT * FROM section;"
# It went in. Nothing in the file remembered.

# --- 10. The measured version ---------------------------------------
python code/m08_constraints.py
\end{lstlisting}

\smallskip
\textbf{What to notice.} Step 3 is the one that makes the point. The constraint
is \emph{there} --- \texttt{PRAGMA foreign\_key\_list} prints it, the schema
stores it, any tool reading the database would report it --- and step 4 walks
straight past it. From now on, every script you write in this course opens
with \texttt{PRAGMA foreign\_keys = ON;}, and \texttt{code/run\_all.py}
checks that they all do.
\end{lab}

\begin{checkpoint}
\begin{tightnum}
  \item You are handed a SQLite database and told its foreign keys are
        enforced. List the commands you would run to find out whether that is
        true \emph{now}, and whether it has ever been true, and say what each
        result would and would not establish.
  \item For each of \ent{section}~$\rightarrow$~\ent{course},
        \ent{section}~$\rightarrow$~\ent{instructor} and
        \ent{enrolment}~$\rightarrow$~\ent{student}, choose an
        \sqlkw{ON DELETE} action and justify it in one sentence. At least one
        of your three must not be \sqlkw{CASCADE}.
  \item A table is declared \texttt{roll\_no TEXT PRIMARY KEY}. Explain what
        value it will accept that it should not, which integrity constraint
        that violates, and give the two different one-word changes that each
        fix it.
\end{tightnum}
\end{checkpoint}

\begin{chaptersummary}
\begin{tight}
  \item The model gives three constraints: \textbf{domain} (values come from
        the right set), \textbf{entity integrity} (no part of a primary key is
        null), \textbf{referential integrity} (a reference points at something
        that exists). \textbf{Semantic} constraints come from the university.
  \item A \textbf{foreign key} must match a primary key in the referenced
        relation, \emph{or} be wholly null --- and the null case is
        \chref{cardinality}'s partial participation.
  \item Foreign keys are not invented: a 1:N relationship puts one on the many
        side, total participation makes it \sqlkw{NOT NULL}, and a weak entity
        makes it part of its own key.
  \item \textbf{Referential actions}: \sqlkw{RESTRICT} by default,
        \sqlkw{CASCADE} only to a weak entity, \sqlkw{SET NULL} when the child
        survives the parent.
  \item \textbf{Measured: \sqlkw{PRAGMA} foreign\_keys is 0 by default.} A
        correct \sqlkw{REFERENCES} clause is parsed, stored, reported --- and
        enforces nothing. The orphan went in without an error.
  \item The setting is \textbf{per connection}. Connection B, on the same
        file, accepted what connection A refused.
  \item Enforcement \textbf{never looks backwards}: orphans already present
        stay. \texttt{PRAGMA foreign\_key\_check} finds them.
  \item \textbf{Measured: \texttt{TEXT PRIMARY KEY} stores a null.}
        \sqlkw{STRICT} or \sqlkw{NOT NULL} fixes it; \texttt{INTEGER PRIMARY
        KEY} only appears to, because it substitutes a rowid.
  \item \textbf{Measured: enforcement costs +24\%} on bulk inserts --- outside
        the run-to-run spread, and cheap for what it buys. A single run
        reported 3\%, which is why the harness repeats.
  \item The general lesson: \textbf{a constraint has two lives, one in the
        schema and one at run time, and only the second protects anything.}
\end{tight}
\end{chaptersummary}

\begin{reviewq}
  \item \textbf{(MCQ)} Entity integrity is violated when: (a)~a foreign key
        points nowhere; (b)~part of a primary key is null; (c)~a value is
        outside its domain; (d)~two tables have the same name.
  \item \textbf{(MCQ)} In SQLite, \sqlkw{PRAGMA} \texttt{foreign\_keys}
        defaults to: (a)~\sqlkw{ON} for \sqlkw{STRICT} tables;
        (b)~\sqlkw{ON}; (c)~\sqlkw{OFF}; (d)~whatever it was last set to in
        that file.
  \item \textbf{(Short)} Give a situation in which \sqlkw{ON DELETE SET NULL}
        is the right action and \sqlkw{CASCADE} would lose data you need.
  \item \textbf{(Short)} Why is a foreign key allowed to be null at all? Answer
        in terms of \chref{cardinality}.
  \item \textbf{(Short)} You turn on foreign key enforcement for a database
        that has been running for a year. What does it protect you from, and
        what does it not?
  \item \textbf{(Applied)} Write the \sqlkw{CREATE TABLE} statements for
        \ent{instructor} and \ent{enrolment}, with every domain, entity
        and referential constraint declared and a referential action chosen
        for each foreign key. Beside each constraint, name the decision from
        \chref{ermodel} or \chref{cardinality} that it implements.
  \item \textbf{(Applied)} A colleague argues that foreign keys should be left
        off because they cost 24\% on bulk inserts and ``the application
        checks anyway''. Write the reply. Use the measured figures, the
        per-connection result, and \chref{whynotfiles}'s argument, and say
        under what single condition their position would actually be
        defensible.
\end{reviewq}


\part{From Diagram to Schema}
%=====================================================================
% PART III OPENER
%=====================================================================
\thispagestyle{plain}
\structuralfigures{P3.}

\begin{center}
\begin{tcolorbox}[enhanced, width=0.92\textwidth, colback=greenhl!5,
  colframe=greenhl, boxrule=0pt, leftrule=5pt, arc=3pt, top=8pt, bottom=8pt]
\large\itshape\color{greenhl!60!black}
A diagram is a drawing until a database agrees to it. Part~III is the shortest in the book and the one with the most permanent consequences: the rules for turning a design into tables, the SQL that creates them, and what to do a fortnight later when the design turns out to have been slightly wrong and the tables are already full.
\end{tcolorbox}
\end{center}

\vspace{6pt}

\dsfigh{%
\begin{tikzpicture}[
  cbox/.style={rounded corners=3pt, draw=greenhl, line width=1pt, fill=greenhl!7,
               text=greenhl!55!black, font=\small\bfseries, align=center,
               minimum width=4.0cm, minimum height=1.0cm},
  hbox/.style={rounded corners=3pt, draw=greenhl, line width=1.8pt, fill=greenhl!16,
               text=greenhl!55!black, font=\small\bfseries, align=center,
               minimum width=4.0cm, minimum height=1.0cm},
  arr/.style={-{Stealth[length=2.4mm]}, greenhl!60, line width=1pt},
  note/.style={font=\scriptsize\itshape, text=greenhl!55!black, align=center,
               fill=white, inner sep=2pt}
]
  \node[cbox] (c0) at (0.00,0) {Ch.\ 9\\Mapping};
  \node[hbox] (c1) at (5.20,0) {Ch.\ 10\\Creating the Schema};
  \node[cbox] (c2) at (10.40,0) {Ch.\ 11\\Changing It};
  \draw[arr] (c0) -- (c1);
  \draw[arr] (c1) -- (c2);
  \node[note] at (2.60,-1.0) {diagram to tables, by rule};
  \node[note] at (7.80,-1.0) {tables that exist};
  \node[font=\scriptsize\itshape, text=accent, align=center] at (5.20,1.15)
       {the chapter in which a column's declared type turns out to be advice};
\end{tikzpicture}}%
{Reading path through Part~III. Chapter 9 is mechanical and should be:
mapping is a set of rules you apply rather than judgements you make. The
judgement was Part~II and the consequences are Part~V.}{part3map}

\newpage

%=====================================================================
\chapter{Mapping an ER Diagram to Tables}\label{ch:mapping}
%=====================================================================
\locator{3}{Part III --- From Diagram to Schema}

\begin{outcomes}
By the end of this chapter you will be able to:
\begin{tight}
  \item \textbf{Apply} the seven mapping rules, in order, to any ER diagram.
  \item \textbf{Map} a 1:N relationship by placing a foreign key, and
        \textbf{say} which side it goes on and why.
  \item \textbf{Resolve} an M:N relationship into a table, and
        \textbf{determine} its primary key.
  \item \textbf{Map} a weak entity, a multi-valued attribute and a recursive
        relationship, each of which has its own rule.
  \item \textbf{Check} a finished schema back against the diagram it came
        from, constraint by constraint.
\end{tight}
\end{outcomes}

\begin{prereq}
\chref{cardinality}: the finished, constrained diagram --- you will map the one
you drew. \chref{constraints} for foreign keys and referential actions, which
are what the rules produce.
\end{prereq}

\begin{keyterms}
mapping $\bullet$ logical design $\bullet$ strong entity $\bullet$ weak entity
$\bullet$ relationship table $\bullet$ junction table $\bullet$ composite
primary key $\bullet$ foreign key placement $\bullet$ resolving a
many-to-many $\bullet$ surrogate key $\bullet$ recursive relationship
$\bullet$ ternary relationship $\bullet$ flattening a composite attribute
\end{keyterms}

\section{The Part of Design That Is Mechanical}

Everything in Part~II was judgement: whether a noun is an entity, whether
participation is really total, which candidate key to choose. Those decisions
were expensive and nobody could make them for you.

This chapter is the opposite. Given a finished diagram, turning it into tables
is a \emph{procedure}. Seven rules, applied in order, and two people who
follow them get the same schema. That is the point of having spent four
chapters on the diagram: the hard thinking is over, and what remains can be
checked.

\begin{conceptbox}[The seven rules, in the order they must be applied]
\begin{tabular}{@{}R{0.7cm}L{5.3cm}L{8.0cm}@{}}
\textbf{1} & Strong entity type & one table; its key becomes the primary key \\
\textbf{2} & Composite attribute & store the \emph{parts}, not the whole \\
\textbf{3} & Multi-valued attribute & \textbf{a new table}: owner's key +
  the value, both in the primary key \\
\textbf{4} & 1:1 relationship & a foreign key on the side with \emph{total}
  participation \\
\textbf{5} & 1:N relationship & a foreign key on the \textbf{many} side \\
\textbf{6} & M:N relationship & \textbf{a new table}: both keys, together the
  primary key, plus the relationship's own attributes \\
\textbf{7} & Weak entity type & a table whose primary key is the owner's key
  \emph{plus} the partial key \\
\end{tabular}

\smallskip
\textbf{Order matters.} Rules 4--7 place foreign keys, and a foreign key can
only point at a primary key that already exists. Do all the entities first.
\end{conceptbox}

\section{Rules 1 to 3: the Entities}

\begin{examplebox}[Rule 1 --- a strong entity becomes a table]
\ent{student}, with key \texttt{roll\_no}, becomes

\medskip
\relation{student}{\relkey{roll\_no}, name, programme, semester}

\medskip
\begin{lstlisting}[language=sql]
CREATE TABLE student (
  roll_no   TEXT    PRIMARY KEY NOT NULL,
  name      TEXT    NOT NULL,
  programme TEXT    NOT NULL,
  semester  INTEGER NOT NULL CHECK (semester BETWEEN 1 AND 8)
) STRICT;
\end{lstlisting}

The \sqlkw{CHECK} is the domain constraint from \chref{constraints}, and it
came from the ER diagram's domain, not from inspiration.
\end{examplebox}

\begin{examplebox}[Rules 2 and 3 --- the two attribute kinds that will not fit]
\textbf{Rule 2, composite.} \texttt{address} = house, street, city. Store the
parts:

\begin{lstlisting}[language=sql]
  addr_house  TEXT,
  addr_street TEXT,
  addr_city   TEXT,
\end{lstlisting}

Three columns, one for each leaf of the ER tree. Never one \texttt{address}
column --- \chref{ermodel} explained what that costs, and the cost is paid at
the moment somebody asks ``how many students live in Bahawalpur''.

\medskip
\textbf{Rule 3, multi-valued.} \texttt{phone} cannot be a column at all. It
becomes its own table:

\medskip
\relation{student\_phone}{\relkey{roll\_no}, \relkey{phone}}

\begin{lstlisting}[language=sql]
CREATE TABLE student_phone (
  roll_no TEXT NOT NULL REFERENCES student(roll_no) ON DELETE CASCADE,
  phone   TEXT NOT NULL,
  PRIMARY KEY (roll_no, phone)
) STRICT;
\end{lstlisting}

\smallskip
\textbf{Both attributes are in the primary key, and that is the rule, not a
convenience.} It says a student may have many numbers and may not have the
same number recorded twice. Make \texttt{roll\_no} alone the key and each
student gets one phone --- which is where you started. Leave the key off
altogether and \chref{relational}'s step 2 happens: the same number goes in
twice and no \sqlkw{DELETE} can remove one of them.

\smallskip
The cascade is right here because a phone number belongs to a student in
exactly the way \chref{constraints} requires: it cannot exist without one.
\end{examplebox}

\section{Rules 4 to 6: the Relationships}

\begin{definitionbox}[Where a foreign key goes, and why there is no choice]
In a \textbf{1:N} relationship the foreign key goes on the \textbf{many} side,
always. The reason is the atomicity of \chref{relational}: a course belongs to
one department, so one value fits in one column. A department has many courses,
and many values do not.

\smallskip
In a \textbf{1:1} relationship either side would work, so there is a choice and
a rule for making it: \textbf{put it on the side with total participation}, so
that it can be declared \sqlkw{NOT NULL}. Put it on the partial side and most
of the column is null.

\smallskip
In an \textbf{M:N} relationship \emph{neither} side works, and that is why the
relationship needs a table of its own.
\end{definitionbox}

\begin{worked}[Rule 6 --- resolving \ent{enrols}, and finding its key]
\ent{student} (0,N) --- \ent{enrols} --- (0,N) \ent{section}, carrying
\texttt{grade}.

\medskip
\textbf{Step 1 --- why no foreign key will do.} Put \texttt{section\_id} on
\ent{student} and a student takes one section. Put \texttt{roll\_no} on
\ent{section} and a section holds one student. Both are wrong, and neither is
fixable by choosing the other side.

\medskip
\textbf{Step 2 --- make a table.} It gets the primary key of each participant:

\medskip
\relation{enrolment}{\relkey{roll\_no}, \relkey{course\_code},
\relkey{semester}, \relkey{section\_no}, grade}

\medskip
Four key attributes, not two --- because \ent{section} is weak
(\chref{cardinality}) and its own key is already three attributes wide. A weak
entity on either end of an M:N relationship makes the junction table wide, and
that is the honest consequence of the design rather than a mistake.

\medskip
\textbf{Step 3 --- the primary key is the combination.} All four together.
This states exactly one thing, and it is the thing the university means:
\emph{a student may appear in a section only once}.

\medskip
\textbf{Step 4 --- the relationship's attribute comes too.} \texttt{grade} is
not in the key. It is a fact \emph{about} the enrolment, nullable because the
semester is not over.

\medskip
\begin{lstlisting}[language=sql]
CREATE TABLE enrolment (
  roll_no     TEXT    NOT NULL REFERENCES student(roll_no),
  course_code TEXT    NOT NULL,
  semester    TEXT    NOT NULL,
  section_no  INTEGER NOT NULL,
  grade       TEXT    CHECK (grade IN ('A','B','C','D','F')),
  PRIMARY KEY (roll_no, course_code, semester, section_no),
  FOREIGN KEY (course_code, semester, section_no)
    REFERENCES section(course_code, semester, section_no)
) STRICT;
\end{lstlisting}

\smallskip
\textbf{Notice the second foreign key is written separately.} A foreign key of
three columns cannot be declared inline against one column; it needs its own
\sqlkw{FOREIGN KEY} clause naming all three, in the same order as the parent's
primary key. This is where a wide weak-entity key is actually felt, and it is
the strongest practical argument for giving \ent{section} a surrogate
\texttt{section\_id} instead --- an argument \chref{evolution} makes with a
measurement rather than with taste.
\end{worked}

\begin{alertbox}[The M:N table is not ``just a link table''; it is where the facts live]
Students meet this table and call it plumbing. It is the most informative
table in the database:

\begin{tight}
  \item \texttt{grade} is there, and nowhere else. Every transcript, every GPA,
        every pass rate is a query over this table.
  \item It is the only place that records \emph{when} a student took something.
  \item It is almost always the largest table --- 4{,}000 students $\times$ 5
        sections $\times$ 8 semesters is 160{,}000 rows against a few thousand
        everywhere else, which is why \chref{indexes} measures this table and
        not the others.
\end{tight}

\smallskip
If an M:N table has no attributes of its own today, it still will: ``when did
they enrol'', ``who approved it'', ``was it a repeat''. Design it as a first
class table, because it is one.
\end{alertbox}

\dsfig{%
\begin{tikzpicture}[
  ent/.style={draw=primary, line width=1pt, fill=primary!8,
              text=primary!60!black, font=\scriptsize\bfseries,
              minimum width=1.8cm, minimum height=0.62cm},
  weak/.style={ent, draw=purplehl, fill=purplehl!8, text=purplehl!65!black,
               double, double distance=1.1pt},
  rel/.style={draw=greenhl, line width=0.9pt, fill=greenhl!8,
              text=greenhl!55!black, font=\scriptsize\bfseries,
              diamond, aspect=2.2, inner sep=0.5pt},
  tb/.style={rounded corners=2pt, draw=#1, line width=0.9pt, fill=#1!6,
             text=#1!60!black, font=\scriptsize, align=left,
             inner sep=3pt, text width=5.2cm},
  ln/.style={draw=neutral!55, line width=0.7pt},
  rr/.style={-{Stealth[length=2.2mm]}, accent!75, line width=1pt},
  rl/.style={font=\scriptsize\bfseries, text=accent},
  hd/.style={font=\scriptsize\bfseries, text=primary}
]
  \node[hd] at (-3.4,2.5) {the diagram};
  \node[hd] at ( 4.4,2.5) {the tables};

  \node[ent]  (stu) at (-5.3, 1.3) {STUDENT};
  \node[rel]  (enr) at (-3.4, 1.3) {ENROLS};
  \node[weak] (sec) at (-1.5, 1.3) {SECTION};
  % Double diamond: HAS is the identifying relationship, and drawing it as an
  % ordinary one here contradicted Figure 6.1 two chapters earlier.
  \node[rel, draw=purplehl, fill=purplehl!8, text=purplehl!65!black,
        double, double distance=1.1pt] (has) at (-1.5,-0.3) {HAS};
  \node[ent]  (crs) at (-1.5,-1.6) {COURSE};
  \draw[ln] (stu) -- (enr);  \draw[ln] (enr) -- (sec);
  \draw[ln, double, double distance=1pt] (sec) -- (has);
  \draw[ln] (has) -- (crs);

  \node[tb=primary]  (t1) at (4.4, 1.75)
       {\textbf{student}(\relkey{roll\_no}, name, \dots)};
  \node[tb=accent]   (t2) at (4.4, 0.75)
       {\textbf{enrolment}(\relkey{roll\_no}, \relkey{course\_code},
        \relkey{semester}, \relkey{section\_no}, grade)};
  \node[tb=purplehl] (t3) at (4.4,-0.55)
       {\textbf{section}(\relkey{course\_code}, \relkey{semester},
        \relkey{section\_no}, room)};
  \node[tb=primary]  (t4) at (4.4,-1.6)
       {\textbf{course}(\relkey{course\_code}, title, \dots)};

  % Each label rides its OWN curve with pos=, rather than floating between
  % two of them: the first draft put "rule 6" in the gap where the enrolment
  % and section arrows cross, and it was impossible to tell which it named.
  \draw[rr] (stu.east) ..controls (-4.2,2.6) and (0.6,2.9)..
       node[rl, above, pos=0.55] {rule 1} (t1.west);
  % These two curves run close together in the middle of the figure, so each
  % label goes where its OWN curve is alone: rule 6 late, on the flat approach
  % to the enrolment table; rule 7 early, just as it leaves SECTION. Placed at
  % mid-curve they printed on each other's arrows.
  \draw[rr] (enr.south) ..controls (-3.4,0.2) and (0.6,0.75)..
       node[rl, above, pos=0.80] {rule 6} (t2.west);
  \draw[rr] (sec.east) ..controls (0.0,1.3) and (0.6,-0.2)..
       node[rl, right, pos=0.18] {rule 7} (t3.west);
  \draw[rr] (crs.east) ..controls (0.0,-1.6) and (0.6,-1.6)..
       node[rl, above, pos=0.5] {rule 1} (t4.west);
\end{tikzpicture}}%
{Four of the seven rules on the part of the diagram that uses them. Rule~1
makes a table per box; rule~6 turns the diamond into a table of its own;
rule~7 gives the weak entity its owner's key. The arrows are the whole of
logical design.}{mapping}

\section{Rule 7, and the Two Awkward Cases}

\begin{examplebox}[Rule 7 --- a weak entity]
\ent{section} is identified by its owner plus its discriminator, so its primary
key is the owner's key plus the partial key, and the owner's key is
\emph{simultaneously} a foreign key:

\begin{lstlisting}[language=sql]
CREATE TABLE section (
  course_code TEXT    NOT NULL
              REFERENCES course(course_code) ON DELETE CASCADE,
  semester    TEXT    NOT NULL,
  section_no  INTEGER NOT NULL CHECK (section_no > 0),
  room        TEXT,
  emp_id      TEXT REFERENCES instructor(emp_id) ON DELETE SET NULL,
  PRIMARY KEY (course_code, semester, section_no)
) STRICT;
\end{lstlisting}

\smallskip
Three things are true of \texttt{course\_code} at once: it is part of the
primary key, it is a foreign key, and it is \sqlkw{NOT NULL}. All three come
straight off the diagram --- the double line in Figure~\ref{fig:ercomplete}
said the participation was total, and \sqlkw{CASCADE} is correct here for the
reason \chref{constraints} gave and nowhere else in this schema.
\end{examplebox}

\begin{conceptbox}[The recursive relationship, and the ternary one]
Two cases the seven rules cover, but which look as though they need an eighth.

\medskip
\textbf{Recursive (\ent{course} \ent{requires} \ent{course}, M:N).} Rule 6
applies unchanged. The new table takes the key of each participant --- and
both participants are \ent{course}, so it takes \texttt{course\_code}
\emph{twice}, under the two role names:

\begin{lstlisting}[language=sql]
CREATE TABLE prerequisite (
  course_code TEXT NOT NULL REFERENCES course(course_code),
  requires    TEXT NOT NULL REFERENCES course(course_code),
  PRIMARY KEY (course_code, requires),
  CHECK (course_code <> requires)          -- nothing requires itself
) STRICT;
\end{lstlisting}

The role names from \chref{ermodel} become the column names. Without them you
would have two columns called \texttt{course\_code} and no legal way to name
either.

\medskip
\textbf{Ternary.} A relationship among three entity types becomes one table
with three foreign keys, and its primary key is normally all three. Do
\emph{not} map it as three binary relationships: ``student, section,
instructor'' as three pairs cannot record which instructor taught which
student in which section, only that each pair occurred.
\end{conceptbox}

\begin{pitfall}
\begin{tight}
  \item \textbf{Putting the foreign key on the ``one'' side of a 1:N.} One
        column cannot hold many values. This is the single most common error.
  \item \textbf{Mapping M:N with a foreign key.} It cannot be done; make the
        table.
  \item \textbf{Giving the M:N table a surrogate key and nothing else.} An
        \texttt{id} column makes rows addressable; without the composite
        \sqlkw{UNIQUE} the same student enrols twice.
  \item \textbf{Making the multi-valued table's key just the owner.} Then the
        owner has one value again, which is what you were escaping.
  \item \textbf{Storing a composite attribute whole.} Three columns, one per
        leaf.
  \item \textbf{Declaring a multi-column foreign key inline.} It needs its own
        \sqlkw{FOREIGN KEY} clause, with the columns in the parent's key order.
  \item \textbf{Mapping a ternary relationship as three binary ones.} It loses
        exactly the fact the relationship existed to record.
\end{tight}
\end{pitfall}

\begin{lab}[Turn your diagram into the database you will use all term]
This lab produces \texttt{code/schema/v1\_create.sql}, and every later chapter
runs against what you build here.

\begin{lstlisting}[language=bash]
# --- 1. Start the file ----------------------------------------------
# Open code/schema/v1_create.sql in an editor. First line, always:
#   PRAGMA foreign_keys = ON;
\end{lstlisting}
\begin{lstlisting}[language=sql]
-- --- 2. Rule 1: one table per strong entity ---------------------
--     department, course, instructor, student.
--     Primary key first, NOT NULL, STRICT on every table.

-- --- 3. Rules 2 and 3: the awkward attributes -------------------
--     Flatten the composite one into its parts.
--     Give the multi-valued one its own table, with BOTH
--     attributes in the primary key.

-- --- 4. Rule 5: the 1:N foreign keys ----------------------------
--     course -> department.  Which side, and NOT NULL or not?
--     Your Chapter 6 diagram already decided both.

-- --- 5. Rule 7: the weak entity ---------------------------------
--     section. Owner's key plus discriminator, and the owner's
--     key is a foreign key at the same time.

-- --- 6. Rule 6: the M:N tables ----------------------------------
--     enrolment, and prerequisite (the recursive one).
--     Composite primary keys. Relationship attributes included.

-- --- 7. A referential action for every foreign key --------------
--     There are at least six. CASCADE only to a weak entity.
--     Write a one-line comment on each saying why.
\end{lstlisting}
\begin{lstlisting}[language=bash]
# --- 8. Create it ---------------------------------------------------
rm -f university.db
sqlite3 university.db ".read code/schema/v1_create.sql"
sqlite3 university.db ".tables"
sqlite3 university.db ".schema section"

# --- 9. Check it against the diagram, constraint by constraint ------
sqlite3 university.db "PRAGMA foreign_key_list(section);"
sqlite3 university.db "PRAGMA foreign_key_list(enrolment);"
sqlite3 university.db "PRAGMA table_info(course);"
# For EVERY line of Figure 6.1, find the line of SQL that implements
# it. Any line of the diagram with no SQL is a mapping you skipped.

# --- 10. Prove the M:N key does its job ------------------------------
sqlite3 university.db "PRAGMA foreign_keys=ON;
  INSERT INTO student (roll_no, name, programme, semester)
    VALUES ('2021-IT-001','Ayesha Khan','BSIT',2);
  INSERT INTO department VALUES ('CS','Computer Science');
  INSERT INTO course VALUES ('CS-103','Database Systems',4,'CS');
  INSERT INTO section VALUES ('CS-103','FALL2026',1,'R7',NULL);
  INSERT INTO enrolment VALUES
    ('2021-IT-001','CS-103','FALL2026',1,NULL);
  INSERT INTO enrolment VALUES
    ('2021-IT-001','CS-103','FALL2026',1,NULL);"
# The second one must fail. Record the message.
\end{lstlisting}

\smallskip
\textbf{What to notice.} Step 9 is the check that makes mapping a procedure
rather than a craft: every line of the diagram should have a line of SQL, and
every line of SQL should have a line of the diagram. Lines of SQL with no
diagram behind them are decisions you made without noticing, and they are
exactly the ones that turn out to be wrong in week twelve.
\end{lab}

\begin{checkpoint}
\begin{tightnum}
  \item A 1:1 relationship joins \ent{department} and \ent{instructor} (the
        head of department). Say which table gets the foreign key, justify the
        choice from the participation constraints, and say what the column
        would look like if you chose the other side.
  \item \ent{enrolment}'s primary key has four attributes. Explain what fact
        about the university each one contributes, and say precisely what
        becomes possible if you remove \texttt{semester} from the key.
  \item Map the recursive relationship \ent{course} \ent{requires}
        \ent{course}. Give the table, its primary key, its two foreign keys,
        and one \sqlkw{CHECK} constraint that the ER diagram does not state but
        that the university certainly means.
\end{tightnum}
\end{checkpoint}

\begin{chaptersummary}
\begin{tight}
  \item Mapping is \textbf{mechanical}. Seven rules, applied in order, and two
        people get the same schema --- which is what four chapters of diagram
        bought.
  \item \textbf{Entities first} (rules 1--3), then relationships (4--7): a
        foreign key can only point at a primary key that already exists.
  \item A \textbf{composite} attribute is stored as its parts; a
        \textbf{multi-valued} attribute becomes a table whose primary key is
        \emph{both} columns.
  \item \textbf{1:N} puts the foreign key on the \textbf{many} side, because
        one column holds one value. \textbf{1:1} puts it on the \emph{totally
        participating} side so it can be \sqlkw{NOT NULL}.
  \item \textbf{M:N} becomes a table carrying both keys plus the
        relationship's own attributes; the primary key is the combination, and
        it is what forbids enrolling twice.
  \item The M:N table is not plumbing. It holds the grades, it records the
        time, and it is the biggest table in the database.
  \item A \textbf{weak entity} becomes a table whose key is the owner's key
        plus the partial key, with the owner's key simultaneously a foreign
        key and \sqlkw{NOT NULL}.
  \item A \textbf{recursive} M:N uses rule 6 unchanged, taking the same key
        twice under the ER diagram's \textbf{role names}.
  \item A \textbf{ternary} relationship is one table with three foreign keys,
        never three binary relationships --- that loses the fact it existed to
        record.
  \item A multi-column foreign key needs its own \sqlkw{FOREIGN KEY} clause,
        with columns in the parent key's order.
  \item Check the schema back: every line of the diagram should have a line of
        SQL, and every line of SQL a line of the diagram.
\end{tight}
\end{chaptersummary}

\begin{reviewq}
  \item \textbf{(MCQ)} In a 1:N relationship the foreign key is placed:
        (a)~on the one side; (b)~on the many side; (c)~in a new table;
        (d)~on either side.
  \item \textbf{(MCQ)} A multi-valued attribute maps to: (a)~several columns;
        (b)~a comma-separated column; (c)~a new table keyed on owner and
        value; (d)~a \sqlkw{CHECK} constraint.
  \item \textbf{(Short)} Why can an M:N relationship not be mapped with a
        foreign key? Answer using a property of relations from
        \chref{relational}.
  \item \textbf{(Short)} A 1:1 relationship is partial on both sides. Where do
        you put the foreign key, and what is the consequence either way?
  \item \textbf{(Short)} What does the composite primary key of
        \ent{enrolment} forbid, and what would be permitted without it?
  \item \textbf{(Applied)} Map this diagram completely: a \ent{library} has
        many \ent{book}s; each book has many \ent{copy}s; \ent{member}s
        \ent{borrow} copies and a borrowing records a date out and a date
        back. Give every table, every key, every foreign key and a referential
        action for each, and name the rule that produced each table.
  \item \textbf{(Applied)} Take your own \texttt{v1\_create.sql} and go through
        Figure~\ref{fig:ercomplete} line by line. Produce a two-column table:
        the diagram constraint, and the SQL that implements it. Then list
        anything in either column with no partner, and for each one say
        whether the diagram or the schema is wrong.
\end{reviewq}

%=====================================================================
\chapter{SQL: Creating the Schema}\label{ch:ddl}
%=====================================================================
\locator{3}{Part III --- From Diagram to Schema}

\begin{outcomes}
By the end of this chapter you will be able to:
\begin{tight}
  \item \textbf{Write} a complete \sqlkw{CREATE TABLE} statement with column
        types, key, nullability, defaults and \sqlkw{CHECK} constraints.
  \item \textbf{Explain} what type affinity is, and \textbf{predict} what
        SQLite stores for a given value in a given declared type.
  \item \textbf{Demonstrate} that an \texttt{INTEGER} column will hold text,
        and that \sqlkw{STRICT} stops it.
  \item \textbf{Trace} a wrong query result back to a value of the wrong type.
  \item \textbf{Choose} between a column constraint and a table constraint, and
        \textbf{name} a case where only the second will do.
\end{tight}
\end{outcomes}

\begin{prereq}
\chref{mapping}: the tables and keys you are now declaring.
\chref{constraints} for \sqlkw{PRAGMA} \texttt{foreign\_keys} and \sqlkw{NOT
NULL}, both of which appear again here for a second reason. You need
\texttt{code/m10\_ddl.py}.
\end{prereq}

\begin{keyterms}
DDL $\bullet$ \sqlkw{CREATE TABLE} $\bullet$
storage class $\bullet$ type affinity $\bullet$ declared type $\bullet$
\sqlkw{STRICT} $\bullet$ \texttt{typeof()} $\bullet$ column constraint
$\bullet$ table constraint $\bullet$ \sqlkw{CHECK} $\bullet$ \sqlkw{DEFAULT}
$\bullet$ \sqlkw{UNIQUE} $\bullet$ \sqlkw{AUTOINCREMENT} $\bullet$ rowid
$\bullet$ generated column
\end{keyterms}

\section{The Statement, in Full}

\sqlkw{CREATE TABLE} is the whole of data definition that this course needs.
Everything in it is a promise, and the order of the clauses is fixed.

\begin{definitionbox}[The anatomy of a \sqlkw{CREATE TABLE}]
\begin{lstlisting}[language=sql]
CREATE TABLE student (
  roll_no   TEXT    PRIMARY KEY NOT NULL,     -- column constraints
  name      TEXT    NOT NULL,
  programme TEXT    NOT NULL DEFAULT 'BSIT',
  semester  INTEGER NOT NULL CHECK (semester BETWEEN 1 AND 8),
  cnic      TEXT    UNIQUE,
  CHECK (length(roll_no) = 11)                -- a table constraint
) STRICT;
\end{lstlisting}

\smallskip
A \term{column constraint} follows its column and can mention only that
column. A \term{table constraint} is listed after all the columns and may
mention several. \sqlkw{PRIMARY KEY} on more than one column, and any
\sqlkw{CHECK} comparing two columns, must be table constraints --- there is no
other way to write them.
\end{definitionbox}

\begin{examplebox}[The five clauses, and what each one is for]
\begin{tabular}{@{}L{3.0cm}L{4.5cm}L{6.5cm}@{}}
\textbf{Clause} & \textbf{Promises} & \textbf{Comes from} \\
\sqlkw{PRIMARY KEY} & these values identify a row & the key chosen in
  \chref{ermodel} \\
\sqlkw{NOT NULL} & a value is always present & total participation
  (\chref{cardinality}), or entity integrity \\
\sqlkw{UNIQUE} & no two rows agree here & an \emph{alternate} key \\
\sqlkw{CHECK} & every value satisfies this & the attribute's domain
  (\chref{constraints}) \\
\sqlkw{DEFAULT} & what to use if none given & a university policy, not a
  guess \\
\end{tabular}

\smallskip
\textbf{\sqlkw{DEFAULT} is the one to be careful with.} It turns a missing
value into a real one, so a row that \emph{should} have been rejected is
accepted with a plausible value instead. Use it where the university genuinely
has a default --- \texttt{semester} is 1 for a new admission --- and not to
avoid thinking about nulls.
\end{examplebox}

\section{Type Affinity, and What It Costs}

Every other relational system treats a declared type as a rule. SQLite treats
it as a \emph{preference}, and the difference is the subject of this chapter.

\begin{definitionbox}[Storage class, declared type, affinity]
A \term{storage class} is what a value actually \emph{is} in the file: null,
integer, real, text or blob. Every value has one, and \texttt{typeof()} reports
it.

\smallskip
A \term{declared type} is what you wrote in \sqlkw{CREATE TABLE}.

\smallskip
\term{Type affinity} is the rule connecting them: on insert, SQLite
\emph{tries} to convert the value to the column's declared type. If it can, it
does. \textbf{If it cannot, it stores the value unchanged.} No error.

\smallskip
\sqlkw{STRICT}, added to SQLite in 2021, changes the last sentence: the value
is refused. A \sqlkw{STRICT} table may only use the types \texttt{INT},
\texttt{INTEGER}, \texttt{REAL}, \texttt{TEXT}, \texttt{BLOB} and
\texttt{ANY}.
\end{definitionbox}

\dsfig{%
\begin{tikzpicture}[
  val/.style={rounded corners=2pt, draw=secondary, line width=0.9pt,
              fill=secondary!8, text=secondary!60!black,
              font=\scriptsize\ttfamily, minimum width=1.7cm,
              minimum height=0.6cm},
  dec/.style={draw=primary, line width=1pt, fill=primary!7,
              text=primary!60!black, font=\scriptsize\bfseries,
              align=center, text width=3.1cm, minimum height=1.0cm,
              rounded corners=3pt},
  good/.style={rounded corners=2pt, draw=greenhl, line width=0.9pt,
               fill=greenhl!8, text=greenhl!55!black,
               font=\scriptsize\bfseries, align=center, text width=3.0cm,
               minimum height=0.75cm},
  bad/.style={good, draw=accent, fill=accent!8, text=accent!75!black},
  strict/.style={good, draw=purplehl, fill=purplehl!8,
                 text=purplehl!65!black},
  arr/.style={-{Stealth[length=2.2mm]}, primary!55, line width=0.9pt},
  tag/.style={font=\scriptsize\bfseries, text=primary!70!black},
  nt/.style={font=\scriptsize\itshape, text=neutral, align=left}
]
  \node[val] (v) at (-6.0,0) {'hello'};
  \node[nt, anchor=south, text width=2.6cm, align=center]
       at (-6.0,0.45) {a value, on its way into\\\texttt{c INTEGER}};

  \node[dec] (d) at (-2.4,0) {can it be converted\\to the column's\\declared type?};
  \draw[arr] (v) -- (d);

  \node[good]   (yes) at ( 2.0, 1.5) {converted\\{\scriptsize\mdseries
       \texttt{'42'} $\rightarrow$ integer 42}};
  \node[bad]    (no)  at ( 2.0, 0.0) {\textbf{stored unchanged}\\{\scriptsize
       \mdseries \texttt{'hello'} stays text}};
  \node[strict] (st)  at ( 2.0,-1.6) {refused\\{\scriptsize\mdseries
       \texttt{cannot store TEXT}}};

  \draw[arr] (d) -- node[tag, above, pos=0.6] {yes} (yes);
  \draw[arr] (d) -- node[tag, above, pos=0.62] {no} (no);
  % pos=0.33, not 0.6: further along the line the label landed on the box it
  % points at, and the arrowhead disappeared behind its own caption.
  \draw[arr, purplehl!60] (d) -- node[tag, below, pos=0.33,
       text=purplehl!70!black] {no, and \sqlkw{STRICT}} (st);

  \node[nt, anchor=west, text width=4.2cm, text=greenhl!55!black]
       at (3.9, 1.5) {silently succeeds};
  \node[nt, anchor=west, text width=4.4cm, text=accent!75!black]
       at (3.9, 0.0) {\textbf{silently fails.} No error, and the column
        looks the same as the row above it};
  \node[nt, anchor=west, text width=4.2cm, text=purplehl!65!black]
       at (3.9,-1.6) {the only outcome you can reason about};
\end{tikzpicture}}%
{Type affinity in one decision. The middle branch is the whole problem: it is
not an error, it is not reported, and a column containing both of the top two
outcomes looks exactly like a column containing only the first.}{affinity}

\begin{measured}[A column declared \texttt{INTEGER} that holds the word ``absent'']
\textsf{\footnotesize \texttt{code/m10\_ddl.py}, SQLite 3.51.2, Intel i5-6600,
16\,GB, Windows 11.}

\medskip
\textbf{1. One value, five declared types --- what is actually stored}

\begin{tabular}{@{}L{2.8cm}L{2.0cm}L{1.8cm}L{1.8cm}L{2.2cm}L{2.2cm}@{}}
\textbf{inserted} & \texttt{INTEGER} & \texttt{TEXT} & \texttt{REAL} &
\texttt{NUMERIC} & \textbf{(no type)} \\
\texttt{42} & integer & text & real & integer & integer \\
\texttt{'42'} & \textbf{integer} & text & real & integer & text \\
\texttt{'hello'} & \textbf{text} & text & \textbf{text} & \textbf{text} & text \\
\texttt{3.7} & \textbf{real} & text & real & real & real \\
\sqlkw{NULL} & null & null & null & null & null \\
\end{tabular}

\smallskip
Read the \texttt{'hello'} row. Five columns, five different declared types,
and every one of them stored the string.

\medskip
\textbf{2. The same column, one keyword different}

\begin{tabular}{@{}L{3.4cm}L{4.4cm}L{6.0cm}@{}}
\textbf{inserted into \texttt{c INTEGER}} & \textbf{plain} & \textbf{\sqlkw{STRICT}} \\
\texttt{42}      & stored as integer & stored as integer \\
\texttt{'42'}    & stored as integer & stored as integer \\
\texttt{'hello'} & \textbf{stored as text} & refused: \texttt{cannot store TEXT value in INT} \\
\texttt{3.7}     & \textbf{stored as real} & refused: \texttt{cannot store REAL value in INT} \\
\sqlkw{NULL}     & stored as null & stored as null \\
\end{tabular}

\medskip
\textbf{3. What it does to the answers}

Marks entered as \texttt{A=95, B=72, C='absent', D=88, E='100'}:

\begin{tabular}{@{}L{5.6cm}R{3.2cm}L{5.4cm}@{}}
\sqlkw{ORDER BY} \texttt{score} & \texttt{absent} last &
  text sorts \emph{above} every number \\
\sqlkw{MAX}(\texttt{score}) & \texttt{'absent'} &
  so C is top of the class \\
\sqlkw{AVG}(\texttt{score}) as written & \textbf{71.00} &
  \texttt{'absent'} read as 0, still divided by 5 \\
\sqlkw{AVG} over the numeric rows & 88.75 & the honest figure \\
\textbf{the reported average is} & \textbf{20\% out} & and nothing says so \\
\end{tabular}

\smallskip
Note also that \texttt{'100'} \emph{was} converted, to the integer 100.
Affinity succeeded silently there and failed silently on \texttt{'absent'},
and the column looks identical either way.

\medskip
\textbf{4. What \sqlkw{STRICT} costs} (50{,}000 rows, 5 runs)

\begin{tabular}{@{}L{5.0cm}R{2.6cm}R{2.4cm}R{2.4cm}@{}}
& \textbf{median} & \textbf{min} & \textbf{max} \\
plain          & 67.9\,ms & 66.2 & 78.5 \\
\sqlkw{STRICT} & 67.9\,ms & 66.8 & 72.8 \\
\textbf{difference} & \textbf{0.0\,ms} & \multicolumn{2}{l}{\textbf{(+0\%)}} \\
\end{tabular}

\smallskip
\textbf{The two ranges overlap}, so the difference is noise rather than a
cost --- which is the opposite of \chref{constraints}'s foreign-key figure,
where they did not overlap and the 24\% was real. \sqlkw{STRICT} is free: it
is checking a type it already knows.

\smallskip
\textbf{So the whole argument is one-sided.} The protection costs nothing
measurable and the absence of it silently moved a class average by twenty per
cent. Every table in this course is \sqlkw{STRICT}.
\end{measured}

\begin{alertbox}[Why SQLite is like this, and why that does not help you]
This is not a bug and it will not be fixed. SQLite was designed to be
embedded, tolerant and backward-compatible for decades; rejecting a value that
an older version accepted would break databases in the field. Affinity is the
documented, deliberate behaviour, and \sqlkw{STRICT} was added as an opt-in
twenty-one years later.

\smallskip
Knowing the reason does not change what you must do:

\begin{tight}
  \item \textbf{Write \sqlkw{STRICT} on every table.} There is no measured
        reason not to.
  \item \textbf{Use \texttt{typeof()} when a result looks wrong.} It is the
        first diagnostic, and the one nobody reaches for.
  \item \textbf{Be most suspicious of data you imported.} Values typed by hand
        into \sqlkw{INSERT} statements are usually the right type; values read
        from a CSV are all text until something converts them, and
        \chref{evolution} is where that bites.
  \item \textbf{Do not assume other systems behave this way.} PostgreSQL and
        MySQL reject \texttt{'hello'} in an integer column outright. This
        chapter's behaviour is SQLite's, and \chref{project} is where the
        comparison is drawn properly.
\end{tight}
\end{alertbox}

\section{Three Smaller Things That Matter}

\begin{conceptbox}[Rowid, \sqlkw{AUTOINCREMENT}, and generated columns]
\textbf{Every ordinary SQLite table has a hidden \texttt{rowid}} --- a 64-bit
integer key the system maintains. Declaring a column \texttt{INTEGER PRIMARY
KEY} makes that column \emph{be} the rowid rather than sit beside it, which is
why \chref{constraints} measured it substituting an integer for a null.

\smallskip
\textbf{\sqlkw{AUTOINCREMENT} is almost never what you want.} Without it,
SQLite reuses the largest rowid when rows are deleted; with it, it never does,
at the cost of an extra table and a write per insert. Use it only when a
reused identifier would be a real problem --- an invoice number, not a
\texttt{section\_id}.

\smallskip
\textbf{A generated column} is computed from the others and never stored
wrongly, which makes it the right home for \chref{ermodel}'s derived
attributes:

\begin{lstlisting}[language=sql]
  date_of_birth TEXT NOT NULL,
  age INTEGER GENERATED ALWAYS AS
      (CAST((julianday('now') - julianday(date_of_birth)) / 365.25
            AS INTEGER)) VIRTUAL
\end{lstlisting}

\sqlkw{VIRTUAL} computes it on every read and stores nothing;
\sqlkw{STORED} writes it to disk. For \texttt{age}, \sqlkw{VIRTUAL} is the
only correct choice: a stored one would be wrong by morning, which was exactly
\chref{ermodel}'s objection to storing it at all.
\end{conceptbox}

\begin{pitfall}
\begin{tight}
  \item \textbf{Trusting a declared type.} Without \sqlkw{STRICT} it is a
        preference. Measured above.
  \item \textbf{Using \texttt{VARCHAR(50)}.} SQLite accepts the words and
        ignores the 50 entirely. Write \texttt{TEXT}, and enforce a length
        with \sqlkw{CHECK} if you mean one.
  \item \textbf{Declaring a composite primary key as two column
        constraints.} \texttt{a INTEGER PRIMARY KEY, b TEXT PRIMARY KEY} is an
        error; it must be one table constraint.
  \item \textbf{Using \sqlkw{DEFAULT} to dodge a null.} It replaces a
        rejection with a plausible wrong value.
  \item \textbf{Reaching for \sqlkw{AUTOINCREMENT} by habit.} It costs a write
        per insert and solves a problem you probably do not have.
  \item \textbf{Storing a derived value.} Use a \sqlkw{VIRTUAL} generated
        column.
  \item \textbf{Forgetting that \sqlkw{CHECK} is not re-run on existing rows.}
        \chref{evolution} is about exactly that.
\end{tight}
\end{pitfall}

\begin{lab}[Store the wrong type on purpose, then make it impossible]
\begin{lstlisting}[language=bash]
# --- 1. A scratch database; keep university.db clean ----------------
sqlite3 types.db
\end{lstlisting}
\begin{lstlisting}[language=sql]
-- --- 2. A table that looks completely ordinary ------------------
.headers on
.mode box
CREATE TABLE marks (roll TEXT PRIMARY KEY, score INTEGER NOT NULL);

-- --- 3. Put four marks in, one of them the word 'absent' --------
INSERT INTO marks VALUES ('A',95),('B',72),('C','absent'),('D',88);
SELECT roll, score, typeof(score) FROM marks;
-- One row is text. Nothing warned you. The column says INTEGER.

-- --- 4. Ask three ordinary questions ----------------------------
SELECT roll, score FROM marks ORDER BY score DESC;
SELECT MAX(score) FROM marks;
SELECT AVG(score) FROM marks;
-- Write all three answers down. Then work out, by hand, what each
-- one SHOULD be. Two of the three are wrong.

-- --- 5. Find the bad row, which is the real skill ---------------
SELECT roll, score FROM marks
 WHERE typeof(score) NOT IN ('integer','real');
-- Learn this query. It is how you audit a column you inherited.

-- --- 6. The honest average --------------------------------------
SELECT AVG(score) FROM marks
 WHERE typeof(score) IN ('integer','real');
-- Compare with step 4. How far out was it, as a percentage?

-- --- 7. Make it impossible --------------------------------------
CREATE TABLE marks2 (
  roll  TEXT    PRIMARY KEY NOT NULL,
  score INTEGER NOT NULL CHECK (score BETWEEN 0 AND 100)
) STRICT;
INSERT INTO marks2 VALUES ('A',95);
INSERT INTO marks2 VALUES ('C','absent');
INSERT INTO marks2 VALUES ('E',150);
-- Two refusals, two different messages. Record both, and say
-- which constraint produced each.

-- --- 8. Prove VARCHAR(50) is not a length limit -----------------
CREATE TABLE v (s VARCHAR(5));
INSERT INTO v VALUES ('this string is far longer than five');
SELECT length(s) FROM v;
-- Say in one sentence what VARCHAR(5) accomplished.
\end{lstlisting}
\begin{lstlisting}[language=bash]
# --- 9. The measured version ----------------------------------------
python code/m10_ddl.py

# --- 10. Apply it to your own schema ---------------------------------
# Every table in code/schema/v1_create.sql must end in STRICT.
# Check, and then prove it:
sqlite3 university.db "SELECT name, sql FROM sqlite_master
  WHERE type='table' AND sql NOT LIKE '%STRICT%';"
# That query should return nothing.
\end{lstlisting}

\smallskip
\textbf{What to notice.} Step 4 is the point of the chapter: three ordinary
questions, no errors, two wrong answers. Step 5 is the skill --- a
\texttt{typeof()} audit is the only way to find this in a database you did not
build, and it is worth running on any data you import. Step 10 is the
protection, and it costs nothing measurable.
\end{lab}

\begin{checkpoint}
\begin{tightnum}
  \item A column is declared \texttt{INTEGER} and \sqlkw{AVG} over it returns
        71 when the true average of the real marks is 88.75. Explain exactly
        how the database arrived at 71, and give the two queries you would run
        to prove your explanation.
  \item State the difference between a column constraint and a table
        constraint, and give two constraints from your own
        \texttt{v1\_create.sql} that \emph{must} be table constraints, saying
        why each one cannot be written as a column constraint.
  \item \sqlkw{STRICT} costs nothing measurable and prevents a class of silent
        error. Give the one argument somebody could still make against putting
        it on every table, and say whether you accept it.
\end{tightnum}
\end{checkpoint}

\begin{chaptersummary}
\begin{tight}
  \item \sqlkw{CREATE TABLE} carries five kinds of promise: \sqlkw{PRIMARY
        KEY}, \sqlkw{NOT NULL}, \sqlkw{UNIQUE}, \sqlkw{CHECK} and
        \sqlkw{DEFAULT}, each traceable to a decision in Part~II.
  \item A \textbf{column constraint} mentions one column; a \textbf{table
        constraint} may mention several. Composite keys and cross-column
        \sqlkw{CHECK}s must be table constraints.
  \item A \textbf{storage class} is what a value is; a \textbf{declared type}
        is what you asked for; \textbf{affinity} converts if it can and
        \textbf{stores the value unchanged if it cannot}.
  \item \textbf{Measured:} \texttt{'hello'} went into columns declared
        \texttt{INTEGER}, \texttt{REAL} and \texttt{NUMERIC}, unchanged and
        without an error.
  \item \textbf{Measured:} on a mixed column, \sqlkw{ORDER BY} sorts text above
        every number, \sqlkw{MAX} returns \texttt{'absent'}, and \sqlkw{AVG}
        reads it as 0 while still dividing by the full count --- a class
        average reported \textbf{20\% low}, silently.
  \item \texttt{'100'} \emph{was} converted to 100. Affinity succeeds and fails
        silently, and the column looks the same either way.
  \item \textbf{Measured: \sqlkw{STRICT} costs nothing} --- the two timing
        ranges overlap, unlike \chref{constraints}'s foreign-key figure. The
        argument for it is one-sided.
  \item \texttt{typeof()} is the first diagnostic when a result looks wrong,
        and the one nobody reaches for. Be most suspicious of imported data.
  \item \texttt{VARCHAR(50)} has no length limit in SQLite. Use \texttt{TEXT}
        plus a \sqlkw{CHECK}.
  \item A \textbf{\sqlkw{VIRTUAL} generated column} is the right home for a
        derived attribute: computed on read, never stale.
\end{tight}
\end{chaptersummary}

\begin{reviewq}
  \item \textbf{(MCQ)} Inserting \texttt{'hello'} into a column declared
        \texttt{INTEGER} in a non-\sqlkw{STRICT} table: (a)~raises an error;
        (b)~stores 0; (c)~stores the text unchanged; (d)~stores null.
  \item \textbf{(MCQ)} \texttt{VARCHAR(5)} in SQLite limits the column to:
        (a)~5 characters; (b)~5 bytes; (c)~nothing; (d)~5 characters only in
        \sqlkw{STRICT} tables.
  \item \textbf{(Short)} Why does \texttt{'42'} become an integer while
        \texttt{'hello'} stays text? State the rule in one sentence.
  \item \textbf{(Short)} Give a constraint that can only be written as a table
        constraint, and say why.
  \item \textbf{(Short)} \sqlkw{STRICT} showed no measurable cost while foreign
        key enforcement showed 24\%. Explain why the two results differ.
  \item \textbf{(Applied)} You are given a table of 50{,}000 imported rows with
        columns \texttt{roll\_no}, \texttt{score} and \texttt{date\_sat}.
        Write the queries that would tell you whether any value is of the wrong
        type, decide which results you would report as faults, and describe
        how you would repair the table without losing a row.
  \item \textbf{(Applied)} Rewrite one table of your \texttt{v1\_create.sql} to
        add a \sqlkw{DEFAULT}, a cross-column \sqlkw{CHECK} and a
        \sqlkw{VIRTUAL} generated column. For each, say what fact about the
        university it records and what would go wrong if it were omitted ---
        and for the \sqlkw{DEFAULT}, say what it would wrongly accept.
\end{reviewq}

%=====================================================================
\chapter{Changing a Schema That Is Already Full}\label{ch:evolution}
%=====================================================================
\locator{3}{Part III --- From Diagram to Schema}

\begin{outcomes}
By the end of this chapter you will be able to:
\begin{tight}
  \item \textbf{Classify} a proposed schema change as cheap, expensive or
        impossible, and \textbf{say} what decides which.
  \item \textbf{Use} the four forms of \sqlkw{ALTER TABLE} SQLite supports, and
        \textbf{state} the limits of each.
  \item \textbf{Perform} the twelve-step rebuild for a change
        \sqlkw{ALTER TABLE} cannot make.
  \item \textbf{Write} a migration as a numbered, repeatable script, and
        \textbf{record} the schema version in the database itself.
  \item \textbf{Backfill} a new \sqlkw{NOT NULL} column on a table that already
        has rows.
\end{tight}
\end{outcomes}

\begin{prereq}
\chref{ddl}: \sqlkw{CREATE TABLE} and \sqlkw{STRICT}. \chref{mapping} for the
schema you are about to change --- this chapter modifies the
\texttt{university.db} you built there, and produces
\texttt{v2\_constraints.sql}.
\end{prereq}

\begin{keyterms}
schema evolution $\bullet$ migration $\bullet$ \sqlkw{ALTER TABLE} $\bullet$
\sqlkw{ADD COLUMN} $\bullet$ \sqlkw{RENAME TO} $\bullet$ \sqlkw{RENAME COLUMN}
$\bullet$ \sqlkw{DROP COLUMN} $\bullet$ the twelve-step rebuild $\bullet$
backfill $\bullet$ \sqlkw{PRAGMA} \texttt{user\_version} $\bullet$
forward migration $\bullet$ rollback
\end{keyterms}

\section{The Chapter the Outline Does Not Name}

HEC's course content lists nothing about changing a schema --- \term{schema
evolution}, in the textbooks' phrase. It is here anyway,
for a reason you can check against your own timetable: this course carries one
database through sixteen weeks. In week~5 it had seven tables; by
\chref{normalforms} it has more, by \chref{views} it has views and triggers,
and by \chref{indexes} it has indexes. Every one of those is a schema change
made to a database that already has rows in it.

A student who can only create a schema from scratch has to start again every
time the design improves --- and in a real system, starting again means losing
the data.

\begin{conceptbox}[Three kinds of change, and what decides the cost]
\begin{tabular}{@{}L{3.6cm}L{4.6cm}L{5.8cm}@{}}
\textbf{Kind} & \textbf{Examples} & \textbf{Cost} \\
\textbf{Cheap} & add a nullable column; rename; add an index &
  a change to the \emph{schema only}. The rows are untouched. \\
\textbf{Expensive} & change a type; add a constraint; split a table &
  every row must be read and rewritten. \\
\textbf{Impossible} & recover a fact nobody recorded &
  no amount of work helps; only a new requirement and future data. \\
\end{tabular}

\smallskip
\textbf{What decides it is whether existing rows must change.} Adding a
nullable column does not: every existing row's new value is null, and null
needs no writing. Adding \sqlkw{NOT NULL} does, because every existing row
needs a value that is not there.

\smallskip
The third row is the one worth remembering. \chref{ermodel} said a
\texttt{address} stored whole could never answer ``who lives in Bahawalpur''.
That is not expensive --- the city was never recorded separately and splitting
the column cannot invent it. Design errors in Part~II become \emph{impossible}
here, which is why Part~II took four chapters.
\end{conceptbox}

\section{What \sqlkw{ALTER TABLE} Can Do}

\begin{definitionbox}[The four forms, and the limits of each]
\begin{tight}
  \item \sqlkw{ALTER TABLE} $t$ \sqlkw{RENAME TO} $u$ --- renames the table,
        and since 3.25 updates every reference to it in other tables' foreign
        keys, views and triggers.
  \item \sqlkw{ALTER TABLE} $t$ \sqlkw{RENAME COLUMN} $a$ \sqlkw{TO} $b$ ---
        renames a column, updating references.
  \item \sqlkw{ALTER TABLE} $t$ \sqlkw{ADD COLUMN} $c$ \dots\ --- adds a
        column at the \emph{end}. It may not be \sqlkw{PRIMARY KEY} or
        \sqlkw{UNIQUE}; it may not have a non-constant default; and it may be
        \sqlkw{NOT NULL} \textbf{only if it has a default}.
  \item \sqlkw{ALTER TABLE} $t$ \sqlkw{DROP COLUMN} $c$ --- removes a column,
        provided nothing depends on it: not a primary key, not indexed, not
        mentioned in a constraint, view or trigger.
\end{tight}

\smallskip
That is the complete list. \textbf{There is no \sqlkw{ALTER COLUMN}}, so a
type cannot be changed, a constraint cannot be added, and a column cannot be
moved --- not by \sqlkw{ALTER TABLE}.
\end{definitionbox}

\begin{examplebox}[The three cheap changes, and the one that catches people]
\begin{lstlisting}[language=sql]
-- Cheap: a nullable column. Instant on any number of rows.
ALTER TABLE student ADD COLUMN email TEXT;

-- Cheap: renames.
ALTER TABLE student RENAME COLUMN cnic TO national_id;

-- Refused: NOT NULL with no default, on a table with rows.
ALTER TABLE student ADD COLUMN email TEXT NOT NULL;
--   Cannot add a NOT NULL column with default value NULL
\end{lstlisting}

\smallskip
\textbf{The third is not an arbitrary restriction.} Every existing row would
need a value for \texttt{email} and there is none, so the table would be in a
state its own schema forbids the instant the statement finished. SQLite
refuses rather than produce that.

\smallskip
The two honest ways out are both in the next section, and the dishonest one is
\sqlkw{DEFAULT} \texttt{''} --- which satisfies the constraint with a value
that means nothing, and is \chref{relational}'s magic value wearing a new hat.
\end{examplebox}

\section{The Twelve-Step Rebuild}

For everything else --- changing a type, adding a \sqlkw{CHECK}, adding
\sqlkw{STRICT}, reordering columns --- there is one procedure, and it is in
SQLite's own documentation. It is worth doing by hand once.

\begin{worked}[Adding a constraint to a table that is already full]
The requirement: \texttt{student.semester} must be between 1 and 8. The table
exists, has 4{,}000 rows, and has no such \sqlkw{CHECK}.

\medskip
\textbf{Step 0 --- find out whether the data already satisfies it.} Before
anything else:
\begin{lstlisting}[language=sql]
SELECT roll_no, semester FROM student
 WHERE semester NOT BETWEEN 1 AND 8 OR semester IS NULL;
\end{lstlisting}
If this returns rows, the migration will fail at step 5 and you need a
decision from the university first --- not a cleverer script. \textbf{This is
the step people skip, and it is the one that turns a ten-minute migration into
an afternoon.}

\medskip
\textbf{Steps 1--12, as a script:}
\begin{lstlisting}[language=sql]
PRAGMA foreign_keys = OFF;          -- 1. and NOT inside the transaction
BEGIN;                              -- 2. all of it, or none of it

CREATE TABLE student_new (          -- 3. the table you wish you had
  roll_no   TEXT    PRIMARY KEY NOT NULL,
  name      TEXT    NOT NULL,
  programme TEXT    NOT NULL,
  semester  INTEGER NOT NULL CHECK (semester BETWEEN 1 AND 8),
  cnic      TEXT    UNIQUE
) STRICT;

INSERT INTO student_new                -- 4. copy, naming every column
  SELECT roll_no, name, programme, semester, cnic FROM student;

DROP TABLE student;                    -- 5. the old one
ALTER TABLE student_new RENAME TO student;   -- 6. take its name

CREATE INDEX idx_student_name ON student(name);  -- 7. rebuild indexes
-- 8. recreate views and triggers that referenced it

PRAGMA foreign_key_check;           -- 9. did anything become an orphan?
COMMIT;                             -- 10.
PRAGMA foreign_keys = ON;           -- 11.
PRAGMA integrity_check;             -- 12.
\end{lstlisting}

\medskip
\textbf{Four things about this that are not obvious.}
\begin{tight}
  \item \textbf{\sqlkw{PRAGMA} foreign\_keys goes \emph{outside} the
        transaction}, both times. Inside one it is silently ignored --- which
        is \chref{constraints}'s pitfall arriving in a place where the
        consequence is a dropped foreign key rather than a rejected insert.
  \item \textbf{Name every column in the \sqlkw{INSERT}.} \sqlkw{SELECT *}
        depends on column order, which is exactly what you may be changing.
  \item \textbf{Indexes, views and triggers do not survive \sqlkw{DROP
        TABLE}.} Step 7 is not optional and is the most commonly forgotten.
  \item \textbf{It is $O(n)$.} Every row is read and written. On this
        database that is instant; on a real one it is the difference between a
        maintenance window and a Tuesday.
\end{tight}
\end{worked}

\dsfig{%
\begin{tikzpicture}[
  box/.style={rounded corners=3pt, draw=#1, line width=1pt, fill=#1!7,
              text=#1!60!black, font=\scriptsize\bfseries, align=center,
              text width=2.9cm, minimum height=0.8cm},
  arr/.style={-{Stealth[length=2.2mm]}, primary!55, line width=0.9pt},
  q/.style={font=\scriptsize\itshape, text=neutral, align=center},
  tag/.style={font=\scriptsize\bfseries, align=center}
]
  \node[box=primary] (start) at (0,2.6) {the change\\you want};
  \node[q, text width=4.6cm] (q1) at (0,1.2)
       {does any \emph{existing row} have to change?};
  \node[box=greenhl] (cheap) at (-4.6,-0.3) {ALTER TABLE\\{\scriptsize
       \mdseries schema only, instant}};
  \node[q, text width=4.4cm] (q2) at (2.2,-0.3)
       {is the information already in the database?};
  \node[box=goldhl] (reb) at (0.4,-2.1) {the 12-step\\rebuild};
  \node[box=accent]  (imp) at (5.0,-2.1) {\textbf{impossible}};

  \draw[arr] (start) -- (q1);
  \draw[arr] (q1) -- node[tag, above, text=greenhl!60!black, pos=0.55] {no}
             (cheap);
  \draw[arr] (q1) -- node[tag, above right, text=neutral, pos=0.5] {yes}
             (q2);
  \draw[arr] (q2) -- node[tag, left, text=goldhl!70!black] {yes} (reb);
  \draw[arr] (q2) -- node[tag, right, text=accent] {no} (imp);

  \node[q, anchor=north, text width=3.6cm, text=greenhl!60!black]
       at (-4.6,-1.1) {add a nullable column, rename, add an index};
  \node[q, anchor=north, text width=3.6cm, text=goldhl!70!black]
       at (0.4,-2.9) {change a type, add a \sqlkw{CHECK}, split a table,
        backfill};
  \node[q, anchor=north, text width=3.6cm, text=accent]
       at (5.0,-2.9) {recover a fact nobody ever recorded --- a Part~II error};
\end{tikzpicture}}%
{Two questions decide what a schema change costs. The right-hand branch is the
one Part~II existed to prevent: no procedure recovers information that was
never stored.}{evolution}

\section{Migrations as Files}

\begin{definitionbox}[A migration]
A \term{migration} is a numbered script that moves the database from one schema
version to the next. It is kept in the repository beside the code, it is run
once per database, and it is never edited after it has been run anywhere.

\smallskip
\sqlkw{PRAGMA} \texttt{user\_version} is a 32-bit integer SQLite stores in the
database file and otherwise ignores completely. It is yours to use, and it is
how a script knows which migrations a given file has already had.
\end{definitionbox}

\begin{examplebox}[The pattern, and it is four lines]
\begin{lstlisting}[language=sql]
-- v2_constraints.sql -- run only against a version 1 database
PRAGMA foreign_keys = OFF;
BEGIN;

-- ... the changes ...

PRAGMA user_version = 2;      -- the last thing before COMMIT
COMMIT;
PRAGMA foreign_keys = ON;
\end{lstlisting}

and the check before running it:
\begin{lstlisting}[language=bash]
sqlite3 university.db "PRAGMA user_version;"
# 1  -> run v2.   2 -> already applied, do nothing.
\end{lstlisting}

\smallskip
\textbf{Why the version is set inside the transaction.} If the migration
fails half-way, the rollback takes the version number with it, and the
database is honestly still at version 1. Set it afterwards and a failed
migration leaves a database that claims to be something it is not --- which is
worse than the failure.
\end{examplebox}

\begin{alertbox}[Four rules that are about people, not about SQL]
\begin{tight}
  \item \textbf{Never edit a migration that has run anywhere.} Somebody else's
        database already has the old version applied; changing the file means
        two databases with the same version number and different schemas, and
        nothing will ever tell you.
  \item \textbf{Take a backup first, and check that you can restore it.} For
        SQLite this is \texttt{copy university.db university.db.bak} and then
        actually opening the copy. An untested backup is a hope.
  \item \textbf{Make the change in two deployments when a column is going
        away.} Stop writing it, ship, confirm nothing broke, then drop it.
        Dropping and deploying together means the rollback is also a
        migration.
  \item \textbf{Write the \emph{reverse} migration, even if you never run
        it.} Writing it is how you find out that the change is irreversible
        --- and finding that out before you run it is the entire point.
  \end{tight}
\end{alertbox}

\begin{pitfall}
\begin{tight}
  \item \textbf{Adding \sqlkw{NOT NULL} without a default to a full table.}
        Refused, and correctly: every existing row would violate it.
  \item \textbf{Satisfying it with \sqlkw{DEFAULT} \texttt{''}.} A magic value
        that will be counted and compared as though it meant something.
  \item \textbf{Setting \sqlkw{PRAGMA} foreign\_keys inside the transaction.}
        Silently ignored, and in a rebuild the consequence is a dropped
        foreign key.
  \item \textbf{\sqlkw{INSERT} \dots\ \sqlkw{SELECT *} in a rebuild.} It
        depends on the column order you are changing.
  \item \textbf{Forgetting to recreate indexes, views and triggers.} They go
        with the dropped table and nothing reminds you.
  \item \textbf{Skipping step 0.} Check the data satisfies the new constraint
        \emph{before} writing the migration.
  \item \textbf{Editing a migration after it has run.} Two databases, one
        version number, different schemas.
\end{tight}
\end{pitfall}

\begin{lab}[Migrate the database you built in Chapter 9]
This lab produces \texttt{code/schema/v2\_constraints.sql}.

\begin{lstlisting}[language=bash]
# --- 1. Back it up, and prove the backup opens ----------------------
cp university.db university.db.bak
sqlite3 university.db.bak ".tables"
sqlite3 university.db "PRAGMA user_version;"     # 0 or 1
\end{lstlisting}
\begin{lstlisting}[language=sql]
-- --- 2. The cheap change: a nullable column ---------------------
ALTER TABLE student ADD COLUMN email TEXT;
SELECT roll_no, email FROM student LIMIT 3;
-- Instant, whatever the row count. Nothing was rewritten.

-- --- 3. The same column, NOT NULL -------------------------------
ALTER TABLE student ADD COLUMN email2 TEXT NOT NULL;
-- Record the error. Explain it in terms of the existing rows.

-- --- 4. Step 0 of any rebuild: does the data already comply? ----
SELECT COUNT(*) FROM student WHERE semester NOT BETWEEN 1 AND 8;
-- If this is not 0, stop and decide what the right value is.

-- --- 5. Backfill, so NOT NULL becomes possible ------------------
UPDATE student SET email = roll_no || '@iub.edu.pk' WHERE email IS NULL;
SELECT COUNT(*) FROM student WHERE email IS NULL;
-- Now every row has a value, and only now is the constraint addable.
\end{lstlisting}
\begin{lstlisting}[language=bash]
# --- 6. Write the migration -----------------------------------------
# Put steps 1-12 into code/schema/v2_constraints.sql, adding to the
# student table:  email TEXT NOT NULL, and a CHECK that semester is
# between 1 and 8. Set PRAGMA user_version = 2 before COMMIT.
\end{lstlisting}
\begin{lstlisting}[language=sql]
-- --- 7. Run it --------------------------------------------------
.read code/schema/v2_constraints.sql
PRAGMA user_version;
.schema student
PRAGMA integrity_check;
PRAGMA foreign_key_check;

-- --- 8. Prove the new constraint is live ------------------------
INSERT INTO student VALUES
  ('2021-IT-999','Test','BSIT',99,NULL,NULL,NULL,NULL,'x@y.z');
-- Must fail. Which constraint caught it?

-- --- 9. Prove the old data survived -----------------------------
SELECT COUNT(*) FROM student;
-- Compare with the count before step 6. They must match.

-- --- 10. Write the reverse migration ----------------------------
-- Write v2_rollback.sql, which undoes step 6. Do not run it.
-- Then answer: is it a true reverse? What cannot be undone?
\end{lstlisting}

\smallskip
\textbf{What to notice.} Step 10 is the exercise. The forward migration added a
\sqlkw{CHECK} and a column; the reverse can drop both --- but the
\emph{backfilled email addresses} are now indistinguishable from real ones,
and no rollback can tell them apart. \textbf{A migration that writes data is
not reversible, whatever the script says}, and knowing which of your
migrations are in that class is the whole of the skill.
\end{lab}

\begin{checkpoint}
\begin{tightnum}
  \item Classify each as cheap, expensive or impossible, with the reason:
        adding a nullable \texttt{email}; adding \texttt{email} \sqlkw{NOT
        NULL}; changing \texttt{semester} from \texttt{TEXT} to
        \texttt{INTEGER}; splitting \texttt{address} into house, street and
        city on a table where only the whole address was ever stored.
  \item A migration adds a \sqlkw{NOT NULL} column by backfilling it with a
        computed value. Explain why the rollback script is not a true reverse,
        and say what you would record at migration time to make it one.
  \item \sqlkw{PRAGMA} \texttt{user\_version} is set inside the transaction
        rather than after it. Say what goes wrong if it is set after, and
        describe the state the database would be left in.
\end{tightnum}
\end{checkpoint}

\begin{chaptersummary}
\begin{tight}
  \item A schema change is \textbf{cheap} if no existing row must change,
        \textbf{expensive} if every row must be rewritten, and
        \textbf{impossible} if it needs information nobody recorded.
  \item The third class is a Part~II design error arriving late. No procedure
        recovers a fact that was never stored.
  \item SQLite's \sqlkw{ALTER TABLE} does four things: \sqlkw{RENAME TO},
        \sqlkw{RENAME COLUMN}, \sqlkw{ADD COLUMN}, \sqlkw{DROP COLUMN}.
        \textbf{There is no \sqlkw{ALTER COLUMN}}, so types and constraints
        cannot be changed by it.
  \item \sqlkw{ADD COLUMN} may be \sqlkw{NOT NULL} only with a default,
        because every existing row needs a value --- and a default of
        \texttt{''} is a magic value, not a solution.
  \item Everything else is the \textbf{twelve-step rebuild}: new table, copy
        naming every column, drop, rename, \textbf{rebuild indexes, views and
        triggers}, check, commit.
  \item \sqlkw{PRAGMA} \texttt{foreign\_keys} goes \emph{outside} the
        transaction. Inside one it is ignored, and in a rebuild that costs you
        a foreign key.
  \item A \textbf{migration} is a numbered script, kept in the repository, run
        once, and \textbf{never edited after it has run anywhere}.
  \item \sqlkw{PRAGMA} \texttt{user\_version} records which migrations a file
        has had. Set it \emph{inside} the transaction, so a failed migration
        leaves an honest version number.
  \item Back up first and open the backup. Drop a column in a second
        deployment. Write the reverse migration even if you never run it.
  \item \textbf{A migration that writes data is not reversible}, whatever the
        rollback script claims --- backfilled values cannot be told from real
        ones afterwards.
\end{tight}
\end{chaptersummary}

\begin{reviewq}
  \item \textbf{(MCQ)} Which of these can SQLite's \sqlkw{ALTER TABLE} do?
        (a)~change a column's type; (b)~add a \sqlkw{CHECK} constraint;
        (c)~rename a column; (d)~reorder columns.
  \item \textbf{(MCQ)} \sqlkw{ADD COLUMN} \texttt{x TEXT NOT NULL} on a table
        with rows is: (a)~instant; (b)~refused unless a default is given;
        (c)~refused always; (d)~allowed but leaves nulls.
  \item \textbf{(Short)} Why must \sqlkw{PRAGMA} \texttt{foreign\_keys = OFF}
        appear before \sqlkw{BEGIN} rather than after it?
  \item \textbf{(Short)} Name two things that are lost by \sqlkw{DROP TABLE}
        during a rebuild and must be recreated explicitly.
  \item \textbf{(Short)} Give a schema change that is impossible, and say what
        decision in Part~II would have prevented it.
  \item \textbf{(Applied)} Write the full twelve-step migration that changes
        \texttt{enrolment.grade} from an unconstrained \texttt{TEXT} column to
        one with \sqlkw{CHECK} \texttt{(grade IN ('A','B','C','D','F'))}.
        Include step 0, say what you would do if step 0 returned 40 rows with
        a grade of \texttt{'I'}, and state which parts of your script are not
        reversible.
  \item \textbf{(Applied)} Your university decides, in week 12, that an
        instructor may co-teach: \ent{teaches} becomes M:N. Describe the
        complete migration --- new table, backfill from the existing
        \texttt{section.emp\_id}, what to do with the old column, and in how
        many deployments you would ship it. Then say which single step is
        irreversible and why.
\end{reviewq}


\part{Asking Questions}
%=====================================================================
% PART IV OPENER
%=====================================================================
\thispagestyle{plain}
\structuralfigures{P4.}

\begin{center}
\begin{tcolorbox}[enhanced, width=0.92\textwidth, colback=purplehl!5,
  colframe=purplehl, boxrule=0pt, leftrule=5pt, arc=3pt, top=8pt, bottom=8pt]
\large\itshape\color{purplehl!60!black}
This is the longest part of the course and the one employers mean when they ask whether you know databases. It begins with the algebra, because SQL is a notation for it and a reader who has seen the algebra first can tell what a query is doing from its shape. Then five chapters of asking the database questions --- and one recurring difficulty, which is that the answer to a question about data you do not have is not ``none''.
\end{tcolorbox}
\end{center}

\vspace{6pt}

\dsfigh{%
\begin{tikzpicture}[
  cbox/.style={rounded corners=3pt, draw=purplehl, line width=1pt, fill=purplehl!7,
               text=purplehl!55!black, font=\small\bfseries, align=center,
               minimum width=2.8cm, minimum height=1.0cm},
  hbox/.style={rounded corners=3pt, draw=purplehl, line width=1.8pt, fill=purplehl!16,
               text=purplehl!55!black, font=\small\bfseries, align=center,
               minimum width=2.8cm, minimum height=1.0cm},
  arr/.style={-{Stealth[length=2.4mm]}, purplehl!60, line width=1pt},
  note/.style={font=\scriptsize\itshape, text=purplehl!55!black, align=center,
               fill=white, inner sep=2pt}
]
  \node[cbox] (c0) at (0.00,0)  {Ch.\ 12\\Algebra};
  \node[cbox] (c1) at (3.40,0)  {Ch.\ 13\\Joins};
  \node[hbox] (c2) at (6.80,0)  {Ch.\ 14\\Selecting};
  \node[cbox] (c3) at (10.20,0) {Ch.\ 15\\Sub-queries};
  \node[cbox] (c4) at (13.60,0) {Ch.\ 16\\Grouping};
  \draw[arr] (c0) -- (c1);
  \draw[arr] (c1) -- (c2);
  \draw[arr] (c2) -- (c3);
  \draw[arr] (c3) -- (c4);
  \node[note] at (1.70,-1.0) {the theory};
  \node[note] at (8.50,-1.0) {and the trouble with NULL};
  \node[note] at (11.90,-1.0) {questions inside questions};
  \node[font=\scriptsize\itshape, text=accent, align=center] at (6.80,1.15)
       {three-valued logic: the one idea in this part that is genuinely
        counter-intuitive};
\end{tikzpicture}}%
{Reading path through Part~IV. Chapter 12 is the only one here with no SQL in
it, and it is the one that makes the other four make sense.}{part4map}

\newpage

%=====================================================================
\chapter{Relational Algebra}\label{ch:algebra}
%=====================================================================
\locator{4}{Part IV --- Asking Questions}

\begin{outcomes}
By the end of this chapter you will be able to:
\begin{tight}
  \item \textbf{Apply} selection, projection, Cartesian product, union,
        difference and rename to small relations by hand.
  \item \textbf{Write} an algebra expression for a question stated in English.
  \item \textbf{Translate} between an algebra expression and the
        \sqlkw{SELECT} statement that computes it.
  \item \textbf{Explain} why the output of every operator is a relation, and
        what that makes possible.
  \item \textbf{State} which operators need union-compatible inputs, and
        \textbf{check} that two relations are.
\end{tight}
\end{outcomes}

\begin{prereq}
\chref{relational}: a relation is a set of tuples. Appendix~B for set union,
intersection and difference --- if those three words are not routine, read it
first; it is four pages. This chapter is the first in the course with a
symbol on most pages.
\end{prereq}

\begin{keyterms}
relational algebra $\bullet$ operator $\bullet$ closure $\bullet$ selection
($\sigma$) $\bullet$ projection ($\pi$) $\bullet$ Cartesian product
($\times$) $\bullet$ union ($\cup$) $\bullet$ set difference ($-$) $\bullet$
intersection ($\cap$) $\bullet$ rename ($\rho$) $\bullet$ union-compatible
$\bullet$ predicate $\bullet$ query tree $\bullet$ procedural $\bullet$
declarative
\end{keyterms}

\section{Why a Course on SQL Starts With Algebra}

SQL is the language you will use. Relational algebra is the language the
\emph{database} uses --- the optimiser of \chref{whodoeswhat} turns your
\sqlkw{SELECT} into an algebra expression, rearranges it, and runs that.
Knowing the algebra is knowing what your query becomes.

It also answers a question that SQL hides: why can a query be used inside
another query? The answer is one property, and it is the first definition.

\begin{definitionbox}[Relational algebra, and closure]
The \term{relational algebra} is a set of operators. Each takes one or two
relations as input and produces \textbf{a relation} as output.

\smallskip
That last fact is \term{closure}: because every result is a relation, every
result can be the input to another operator. Expressions nest without limit,
exactly as $(3+4)\times 2$ does in arithmetic --- and for the same reason, which
is that $3+4$ is a number.

\smallskip
Six operators are fundamental: $\sigma$, $\pi$, $\times$, $\cup$, $-$ and
$\rho$. Every other operator, including every kind of join, can be written in
terms of these six.
\end{definitionbox}

\section{The Two Unary Operators}

\begin{definitionbox}[Selection and projection]
\term{Selection} $\sigma_{p}(R)$ returns the tuples of $R$ for which the
predicate $p$ is true. It chooses \textbf{rows}. The result has the same
attributes as $R$ and at most as many tuples.

\smallskip
\term{Projection} $\pi_{A_1,\dots,A_k}(R)$ returns the tuples of $R$
restricted to the attributes $A_1,\dots,A_k$. It chooses \textbf{columns}.
Because the result is a set, \textbf{duplicates are removed} --- so it may have
fewer tuples than $R$.
\end{definitionbox}

\dsfig{%
\begin{tikzpicture}[
  cell/.style={draw=primary!35, line width=0.5pt, minimum width=1.55cm,
               minimum height=0.42cm, font=\scriptsize\ttfamily,
               text=primary!65!black, inner sep=0pt},
  hd/.style={cell, fill=primary, text=white, font=\scriptsize\bfseries},
  sel/.style={cell, fill=greenhl!18},
  prj/.style={cell, fill=secondary!18},
  ttl/.style={font=\scriptsize\bfseries, text=#1, anchor=south},
  nt/.style={font=\scriptsize\itshape, text=neutral, align=center}
]
  % --- selection: rows
  \node[ttl=greenhl!60!black] at (-3.9,0.55)
       {$\sigma_{\,\mathit{programme}=\mathrm{'BSIT'}}(\mathit{student})$};
  \foreach \c/\h [count=\i] in {roll\_no/r, name/n, programme/p}{
    \node[hd] at (-5.45+\i*1.55-1.55,0.21) {\c};}
  \foreach \r/\a/\b/\c/\s in {1/IT-001/Ayesha/BSIT/sel, 2/IT-002/Bilal/BSIT/sel,
                              3/CS-007/Hira/BSCS/cell, 4/IT-014/Saad/BSIT/sel}{
    \node[\s] at (-5.45,-0.42*\r+0.21) {\a};
    \node[\s] at (-3.90,-0.42*\r+0.21) {\b};
    \node[\s] at (-2.35,-0.42*\r+0.21) {\c};}
  \node[nt, text width=4.6cm] at (-3.9,-2.25)
       {chooses \textbf{rows}; every column survives};

  % --- projection: columns
  \node[ttl=secondary!70!black] at (3.1,0.55)
       {$\pi_{\,\mathit{programme}}(\mathit{student})$};
  \foreach \c [count=\i] in {roll\_no, name, programme}{
    \node[hd] at (1.55+\i*1.55-1.55,0.21) {\c};}
  \foreach \r/\a/\b/\c in {1/IT-001/Ayesha/BSIT, 2/IT-002/Bilal/BSIT,
                           3/CS-007/Hira/BSCS, 4/IT-014/Saad/BSIT}{
    \node[cell] at (1.55,-0.42*\r+0.21) {\a};
    \node[cell] at (3.10,-0.42*\r+0.21) {\b};
    \node[prj]  at (4.65,-0.42*\r+0.21) {\c};}
  \draw[-{Stealth[length=2mm]}, secondary!70, line width=0.9pt]
       (5.55,-0.6) -- (6.35,-0.6);
  \node[hd]   at (7.2,-0.21) {programme};
  \node[prj]  at (7.2,-0.63) {BSIT};
  \node[prj]  at (7.2,-1.05) {BSCS};
  \node[nt, text width=4.8cm] at (3.1,-2.25)
       {chooses \textbf{columns}, then removes duplicates: four rows in, two
        out};
\end{tikzpicture}}%
{The two unary operators. Selection cuts horizontally and never loses a
column; projection cuts vertically and can lose rows, because three
\texttt{BSIT} values become one.}{selproj}

\begin{examplebox}[Both operators, and their SQL]
\begin{tabular}{@{}L{5.6cm}L{8.4cm}@{}}
\textbf{Algebra} & \textbf{SQL} \\
$\sigma_{\mathit{semester}=2}(\mathit{student})$ &
  \texttt{SELECT * FROM student WHERE semester = 2;} \\
$\pi_{\mathit{name},\,\mathit{programme}}(\mathit{student})$ &
  \texttt{SELECT DISTINCT name, programme FROM student;} \\
$\pi_{\mathit{name}}(\sigma_{\mathit{semester}=2}(\mathit{student}))$ &
  \texttt{SELECT DISTINCT name FROM student WHERE semester = 2;} \\
\end{tabular}

\smallskip
\textbf{Look at the word \sqlkw{DISTINCT}.} It is in all three projection rows
because algebra projection removes duplicates and SQL's does not. SQL returns
a \emph{bag} --- duplicates kept --- unless told otherwise, which is
\chref{relational}'s ``a table is not a relation'' showing up in the query
language. Without \sqlkw{DISTINCT}, the second query returns 4{,}000 rows; with
it, one per distinct name and programme pair.

\smallskip
\textbf{And notice the nesting in the third row.} $\sigma$ inside, $\pi$
outside, and SQL writes them in the opposite order from how they run: the
\sqlkw{SELECT} list is applied last, after \sqlkw{WHERE}. Reading SQL in
execution order --- \sqlkw{FROM}, then \sqlkw{WHERE}, then \sqlkw{SELECT} --- is
the single most useful habit this chapter can give you.
\end{examplebox}

\section{The Binary Operators}

\begin{definitionbox}[Product, union, difference, intersection]
The \term{Cartesian product} $R \times S$ pairs every tuple of $R$ with every
tuple of $S$. If $R$ has $m$ tuples and $S$ has $n$, the result has
$m \times n$ tuples and the attributes of both.

\smallskip
\term{Union} $R \cup S$, \term{difference} $R - S$ and \term{intersection}
$R \cap S$ are the set operations of Appendix~B, applied to tuples. All three
require $R$ and $S$ to be \term{union-compatible}: the same number of
attributes, with compatible domains, in the same order.

\smallskip
Intersection is not fundamental: $R \cap S = R - (R - S)$.
\end{definitionbox}

\begin{worked}[How big is a product, and why that matters]
On the seeded database that \texttt{code/schema/seed.py} builds,
\textit{student} has 4{,}000 tuples and \textit{enrolment} has 60{,}037.

\medskip
\textbf{Step 1 --- the size.} $|\mathit{student} \times \mathit{enrolment}|
= 4{,}000 \times 60{,}037 = 240{,}148{,}000$ tuples.

\medskip
\textbf{Step 2 --- how many of them mean anything.} A pair is meaningful only
when the student in it \emph{is} the student the enrolment is about. Each
enrolment has exactly one such student, so the meaningful pairs number exactly
60{,}037.

\medskip
\textbf{Step 3 --- the proportion.} $60{,}037 / 240{,}148{,}000 = 0.025\%$.
One pair in four thousand is true; the other 3{,}999 pair a student with an
enrolment that belongs to somebody else.

\medskip
\textbf{Step 4 --- the fix, in algebra.} Select the meaningful ones:
\[
  \sigma_{\mathit{student.roll\_no} = \mathit{enrolment.roll\_no}}
  (\mathit{student} \times \mathit{enrolment})
\]
That expression \emph{is} a join, and \chref{joins} gives it its own symbol.
Every join in SQL is a product followed by a selection --- in meaning. In
execution it is never computed that way, because 240 million intermediate
tuples to keep 60 thousand is exactly the plan an optimiser exists to avoid,
and \chref{joins} measures how it avoids it.
\end{worked}

\begin{alertbox}[Union-compatibility is checked by position, and that is a trap]
\begin{lstlisting}[language=sql]
SELECT roll_no, name FROM student
UNION
SELECT emp_id,  name FROM instructor;
\end{lstlisting}

This runs. Two attributes each side, both text: union-compatible. The result
is a list of people --- students and staff --- and that is what was meant.

\begin{lstlisting}[language=sql]
SELECT name,    roll_no FROM student
UNION
SELECT emp_id,  name    FROM instructor;
\end{lstlisting}

This \emph{also} runs, and produces a column that holds student names and
employee numbers, and a second column holding roll numbers and instructor
names. SQL checks that the count and types match \emph{by position}; it does
not check that the meanings do, because it cannot.

\smallskip
\textbf{The rule:} in a \sqlkw{UNION}, \sqlkw{EXCEPT} or \sqlkw{INTERSECT},
list the columns in the same order on both sides and give them the same
meaning. The database will hold you to the first half only.
\end{alertbox}

\section{Rename, and a Question Written Three Ways}

\begin{definitionbox}[Rename]
\term{Rename} $\rho_{S}(R)$ gives the relation $R$ the name $S$;
$\rho_{S(B_1,\dots,B_n)}(R)$ also renames its attributes.

\smallskip
It is needed whenever a relation appears twice in one expression --- the
recursive relationship of \chref{cardinality} --- because otherwise
$\mathit{course.course\_code}$ would name two different things. In SQL it is
the \emph{alias}: \texttt{FROM course AS c1, course AS c2}.
\end{definitionbox}

\begin{worked}[``Which students have taken a course but never received an A?'']
\textbf{Step 1 --- in English, the two sets.} Students who have taken
something, minus students who have an A somewhere.

\medskip
\textbf{Step 2 --- each set in algebra.}
\begin{align*}
  T &= \pi_{\mathit{roll\_no}}(\mathit{enrolment}) \\
  A &= \pi_{\mathit{roll\_no}}(\sigma_{\mathit{grade}=\mathrm{'A'}}
        (\mathit{enrolment}))
\end{align*}

\textbf{Step 3 --- combine.} $T - A$. Both are projections onto
\textit{roll\_no}, so they are union-compatible by construction.

\medskip
\textbf{Step 4 --- the SQL.}
\begin{lstlisting}[language=sql]
SELECT roll_no FROM enrolment
EXCEPT
SELECT roll_no FROM enrolment WHERE grade = 'A';
\end{lstlisting}

\medskip
\textbf{Step 5 --- check it against a wrong answer.} The tempting version is
\[
  \pi_{\mathit{roll\_no}}(\sigma_{\mathit{grade} \neq \mathrm{'A'}}
  (\mathit{enrolment}))
\]
--- ``students with a grade that is not A''. That includes every student who
has one B, even if they also have five A's. It answers a different question
and it runs without complaint. On the seeded database the right expression
returns \textbf{551} students; the tempting one returns \textbf{4{,}000} ---
every student in the university, because every one of them has at least one
grade that is not an A.

\smallskip
\textbf{The lesson.} ``Never'' is a difference, not a selection. Questions
with \emph{no}, \emph{never}, \emph{none} or \emph{only} in them almost
always need $-$ or its SQL equivalents, and \chref{subqueries} returns to
exactly this, with \sqlkw{NOT EXISTS} and a measurement.
\end{worked}

\begin{conceptbox}[Algebra is procedural; SQL is declarative --- and why both exist]
An algebra expression says \emph{how}: first select, then project, then take
the difference. Its nesting is an order of operations.

\smallskip
SQL says \emph{what}: the rows satisfying this condition, from these tables.
It deliberately does not say in what order.

\smallskip
The optimiser sits between them. It turns your SQL into one algebra
expression, then uses algebraic laws --- $\sigma$ before $\times$ rather than
after, $\pi$ as early as possible --- to rewrite it into a cheaper one that
computes the same relation. Those laws are true \emph{only because} every
operator is closed and every intermediate result is a relation.

\smallskip
So closure, the first definition in this chapter, is not a curiosity. It is
the reason a database can choose a plan at all, and \chref{indexes} is about
what happens when it chooses badly.
\end{conceptbox}

\begin{pitfall}
\begin{tight}
  \item \textbf{Forgetting \sqlkw{DISTINCT}.} Algebra $\pi$ removes
        duplicates; SQL's \sqlkw{SELECT} list does not.
  \item \textbf{Reading SQL in written order.} It runs \sqlkw{FROM},
        \sqlkw{WHERE}, \sqlkw{SELECT}. Read it that way.
  \item \textbf{Writing ``never'' as a selection.} $\sigma_{\neq}$ answers
        ``at least once something else''. ``Never'' is a difference.
  \item \textbf{Trusting \sqlkw{UNION} to check meaning.} It checks count and
        type by position. Order the columns yourself.
  \item \textbf{Forgetting the size of a product.} $m \times n$. Four thousand
        by sixty thousand is 240 million.
  \item \textbf{Using a relation twice without renaming it.} The attribute
        names become ambiguous. Use an alias.
\end{tight}
\end{pitfall}

\begin{lab}[Compute five algebra expressions, and catch one wrong answer]
\begin{lstlisting}[language=bash]
# --- 1. A populated database, the same rows as the handout ---------
python code/schema/seed.py university.db
sqlite3 university.db
\end{lstlisting}
\begin{lstlisting}[language=sql]
-- --- 2. Selection and projection -------------------------------
.headers on
.mode box
SELECT COUNT(*) FROM student WHERE semester = 2;
SELECT COUNT(*) FROM (SELECT programme FROM student);
SELECT COUNT(*) FROM (SELECT DISTINCT programme FROM student);
-- The last two are both "projection onto programme". Why do they
-- give different numbers, and which one is the algebra's answer?

-- --- 3. The size of a product, without computing it -------------
SELECT (SELECT COUNT(*) FROM student) * (SELECT COUNT(*) FROM enrolment);
-- Do NOT run SELECT COUNT(*) FROM student, enrolment on this
-- database unless you are prepared to wait. Say why.

-- --- 4. A product made meaningful: the selection that is a join -
SELECT COUNT(*) FROM student s, enrolment e
 WHERE s.roll_no = e.roll_no;
-- Compare with the size of enrolment. Explain the match.

-- --- 5. Union, and the positional trap --------------------------
SELECT roll_no AS id, name FROM student
UNION
SELECT emp_id, name FROM instructor
LIMIT 5;
SELECT name, roll_no FROM student
UNION
SELECT emp_id, name FROM instructor
LIMIT 5;
-- Both run. Describe precisely what is wrong with the second.

-- --- 6. "Never an A" -- the right way ---------------------------
SELECT COUNT(*) FROM (
  SELECT roll_no FROM enrolment
  EXCEPT
  SELECT roll_no FROM enrolment WHERE grade = 'A');

-- --- 7. "Never an A" -- the tempting way ------------------------
SELECT COUNT(DISTINCT roll_no) FROM enrolment WHERE grade <> 'A';
-- Write both counts down. One of them is every student in the
-- university. Which, and why? This is the worked example, run.

-- --- 8. Rename: the recursive relationship ----------------------
SELECT c1.title AS course, c2.title AS requires
  FROM prerequisite p
  JOIN course c1 ON c1.course_code = p.course_code
  JOIN course c2 ON c2.course_code = p.requires
 LIMIT 5;
-- course appears twice. Remove the aliases and record what fails.
\end{lstlisting}

\smallskip
\textbf{What to notice.} Steps 6 and 7 differ by one idea --- 551 against
4{,}000 on the seeded data --- and neither one raises an error. Every query you
write from here should be checked against a hand-worked expectation, because
a confident wrong answer looks identical to a confident right one.
\end{lab}

\begin{checkpoint}
\begin{tightnum}
  \item $\pi_{\mathit{programme}}(\mathit{student})$ on 4{,}000 students
        returns 4 tuples, and \texttt{SELECT programme FROM student} returns
        4{,}000 rows. Explain the difference in terms of the definition of a
        relation, and say which keyword reconciles them.
  \item Write, in algebra and in SQL, ``students in semester 2 who are
        enrolled in nothing in FALL2026''. Say which operator expresses
        ``nothing'', and give the tempting wrong version and what it returns.
  \item Every join is a selection over a product. Using the seeded sizes,
        explain why no database computes a join that way, and say which
        property of the algebra lets the optimiser compute something else
        instead and still be correct.
\end{tightnum}
\end{checkpoint}

\begin{chaptersummary}
\begin{tight}
  \item Relational algebra is a set of operators on relations whose every
        result is \textbf{a relation}. That is \textbf{closure}, and it is why
        queries nest.
  \item Six operators are fundamental: \textbf{selection} $\sigma$ (rows),
        \textbf{projection} $\pi$ (columns), \textbf{product} $\times$,
        \textbf{union} $\cup$, \textbf{difference} $-$, \textbf{rename}
        $\rho$. Intersection and every join are built from them.
  \item Algebra projection \textbf{removes duplicates}; SQL's does not unless
        you write \sqlkw{DISTINCT}. A table is a bag; a relation is a set.
  \item SQL runs \sqlkw{FROM}, then \sqlkw{WHERE}, then \sqlkw{SELECT} --- the
        opposite of the order it is written in.
  \item $|R \times S| = |R| \times |S|$. On the seeded data that is 240
        million tuples, of which 0.025\% mean anything. A join is a selection
        over a product \emph{in meaning}, never in execution.
  \item \sqlkw{UNION}, \sqlkw{EXCEPT} and \sqlkw{INTERSECT} check
        union-compatibility \textbf{by position and type}, never by meaning.
  \item ``Never'', ``none'' and ``only'' are \textbf{differences}, not
        selections. The selection version runs and answers a different
        question.
  \item \textbf{Rename} is the SQL alias, and it is required whenever a
        relation appears twice in one query.
  \item Algebra is procedural, SQL declarative; the optimiser turns one into
        the other and rewrites it using laws that hold only because of
        closure.
\end{tight}
\end{chaptersummary}

\begin{reviewq}
  \item \textbf{(MCQ)} Which operator can return fewer tuples than its input
        even when no condition is given? (a)~selection; (b)~projection;
        (c)~rename; (d)~Cartesian product.
  \item \textbf{(MCQ)} $|R| = 30$ and $|S| = 50$. $|R \times S|$ is:
        (a)~80; (b)~30; (c)~1{,}500; (d)~it depends on the keys.
  \item \textbf{(Short)} Define union-compatibility and give two relations
        from the university schema that are union-compatible but whose union
        would be meaningless.
  \item \textbf{(Short)} Show that intersection can be written using only
        difference.
  \item \textbf{(Short)} Why does SQL need \sqlkw{DISTINCT} when relational
        algebra does not?
  \item \textbf{(Applied)} Write in algebra and then in SQL: (a)~the names of
        instructors in the CS department; (b)~course codes that are a
        prerequisite for something but require nothing themselves;
        (c)~students who have received both an A and an F. For each, say which
        operators you used and why.
  \item \textbf{(Applied)} Draw the query tree for
        $\pi_{\mathit{name}}(\sigma_{\mathit{grade}=\mathrm{'A'}}
        (\mathit{student} \times \mathit{enrolment}))$ with the join condition
        added. Then draw an equivalent tree that is cheaper, say which
        algebraic law justifies each change, and estimate from the seeded
        table sizes how many intermediate tuples each tree produces.
\end{reviewq}

%=====================================================================
\chapter{Joins: the Algebra and the SQL}\label{ch:joins}
%=====================================================================
\locator{4}{Part IV --- Asking Questions}

\begin{outcomes}
By the end of this chapter you will be able to:
\begin{tight}
  \item \textbf{Define} theta join, equijoin and natural join in algebra, and
        \textbf{write} each in SQL.
  \item \textbf{Distinguish} inner from outer joins, and \textbf{predict}
        which rows an inner join will drop.
  \item \textbf{Explain} why \sqlkw{NATURAL JOIN} is dangerous, using a case
        from this course's own schema.
  \item \textbf{Read} an \sqlkw{EXPLAIN QUERY PLAN} for a join and
        \textbf{say} whether it is using an index.
  \item \textbf{Account for} a hundred-fold difference between two joins that
        return identical rows.
\end{tight}
\end{outcomes}

\begin{prereq}
\chref{algebra}: a join is a selection over a product. \chref{relational} for
nulls, whose consequence for joins this chapter measures. You need
\texttt{code/m13\_joins.py} and a seeded \texttt{university.db}.
\end{prereq}

\begin{keyterms}
join $\bullet$ theta join $\bullet$ equijoin $\bullet$ natural join ($\bowtie$)
$\bullet$ inner join $\bullet$ outer join $\bullet$ left outer join $\bullet$
right outer join $\bullet$ full outer join $\bullet$ self-join $\bullet$
\sqlkw{USING} $\bullet$ \sqlkw{ON} $\bullet$ nested-loop join $\bullet$
automatic index $\bullet$ \sqlkw{EXPLAIN QUERY PLAN} $\bullet$ \sqlkw{SCAN}
$\bullet$ \sqlkw{SEARCH}
\end{keyterms}

\section{One Operator, Three Strengths}

\chref{algebra} ended on the expression
$\sigma_{\mathit{s.roll\_no}=\mathit{e.roll\_no}}(\mathit{s} \times
\mathit{e})$. It occurs so often that it has its own symbol, and three
variants that differ only in how much the condition is allowed to say.

\begin{definitionbox}[Theta join, equijoin, natural join]
A \term{theta join} $R \bowtie_{\theta} S$ is $\sigma_{\theta}(R \times S)$
for any condition $\theta$ comparing attributes of $R$ with attributes of $S$
using $=, <, \le, >, \ge$ or $\neq$.

\smallskip
An \term{equijoin} is a theta join whose condition uses only $=$. Almost every
join you will write is one.

\smallskip
A \term{natural join} $R \bowtie S$ is an equijoin on \emph{every attribute
the two relations share by name}, keeping one copy of each. It has no
condition at all --- the condition is whatever the names happen to be.
\end{definitionbox}

\begin{examplebox}[The same join, four ways in SQL]
\begin{lstlisting}[language=sql]
-- 1. the product, then the selection: chapter 12 exactly
SELECT * FROM student s, enrolment e WHERE s.roll_no = e.roll_no;

-- 2. explicit, and the form this course uses
SELECT * FROM student s JOIN enrolment e ON s.roll_no = e.roll_no;

-- 3. when the column has the same name on both sides
SELECT * FROM student JOIN enrolment USING (roll_no);

-- 4. natural: no condition written at all
SELECT * FROM student NATURAL JOIN enrolment;
\end{lstlisting}

\smallskip
The first three are one query to the optimiser --- the measurement below shows
the same plan for each. The fourth is not, on this schema, and the reason is
the most useful thing in this chapter.
\end{examplebox}

\section{What a Join Drops}

\begin{definitionbox}[Inner and outer joins]
An \term{inner join} returns only the pairs that satisfy the condition. A row
with no partner on the other side does not appear.

\smallskip
A \term{left outer join} of $R$ with $S$ returns every row of $R$; where
there is no partner in $S$, the $S$ attributes are null. A \term{right outer
join} keeps every row of $S$; a \term{full outer join} keeps every row of
both.

\smallskip
SQLite supports all three. \sqlkw{LEFT JOIN} is the one you will need, and the
question that decides it is always the same: \emph{must a row with no partner
still appear?}
\end{definitionbox}

\dsfig{%
\begin{tikzpicture}[
  cell/.style={draw=primary!35, line width=0.5pt, minimum width=1.5cm,
               minimum height=0.40cm, font=\scriptsize\ttfamily,
               text=primary!65!black, inner sep=0pt},
  hd/.style={cell, fill=primary, text=white, font=\scriptsize\bfseries},
  gone/.style={cell, fill=accent!12, text=accent},
  nul/.style={cell, text=neutral},
  ttl/.style={font=\scriptsize\bfseries, text=#1, anchor=south},
  nt/.style={font=\scriptsize\itshape, text=neutral, align=center}
]
  % section (left table)
  \node[ttl=primary] at (-5.4,0.5) {section};
  \foreach \c [count=\i] in {course, emp\_id}{
    \node[hd] at (-6.95+\i*1.5,0.2) {\c};}
  \foreach \r/\a/\b/\s in {1/CS-103/E-114/cell, 2/CS-201/E-120/cell,
                           3/CS-305/NULL/gone}{
    \node[\s] at (-5.45,-0.4*\r+0.2) {\a};
    \node[\s] at (-3.95,-0.4*\r+0.2) {\b};}

  % inner join result
  \node[ttl=greenhl!60!black] at (0.75,0.5) {\sqlkw{JOIN} instructor};
  \foreach \c [count=\i] in {course, emp\_id, name}{
    \node[hd] at (-1.5+\i*1.5,0.2) {\c};}
  \foreach \r/\a/\b/\c in {1/CS-103/E-114/Malik, 2/CS-201/E-120/Raza}{
    \node[cell] at (0.0,-0.4*\r+0.2) {\a};
    \node[cell] at (1.5,-0.4*\r+0.2) {\b};
    \node[cell] at (3.0,-0.4*\r+0.2) {\c};}
  \node[nt, text=accent, text width=4.4cm] at (1.5,-1.35)
       {CS-305 is \textbf{gone}: its emp\_id is null and null equals
        nothing};

  % left join result
  \node[ttl=secondary!70!black] at (6.95,0.5) {\sqlkw{LEFT JOIN} instructor};
  \foreach \c [count=\i] in {course, emp\_id, name}{
    \node[hd] at (4.7+\i*1.5,0.2) {\c};}
  \foreach \r/\a/\b/\c/\s in {1/CS-103/E-114/Malik/cell,
                              2/CS-201/E-120/Raza/cell,
                              3/CS-305/NULL/NULL/nul}{
    \node[\s] at (6.2,-0.4*\r+0.2) {\a};
    \node[\s] at (7.7,-0.4*\r+0.2) {\b};
    \node[\s] at (9.2,-0.4*\r+0.2) {\c};}
  \node[nt, text width=4.4cm] at (7.7,-1.55)
       {every section survives; the missing partner is filled with null};
\end{tikzpicture}}%
{The section nobody has been assigned to teach. An inner join drops it without
comment; a left join keeps it and marks what is missing. On the seeded
database this is 73 sections out of 1{,}500.}{innerleft}

\section{Measured: Plans, Nulls and the Cost of a Missing Index}

\begin{measured}[Three joins that return the same rows and one that returns none]
\textsf{\footnotesize \texttt{code/m13\_joins.py}, seeded database (4{,}000
students, 60{,}037 enrolments), SQLite 3.51.2, Intel i5-6600, Windows 11.}

\medskip
\textbf{1. Four spellings of one inner join}

\begin{tabular}{@{}L{4.6cm}R{1.9cm}R{1.7cm}L{5.6cm}@{}}
\textbf{written as} & \textbf{rows} & \textbf{ms} & \textbf{plan} \\
comma + \sqlkw{WHERE}  & 60{,}037 & 29.1 & \texttt{SCAN e} + \texttt{SEARCH s} \\
\sqlkw{JOIN} \dots\ \sqlkw{ON} & 60{,}037 & 30.4 & identical \\
\sqlkw{JOIN} \dots\ \sqlkw{USING} & 60{,}037 & 30.1 & identical \\
\sqlkw{NATURAL JOIN} & \textbf{0} & 43.7 & \\
\end{tabular}

\smallskip
The two tables share \textbf{two} attribute names, not one:
\texttt{roll\_no}, and \texttt{semester} --- which is an \texttt{INTEGER} year
of study (e.g.\ \texttt{8}) in \textit{student} and a \texttt{TEXT} term label
(e.g.\ \texttt{'FALL2025'}) in \textit{enrolment}. \sqlkw{NATURAL JOIN} matched
on both, demanded a year equal a term label, and returned nothing. \textbf{No
error, on the schema exactly as \chref{mapping} built it.}

\medskip
\textbf{2. Inner against left join}

\begin{tabular}{@{}L{9.4cm}R{2.4cm}@{}}
sections in the database & 1{,}500 \\
sections with no instructor yet (\texttt{emp\_id} null) & 73 \\
timetable as an inner join & 1{,}427 \\
timetable as a left join & 1{,}500 \\
\end{tabular}

\smallskip
The 73 missing from the inner join are exactly the sections nobody has been
assigned to --- the rows a timetable most needs to show.

\medskip
\textbf{3. Students who share a name: a join on an unindexed column}

\begin{tabular}{@{}R{1.6cm}R{2.0cm}R{1.9cm}R{2.4cm}R{2.0cm}R{1.8cm}@{}}
\textbf{students} & \textbf{pairs found} & \textbf{as written}
& \textbf{automatic index off} & \textbf{real index} & \textbf{slowdown} \\
1{,}000 & 1{,}785  & 3.2\,ms  & 235.1\,ms  & 2.9\,ms  & 82$\times$ \\
2{,}000 & 7{,}109  & 8.6\,ms  & 1{,}088.1\,ms & 11.1\,ms & 98$\times$ \\
4{,}000 & 28{,}517 & 25.1\,ms & 4{,}170.1\,ms & 38.7\,ms & \textbf{108$\times$} \\
\end{tabular}

\smallskip
Without an index the join is a \term{nested-loop join}: for each row of
one side, every row of the other is examined --- two loops, one inside the
other, which is where the $n^2$ comes from.

\smallskip
\textbf{Compare across a row, not down a column.} The pairs found grow about
four-fold per doubling for \emph{every} method, because names come from a
fixed pool and the number of real matches is itself quadratic --- so growth
ratios here mix up how much the database \emph{looked at} with how much it
\emph{returned}. Across a row the output is identical, and the difference is
pure search: without an index, all $4000 \times 3999 / 2 \approx 8$ million
pairs are examined to find 28 thousand. The slowdown itself rises with size,
82 to 108, which is the real signature of a quadratic method.

\medskip
\textbf{4. Why ``as written'' is fast anyway}

\begin{tabular}{@{}L{4.8cm}L{9.2cm}@{}}
automatic index on (default) & \texttt{SEARCH b USING AUTOMATIC COVERING INDEX (name=?)} \\
automatic index off & \texttt{SEARCH b USING INDEX \dots (roll\_no>?)} \\
\texttt{CREATE INDEX} on name & \texttt{SEARCH b USING INDEX idx\_name (name=?)} \\
\end{tabular}

\smallskip
SQLite noticed the missing index, \textbf{built one for this single query,
used it, and threw it away}. You pay to build an index on every run and never
get to keep it.
\end{measured}

\begin{alertbox}[What the measurement means for the queries you write]
\begin{tight}
  \item \textbf{Never use \sqlkw{NATURAL JOIN}.} Its condition is decided by
        column names, and column names change for reasons that have nothing to
        do with the join. This course's own schema has the trap in it, and it
        was not put there on purpose --- a year of study and a term are both
        reasonably called ``semester''. Write \sqlkw{ON} or \sqlkw{USING} and
        say what you mean.
  \item \textbf{Ask, for every join, whether an unmatched row must still
        appear.} If yes, it is a \sqlkw{LEFT JOIN}. The inner join's silence
        about the 73 is not a bug in SQL; it is the definition, applied to a
        null foreign key that \chref{cardinality} deliberately allowed.
  \item \textbf{Read the plan.} \sqlkw{AUTOMATIC INDEX} is the optimiser
        telling you an index is missing, in plain words. It is the cheapest
        performance diagnosis there is, and \chref{indexes} is about acting on
        it.
\end{tight}
\end{alertbox}

\section{Joining a Table to Itself}

\begin{examplebox}[The self-join, and the condition that halves it]
``Pairs of students who share a name'' joins \textit{student} to itself, which
needs two aliases --- \chref{algebra}'s rename:

\begin{lstlisting}[language=sql]
SELECT a.roll_no, b.roll_no, a.name
  FROM student a
  JOIN student b ON a.name = b.name
                AND a.roll_no < b.roll_no;
\end{lstlisting}

\smallskip
\textbf{The second condition does two jobs.} Without it every student matches
themselves --- Ayesha shares a name with Ayesha --- and every real pair appears
twice, once each way round. \texttt{a.roll\_no < b.roll\_no} removes both: no
row is paired with itself, and each pair appears in exactly one order. Use
$<$, never $\neq$, which would keep both orders.

\smallskip
The prerequisite query in \chref{algebra}'s laboratory was a self-join too, of
\textit{course}, through \textit{prerequisite}. Every recursive relationship in
an ER diagram becomes a self-join in SQL.
\end{examplebox}

\begin{pitfall}
\begin{tight}
  \item \textbf{\sqlkw{NATURAL JOIN}.} It joined on two columns, not one, and
        returned nothing. Measured.
  \item \textbf{An inner join where a left join was meant.} The unmatched rows
        vanish silently --- 73 sections from a timetable.
  \item \textbf{A condition on the right table in \sqlkw{WHERE} after a
        \sqlkw{LEFT JOIN}.} \texttt{WHERE i.name = 'Malik'} discards the
        null-extended rows and turns it back into an inner join. Put such
        conditions in \sqlkw{ON}.
  \item \textbf{Forgetting the join condition.} You get the product: 240
        million rows on the seeded data.
  \item \textbf{A self-join with $\neq$.} Every pair twice. Use $<$.
  \item \textbf{Ignoring \sqlkw{AUTOMATIC INDEX} in a plan.} It is a missing
        index announcing itself.
\end{tight}
\end{pitfall}

\begin{lab}[Lose 73 sections, find them, and make a join a hundred times faster]
\begin{lstlisting}[language=bash]
# --- 1. The seeded database ----------------------------------------
python code/schema/seed.py university.db
sqlite3 university.db
\end{lstlisting}
\begin{lstlisting}[language=sql]
-- --- 2. The timetable, as an inner join -------------------------
.headers on
.mode box
SELECT COUNT(*) FROM section s JOIN instructor i ON s.emp_id = i.emp_id;
SELECT COUNT(*) FROM section;
-- The two numbers differ. By how much, and which sections are they?

-- --- 3. Find them -----------------------------------------------
SELECT s.course_code, s.section_no, s.semester
  FROM section s LEFT JOIN instructor i ON s.emp_id = i.emp_id
 WHERE i.emp_id IS NULL
 LIMIT 10;
-- This is the anti-join idiom: LEFT JOIN, then keep the rows that
-- found no partner. Learn it; Chapter 15 measures it.

-- --- 4. Break the left join by mistake --------------------------
SELECT COUNT(*) FROM section s LEFT JOIN instructor i
    ON s.emp_id = i.emp_id
 WHERE i.designation = 'Professor';
SELECT COUNT(*) FROM section s LEFT JOIN instructor i
    ON s.emp_id = i.emp_id AND i.designation = 'Professor';
-- Same condition, two places, two very different counts. Explain.

-- --- 5. The natural join, on the schema you built ---------------
SELECT COUNT(*) FROM student NATURAL JOIN enrolment;
-- Zero. Find out why:
PRAGMA table_info(student);
PRAGMA table_info(enrolment);
-- Name the second shared column, and its two meanings.

-- --- 6. Read a plan ---------------------------------------------
EXPLAIN QUERY PLAN
SELECT COUNT(*) FROM student a JOIN student b
    ON a.name = b.name AND a.roll_no < b.roll_no;
-- Find the word AUTOMATIC. Say what SQLite just did.

-- --- 7. Give it the index it keeps building ---------------------
CREATE INDEX idx_student_name ON student(name);
EXPLAIN QUERY PLAN
SELECT COUNT(*) FROM student a JOIN student b
    ON a.name = b.name AND a.roll_no < b.roll_no;
\end{lstlisting}
\begin{lstlisting}[language=bash]
# --- 8. The measured version ----------------------------------------
python code/m13_joins.py
\end{lstlisting}

\smallskip
\textbf{What to notice.} Step 4 is the one that catches people in examinations
and in production alike. A condition in \sqlkw{WHERE} is applied \emph{after}
the outer join has produced its null-extended rows, and since a null
\texttt{designation} is not \texttt{'Professor'}, they are all discarded ---
the outer join has silently become an inner one.
\end{lab}

\begin{checkpoint}
\begin{tightnum}
  \item \sqlkw{NATURAL JOIN} of \textit{student} and \textit{enrolment}
        returns 0 rows on this course's schema. Explain exactly why, say what
        you would change in the query, and say what you would change in the
        schema if you could change only one column name.
  \item An inner join of \textit{section} to \textit{instructor} returns 1{,}427
        rows and the table has 1{,}500. Identify the missing 73 by a property
        they share, write the query that lists them, and say which design
        decision in \chref{cardinality} made them possible.
  \item Two joins return the same 28{,}517 rows, one in 38.7\,ms and one in
        4{,}170\,ms. Explain the difference using the number of pairs each one
        examined, and say why comparing growth ratios down the columns would
        have given a misleading answer.
\end{tightnum}
\end{checkpoint}

\begin{chaptersummary}
\begin{tight}
  \item A \textbf{theta join} is $\sigma_\theta(R \times S)$; an
        \textbf{equijoin} uses only $=$; a \textbf{natural join} equijoins on
        \emph{every} shared attribute name.
  \item Comma-plus-\sqlkw{WHERE}, \sqlkw{JOIN \dots ON} and \sqlkw{JOIN \dots
        USING} produce the same plan. Use \sqlkw{ON}.
  \item \textbf{Measured: \sqlkw{NATURAL JOIN} returned 0 rows} on this
        course's schema, because \textit{student} and \textit{enrolment} both
        have a \texttt{semester} --- a year of study in one, a term label in
        the other. Never use it.
  \item An \textbf{inner join} drops rows with no partner; a \textbf{left
        join} keeps them with nulls. \textbf{Measured: 73 of 1{,}500 sections}
        vanished from the inner-join timetable --- the ones with no instructor.
  \item A condition on the right-hand table belongs in \sqlkw{ON}. In
        \sqlkw{WHERE} it turns a left join back into an inner one.
  \item A \textbf{self-join} needs two aliases, and $<$ rather than $\neq$ to
        avoid both self-pairs and duplicate pairs.
  \item \textbf{Measured: a join on an unindexed column was 108$\times$ slower}
        at 4{,}000 students than the same join with an index --- 8 million
        pairs examined against 28 thousand found.
  \item Compare methods \textbf{at the same size}. Here the output itself grew
        quadratically, so growth ratios mixed search cost with output size.
  \item SQLite builds an \textbf{automatic index} for a single query when one
        is missing, and discards it. \sqlkw{AUTOMATIC} in a plan is a missing
        index announcing itself.
\end{tight}
\end{chaptersummary}

\begin{reviewq}
  \item \textbf{(MCQ)} A natural join of $R$ and $S$ is an equijoin on:
        (a)~their primary keys; (b)~their foreign keys; (c)~every attribute
        they share by name; (d)~the first column of each.
  \item \textbf{(MCQ)} A section whose \texttt{emp\_id} is null, joined to
        \textit{instructor} with an inner join, appears: (a)~once, with nulls;
        (b)~not at all; (c)~once per instructor; (d)~with an error.
  \item \textbf{(Short)} Why does \texttt{a.roll\_no < b.roll\_no} in a
        self-join give each pair exactly once, while $\neq$ gives it twice?
  \item \textbf{(Short)} What does \sqlkw{AUTOMATIC COVERING INDEX} in a query
        plan tell you, and what should you do about it?
  \item \textbf{(Short)} Give a query in the university schema that needs a
        \sqlkw{LEFT JOIN}, and say what an inner join would silently omit.
  \item \textbf{(Applied)} Write the query that lists every course with its
        number of sections in FALL2026, including courses with none. (On the
        seeded database every course has one, so first insert a new course
        with no sections, or there is nothing to omit.) Show the inner-join
        version that gets it wrong, say which courses it omits, and then give
        the correct version and explain where its condition on semester must
        go and why.
  \item \textbf{(Applied)} The measured slowdown was 82$\times$, 98$\times$ and
        108$\times$ at 1{,}000, 2{,}000 and 4{,}000 students. Predict it at
        8{,}000, show your reasoning from the number of pairs examined and
        found, and explain why the slowdown is rising rather than constant.
\end{reviewq}

%=====================================================================
\chapter{Selecting Rows, and the Trouble with NULL}\label{ch:select}
%=====================================================================
\locator{4}{Part IV --- Asking Questions}

\begin{outcomes}
By the end of this chapter you will be able to:
\begin{tight}
  \item \textbf{Write} a \sqlkw{SELECT} with \sqlkw{WHERE}, \sqlkw{ORDER BY}
        and \sqlkw{LIMIT}, and \textbf{state} the order in which its clauses
        run.
  \item \textbf{Use} \sqlkw{LIKE}, \sqlkw{IN}, \sqlkw{BETWEEN} and
        \sqlkw{CASE} correctly, including at their boundaries.
  \item \textbf{Evaluate} a condition in three-valued logic, and
        \textbf{predict} which rows a condition involving null will return.
  \item \textbf{Explain} why \sqlkw{NOT IN} can return nothing, and
        \textbf{replace} it with a form that cannot.
  \item \textbf{Account for} every row of a table when a condition and its
        negation together do not cover it.
\end{tight}
\end{outcomes}

\begin{prereq}
\chref{relational}: what a null is. \chref{joins} for the 73 sections a null
foreign key dropped. A seeded \texttt{university.db}; every number in this
chapter is reproduced by its laboratory.
\end{prereq}

\begin{keyterms}
\sqlkw{SELECT} $\bullet$ \sqlkw{WHERE} $\bullet$ \sqlkw{ORDER BY} $\bullet$
\sqlkw{LIMIT} $\bullet$ \sqlkw{OFFSET} $\bullet$
\sqlkw{LIKE} $\bullet$ wildcard $\bullet$ \sqlkw{IN} $\bullet$
\sqlkw{BETWEEN} $\bullet$ \sqlkw{CASE} $\bullet$ three-valued logic $\bullet$
\unk{} $\bullet$ \sqlkw{IS NULL} $\bullet$ \sqlkw{COALESCE}
$\bullet$ \sqlkw{NULLIF} $\bullet$ \sqlkw{NOT IN} $\bullet$ \sqlkw{NOT EXISTS}
\end{keyterms}

\section{The Statement, and the Order It Runs In}

\begin{definitionbox}[\sqlkw{SELECT}, in execution order]
\begin{lstlisting}[language=sql]
SELECT   name, semester          -- 4. keep these columns
FROM     student                 -- 1. start from this relation
WHERE    programme = 'BSIT'      -- 2. keep these rows
ORDER BY semester DESC, name     -- 5. sort what is left
LIMIT    10 OFFSET 20;           -- 6. then take a slice
\end{lstlisting}

\smallskip
Step 3, between \sqlkw{WHERE} and \sqlkw{SELECT}, is \sqlkw{GROUP BY}, and
\chref{grouping} adds it. The order matters in one practical way: a name given
in the \sqlkw{SELECT} list with \sqlkw{AS} \emph{does not exist yet} when
\sqlkw{WHERE} runs, so it cannot be used there --- but it can be used in
\sqlkw{ORDER BY}, which runs after.
\end{definitionbox}

\begin{examplebox}[The four predicates, and where each one surprises]
\begin{tabular}{@{}L{4.6cm}L{9.4cm}@{}}
\texttt{name LIKE 'A\%'} & the \term{wildcards}: \texttt{\%} is any run of characters,
  \texttt{\_} exactly one. In SQLite \sqlkw{LIKE} is \textbf{case-insensitive}
  for ASCII letters, so this also matches \texttt{'ayesha'}. \\
\texttt{programme IN ('BSIT','BSCS')} & shorthand for two \sqlkw{OR}s. Fine
  --- until the list contains a null, below. \\
\texttt{semester BETWEEN 2 AND 4} & \textbf{inclusive at both ends}: 2, 3 and
  4. \\
\texttt{CASE WHEN g='A' THEN 4 \dots\ END} & evaluated top to bottom; the
  first true branch wins; with no \sqlkw{ELSE}, the result is null. \\
\end{tabular}

\smallskip
\textbf{And one that is not a predicate at all.} \texttt{LIMIT 10} without
\sqlkw{ORDER BY} returns \emph{some} ten rows. Which ten depends on the plan,
and \chref{relational} said the plan is allowed to change. A ``top ten'' with
no \sqlkw{ORDER BY} is ten rows chosen by the optimiser.
\end{examplebox}

\section{Three-Valued Logic}

Every condition so far has been true or false. As soon as a null is involved
there is a third answer, and it is the source of most quietly wrong results in
SQL.

\begin{definitionbox}[\unk{}]
Any comparison with a null --- \texttt{=}, \texttt{<>}, \texttt{<},
\texttt{>}, even \texttt{NULL = NULL} --- evaluates to \unk{}, not to
true or false. It means ``the value is missing, so I cannot say''.

\smallskip
\sqlkw{WHERE} keeps a row \textbf{only if its condition is true}. False and
\unk{} are both discarded.

\smallskip
\sqlkw{IS NULL} and \sqlkw{IS NOT NULL} are the only tests that return true
or false for a null. They exist because \texttt{= NULL} cannot work.
\end{definitionbox}

\begin{conceptbox}[The truth tables, and the one row to remember]
\begin{tabular}{@{}L{1.9cm}L{1.9cm}|L{1.9cm}L{1.9cm}|L{2.2cm}L{2.4cm}@{}}
\textbf{\sqlkw{AND}} & & \textbf{\sqlkw{OR}} & & \textbf{\sqlkw{NOT}} & \\
T \sqlkw{AND} U & = U & T \sqlkw{OR} U & = \textbf{T} & \sqlkw{NOT} T & = F \\
F \sqlkw{AND} U & = \textbf{F} & F \sqlkw{OR} U & = U & \sqlkw{NOT} F & = T \\
U \sqlkw{AND} U & = U & U \sqlkw{OR} U & = U & \sqlkw{NOT} U & = \textbf{U} \\
\end{tabular}

\smallskip
Read the bold entries. \sqlkw{AND} with a false is false whatever the other
side; \sqlkw{OR} with a true is true whatever the other side --- those are the
cases where a missing value genuinely does not matter.

\smallskip
\textbf{The row to remember is the last one: \sqlkw{NOT} \unk{} is
\unk{}.} Negating a condition does \emph{not} give you the other
rows. It gives you the other rows \emph{that are not null}. A condition and its
negation together do not cover a table that contains nulls --- and the next box
counts exactly how many rows fall through the gap.
\end{conceptbox}

\begin{worked}[Grade A, not grade A, and the 19{,}992 rows that are neither]
The seeded \textit{enrolment} table has 60{,}037 rows. Three queries:

\medskip
\begin{tabular}{@{}L{7.2cm}R{2.6cm}@{}}
\texttt{WHERE grade = 'A'}       & 7{,}268 \\
\texttt{WHERE grade <> 'A'}      & 32{,}777 \\
\texttt{WHERE NOT (grade = 'A')} & 32{,}777 \\
\end{tabular}

\medskip
\textbf{Step 1 --- the sum.} $7{,}268 + 32{,}777 = 40{,}045$. The table has
60{,}037. \textbf{19{,}992 rows are in neither answer.}

\medskip
\textbf{Step 2 --- which rows.} The FALL2026 enrolments, whose semester is not
over: \texttt{grade} is null. For those rows \texttt{grade = 'A'} is
\unk{}, so is \texttt{grade <> 'A'}, and so --- by the last row of
the truth table --- is \sqlkw{NOT} \texttt{(grade = 'A')}. Every one of the
three conditions discards them.

\medskip
\textbf{Step 3 --- the check.}
\begin{lstlisting}[language=sql]
SELECT COUNT(*) FROM enrolment WHERE grade IS NULL;    -- 19992
\end{lstlisting}
$7{,}268 + 32{,}777 + 19{,}992 = 60{,}037$. Every row accounted for.

\medskip
\textbf{Step 4 --- what goes wrong in practice.} A report of ``students who
did not get an A'' built on \texttt{grade <> 'A'} omits a third of all
enrolments without a word. Whether it \emph{should} include them is a question
for the examinations office --- an ungraded course is not a failure --- but
the query decided it silently, and decided it the way most people would not.

\medskip
\textbf{The rule.} When a column can be null, write the null case
explicitly, whichever way you want it:
\begin{lstlisting}[language=sql]
WHERE grade <> 'A' OR grade IS NULL     -- include ungraded
WHERE grade <> 'A'                      -- exclude them, and say so
\end{lstlisting}
\end{worked}

\dsfig{%
\begin{tikzpicture}[
  seg/.style={draw=white, line width=1pt, minimum height=0.9cm,
              font=\scriptsize\bfseries, text=white, inner sep=0pt},
  br/.style={draw=#1, line width=0.9pt,
             decorate, decoration={brace, amplitude=5pt}},
  lb/.style={font=\scriptsize, align=center, text=#1}
]
  % 60,037 rows drawn to scale across 13 cm. The widths are computed
  % here, not with pgfmath: the first draft used macros named a, b and c
  % -- which shadow TeX's own accent commands -- and a 13/60037 scale
  % factor, and three 'Dimension too large' errors followed.
  %   13 x 7268/60037 = 1.5738    13 x 32777/60037 = 7.0973
  %   13 x 19992/60037 = 4.3289
  \def\wA{1.5738}\def\wN{7.0973}\def\wU{4.3289}
  \node[seg, fill=greenhl, minimum width=\wA cm, anchor=west]
       (A) at (0,0) {A};
  \node[seg, fill=secondary, minimum width=\wN cm, anchor=west]
       (N) at (\wA,0) {B, C, D, F};
  \node[seg, fill=neutral!70, minimum width=\wU cm, anchor=west]
       (U) at (\wA+\wN,0) {NULL --- not yet graded};

  \draw[br=greenhl] (0,0.55) -- (\wA,0.55);
  \node[lb=greenhl!60!black] at (\wA/2,1.05)
       {\texttt{grade = 'A'}\\7{,}268};
  \draw[br=secondary] (\wA,0.55) -- (\wA+\wN,0.55);
  \node[lb=secondary!70!black] at (\wA+\wN/2,1.05)
       {\texttt{grade <> 'A'} \ and \ \texttt{NOT (grade = 'A')}\\32{,}777};
  \draw[br=accent] (\wA+\wN+\wU,-0.55) -- (\wA+\wN,-0.55);
  \node[lb=accent] at (\wA+\wN+\wU/2,-1.15)
       {\textbf{caught by neither}\\19{,}992 --- \sqlkw{IS NULL} only};
  \node[lb=neutral, anchor=east] at (-0.15,0) {60{,}037\\rows};
\end{tikzpicture}}%
{The seeded \textit{enrolment} table, drawn to scale. A condition and its
negation cover the green and the blue. The grey third is \unk{} to
both, and \sqlkw{WHERE} discards it either way.}{nullgap}

\section{\sqlkw{NOT IN}, Which Can Return Nothing}

\begin{alertbox}[Three new instructors, and a query that cannot find them]
Hire three instructors who are not yet teaching anything. Then ask which
instructors teach no section --- the obvious way:

\begin{lstlisting}[language=sql]
SELECT COUNT(*) FROM instructor
 WHERE emp_id NOT IN (SELECT emp_id FROM section);
\end{lstlisting}

On the seeded database, with the three added, this returns \textbf{0}.

\begin{lstlisting}[language=sql]
SELECT COUNT(*) FROM instructor i
 WHERE NOT EXISTS (SELECT 1 FROM section s WHERE s.emp_id = i.emp_id);
\end{lstlisting}

This returns \textbf{3}.

\medskip
\textbf{Why.} \texttt{x NOT IN (a, b, c)} means \texttt{x <> a AND x <> b AND
x <> c}. The sub-query's list contains the 73 nulls from the unassigned
sections, so the expansion contains \texttt{x <> NULL} --- which is
\unk{}. By the truth table, \texttt{T AND U = U}. For every
instructor, including the three new ones, the whole condition is at best
\unk{}, and \sqlkw{WHERE} discards \unk{}. \textbf{One null
in the list empties the result.}

\medskip
\textbf{The fix is either of these:}
\begin{tight}
  \item use \sqlkw{NOT EXISTS}, which tests for the \emph{existence} of a row
        rather than comparing with a value, and is unaffected by nulls; or
  \item filter the nulls out of the list:
        \texttt{\dots\ NOT IN (SELECT emp\_id FROM section WHERE emp\_id IS NOT
        NULL)}, which also returns 3.
\end{tight}

\smallskip
\textbf{This course uses \sqlkw{NOT EXISTS}.} The second fix works, but it has
to be remembered every time, and it will be forgotten the first time the
column becomes nullable --- which is a schema change somebody else will make.
\chref{subqueries} measures the two forms against each other.
\end{alertbox}

\begin{examplebox}[\sqlkw{COALESCE} and \sqlkw{NULLIF}: deciding what null means]
\sqlkw{COALESCE}(a, b, c) returns the first argument that is not null. It is
how you \emph{decide}, in one place, what a null should be treated as:

\begin{lstlisting}[language=sql]
SELECT roll_no, COALESCE(cnic, 'not on file') AS cnic FROM student;
SELECT COALESCE(grade, 'in progress') AS grade FROM enrolment;
\end{lstlisting}

\sqlkw{NULLIF}(a, b) is the reverse: it returns null if $a = b$, otherwise
$a$. Its main use is cleaning up \chref{relational}'s magic values ---
\texttt{NULLIF(cnic, 'N/A')} turns an invented marker back into the null it
should have been.

\smallskip
\textbf{Use \sqlkw{COALESCE} for display, not for arithmetic.} Writing
\texttt{AVG(COALESCE(score, 0))} decides that a missing mark is a zero, which
is \chref{ddl}'s 20\% error made on purpose.
\end{examplebox}

\begin{conceptbox}[Where nulls go in \sqlkw{ORDER BY}]
SQLite sorts nulls \textbf{first} in ascending order and last in descending,
treating null as smaller than every value. Other systems differ: PostgreSQL
puts them last in ascending order. SQLite since 3.30 accepts the standard
clause that makes it explicit:

\begin{lstlisting}[language=sql]
SELECT roll_no, grade FROM enrolment ORDER BY grade NULLS LAST;
\end{lstlisting}

Write \sqlkw{NULLS FIRST} or \sqlkw{NULLS LAST} whenever the column can be
null and the order matters, because ``what this database happens to do'' is
not the same on the next one.
\end{conceptbox}

\begin{pitfall}
\begin{tight}
  \item \textbf{\texttt{WHERE x = NULL}.} Always \unk{}, so always
        empty. Use \sqlkw{IS NULL}.
  \item \textbf{Assuming a condition and its negation cover the table.} They
        miss the nulls --- 19{,}992 of 60{,}037 enrolments.
  \item \textbf{\sqlkw{NOT IN} against a nullable column.} One null in the
        list and the result is empty. Use \sqlkw{NOT EXISTS}.
  \item \textbf{Using a column alias in \sqlkw{WHERE}.} It does not exist yet.
  \item \textbf{\sqlkw{LIMIT} without \sqlkw{ORDER BY}.} Some rows, chosen by
        the plan.
  \item \textbf{Forgetting \sqlkw{BETWEEN} is inclusive.} Both ends.
  \item \textbf{Expecting \sqlkw{LIKE} to be case-sensitive.} In SQLite it is
        not, for ASCII.
  \item \textbf{\sqlkw{COALESCE} to zero inside an average.} A decision that a
        missing mark is a fail, made silently.
\end{tight}
\end{pitfall}

\begin{lab}[Lose 19{,}992 rows and three instructors, and find them all]
\begin{lstlisting}[language=bash]
# --- 1. A fresh seeded database --------------------------------------
python code/schema/seed.py university.db
sqlite3 university.db
\end{lstlisting}
\begin{lstlisting}[language=sql]
-- --- 2. A condition and its negation ----------------------------
.headers on
.mode box
SELECT COUNT(*) FROM enrolment WHERE grade = 'A';
SELECT COUNT(*) FROM enrolment WHERE grade <> 'A';
SELECT COUNT(*) FROM enrolment;
-- Add the first two. Compare with the third. Explain the gap.

-- --- 3. Account for every row -----------------------------------
SELECT COUNT(*) FROM enrolment WHERE grade IS NULL;
SELECT COUNT(*) FROM enrolment WHERE grade = NULL;
-- Two queries that look alike. One finds the gap; one finds nothing.

-- --- 4. Hire three people who teach nothing yet -----------------
PRAGMA foreign_keys = ON;
INSERT INTO instructor VALUES
  ('E-9001','New Hire One','Lecturer','CS'),
  ('E-9002','New Hire Two','Lecturer','IT'),
  ('E-9003','New Hire Three','Lecturer','SE');

-- --- 5. Look for them with NOT IN -------------------------------
SELECT COUNT(*) FROM instructor
 WHERE emp_id NOT IN (SELECT emp_id FROM section);
-- There are three. How many did it find?

-- --- 6. Look for them with NOT EXISTS ---------------------------
SELECT name FROM instructor i
 WHERE NOT EXISTS (SELECT 1 FROM section s WHERE s.emp_id = i.emp_id);

-- --- 7. Prove the nulls are to blame ----------------------------
SELECT COUNT(*) FROM section WHERE emp_id IS NULL;
SELECT COUNT(*) FROM instructor
 WHERE emp_id NOT IN (SELECT emp_id FROM section
                       WHERE emp_id IS NOT NULL);

-- --- 8. Decide what null means, once ----------------------------
SELECT COALESCE(grade, 'in progress') AS g, COUNT(*)
  FROM enrolment GROUP BY g ORDER BY g;
SELECT roll_no, grade FROM enrolment ORDER BY grade NULLS LAST LIMIT 3;
\end{lstlisting}

\smallskip
\textbf{What to notice.} Step 5 against step 6 is the chapter in one
comparison: the same question, two spellings, 0 against 3. Neither raised an
error, and the wrong one is the one most people write first. If you remember
one rule from Part~IV, make it this: \textbf{never \sqlkw{NOT IN} a column
that can be null.}
\end{lab}

\begin{checkpoint}
\begin{tightnum}
  \item \texttt{grade = 'A'} returns 7{,}268 rows and \texttt{grade <> 'A'}
        returns 32{,}777, on a table of 60{,}037. Account for the remaining
        rows using the truth table, and write the single query that returns
        every row that is \emph{not} an A, including the ungraded ones.
  \item Explain, step by step from the expansion of \sqlkw{NOT IN} into
        \sqlkw{AND}s, why one null in the sub-query's result makes the outer
        query return no rows at all.
  \item A colleague writes \texttt{AVG(COALESCE(score, 0))} ``so that the
        average includes everybody''. Say what decision this makes about
        missing marks, show with a three-row example how it changes the
        average, and say what you would write instead and why.
\end{tightnum}
\end{checkpoint}

\begin{chaptersummary}
\begin{tight}
  \item \sqlkw{SELECT} runs \sqlkw{FROM}, \sqlkw{WHERE}, \sqlkw{GROUP BY},
        \sqlkw{SELECT}, \sqlkw{ORDER BY}, \sqlkw{LIMIT}. An alias does not
        exist in \sqlkw{WHERE}; it does in \sqlkw{ORDER BY}.
  \item \sqlkw{BETWEEN} is inclusive; SQLite's \sqlkw{LIKE} is
        case-insensitive for ASCII; \sqlkw{LIMIT} without \sqlkw{ORDER BY}
        returns rows the plan chose.
  \item Any comparison with a null is \textbf{\unk{}}.
        \sqlkw{WHERE} keeps only \emph{true}, so \unk{} rows vanish.
  \item \textbf{\sqlkw{NOT} \unk{} is \unk{}}: a condition
        and its negation do not cover a table with nulls.
  \item \textbf{On the seeded data:} \texttt{= 'A'} found 7{,}268,
        \texttt{<> 'A'} found 32{,}777, and \textbf{19{,}992 ungraded rows were
        in neither}.
  \item \texttt{= NULL} is always empty. \sqlkw{IS NULL} is the only test.
  \item \textbf{\sqlkw{NOT IN} returned 0 where \sqlkw{NOT EXISTS} returned 3},
        because one null in the list makes every comparison \unk{}.
        Never \sqlkw{NOT IN} a nullable column.
  \item \sqlkw{COALESCE} decides what null means; use it for display, never to
        turn a missing mark into a zero.
  \item SQLite sorts nulls first ascending; say \sqlkw{NULLS LAST} when it
        matters, because other systems differ.
\end{tight}
\end{chaptersummary}

\begin{reviewq}
  \item \textbf{(MCQ)} \texttt{NULL = NULL} evaluates to: (a)~true;
        (b)~false; (c)~\unk{}; (d)~an error.
  \item \textbf{(MCQ)} \texttt{semester BETWEEN 2 AND 4} includes:
        (a)~3 only; (b)~2 and 3; (c)~3 and 4; (d)~2, 3 and 4.
  \item \textbf{(Short)} Why can a column alias be used in \sqlkw{ORDER BY}
        but not in \sqlkw{WHERE}?
  \item \textbf{(Short)} Evaluate \texttt{F OR U}, \texttt{T AND U} and
        \sqlkw{NOT} \texttt{U}, and say which one causes the most errors in
        practice.
  \item \textbf{(Short)} Give the two fixes for a \sqlkw{NOT IN} against a
        nullable column, and say why this course prefers one of them.
  \item \textbf{(Applied)} Write a query listing every student with a label
        computed by \sqlkw{CASE}: \texttt{'senior'} for semesters 7--8,
        \texttt{'mid'} for 3--6, \texttt{'new'} for 1--2. Then say what your
        query does with a null semester, and change it so that a null is
        labelled \texttt{'unknown'} explicitly.
  \item \textbf{(Applied)} You are asked for ``students who have never failed
        a course''. Write it three ways --- with \sqlkw{NOT IN}, with
        \sqlkw{NOT EXISTS}, and with \sqlkw{EXCEPT} --- run all three on the
        seeded database, and explain any difference between their counts in
        terms of the null grades of FALL2026.
\end{reviewq}

%=====================================================================
\chapter{Sub-queries, \texorpdfstring{\sqlkw{EXISTS}}{EXISTS} and Set Operations}\label{ch:subqueries}
%=====================================================================
\locator{4}{Part IV --- Asking Questions}

\begin{outcomes}
By the end of this chapter you will be able to:
\begin{tight}
  \item \textbf{Write} a sub-query in the \sqlkw{WHERE}, \sqlkw{FROM} and
        \sqlkw{SELECT} clauses, and \textbf{say} what each one must return.
  \item \textbf{Distinguish} a correlated from an uncorrelated sub-query by
        reading it.
  \item \textbf{Write} a semi-join and an anti-join in several forms, and
        \textbf{choose} one that is correct whatever the nulls.
  \item \textbf{Explain} what makes a correlated sub-query expensive, and
        \textbf{show} that it is not the correlation.
  \item \textbf{Use} a common table expression to give a sub-query a name.
\end{tight}
\end{outcomes}

\begin{prereq}
\chref{select}, and especially its \sqlkw{NOT IN} alert --- this chapter
measures the four alternatives. \chref{joins} for the anti-join idiom. You
need \texttt{code/m15\_subqueries.py}.
\end{prereq}

\begin{keyterms}
sub-query $\bullet$ nested query $\bullet$ scalar sub-query $\bullet$
correlated sub-query $\bullet$ uncorrelated sub-query $\bullet$ semi-join
$\bullet$ anti-join $\bullet$ \sqlkw{IN} $\bullet$ \sqlkw{EXISTS} $\bullet$
\sqlkw{ANY} $\bullet$ \sqlkw{ALL} $\bullet$ derived table $\bullet$ common
table expression $\bullet$ \sqlkw{WITH} $\bullet$ \sqlkw{UNION} $\bullet$
\sqlkw{UNION ALL} $\bullet$ \sqlkw{EXCEPT}
\end{keyterms}

\section{A Query Inside a Query}

\chref{algebra}'s closure property --- every result is a relation --- is what
makes this chapter possible. A \sqlkw{SELECT} produces a relation, so a
\sqlkw{SELECT} can appear wherever a relation, a list or a single value is
expected.

\begin{definitionbox}[Three places a sub-query can go]
A sub-query is also called a \term{nested query}.

\smallskip
\textbf{In \sqlkw{WHERE}}, after \sqlkw{IN} or \sqlkw{EXISTS}: it supplies a
list of values, or the fact that some row exists.

\smallskip
\textbf{In \sqlkw{FROM}}: it supplies a whole relation, called a
\term{derived table}, which must be given an alias.

\smallskip
\textbf{In the \sqlkw{SELECT} list}, or in a comparison: it must be a
\term{scalar sub-query}, returning exactly one column and at most one row. If
it returns no rows its value is null; if it returns two, SQLite silently uses
the first --- other systems raise an error.
\end{definitionbox}

\begin{definitionbox}[Correlated and uncorrelated]
An \term{uncorrelated} sub-query mentions only its own tables. It could be run
on its own, and in principle it is run once.

\smallskip
A \term{correlated} sub-query mentions a column of the \emph{outer} query, so
its answer depends on which outer row is being considered. In principle it is
run once per outer row.

\begin{lstlisting}[language=sql]
-- uncorrelated: the inner query stands alone
WHERE roll_no IN (SELECT roll_no FROM enrolment WHERE grade = 'A')

-- correlated: s.roll_no belongs to the OUTER query
WHERE EXISTS (SELECT 1 FROM enrolment e
               WHERE e.roll_no = s.roll_no AND e.grade = 'A')
\end{lstlisting}

\smallskip
``In principle'' is doing work in both definitions. The optimiser is free to
run either one however it likes, and the measurement below shows it doing so.
\end{definitionbox}

\section{Semi-joins and Anti-joins}

\begin{definitionbox}[Semi-join, anti-join]
A \term{semi-join} returns the rows of $R$ that have \emph{at least one}
partner in $S$ --- each row once, however many partners it has. ``Students with
at least one A.''

\smallskip
An \term{anti-join} returns the rows of $R$ that have \emph{no} partner in
$S$. ``Instructors who teach nothing.'' It is \chref{algebra}'s difference,
and ``never'', ``none'' and ``no'' always mean one.
\end{definitionbox}

\begin{measured}[Seven spellings, one wrong answer, and correlation that was not the problem]
\textsf{\footnotesize \texttt{code/m15\_subqueries.py}, seeded database,
SQLite 3.51.2, Intel i5-6600, Windows 11. Best of three runs.}

\medskip
\textbf{1. A semi-join, three ways} --- students with at least one A

\begin{tabular}{@{}L{5.0cm}R{2.0cm}R{1.9cm}L{5.0cm}@{}}
\textbf{written as} & \textbf{answer} & \textbf{ms} & \\
\sqlkw{IN} (uncorrelated)   & 3{,}449 & 20.3 & \\
\sqlkw{EXISTS} (correlated) & 3{,}449 & 40.1 & twice the \sqlkw{IN} \\
\sqlkw{JOIN} + \sqlkw{DISTINCT} & 3{,}449 & 21.4 & \\
\end{tabular}

\medskip
\textbf{2. An anti-join, four ways} --- instructors who teach nothing, after
three are hired who teach nothing yet; \texttt{section.emp\_id} has 73 nulls

\begin{tabular}{@{}L{5.0cm}R{2.0cm}R{1.9cm}L{5.0cm}@{}}
\textbf{written as} & \textbf{answer} & \textbf{ms} & \textbf{right?} \\
\sqlkw{NOT IN}                 & \textbf{0} & 0.76 & \textbf{no} --- one null empties it \\
\sqlkw{NOT EXISTS}             & 3 & 2.18 & yes \\
\sqlkw{LEFT JOIN} \dots\ \sqlkw{IS NULL} & 3 & 2.52 & yes \\
\sqlkw{EXCEPT}                 & 3 & 0.44 & yes \\
\end{tabular}

\smallskip
The fastest wrong answer is no use. \sqlkw{EXCEPT} is a set operation rather
than a comparison, so the null on its right-hand side can only remove a null
on its left --- and \texttt{instructor.emp\_id}, a primary key, has none.

\medskip
\textbf{3. A correlated sub-query in the \sqlkw{SELECT} list} --- enrolments
per course

\begin{tabular}{@{}R{1.7cm}R{2.1cm}R{2.4cm}R{2.1cm}R{2.4cm}R{2.4cm}@{}}
& & \multicolumn{2}{c}{\textbf{no index on course\_code}}
& \multicolumn{2}{c}{\textbf{with the index}} \\
\textbf{students} & \textbf{enrolments} & \textbf{correlated} & \textbf{\sqlkw{GROUP BY}}
& \textbf{correlated} & \textbf{\sqlkw{GROUP BY}} \\
2{,}000 & 30{,}027  & 191.1 & 53.8  & \textbf{2.4}  & 5.4 \\
4{,}000 & 60{,}037  & 517.7 & 121.3 & \textbf{5.1}  & 12.2 \\
8{,}000 & 120{,}134 & 1{,}101.4 & 275.8 & \textbf{10.3} & 23.7 \\
\end{tabular}

\smallskip
All four agree on every count. Times in milliseconds.

\medskip
\textbf{Read part 3 carefully, because it contradicts the folklore.} Without
an index the correlated form is about four times slower than \sqlkw{GROUP BY}:
\textit{enrolment}'s primary key is \texttt{(roll\_no, course\_code, \dots)},
so it cannot be searched by \texttt{course\_code}, and each of the 48 courses
scans all 60 thousand enrolments. \textbf{Add one index and the correlated form
becomes the faster of the two} --- 5.1 against 12.2\,ms --- because it now does
48 index lookups, while \sqlkw{GROUP BY} still reads every enrolment. Both
grow linearly with the data; neither diverges from the other as the table
grows.

\smallskip
\textbf{So the cost was never the correlation.} It was whether the inner query
had an index on the column it is correlated on. The plans say so directly:
\texttt{SCAN e} inside \texttt{CORRELATED SCALAR SUBQUERY} without it,
\texttt{SEARCH e USING \dots\ idx\_enr\_course (course\_code=?)} with it.
\end{measured}

\begin{alertbox}[What to take from the measurement]
\begin{tight}
  \item \textbf{For anti-joins, use \sqlkw{NOT EXISTS}.} It is correct
        whatever the nulls, and its cost is in the same range as the
        alternatives. \sqlkw{EXCEPT} is faster when it fits but returns only
        the compared columns; \sqlkw{LEFT JOIN} \dots\ \sqlkw{IS NULL} is
        correct but easy to break with a condition in the wrong clause
        (\chref{joins}).
  \item \textbf{Do not avoid correlated sub-queries on principle.} Avoid
        unindexed ones. ``Correlated sub-queries are slow'' is a rule of thumb
        learned on databases without the right index, and on this one it was
        backwards by a factor of two.
  \item \textbf{Read the plan before rewriting a query.} \sqlkw{SCAN} inside a
        \sqlkw{CORRELATED} sub-query is the specific thing to look for; it
        means the inner query reads a whole table once per outer row.
\end{tight}
\end{alertbox}

\section{\sqlkw{ANY}, \sqlkw{ALL}, and Naming a Sub-query}

\begin{examplebox}[\sqlkw{ALL} and \sqlkw{ANY}, and the null trap again]
\texttt{x > ALL (sub-query)} is true when $x$ exceeds every value returned;
\texttt{x > ANY (sub-query)} when it exceeds at least one. SQLite does not
support either keyword, which is no loss --- both are clearer as an aggregate:

\begin{lstlisting}[language=sql]
-- "credit hours greater than every MA course"
WHERE credit_hours > (SELECT MAX(credit_hours) FROM course
                       WHERE dept_code = 'MA')
\end{lstlisting}

\smallskip
And \sqlkw{ALL} has \sqlkw{NOT IN}'s trap: \texttt{x > ALL (\dots, NULL)} is
\unk{} for every $x$, because one of its comparisons is. The
aggregate form ignores nulls, which is usually what was meant.
\end{examplebox}

\begin{definitionbox}[Common table expression]
A \term{common table expression} (CTE) names a sub-query before the main query
uses it:

\begin{lstlisting}[language=sql]
WITH a_students AS (
  SELECT DISTINCT roll_no FROM enrolment WHERE grade = 'A'
)
SELECT s.name
  FROM student s JOIN a_students a ON a.roll_no = s.roll_no
 WHERE s.programme = 'BSIT';
\end{lstlisting}

\smallskip
It computes nothing a derived table could not. Its value is that the
sub-query has a \emph{name}, can be read top to bottom, and can be referred to
more than once. \sqlkw{WITH RECURSIVE} --- used in \chref{whodoeswhat}'s
laboratory to generate 300 rows --- is the one form that does something new:
it can walk the prerequisite chain to any depth, which no fixed number of joins
can.
\end{definitionbox}

\begin{worked}[Every prerequisite of a course, to any depth]
\textbf{The question.} Before taking a course, which courses must a student
have passed --- directly, or because a prerequisite has a prerequisite?

\medskip
\textbf{Step 1 --- why a join cannot do it.} One join finds direct
prerequisites. Two joins find their prerequisites. Each extra level needs an
extra join, and nobody knows in advance how deep the chain goes.

\medskip
\textbf{Step 2 --- the recursive CTE.} Start from the direct prerequisites
(the \emph{anchor}), then repeatedly add the prerequisites of whatever is
already found (the \emph{recursive step}), until a step adds nothing:

\begin{lstlisting}[language=sql]
WITH RECURSIVE needs(code) AS (
  SELECT requires FROM prerequisite
   WHERE course_code = 'CS-305'                 -- the anchor
  UNION
  SELECT p.requires FROM prerequisite p
    JOIN needs n ON p.course_code = n.code      -- the step
)
SELECT code FROM needs;
\end{lstlisting}

\medskip
\textbf{Step 3 --- the answer, on the seeded database.} \texttt{CS-305}
directly requires two courses, \texttt{CS-201} and \texttt{CS-258}. The query
returns \textbf{four}: those two, plus \texttt{CS-150} (required by
\texttt{CS-258}) and \texttt{CS-101} (required by \texttt{CS-150}). The chain
is three levels deep --- 305, 258, 150, 101 --- and a query with one join
would have found half of it.

\medskip
\textbf{Step 4 --- why \sqlkw{UNION} and not \sqlkw{UNION ALL}.} \sqlkw{UNION}
discards rows already found, so a course reachable by two paths is listed
once --- and, more importantly, a cycle in the data (A requires B requires A)
terminates instead of running for ever. \textbf{\chref{mapping}'s
\sqlkw{CHECK} stops a course requiring itself, but not a two-step cycle.} Only
\sqlkw{UNION} saves you from that one.
\end{worked}

\dsfig{%
\begin{tikzpicture}[
  crs/.style={rounded corners=3pt, draw=#1, line width=1pt, fill=#1!8,
              text=#1!60!black, font=\scriptsize\bfseries\ttfamily,
              minimum width=1.6cm, minimum height=0.62cm},
  req/.style={-{Stealth[length=2.2mm]}, neutral!70, line width=0.9pt},
  lvl/.style={font=\scriptsize\itshape, text=neutral},
  br/.style={draw=#1, line width=0.9pt,
             decorate, decoration={brace, amplitude=5pt, mirror}}
]
  \node[crs=primary]   (a) at (0,0)    {CS-305};
  \node[crs=greenhl]   (b) at (3.0,0.75) {CS-201};
  \node[crs=greenhl]   (c) at (3.0,-0.75) {CS-258};
  \node[crs=goldhl]    (d) at (6.0,-0.75) {CS-150};
  \node[crs=accent]    (e) at (9.0,-0.75) {CS-101};
  \draw[req] (a) -- (b);
  \draw[req] (a) -- (c);
  \draw[req] (c) -- (d);
  \draw[req] (d) -- (e);
  \node[lvl, anchor=south] at (3.0,1.25) {level 1};
  \node[lvl, anchor=south] at (6.0,-0.15) {level 2};
  \node[lvl, anchor=south] at (9.0,-0.15) {level 3};
  \draw[br=greenhl] (2.1,-1.35) -- (3.9,-1.35);
  \node[lvl, text=greenhl!60!black, anchor=north] at (3.0,-1.6)
       {one join finds these two};
  \draw[br=accent] (2.1,-2.25) -- (9.9,-2.25);
  \node[lvl, text=accent, anchor=north] at (6.0,-2.5)
       {\sqlkw{WITH RECURSIVE} finds all four, however deep};
\end{tikzpicture}}%
{The prerequisites of \texttt{CS-305} on the seeded database. An arrow means
``requires''. A fixed query reaches a fixed depth; the recursive one keeps
going until a step adds nothing new.}{prereqchain}

\begin{pitfall}
\begin{tight}
  \item \textbf{\sqlkw{NOT IN} against a nullable column.} Measured: 0 instead
        of 3.
  \item \textbf{Assuming correlation is the cost.} The missing index is.
        Measured: with the index, the correlated form was the faster.
  \item \textbf{A scalar sub-query that returns two rows.} SQLite quietly uses
        one; other systems refuse.
  \item \textbf{A derived table with no alias.} Required.
  \item \textbf{\sqlkw{UNION} where \sqlkw{UNION ALL} was meant}, or the
        reverse. \sqlkw{UNION} removes duplicates and costs a sort; in a
        recursive CTE it is also what stops a cycle.
  \item \textbf{\sqlkw{x > ALL} over a list with a null.} \unk{}
        for every $x$. Use \sqlkw{MAX}.
\end{tight}
\end{pitfall}

\begin{lab}[Write one question seven ways, and find the one that lies]
\begin{lstlisting}[language=bash]
# --- 1. The seeded database ----------------------------------------
python code/schema/seed.py university.db
sqlite3 university.db
\end{lstlisting}
\begin{lstlisting}[language=sql]
-- --- 2. A semi-join three ways ---------------------------------
.headers on
.mode box
.timer on
SELECT COUNT(*) FROM student WHERE roll_no IN
  (SELECT roll_no FROM enrolment WHERE grade = 'A');
SELECT COUNT(*) FROM student s WHERE EXISTS
  (SELECT 1 FROM enrolment e WHERE e.roll_no = s.roll_no
                               AND e.grade = 'A');
SELECT COUNT(DISTINCT s.roll_no) FROM student s
  JOIN enrolment e ON e.roll_no = s.roll_no WHERE e.grade = 'A';
-- Same answer three times? Record the three timings.

-- --- 3. Hire three people, then an anti-join four ways -----------
PRAGMA foreign_keys = ON;
INSERT INTO instructor VALUES
  ('E-9001','New Hire One','Lecturer','CS'),
  ('E-9002','New Hire Two','Lecturer','IT'),
  ('E-9003','New Hire Three','Lecturer','SE');
SELECT COUNT(*) FROM instructor
 WHERE emp_id NOT IN (SELECT emp_id FROM section);
SELECT COUNT(*) FROM instructor i WHERE NOT EXISTS
  (SELECT 1 FROM section s WHERE s.emp_id = i.emp_id);
SELECT COUNT(*) FROM instructor i LEFT JOIN section s
    ON s.emp_id = i.emp_id WHERE s.emp_id IS NULL;
SELECT COUNT(*) FROM (SELECT emp_id FROM instructor
                      EXCEPT SELECT emp_id FROM section);
-- Three say 3. One says 0. Which, and why?

-- --- 4. A correlated sub-query, and its plan ---------------------
EXPLAIN QUERY PLAN
SELECT c.course_code, (SELECT COUNT(*) FROM enrolment e
                        WHERE e.course_code = c.course_code)
  FROM course c;
-- Find SCAN inside CORRELATED. Then time it:
SELECT c.course_code, (SELECT COUNT(*) FROM enrolment e
                        WHERE e.course_code = c.course_code)
  FROM course c;

-- --- 5. Give the inner query an index, and look again -------------
CREATE INDEX idx_enr_course ON enrolment(course_code);
EXPLAIN QUERY PLAN
SELECT c.course_code, (SELECT COUNT(*) FROM enrolment e
                        WHERE e.course_code = c.course_code)
  FROM course c;
-- SCAN has become SEARCH. Time it again and compare.

-- --- 6. Every prerequisite, to any depth ------------------------
SELECT requires FROM prerequisite WHERE course_code = 'CS-305';
-- The direct prerequisites: two. Now every one, to any depth:
WITH RECURSIVE needs(code) AS (
  SELECT requires FROM prerequisite WHERE course_code = 'CS-305'
  UNION
  SELECT p.requires FROM prerequisite p JOIN needs n
      ON p.course_code = n.code)
SELECT code FROM needs;
-- Four. Draw the chain on paper and check every one.
\end{lstlisting}
\begin{lstlisting}[language=bash]
# --- 7. The measured version ----------------------------------------
python code/m15_subqueries.py
\end{lstlisting}

\smallskip
\textbf{What to notice.} Step 5 is where a rule of thumb meets a measurement.
The correlated query you were told to avoid becomes faster than the
\sqlkw{GROUP BY} you were told to prefer, because of one index and nothing
else. Rules about query \emph{shape} are usually rules about indexes in
disguise, and \chref{indexes} takes the disguise off.
\end{lab}

\begin{checkpoint}
\begin{tightnum}
  \item Classify each sub-query in the laboratory as correlated or
        uncorrelated, and for each correlated one name the outer column it
        refers to.
  \item Of the four anti-join forms, one returned 0 and three returned 3. Say
        which, explain why in terms of three-valued logic, and say why
        \sqlkw{EXCEPT} was not affected by the same nulls.
  \item Without an index the correlated sub-query was four times slower than
        \sqlkw{GROUP BY}; with one it was twice as fast. Explain both results
        in terms of what each query reads, and say what this tells you about
        the advice ``avoid correlated sub-queries''.
\end{tightnum}
\end{checkpoint}

\begin{chaptersummary}
\begin{tight}
  \item A sub-query can go in \sqlkw{WHERE} (a list, or existence),
        \sqlkw{FROM} (a derived table, which needs an alias) or the
        \sqlkw{SELECT} list (a scalar: one column, at most one row).
  \item A \textbf{correlated} sub-query mentions an outer column and, in
        principle, runs per outer row. The optimiser decides what actually
        happens.
  \item A \textbf{semi-join} keeps rows with at least one partner; an
        \textbf{anti-join} keeps rows with none.
  \item \textbf{Measured: the three semi-join forms agreed} (3{,}449);
        \sqlkw{EXISTS} took twice as long as \sqlkw{IN} here.
  \item \textbf{Measured: \sqlkw{NOT IN} returned 0} where \sqlkw{NOT
        EXISTS}, \sqlkw{LEFT JOIN} \dots\ \sqlkw{IS NULL} and \sqlkw{EXCEPT}
        all returned 3. Use \sqlkw{NOT EXISTS}.
  \item \textbf{Measured: without an index on the correlation column, the
        correlated form was 4$\times$ slower than \sqlkw{GROUP BY}; with one
        it was 2$\times$ faster.} The cost was never the correlation.
  \item \sqlkw{SCAN} inside a \sqlkw{CORRELATED} sub-query in a plan means a
        whole table read per outer row. Look for it.
  \item A \textbf{CTE} names a sub-query. \sqlkw{WITH RECURSIVE} walks a chain
        to any depth, and \sqlkw{UNION} in it is what stops a cycle.
  \item \sqlkw{ALL} over a list containing a null is \unk{}; use an
        aggregate.
\end{tight}
\end{chaptersummary}

\begin{reviewq}
  \item \textbf{(MCQ)} A sub-query that mentions a column of the outer query
        is: (a)~scalar; (b)~derived; (c)~correlated; (d)~recursive.
  \item \textbf{(MCQ)} Which anti-join form can return no rows because of a
        single null? (a)~\sqlkw{NOT EXISTS}; (b)~\sqlkw{NOT IN};
        (c)~\sqlkw{EXCEPT}; (d)~\sqlkw{LEFT JOIN} \dots\ \sqlkw{IS NULL}.
  \item \textbf{(Short)} What must a scalar sub-query return, and what does
        SQLite do if it returns more?
  \item \textbf{(Short)} Why does a recursive CTE over \textit{prerequisite}
        need \sqlkw{UNION} rather than \sqlkw{UNION ALL}?
  \item \textbf{(Short)} State in one sentence what decides whether a
        correlated sub-query is expensive.
  \item \textbf{(Applied)} Write ``courses with no enrolments in FALL2026''
        as \sqlkw{NOT EXISTS}, as \sqlkw{LEFT JOIN} \dots\ \sqlkw{IS NULL} and
        as \sqlkw{EXCEPT}. Run all three on the seeded database (insert one
        new course first so the answer is not empty), confirm they agree, and
        read each plan to say which index each one uses.
  \item \textbf{(Applied)} A student's prerequisite chain is to be checked at
        enrolment time. Write the recursive query that lists every course a
        student must have passed before taking a given one, then extend it to
        list only those the student has \emph{not} passed. Say what happens if
        somebody inserts a two-step cycle, and how your query behaves.
\end{reviewq}

%=====================================================================
\chapter{Grouping, Aggregation and Window Functions}\label{ch:grouping}
%=====================================================================
\locator{4}{Part IV --- Asking Questions}

\begin{outcomes}
By the end of this chapter you will be able to:
\begin{tight}
  \item \textbf{Use} \sqlkw{COUNT}, \sqlkw{SUM}, \sqlkw{AVG}, \sqlkw{MIN} and
        \sqlkw{MAX}, and \textbf{predict} what each does with a null.
  \item \textbf{Write} a \sqlkw{GROUP BY} query and \textbf{state} which
        columns may appear in its \sqlkw{SELECT} list.
  \item \textbf{Choose} between \sqlkw{WHERE} and \sqlkw{HAVING} by what the
        condition is about.
  \item \textbf{Write} a window function with \sqlkw{PARTITION BY} and
        \sqlkw{ORDER BY}, and \textbf{explain} how it differs from grouping.
  \item \textbf{Read} a query plan for \sqlkw{TEMP B-TREE}, and \textbf{say}
        what index would remove it.
\end{tight}
\end{outcomes}

\begin{prereq}
\chref{select} for three-valued logic, which the aggregates handle in their
own way. \chref{subqueries} for correlated sub-queries, one of which this
chapter replaces. You need \texttt{code/m16\_grouping.py}.
\end{prereq}

\begin{keyterms}
aggregate function $\bullet$ \sqlkw{COUNT} $\bullet$ \sqlkw{SUM} $\bullet$
\sqlkw{AVG} $\bullet$ \sqlkw{MIN} $\bullet$ \sqlkw{MAX} $\bullet$
\sqlkw{GROUP BY} $\bullet$ group $\bullet$ \sqlkw{HAVING} $\bullet$ bare
column $\bullet$ window function $\bullet$ \sqlkw{OVER} $\bullet$
\sqlkw{PARTITION BY} $\bullet$ \sqlkw{RANK} $\bullet$ \sqlkw{DENSE\_RANK}
$\bullet$ \sqlkw{ROW\_NUMBER} $\bullet$ running total $\bullet$
\sqlkw{TEMP B-TREE}
\end{keyterms}

\section{Collapsing Rows Into One}

Every query so far has returned rows that existed --- possibly filtered,
possibly joined, but each one traceable to a tuple in the database. An
\term{aggregate function} returns a row that existed nowhere: the \emph{average} grade, the
\emph{number} of students.

\begin{definitionbox}[The five aggregates, and what they skip]
\begin{tabular}{@{}L{3.6cm}L{10.4cm}@{}}
\sqlkw{COUNT(*)} & the number of rows. Counts everything, nulls included. \\
\sqlkw{COUNT}$(x)$ & the number of rows where $x$ is \textbf{not null}. \\
\sqlkw{SUM}, \sqlkw{AVG} & over the non-null values only. \sqlkw{AVG}
  divides by the number of non-null values, not by the number of rows. \\
\sqlkw{MIN}, \sqlkw{MAX} & over the non-null values only. \\
\end{tabular}

\smallskip
Every aggregate except \sqlkw{COUNT(*)} \textbf{skips nulls}. Over a set that
is entirely null, \sqlkw{SUM} and \sqlkw{AVG} return null --- not zero ---
and \sqlkw{COUNT} returns 0.
\end{definitionbox}

\section{\sqlkw{GROUP BY}}

\begin{definitionbox}[Grouping]
\sqlkw{GROUP BY} $a$ partitions the rows into groups, one for each distinct
value of $a$, and an aggregate in the \sqlkw{SELECT} list is then computed
\emph{per group}. The result has one row per group.

\smallskip
\textbf{The rule for the \sqlkw{SELECT} list:} every column in it must either
be in the \sqlkw{GROUP BY}, or be inside an aggregate. A column that is neither
--- a \term{bare column} --- has no single value per group, so there is
nothing correct to print.

\smallskip
Null forms a group of its own. All the nulls are grouped together, even though
\texttt{NULL = NULL} is \unk{}: grouping uses ``is the same value''
rather than ``compares equal''.
\end{definitionbox}

\begin{alertbox}[SQLite allows the bare column, and it should not]
\begin{lstlisting}[language=sql]
SELECT programme, name, COUNT(*) FROM student GROUP BY programme;
\end{lstlisting}

PostgreSQL rejects this: \texttt{name} is neither grouped nor aggregated.
SQLite runs it and prints, for each programme, the name of \emph{some} student
in it --- whichever row the engine happened to process last. The output looks
like a meaningful answer, and on each run it may name a different student.

\smallskip
There is one documented exception: with exactly one \sqlkw{MIN} or
\sqlkw{MAX}, SQLite guarantees the bare columns come from the row holding that
minimum or maximum. It is a convenience no other system offers, so code that
depends on it will not move. \textbf{In this course, every non-aggregated
column is in the \sqlkw{GROUP BY}.}
\end{alertbox}

\section{Measured: Nulls, Plans, and a Rank}

\begin{measured}[What aggregates skip, what an index removes, and what a window replaces]
\textsf{\footnotesize \texttt{code/m16\_grouping.py}, seeded database,
SQLite 3.51.2, Intel i5-6600, Windows 11. Best of three runs.}

\medskip
\textbf{1. Aggregates and nulls} --- \textit{enrolment}, 60{,}037 rows

\begin{tabular}{@{}L{6.4cm}R{2.4cm}L{5.4cm}@{}}
\sqlkw{COUNT(*)}      & 60{,}037 & every row \\
\sqlkw{COUNT}(grade)  & 40{,}045 & rows where grade is not null \\
difference            & 19{,}992 & the ungraded FALL2026 rows \\
\sqlkw{AVG}(points)   & \textbf{2.383} & nulls skipped: $\div$ 40{,}045 \\
\sqlkw{AVG}(\sqlkw{COALESCE}(points, 0)) & \textbf{1.590} & nulls as zero:
  $\div$ 60{,}037 \\
\end{tabular}

\smallskip
The second average is \textbf{33\% lower}, and it has decided that an ungraded
course is an F. \sqlkw{GROUP BY} grade puts the 19{,}992 nulls in a group of
their own, beside A (7{,}268), B (12{,}756), C (11{,}283), D (5{,}524) and F
(3{,}214).

\medskip
\textbf{2. \sqlkw{GROUP BY} with and without a supporting index} --- enrolments
per course

\begin{tabular}{@{}R{2.0cm}R{2.4cm}R{2.6cm}R{2.6cm}R{2.4cm}@{}}
\textbf{students} & \textbf{enrolments} & \textbf{no index} & \textbf{with index}
& \textbf{speed-up} \\
4{,}000  & 60{,}037  & 37.0\,ms  & 6.8\,ms  & 5.5$\times$ \\
8{,}000  & 120{,}134 & 86.3\,ms  & 13.9\,ms & 6.2$\times$ \\
16{,}000 & 240{,}025 & 174.2\,ms & 31.3\,ms & 5.6$\times$ \\
\end{tabular}

\smallskip
Without the index the plan says \texttt{USE TEMP B-TREE FOR GROUP BY}: the
rows were sorted into groups before they could be counted. An index on
\texttt{course\_code} is already in that order, so the sort disappears from the
plan --- and five-sixths of the cost goes with it. The speed-up is roughly
constant as the table grows, because both forms are linear and the sort was a
fixed proportion of the work.

\medskip
\textbf{3. The same filter in \sqlkw{WHERE} and in \sqlkw{HAVING}}

\begin{tabular}{@{}L{6.0cm}R{1.5cm}R{1.7cm}L{3.0cm}@{}}
& \textbf{ms} & \textbf{groups} & \\
\texttt{WHERE course\_code LIKE 'CS-\%'}  & 13.1 & 6 & \\
\texttt{HAVING course\_code LIKE 'CS-\%'} & 13.9 & 6 & same answer \\
\end{tabular}

\smallskip
\textbf{The two compiled to the identical program}, 47 bytecode operations
each, with the \sqlkw{LIKE} tested at operation 8, \emph{before} the grouping
sort at operation 12. SQLite moved the \sqlkw{HAVING} condition into
\sqlkw{WHERE} by itself --- which it can do only because the condition mentions
no aggregate. \texttt{HAVING COUNT(*) > 1000} was tested at operation 33,
\emph{after} the sort at operation 9, because a count does not exist until the
grouping is done.

\medskip
\textbf{4. A rank within each programme} --- by number of A grades

\begin{tabular}{@{}R{2.0cm}R{3.4cm}R{3.4cm}R{2.0cm}L{2.4cm}@{}}
\textbf{students} & \textbf{window \sqlkw{RANK()}} & \textbf{correlated}
& \textbf{ratio} & \textbf{agree?} \\
1{,}000 & 16.9\,ms  & 51.6\,ms  & 3$\times$ & yes \\
2{,}000 & 38.4\,ms  & 190.4\,ms & 5$\times$ & yes \\
4{,}000 & 103.2\,ms & 736.7\,ms & \textbf{7$\times$} & yes \\
\end{tabular}

\smallskip
The ratio grows with the data. The correlated form counts, for each student,
how many in the same programme did better --- every student against every
other in the programme, quadratic. The window function sorts each programme
once and reads the rank off the order.
\end{measured}

\begin{alertbox}[What the four results say about how to write these queries]
\begin{tight}
  \item \textbf{Never \sqlkw{COALESCE} to zero inside an aggregate unless a
        missing value really is a zero.} It moved the average by a third.
        That is the same fact \chref{ddl} found with a text value in a numeric
        column, approached from the other side.
  \item \textbf{\sqlkw{TEMP B-TREE} in a plan is a sort you could have avoided.}
        An index on the grouping column supplies the order for free.
        \chref{indexes} shows where the same idea stops working.
  \item \textbf{Put a condition on a column in \sqlkw{WHERE} because that is
        what it means, not because it is faster here.} SQLite fixed the
        \sqlkw{HAVING} version itself; that is an optimisation, not a promise,
        and the next system you use may not make it.
  \item \textbf{For rankings, running totals and ``top $n$ per group'', use a
        window function.} The correlated alternative was seven times slower at
        4{,}000 students and getting worse.
\end{tight}
\end{alertbox}

\section{Window Functions}

\begin{definitionbox}[A window function]
A \term{window function} computes a value for each row from a set of related
rows --- its \emph{window} --- \textbf{without collapsing them}. Grouping turns
many rows into one; a window function keeps every row and adds a column.

\begin{lstlisting}[language=sql]
SELECT roll_no, programme, n_a,
       RANK() OVER (PARTITION BY programme ORDER BY n_a DESC) AS pos
  FROM per_student;
\end{lstlisting}

\smallskip
\sqlkw{PARTITION BY} divides the rows into windows, like \sqlkw{GROUP BY} but
without merging them. \sqlkw{ORDER BY} inside \sqlkw{OVER} orders each window.
The function then runs along the ordered window.
\end{definitionbox}

\dsfig{%
\begin{tikzpicture}[
  cell/.style={draw=primary!35, line width=0.5pt, minimum width=1.45cm,
               minimum height=0.40cm, font=\scriptsize\ttfamily,
               text=primary!65!black, inner sep=0pt},
  hd/.style={cell, fill=primary, text=white, font=\scriptsize\bfseries},
  p1/.style={cell, fill=secondary!14},
  p2/.style={cell, fill=goldhl!16},
  nw/.style={cell, fill=greenhl!20, font=\scriptsize\bfseries\ttfamily,
             text=greenhl!55!black},
  ttl/.style={font=\scriptsize\bfseries, text=#1, anchor=south},
  arr/.style={-{Stealth[length=2.2mm]}, primary!50, line width=0.9pt},
  nt/.style={font=\scriptsize\itshape, text=neutral, align=center}
]
  % source rows
  \node[ttl=primary] at (-5.3,0.5) {per\_student};
  \foreach \c [count=\i] in {prog, n\_a}{\node[hd] at (-6.75+\i*1.45,0.2) {\c};}
  \foreach \r/\p/\n/\s in {1/BSIT/9/p1, 2/BSIT/7/p1, 3/BSIT/5/p1,
                           4/BSCS/8/p2, 5/BSCS/6/p2}{
    \node[\s] at (-5.30,-0.4*\r+0.2) {\p};
    \node[\s] at (-3.85,-0.4*\r+0.2) {\n};}

  % group by
  \draw[arr] (-2.95,-0.25) -- (-1.75,0.15);
  \node[ttl=secondary!70!black] at (0.0,0.5) {\sqlkw{GROUP BY} prog};
  \foreach \c [count=\i] in {prog, MAX(n\_a)}{\node[hd] at (-1.45+\i*1.45,0.2) {\c};}
  \node[p1] at (0.0,-0.2) {BSIT};  \node[p1] at (1.45,-0.2) {9};
  \node[p2] at (0.0,-0.6) {BSCS};  \node[p2] at (1.45,-0.6) {8};
  \node[nt, text width=3.6cm] at (0.72,-1.35)
       {five rows become \textbf{two}: one per group};

  % window
  \draw[arr] (-2.95,-1.15) .. controls (-1.2,-2.4) and (2.5,-2.4) .. (3.4,-1.3);
  \node[ttl=greenhl!60!black] at (5.4,0.5)
       {\sqlkw{RANK() OVER} (\sqlkw{PARTITION BY} prog)};
  \foreach \c [count=\i] in {prog, n\_a, rank}{\node[hd] at (2.95+\i*1.45,0.2) {\c};}
  \foreach \r/\p/\n/\k/\s in {1/BSIT/9/1/p1, 2/BSIT/7/2/p1, 3/BSIT/5/3/p1,
                              4/BSCS/8/1/p2, 5/BSCS/6/2/p2}{
    \node[\s] at (4.40,-0.4*\r+0.2) {\p};
    \node[\s] at (5.85,-0.4*\r+0.2) {\n};
    \node[nw] at (7.30,-0.4*\r+0.2) {\k};}
  \node[nt, text width=4.4cm] at (5.85,-2.35)
       {five rows stay \textbf{five}; each gains a column computed from its
        own partition};
\end{tikzpicture}}%
{The difference in one picture. Grouping collapses each partition into a single
row; a window function keeps every row and adds a value computed from the
partition it belongs to. The colours mark the two partitions.}{groupwindow}

\begin{examplebox}[\sqlkw{ROW\_NUMBER}, \sqlkw{RANK} and \sqlkw{DENSE\_RANK} on a tie]
Four students in one programme with 9, 7, 7 and 5 A grades:

\medskip
\begin{tabular}{@{}R{2.2cm}R{3.2cm}R{2.6cm}R{3.2cm}@{}}
\textbf{A grades} & \sqlkw{ROW\_NUMBER} & \sqlkw{RANK} & \sqlkw{DENSE\_RANK} \\
9 & 1 & 1 & 1 \\
7 & 2 & 2 & 2 \\
7 & 3 & 2 & 2 \\
5 & 4 & \textbf{4} & \textbf{3} \\
\end{tabular}

\smallskip
\sqlkw{ROW\_NUMBER} never ties, so it numbers the two sevens arbitrarily ---
use it only when the order is fully determined. \sqlkw{RANK} gives ties the
same number and \emph{skips} the next, as a sports table does.
\sqlkw{DENSE\_RANK} does not skip. The university's merit list is a policy
decision between the last two, and the query has to make it explicitly.
\end{examplebox}

\begin{worked}[A running total of credit hours, per student, in term order]
\textbf{The question.} For each enrolment of one student, how many credit hours
had they completed up to and including that term?

\medskip
\textbf{Step 1 --- why grouping cannot do it.} \sqlkw{GROUP BY} roll\_no gives
one total per student. The question wants every row kept, each with its own
total so far.

\medskip
\textbf{Step 2 --- the window.}
\begin{lstlisting}[language=sql]
SELECT e.semester, c.course_code, c.credit_hours,
       SUM(c.credit_hours) OVER (
         PARTITION BY e.roll_no
         ORDER BY e.semester
         ROWS BETWEEN UNBOUNDED PRECEDING AND CURRENT ROW
       ) AS hours_so_far
  FROM enrolment e JOIN course c ON c.course_code = e.course_code
 WHERE e.roll_no = (SELECT MIN(roll_no) FROM student);
\end{lstlisting}

\medskip
\textbf{Step 3 --- the frame.} \sqlkw{ROWS BETWEEN UNBOUNDED PRECEDING AND
CURRENT ROW} says ``from the first row of the window to this one''. That is
what makes it a running total rather than a grand total repeated on every row.

\medskip
\textbf{Step 4 --- a trap in the ordering.} \texttt{ORDER BY e.semester} sorts
the text \texttt{'FALL2025'}, \texttt{'FALL2026'}, \texttt{'SPRING2026'}
\emph{alphabetically} --- and so puts Fall 2026 before Spring 2026, which is
wrong. A term label that sorts correctly as text has to be designed to,
\texttt{'2025-2'} rather than \texttt{'FALL2025'}. This is the same column that
broke \sqlkw{NATURAL JOIN} in \chref{joins}, and the second reason to revisit
it when \chref{normalforms} reconsiders the schema.
\end{worked}

\begin{pitfall}
\begin{tight}
  \item \textbf{A bare column in a grouped query.} SQLite prints some row's
        value; other systems refuse.
  \item \textbf{\sqlkw{COALESCE} to zero inside \sqlkw{AVG}.} Measured: a third
        off.
  \item \textbf{Expecting \sqlkw{SUM} of nothing to be 0.} It is null.
        \sqlkw{COALESCE}(\sqlkw{SUM}(x), 0) is the honest fix there.
  \item \textbf{\sqlkw{HAVING} for a condition on a column.} It works here,
        because SQLite rewrites it. Write \sqlkw{WHERE}.
  \item \textbf{A correlated sub-query for a rank.} Measured: seven times
        slower and growing. Use \sqlkw{RANK} \sqlkw{OVER}.
  \item \textbf{\sqlkw{ROW\_NUMBER} where ties are possible.} It breaks them
        arbitrarily.
  \item \textbf{Ordering by a term label stored as text.} It sorts
        alphabetically, which is not chronologically.
\end{tight}
\end{pitfall}

\begin{lab}[Move an average by a third, remove a sort, and rank a programme]
\begin{lstlisting}[language=bash]
# --- 1. The seeded database ----------------------------------------
python code/schema/seed.py university.db
sqlite3 university.db
\end{lstlisting}
\begin{lstlisting}[language=sql]
-- --- 2. Count three ways ----------------------------------------
.headers on
.mode box
.timer on
SELECT COUNT(*), COUNT(grade), COUNT(DISTINCT grade) FROM enrolment;
-- Three numbers. Say what each one counts.

-- --- 3. Two averages, one of them a decision about failure ------
SELECT AVG(CASE grade WHEN 'A' THEN 4 WHEN 'B' THEN 3 WHEN 'C' THEN 2
                      WHEN 'D' THEN 1 WHEN 'F' THEN 0 END) FROM enrolment;
SELECT AVG(COALESCE(CASE grade WHEN 'A' THEN 4 WHEN 'B' THEN 3
              WHEN 'C' THEN 2 WHEN 'D' THEN 1 WHEN 'F' THEN 0 END, 0))
  FROM enrolment;
-- By what percentage do they differ, and what did the second decide?

-- --- 4. The null group -----------------------------------------
SELECT grade, COUNT(*) FROM enrolment GROUP BY grade ORDER BY grade;

-- --- 5. A sort you can see --------------------------------------
EXPLAIN QUERY PLAN
SELECT course_code, COUNT(*) FROM enrolment GROUP BY course_code;
SELECT course_code, COUNT(*) FROM enrolment GROUP BY course_code;
-- Find TEMP B-TREE. Record the time.

-- --- 6. Remove it -----------------------------------------------
CREATE INDEX idx_enr_course ON enrolment(course_code);
EXPLAIN QUERY PLAN
SELECT course_code, COUNT(*) FROM enrolment GROUP BY course_code;
SELECT course_code, COUNT(*) FROM enrolment GROUP BY course_code;
-- What disappeared from the plan, and what happened to the time?

-- --- 7. WHERE against HAVING, and a condition that cannot move ---
SELECT course_code, COUNT(*) FROM enrolment
 GROUP BY course_code HAVING COUNT(*) > 1300;
-- Try to write that one with WHERE instead. Why can you not?

-- --- 8. A rank, with a window ------------------------------------
WITH a AS (SELECT s.roll_no, s.programme, SUM(e.grade = 'A') AS n_a
             FROM student s JOIN enrolment e ON e.roll_no = s.roll_no
            GROUP BY s.roll_no)
SELECT programme, roll_no, n_a,
       RANK()       OVER w AS rnk,
       DENSE_RANK() OVER w AS dense
  FROM a WINDOW w AS (PARTITION BY programme ORDER BY n_a DESC)
 ORDER BY programme, rnk LIMIT 12;
-- Find a tie. Show where RANK and DENSE_RANK disagree.
\end{lstlisting}
\begin{lstlisting}[language=bash]
# --- 9. The measured version ----------------------------------------
python code/m16_grouping.py
\end{lstlisting}

\smallskip
\textbf{What to notice.} Step 7 is the real difference between the two
clauses. A condition on \texttt{course\_code} could go in either, and SQLite
quietly makes them identical. A condition on \sqlkw{COUNT(*)} can only go in
\sqlkw{HAVING}, because when \sqlkw{WHERE} runs the count has not been
computed. Choose by \emph{what the condition is about}, and the clause chooses
itself.
\end{lab}

\begin{checkpoint}
\begin{tightnum}
  \item On the seeded data \sqlkw{AVG}(points) is 2.383 and
        \sqlkw{AVG}(\sqlkw{COALESCE}(points, 0)) is 1.590. Explain both
        numbers in terms of what each one divides by, and say which one the
        examinations office should publish and why.
  \item The \sqlkw{WHERE} and \sqlkw{HAVING} versions of the same filter
        compiled to identical programs. Explain why SQLite was able to do that
        for \texttt{course\_code LIKE 'CS-\%'} and not for \texttt{COUNT(*) >
        1000}, and say why you should write \sqlkw{WHERE} anyway.
  \item Give a question that needs a window function rather than
        \sqlkw{GROUP BY}, write the query, and say what \sqlkw{GROUP BY} would
        have lost.
\end{tightnum}
\end{checkpoint}

\begin{chaptersummary}
\begin{tight}
  \item Every aggregate except \sqlkw{COUNT(*)} \textbf{skips nulls};
        \sqlkw{AVG} divides by the non-null count; \sqlkw{SUM} of nothing is
        null, not 0.
  \item \textbf{Measured:} \sqlkw{AVG} skipping nulls gave 2.383;
        \sqlkw{COALESCE} to zero gave 1.590 --- \textbf{a third lower}, and a
        decision that ungraded means failed.
  \item \sqlkw{GROUP BY} makes one row per group. Every \sqlkw{SELECT} column
        must be grouped or aggregated; SQLite permits a bare column and prints
        an arbitrary row's value.
  \item Null forms its own group.
  \item \textbf{Measured:} an index on the grouping column removed
        \sqlkw{TEMP B-TREE FOR GROUP BY} from the plan and made the query
        \textbf{5.5--6.2$\times$} faster.
  \item \textbf{Measured:} the \sqlkw{WHERE} and \sqlkw{HAVING} forms of a
        column condition \textbf{compiled to the same program} --- SQLite
        moved it. An aggregate condition cannot move. Choose by what the
        condition is about.
  \item A \textbf{window function} adds a value to every row without
        collapsing them. \sqlkw{PARTITION BY} defines the windows,
        \sqlkw{ORDER BY} orders them, the frame bounds them.
  \item \textbf{Measured:} a window \sqlkw{RANK()} was \textbf{7$\times$}
        faster than the correlated equivalent at 4{,}000 students, and the gap
        was widening.
  \item \sqlkw{RANK} skips after a tie, \sqlkw{DENSE\_RANK} does not,
        \sqlkw{ROW\_NUMBER} breaks ties arbitrarily.
  \item A term label stored as \texttt{'FALL2025'} sorts alphabetically, not
        chronologically --- the second fault found in the \texttt{semester}
        column.
\end{tight}
\end{chaptersummary}

\begin{reviewq}
  \item \textbf{(MCQ)} On a column with 10 rows of which 4 are null,
        \sqlkw{COUNT}$(x)$ returns: (a)~10; (b)~6; (c)~4; (d)~null.
  \item \textbf{(MCQ)} A condition on \sqlkw{COUNT(*)} must go in:
        (a)~\sqlkw{WHERE}; (b)~\sqlkw{HAVING}; (c)~\sqlkw{ORDER BY};
        (d)~either.
  \item \textbf{(Short)} Why can a bare column in a grouped query not have a
        correct value, and what does SQLite print instead?
  \item \textbf{(Short)} What does \sqlkw{TEMP B-TREE FOR GROUP BY} in a plan
        mean, and what removes it?
  \item \textbf{(Short)} Give four rows with a tie and show how
        \sqlkw{RANK}, \sqlkw{DENSE\_RANK} and \sqlkw{ROW\_NUMBER} number them.
  \item \textbf{(Applied)} Write the query giving each programme's average
        grade points, number of graded enrolments and number of ungraded ones,
        only for programmes with more than 5{,}000 graded enrolments. State
        for each condition which clause it belongs in, and why.
  \item \textbf{(Applied)} The university wants a merit list: the top three
        students by grade points in each programme, ties included. Write it
        with a window function, say whether you chose \sqlkw{RANK} or
        \sqlkw{DENSE\_RANK} and what the choice means for a student tied in
        third place, and estimate from the measurement how the correlated
        alternative would behave at 16{,}000 students.
\end{reviewq}


\part{Getting the Design Right}
%=====================================================================
% PART V OPENER
%=====================================================================
\thispagestyle{plain}
\structuralfigures{P5.}

\begin{center}
\begin{tcolorbox}[enhanced, width=0.92\textwidth, colback=tealhl!5,
  colframe=tealhl, boxrule=0pt, leftrule=5pt, arc=3pt, top=8pt, bottom=8pt]
\large\itshape\color{tealhl!60!black}
Part~V is the chapter this course exists for, and the three around it. Normalisation has a reputation for being the dull, formal part of databases; it is in fact the only part that tells you your design is wrong \emph{before} the data proves it. The whole of it follows from one idea, met in Chapter~17, which is simpler than its notation suggests.
\end{tcolorbox}
\end{center}

\vspace{6pt}

\dsfigh{%
\begin{tikzpicture}[
  cbox/.style={rounded corners=3pt, draw=tealhl, line width=1pt, fill=tealhl!7,
               text=tealhl!55!black, font=\small\bfseries, align=center,
               minimum width=3.5cm, minimum height=1.0cm},
  hbox/.style={rounded corners=3pt, draw=tealhl, line width=1.8pt, fill=tealhl!16,
               text=tealhl!55!black, font=\small\bfseries, align=center,
               minimum width=3.5cm, minimum height=1.0cm},
  arr/.style={-{Stealth[length=2.4mm]}, tealhl!60, line width=1pt},
  note/.style={font=\scriptsize\itshape, text=tealhl!55!black, align=center,
               fill=white, inner sep=2pt}
]
  \node[cbox] (c0) at (0.00,0)  {Ch.\ 17\\Dependencies};
  \node[hbox] (c1) at (4.30,0)  {Ch.\ 18\\Normal Forms};
  \node[cbox] (c2) at (8.60,0)  {Ch.\ 19\\Decomposition};
  \node[cbox] (c3) at (12.90,0) {Ch.\ 20\\Views and Triggers};
  \draw[arr] (c0) -- (c1);
  \draw[arr] (c1) -- (c2);
  \draw[arr] (c2) -- (c3);
  \node[note] at (2.15,-1.0) {one idea};
  \node[note] at (6.45,-1.0) {four forms from it};
  \node[note] at (10.75,-1.0) {and when to stop};
  \node[font=\scriptsize\itshape, text=accent, align=center] at (4.30,1.15)
       {where the database you built in Part~III turns out to have been
        wrong};
\end{tikzpicture}}%
{Reading path through Part~V. Chapter 18 breaks the running database and
rebuilds it, and Chapter 19 measures what that cost --- because normalising is
a trade and a course that presents it as a pure good has not finished the
argument.}{part5map}

\newpage

%=====================================================================
\chapter{Functional Dependencies}\label{ch:fds}
%=====================================================================
\locator{5}{Part V --- Getting the Design Right}

\begin{outcomes}
By the end of this chapter you will be able to:
\begin{tight}
  \item \textbf{Define} a functional dependency and \textbf{read} one written
        as \fd{X}{Y}.
  \item \textbf{State} Armstrong's three axioms and \textbf{derive} new
        dependencies from given ones.
  \item \textbf{Compute} the closure $X^{+}$ of a set of attributes, step by
        step.
  \item \textbf{Use} attribute closure to test whether a set of attributes is
        a superkey or a candidate key.
  \item \textbf{Reduce} a set of dependencies to a minimal cover.
\end{tight}
\end{outcomes}

\begin{prereq}
\chref{relational} for keys. Appendix~B for sets and the meaning of
``function'' --- the Prerequisite Self-Check's question~12 is the whole
foundation of this chapter, and if it was uncomfortable, read Appendix~B
first, not during.
\end{prereq}

\begin{keyterms}
functional dependency $\bullet$ determinant $\bullet$ dependent $\bullet$
trivial dependency $\bullet$ Armstrong's axioms $\bullet$ reflexivity
$\bullet$ augmentation $\bullet$ transitivity $\bullet$ union rule $\bullet$
decomposition rule $\bullet$ attribute closure $\bullet$ $X^{+}$ $\bullet$
superkey $\bullet$ candidate key $\bullet$ minimal
cover $\bullet$ extraneous attribute $\bullet$ redundant dependency
\end{keyterms}

\section{One Table That Everyone Has Built}

Every university has a spreadsheet like this one. It is the first thing a
department builds when it is asked for ``all the results in one place'', and it
is the table this Part is about:

\medskip
\relation{report}{roll\_no, name, programme, course\_code, title,
credit\_hours, dept\_code, dept\_name, semester, section\_no, room, emp\_id,
instructor, grade}

\medskip
Fourteen columns, one row per enrolment, and on the seeded data 60{,}037 rows.
It answers every question in Part~IV without a single join, which is why
people build it. \chref{normalforms} measures what it costs. This chapter
supplies the vocabulary for saying \emph{why} it costs anything --- and the
vocabulary is one idea.

\begin{definitionbox}[Functional dependency]
A \term{functional dependency} \fd{X}{Y}, read ``$X$ determines $Y$'', holds in
a relation when any two tuples that agree on $X$ must also agree on $Y$.

\smallskip
$X$ is the \term{determinant}; $Y$ is the \term{dependent}. Both are sets of
attributes.

\smallskip
It is the word ``function'' in its mathematical sense, from Appendix~B: give
me the value of $X$ and there is exactly one value of $Y$ it can be paired
with. \fd{roll\_no}{name} says a roll number has exactly one name.
\end{definitionbox}

\begin{examplebox}[The dependencies in \textit{report}, found by asking]
As with keys in \chref{ermodel}, these come from the university, not from the
data. Each one is a rule somebody can confirm:

\medskip
\begin{tabular}{@{}L{6.6cm}L{7.4cm}@{}}
\textbf{Dependency} & \textbf{Because} \\
\fd{roll\_no}{name, programme} & a student has one name and is in one
  programme \\
\fd{course\_code}{title, credit\_hours, dept\_code} & a course is defined
  once \\
\fd{dept\_code}{dept\_name} & a department has one name \\
\fd{course\_code, semester, section\_no}{room, emp\_id} & a section meets in
  one room with one instructor \\
\fd{emp\_id}{instructor} & an employee has one name \\
\fd{roll\_no, course\_code, semester, section\_no}{grade} & one grade per
  enrolment \\
\end{tabular}

\smallskip
Six rules. Notice that each one is a sentence the examinations officer would
agree with immediately --- and that each one says the dependent is stored in
\textit{report} \emph{once per enrolment}, when it is one fact. Ayesha's name
is in the table once for every section she has ever taken.
\end{examplebox}

\begin{alertbox}[An instance can disprove a dependency, never prove one]
\chref{relational} said this about keys, and it is the same fact, because a
key is a dependency (below). On the seeded data, \textit{name} happens to
determine \textit{programme} for some names. That is coincidence, not a rule
--- two students called Hira Nadeem can be in different programmes, and the
next admission may be the one that proves it.

\smallskip
A query \emph{can} disprove a dependency, and it is worth knowing:

\begin{lstlisting}[language=sql]
-- Does course_code -> title hold in this data?
SELECT course_code FROM report
 GROUP BY course_code HAVING COUNT(DISTINCT title) > 1;
-- Any row returned is a counter-example. None returned proves nothing.
\end{lstlisting}
\end{alertbox}

\section{Deriving New Dependencies}

\begin{definitionbox}[Armstrong's axioms]
For sets of attributes $X$, $Y$, $Z$:

\smallskip
\textbf{Reflexivity.} If $Y \subseteq X$ then \fd{X}{Y}. (A set determines its
own parts. Such a dependency is \term{trivial}.)

\smallskip
\textbf{Augmentation.} If \fd{X}{Y} then \fd{XZ}{YZ}. (Adding the same
attributes to both sides keeps it true.)

\smallskip
\textbf{Transitivity.} If \fd{X}{Y} and \fd{Y}{Z} then \fd{X}{Z}.

\smallskip
These three are \emph{sound} (they never derive anything false) and
\emph{complete} (they derive everything that follows). Two derived rules save
time: \textbf{union} --- \fd{X}{Y} and \fd{X}{Z} give \fd{X}{YZ}; and
\textbf{decomposition} --- \fd{X}{YZ} gives \fd{X}{Y} and \fd{X}{Z}.
\end{definitionbox}

\begin{examplebox}[One transitive dependency, and why it is the important one]
From the table:
\begin{tight}
  \item \fd{course\_code}{dept\_code} \hfill (given, by decomposition)
  \item \fd{dept\_code}{dept\_name} \hfill (given)
  \item therefore \fd{course\_code}{dept\_name} \hfill (transitivity)
\end{tight}

\smallskip
Nobody wrote that third rule down, and it is true. It is also the one that
causes trouble: \textit{dept\_name} depends on \textit{course\_code} only
\emph{through} \textit{dept\_code}. Renaming the department means changing
every row of every course it offers --- \chref{normalforms} calls that a
transitive dependency and third normal form forbids it.
\end{examplebox}

\section{Attribute Closure}

Deriving dependencies one axiom at a time is how proofs are written. The
practical question is narrower: \emph{given $X$, what does it determine?} There
is an algorithm, and it is the most useful thing in Part~V.

\begin{definitionbox}[Attribute closure $X^{+}$]
The \term{closure} $X^{+}$ of a set of attributes $X$ under a set of
dependencies $F$ is the set of every attribute that $X$ determines.

\smallskip
\textbf{Algorithm.}
\begin{tightnum}
  \item Start with $X^{+} = X$.
  \item Repeat: for each dependency \fd{A}{B} in $F$, if $A \subseteq X^{+}$,
        add $B$ to $X^{+}$.
  \item Stop when a full pass adds nothing.
\end{tightnum}

\smallskip
It always terminates --- there are finitely many attributes and $X^{+}$ only
grows.
\end{definitionbox}

\begin{worked}[$\{\mathit{course\_code}\}^{+}$, then the key of \textit{report}]
Using the six dependencies from the examplebox, abbreviated
$F$.

\medskip
\textbf{Part A --- what does \textit{course\_code} determine?}

\smallskip
\begin{tabular}{@{}L{1.4cm}L{6.0cm}L{6.6cm}@{}}
\textbf{pass} & \textbf{dependency used} & \textbf{$X^{+}$ is now} \\
start & --- & \{course\_code\} \\
1 & \fd{course\_code}{title, credit\_hours, dept\_code} & \{course\_code,
  title, credit\_hours, dept\_code\} \\
1 & \fd{dept\_code}{dept\_name} --- its left side is now inside & \dots\ plus
  dept\_name \\
2 & nothing new applies & \textbf{stop} \\
\end{tabular}

\smallskip
$\{\mathit{course\_code}\}^{+} = \{$course\_code, title, credit\_hours,
dept\_code, dept\_name$\}$. Five of the fourteen. \textit{course\_code} is not a
key of \textit{report}.

\medskip
\textbf{Part B --- the key.} Try $K = \{$roll\_no, course\_code, semester,
section\_no$\}$:
\begin{tight}
  \item \fd{roll\_no}{name, programme} adds 2;
  \item \fd{course\_code}{title, credit\_hours, dept\_code} adds 3, and then
        \fd{dept\_code}{dept\_name} adds 1;
  \item \fd{course\_code, semester, section\_no}{room, emp\_id} adds 2, and
        then \fd{emp\_id}{instructor} adds 1;
  \item \fd{roll\_no, course\_code, semester, section\_no}{grade} adds 1.
\end{tight}
$4 + 2 + 3 + 1 + 2 + 1 + 1 = 14$. $K^{+}$ is every attribute: \textbf{$K$ is a
superkey}.

\medskip
\textbf{Part C --- is it a candidate key?} Remove each attribute in turn and
recompute. Without \textit{roll\_no}, nothing determines \textit{name}; without
\textit{course\_code}, nothing determines \textit{title}; without
\textit{semester} or \textit{section\_no}, nothing determines \textit{room}. No
proper subset is a superkey, so \textbf{$K$ is a candidate key} --- and it is
exactly \chref{mapping}'s primary key of \textit{enrolment}, found here from
the dependencies alone.
\end{worked}

\dsfig{%
\begin{tikzpicture}[
  atr/.style={rounded corners=2pt, draw=#1, line width=0.8pt, fill=#1!10,
             text=#1!60!black, font=\scriptsize\ttfamily, inner sep=2.5pt,
             minimum height=0.48cm},
  ps/.style={font=\scriptsize\bfseries, text=primary, anchor=east},
  rule/.style={font=\scriptsize\itshape, text=neutral, anchor=west},
  arr/.style={-{Stealth[length=2mm]}, primary!45, line width=0.8pt}
]
  % start
  \node[ps] at (-0.25,1.8) {start};
  \node[atr=primary, anchor=west] (s0) at (0,1.8) {course\_code};

  % pass 1, first rule
  \node[ps] at (-0.25,0.9) {pass 1};
  \node[atr=primary, anchor=west] (a1) at (0,0.9) {course\_code};
  \node[atr=secondary, anchor=west] (a2) at (a1.east) {title};
  \node[atr=secondary, anchor=west] (a3) at (a2.east) {credit\_hours};
  \node[atr=secondary, anchor=west] (a4) at (a3.east) {dept\_code};
  \node[rule] at (7.6,0.9) {\fd{course\_code}{title, credit\_hours,
       dept\_code}};

  % pass 1, second rule, now applicable
  \node[ps] at (-0.25,0.0) {pass 1};
  \node[atr=primary, anchor=west] (b1) at (0,0.0) {course\_code};
  \node[atr=secondary, anchor=west] (b2) at (b1.east) {title};
  \node[atr=secondary, anchor=west] (b3) at (b2.east) {credit\_hours};
  \node[atr=secondary, anchor=west] (b4) at (b3.east) {dept\_code};
  \node[atr=greenhl, anchor=west] (b5) at (b4.east) {dept\_name};
  \node[rule] at (7.6,0.0) {\fd{dept\_code}{dept\_name} --- its left side
       is now inside};

  % pass 2
  \node[ps] at (-0.25,-0.9) {pass 2};
  \node[font=\scriptsize\itshape, text=accent, anchor=west] at (0,-0.9)
       {no dependency adds anything new $\Rightarrow$ \textbf{stop}};

  \draw[arr] (s0.south) ++(0,-0.02) -- ([yshift=1pt]a1.north);
  \draw[arr] (a1.south) ++(0,-0.02) -- ([yshift=1pt]b1.north);
\end{tikzpicture}}%
{$\{\mathit{course\_code}\}^{+}$, one step per row. Each step applies a
dependency whose left side is already inside the set; the second only becomes
applicable because the first added \textit{dept\_code}. The closure stops at
five attributes --- fewer than fourteen, so \textit{course\_code} is not a key
of \textit{report}.}{closure}

\begin{conceptbox}[Why closure is the tool for everything in Part V]
Three questions, one algorithm:

\begin{tight}
  \item \textbf{Is $X$ a superkey?} Compute $X^{+}$; it is if $X^{+}$ is every
        attribute.
  \item \textbf{Does \fd{X}{Y} follow from $F$?} Compute $X^{+}$; it does if
        $Y \subseteq X^{+}$. No axioms needed.
  \item \textbf{Is an attribute in a dependency redundant?} Remove it and
        recompute; if the closure is unchanged, it was. That is the minimal
        cover, next.
\end{tight}

\smallskip
Every normal form in \chref{normalforms} is a test you run with closures, and
\chref{decomposition}'s lossless-join test is one too.
\end{conceptbox}

\section{Minimal Cover}

\begin{definitionbox}[Minimal cover]
A \term{minimal cover} of $F$ is an equivalent set of dependencies --- one with
the same closure --- in which:
\begin{tightnum}
  \item every right-hand side is a single attribute;
  \item no left-hand side has an \term{extraneous attribute} (one that could be
        removed without changing what the dependency implies);
  \item no dependency is \term{redundant} (derivable from the others).
\end{tightnum}

\smallskip
It is the smallest statement of the same rules. \chref{normalforms}'s 3NF
synthesis algorithm starts from one.
\end{definitionbox}

\begin{worked}[A minimal cover, in three steps]
$F = \{$ \fd{A}{BC}, \fd{B}{C}, \fd{AB}{D} $\}$.

\medskip
\textbf{Step 1 --- single right-hand sides.} Decompose:
\fd{A}{B}, \fd{A}{C}, \fd{B}{C}, \fd{AB}{D}.

\medskip
\textbf{Step 2 --- extraneous left-hand attributes.} Only \fd{AB}{D} has two.
Is $B$ extraneous? Compute $\{A\}^{+}$: $A \rightarrow B$ gives $B$, so
$\{A\}^{+} \ni B$, and then \fd{AB}{D} applies. $\{A\}^{+} = \{A,B,C,D\}$,
which contains $D$. So \fd{A}{D} already holds: \textbf{$B$ is extraneous},
and \fd{AB}{D} becomes \fd{A}{D}.

\medskip
\textbf{Step 3 --- redundant dependencies.} Is \fd{A}{C} redundant? Remove it
and compute $\{A\}^{+}$ from the rest: $A \to B$, $B \to C$. It contains $C$.
\textbf{Redundant} --- it follows by transitivity. Remove it.

\medskip
\textbf{Result.} $\{$ \fd{A}{B}, \fd{B}{C}, \fd{A}{D} $\}$. Three dependencies
instead of four, each as small as it can be, and the same closure as the
original. Check: $\{A\}^{+} = \{A,B,C,D\}$ in both.
\end{worked}

\begin{pitfall}
\begin{tight}
  \item \textbf{Reading a dependency off sample data.} Data disproves; it never
        proves. Ask.
  \item \textbf{Reading \fd{X}{Y} backwards.} \fd{roll\_no}{name} does not
        give \fd{name}{roll\_no}; two students share a name.
  \item \textbf{Stopping closure after one pass.} A dependency whose left side
        becomes available in pass 1 must still be applied. Repeat until a pass
        adds nothing.
  \item \textbf{Calling a superkey a candidate key.} Test every proper subset.
  \item \textbf{Splitting a left-hand side.} \fd{AB}{C} does \emph{not} give
        \fd{A}{C}. Only right-hand sides decompose.
  \item \textbf{Forgetting the transitive dependencies nobody wrote down.}
        Closure finds them; that is what it is for.
\end{tight}
\end{pitfall}

\begin{lab}[Build the report table, then find its dependencies and its key]
\begin{lstlisting}[language=bash]
# --- 1. The seeded database ----------------------------------------
python code/schema/seed.py university.db
sqlite3 university.db
\end{lstlisting}
\begin{lstlisting}[language=sql]
-- --- 2. Build the table everybody builds ------------------------
CREATE TABLE report AS
SELECT s.roll_no, s.name, s.programme,
       c.course_code, c.title, c.credit_hours,
       d.dept_code, d.name AS dept_name,
       e.semester, e.section_no, x.room, x.emp_id,
       i.name AS instructor, e.grade
  FROM enrolment e
  JOIN student    s ON s.roll_no     = e.roll_no
  JOIN course     c ON c.course_code = e.course_code
  JOIN department d ON d.dept_code   = c.dept_code
  JOIN section    x ON x.course_code = e.course_code
                   AND x.semester    = e.semester
                   AND x.section_no  = e.section_no
  LEFT JOIN instructor i ON i.emp_id = x.emp_id;
SELECT COUNT(*) FROM report;

-- --- 3. Test a dependency that should hold ----------------------
SELECT course_code FROM report
 GROUP BY course_code HAVING COUNT(DISTINCT title) > 1;
-- No rows: no counter-example. Does that PROVE course_code -> title?

-- --- 4. Test one that should not --------------------------------
SELECT name FROM report
 GROUP BY name HAVING COUNT(DISTINCT roll_no) > 1 LIMIT 5;
-- Counter-examples: name does not determine roll_no. Settled.

-- --- 5. Count the copies of one fact ----------------------------
SELECT roll_no, name, COUNT(*) AS copies FROM report
 GROUP BY roll_no ORDER BY copies DESC LIMIT 3;
-- One student's name, stored this many times. Why?

-- --- 6. The transitive one -------------------------------------
SELECT dept_code, COUNT(*) AS rows_holding_its_name
  FROM report GROUP BY dept_code;
-- To rename one department, how many rows must change?
\end{lstlisting}
\begin{lstlisting}[language=bash]
# --- 7. On paper -----------------------------------------------------
# Compute {course_code}+ , {roll_no}+ and {emp_id}+ by the algorithm,
# one pass per line, from the six dependencies in this chapter.
# Then prove {roll_no, course_code, semester, section_no} is a
# candidate key: superkey first, then every proper subset.
\end{lstlisting}

\smallskip
\textbf{What to notice.} Steps 5 and 6 are the whole motivation for Part~V. One
fact --- a name, a department name --- is stored once for every enrolment it
touches, because the dependency that determines it has a smaller left side than
the table's key. Keep \textit{report}: \chref{normalforms} measures it against
the normalised schema, and \chref{decomposition} measures it again.
\end{lab}

\begin{checkpoint}
\begin{tightnum}
  \item Compute $\{\mathit{emp\_id}\}^{+}$ and
        $\{\mathit{course\_code}, \mathit{semester}, \mathit{section\_no}\}^{+}$
        under the six dependencies of \textit{report}, showing each pass.
  \item Using closure only --- no axioms --- decide whether
        \fd{roll\_no, course\_code}{dept\_name} follows from the six
        dependencies, and show the closure that settles it.
  \item Find a minimal cover of
        $\{$ \fd{A}{B}, \fd{A}{C}, \fd{B}{C}, \fd{AC}{D} $\}$, showing all
        three steps and checking that the result has the same closure of $A$.
\end{tightnum}
\end{checkpoint}

\begin{chaptersummary}
\begin{tight}
  \item \fd{X}{Y} holds when tuples agreeing on $X$ must agree on $Y$. It is a
        function, in Appendix~B's sense, from $X$-values to $Y$-values.
  \item Dependencies come from the university. An instance can disprove one
        with a \sqlkw{GROUP BY} \dots\ \sqlkw{HAVING} \sqlkw{COUNT(DISTINCT)}
        query; it can never prove one.
  \item \textbf{Armstrong's axioms} --- reflexivity, augmentation,
        transitivity --- are sound and complete. Union and decomposition
        follow from them.
  \item Only right-hand sides decompose. \fd{AB}{C} does not give
        \fd{A}{C}.
  \item The \textbf{closure} $X^{+}$ is everything $X$ determines: start with
        $X$, apply every dependency whose left side is inside, repeat until a
        pass adds nothing.
  \item Closure answers three questions: is $X$ a superkey, does a dependency
        follow, is an attribute redundant.
  \item The key of \textit{report} is \{roll\_no, course\_code, semester,
        section\_no\} --- the primary key of \textit{enrolment}, found from
        the dependencies alone.
  \item A \textbf{minimal cover} has single-attribute right sides, no
        extraneous left-side attributes, and no redundant dependencies.
  \item In \textit{report}, every fact whose determinant is smaller than the
        key is stored once per enrolment. That is what \chref{normalforms}
        removes, and measures.
\end{tight}
\end{chaptersummary}

\begin{reviewq}
  \item \textbf{(MCQ)} \fd{AB}{C} and \fd{C}{D} together imply:
        (a)~\fd{A}{D}; (b)~\fd{AB}{D}; (c)~\fd{D}{AB}; (d)~\fd{C}{AB}.
  \item \textbf{(MCQ)} A dependency \fd{X}{Y} with $Y \subseteq X$ is:
        (a)~transitive; (b)~trivial; (c)~redundant; (d)~a key.
  \item \textbf{(Short)} Why can a query find a counter-example to a
        dependency but never prove one holds?
  \item \textbf{(Short)} State the closure algorithm in three lines, and say
        why it must terminate.
  \item \textbf{(Short)} Show by a two-row example that \fd{AB}{C} does not
        imply \fd{A}{C}.
  \item \textbf{(Applied)} A relation \relation{booking}{room, date, slot,
        course\_code, emp\_id, capacity} has the rules: a room has a fixed
        capacity; a room, date and slot is booked for one course; a course has
        one instructor per booking. Write the dependencies, compute the closure
        of \{room, date, slot\}, and find every candidate key.
  \item \textbf{(Applied)} Using the six dependencies of \textit{report}, list
        every attribute whose determinant is a proper subset of the key, and for
        each one compute from the seeded data (step~5 of the laboratory) how
        many copies of a single fact it causes. Which one costs the most rows?
\end{reviewq}

%=====================================================================
\chapter{Normal Forms: 1NF to BCNF}\label{ch:normalforms}
%=====================================================================
\locator{5}{Part V --- Getting the Design Right}

\begin{outcomes}
By the end of this chapter you will be able to:
\begin{tight}
  \item \textbf{Name} the three anomalies --- update, insertion, deletion ---
        and \textbf{show} each one in a real table.
  \item \textbf{State} the conditions for 1NF, 2NF, 3NF and BCNF, and
        \textbf{test} a relation against each using its dependencies.
  \item \textbf{Identify} a partial and a transitive dependency, and
        \textbf{say} which normal form each one violates.
  \item \textbf{Decompose} a relation into BCNF by splitting on a violating
        dependency.
  \item \textbf{Weigh} what normalisation costs in reads against what it saves
        in correctness, using measured figures.
\end{tight}
\end{outcomes}

\begin{prereq}
\chref{fds}: functional dependencies, attribute closure and candidate keys ---
every test in this chapter is a closure. The \textit{report} table from that
chapter's laboratory. You need \texttt{code/m18\_normalforms.py}.
\end{prereq}

\begin{keyterms}
normalisation $\bullet$ anomaly $\bullet$ update anomaly $\bullet$ insertion
anomaly $\bullet$ deletion anomaly $\bullet$ first normal form $\bullet$
atomic $\bullet$ partial dependency $\bullet$ second normal form $\bullet$
transitive dependency $\bullet$ third normal form $\bullet$ prime attribute
$\bullet$ Boyce--Codd normal form $\bullet$ decomposition $\bullet$ redundancy
\end{keyterms}

\section{Three Ways a Wide Table Goes Wrong}

\chref{fds} built \textit{report}: fourteen columns, one row per enrolment,
every fact the university holds about it in one place. It answers questions
without joins. This chapter is about what it does when the data
\emph{changes}, which is the thing it was never designed for.

\begin{definitionbox}[The three anomalies]
An \term{update anomaly}: one fact is stored in many rows, so changing it means
changing all of them --- and changing some of them leaves the table holding two
versions of the truth.

\smallskip
An \term{insertion anomaly}: a fact cannot be recorded without also recording
an unrelated one that does not exist yet.

\smallskip
A \term{deletion anomaly}: deleting one fact removes another that had nowhere
else to live.

\smallskip
All three have one cause. \textbf{A fact whose determinant is smaller than the
key is stored once per key value}, so the table is holding several relations
at once and can only change them together.
\end{definitionbox}

\begin{measured}[The same information, held wide and held normalised]
\textsf{\footnotesize \texttt{code/m18\_normalforms.py}, seeded database
(4{,}000 students), SQLite 3.51.2, Intel i5-6600, Windows 11. Best of three
runs where timed.}

\medskip
\textbf{1. Bytes, and copies of one fact}

\begin{tabular}{@{}L{6.6cm}R{2.0cm}R{2.6cm}L{2.8cm}@{}}
& \textbf{rows} & \textbf{bytes} & \\
\textit{report}, as \sqlkw{CREATE TABLE AS} made it & 60{,}037 & 6{,}860{,}800 & no index \\
\textit{report}, with its key indexed & 60{,}037 & 9{,}289{,}728 & like for like \\
normalised, 6 tables & 65{,}693 & 5{,}574{,}656 & key indexes \\
\end{tabular}

\smallskip
\textbf{Like for like, the wide table is 1.7$\times$ the size.} Compared as
created it looks like 1.2$\times$, because \sqlkw{CREATE TABLE AS} copies rows
and nothing else --- no key, no index --- while every normalised table carries
the index its \sqlkw{PRIMARY KEY} implies. The bytes are not the real cost;
the copies are:

\begin{tabular}{@{}L{6.6cm}R{2.6cm}R{2.6cm}@{}}
\textbf{copies of one fact} & \textbf{\textit{report}} & \textbf{normalised} \\
a student's name & 18 & 1 \\
an instructor's name & 949 & 1 \\
a course's title & 1{,}382 & 1 \\
a department's name & \textbf{7{,}936} & 1 \\
\end{tabular}

\medskip
\textbf{2. The update anomaly} --- rename one department

\begin{tabular}{@{}L{6.6cm}R{2.6cm}R{2.6cm}@{}}
& \textbf{rows changed} & \textbf{ms} \\
\textit{report} & 7{,}869 & 378.9 \\
\textit{department} & 1 & 25.4 \\
\end{tabular}

\smallskip
Most of the 25.4\,ms is the commit, not the change: one row costs about as
much to make durable as it does to write.

\medskip
\textbf{3. A half-finished update} --- the IT department renamed in the
FALL2025 rows only, as a crashed program or a wrong \sqlkw{WHERE} would leave
it

\begin{tabular}{@{}L{6.6cm}R{2.6cm}@{}}
\texttt{'Information Technology'} & 5{,}273 rows \\
\texttt{'Information Tech.'} & 2{,}663 rows \\
\end{tabular}

\smallskip
Two names for one department, in one table, and no constraint anywhere that
could have refused it. In \textit{department} the same mistake is impossible:
there is only one place to write the name.

\medskip
\textbf{4. Insertion and deletion}

\begin{tabular}{@{}L{8.0cm}L{6.0cm}@{}}
a new course, nobody enrolled yet: rows in \textit{report} that can hold it &
  \textbf{0} \\
the same course in \textit{course} & one row, accepted \\
\textit{report} after deleting the 1{,}122 enrolments of PH-350: rows still
  describing the course & \textbf{0} \\
the same course in \textit{course} & still there \\
\end{tabular}

\medskip
\textbf{5. Read cost} --- where the wide table is supposed to win

\begin{tabular}{@{}L{8.0cm}R{2.4cm}L{3.6cm}@{}}
\textbf{(a) one student's transcript} & \textbf{ms} & \\
\textit{report}, no index on roll\_no & 15.35 & \textbf{162$\times$ slower} \\
normalised, one join & 0.09 & \\
\textit{report}, with an index on roll\_no & 0.08 & a tie \\
\textbf{(b) enrolments per department} & & \\
\textit{report}, no joins & 66.51 & \textbf{1.5$\times$ faster} \\
normalised, two joins & 98.34 & \\
\end{tabular}

\smallskip
\textbf{The wide table lost the lookup by two orders of magnitude}, and not
because of a join: the normalised \textit{enrolment} table's primary key
begins with \texttt{roll\_no}, so the transcript was an index search, while
\textit{report} had no index and was read end to end. Given one index the two
tied. It won only the full-scan report, by 1.5$\times$ --- the case where
reading one table instead of joining three is a real saving.
\end{measured}

\begin{alertbox}[So the trade is not ``fast reads against slow writes'']
That is how it is usually described, and the measurement does not support it:

\begin{tight}
  \item \textbf{Lookups were not faster wide.} They were 162$\times$ slower
        until an index was added, and equal after. ``No joins'' is not an
        advantage when the join is an index search.
  \item \textbf{Full scans were faster wide, by 1.5$\times$.} Real, and modest.
  \item \textbf{Writes were 7{,}869 times the work}, and one half-finished
        write left the table permanently self-contradictory.
\end{tight}

\smallskip
The honest summary: normalisation costs a little on reports over every row,
and buys correctness that no amount of care in the application can supply.
\chref{decomposition} asks when the 1.5$\times$ is worth paying for, and the
answer is a measurement, not a preference.
\end{alertbox}

\section{The Normal Forms}

Each normal form forbids one kind of dependency. They are nested: every BCNF
relation is in 3NF, every 3NF relation in 2NF, and every 2NF relation in 1NF.

\begin{definitionbox}[1NF and 2NF]
\textbf{First normal form (1NF):} every attribute value is atomic --- no
lists, no repeating groups, no ``three phone numbers in one field''. It is
\chref{relational}'s fourth property of a relation, and \chref{mapping}'s rule
for multi-valued attributes.

\smallskip
A \term{partial dependency} is one whose determinant is a \emph{proper subset}
of a candidate key: \fd{roll\_no}{name}, when the key is \{roll\_no,
course\_code, semester, section\_no\}.

\smallskip
\textbf{Second normal form (2NF):} 1NF, and no non-prime attribute is partially
dependent on any candidate key. (A \term{prime attribute} is one that belongs
to some candidate key.)
\end{definitionbox}

\begin{definitionbox}[3NF and BCNF]
A \term{transitive dependency} is \fd{K}{Y} holding only because \fd{K}{X} and
\fd{X}{Y}, where $X$ is not a superkey: \fd{course\_code}{dept\_name} through
\textit{dept\_code}.

\smallskip
\textbf{Third normal form (3NF):} for every non-trivial \fd{X}{A}, either $X$
is a superkey, or $A$ is a prime attribute.

\smallskip
\textbf{Boyce--Codd normal form (BCNF):} for every non-trivial \fd{X}{A},
$X$ is a superkey. Full stop --- the prime-attribute exception is gone.

\smallskip
\textbf{The one sentence to remember:} in BCNF, every determinant is a
superkey. Every fact is about the key, the whole key, and nothing but the key.
\end{definitionbox}

\begin{worked}[\textit{report}, from nowhere to BCNF]
Key $K = \{$roll\_no, course\_code, semester, section\_no$\}$ (\chref{fds}).
Test each dependency: is its determinant a superkey?

\medskip
\begin{tabular}{@{}L{6.2cm}L{2.5cm}L{5.3cm}@{}}
\textbf{dependency} & \textbf{superkey?} & \textbf{violates} \\
\fd{roll\_no}{name, programme} & no & 2NF --- partial \\
\fd{course\_code}{title, credit\_hours, dept\_code} & no & 2NF --- partial \\
\fd{dept\_code}{dept\_name} & no & 3NF --- transitive \\
\fd{course\_code, semester, section\_no}{room, emp\_id} & no & 2NF --- partial \\
\fd{emp\_id}{instructor} & no & 3NF --- transitive \\
\fd{$K$}{grade} & \textbf{yes} & nothing \\
\end{tabular}

\medskip
Five violations, one clean dependency. \textbf{The decomposition procedure:}
take a violating \fd{X}{Y}; split $R$ into $R_1 = X \cup Y$ and $R_2 = R - Y$;
repeat on whatever still violates. \textbf{Split on the transitive
dependencies first} --- the box below shows what goes wrong otherwise.

\medskip
\begin{tightnum}
  \item On \fd{dept\_code}{dept\_name}: \relation{department}{\relkey{dept\_code},
        dept\_name}. Remove dept\_name from the rest.
  \item On \fd{emp\_id}{instructor}: \relation{instructor}{\relkey{emp\_id},
        instructor}. Remove instructor.
  \item On \fd{roll\_no}{name, programme}: \relation{student}{\relkey{roll\_no},
        name, programme}. Remove name and programme.
  \item On \fd{course\_code}{title, credit\_hours, dept\_code}:
        \relation{course}{\relkey{course\_code}, title, credit\_hours,
        dept\_code}. Remove title, credit\_hours and dept\_code.
  \item On \fd{course\_code, semester, section\_no}{room, emp\_id}:
        \relation{section}{\relkey{course\_code}, \relkey{semester},
        \relkey{section\_no}, room, emp\_id}. Remove room and emp\_id.
  \item What remains: \relation{enrolment}{\relkey{roll\_no},
        \relkey{course\_code}, \relkey{semester}, \relkey{section\_no},
        grade} --- whose only dependency has the key as its determinant.
\end{tightnum}

\medskip
\textbf{Why the order matters.} Split on \fd{course\_code}{title, credit\_hours,
dept\_code} \emph{first} and $R_2 = R - Y$ removes \textit{dept\_code} from the
remainder --- but leaves \textit{dept\_name} there, with nothing beside it that
determines it except, transitively, \textit{course\_code}. The next split is
then on \fd{course\_code}{dept\_name}, and produces
\relation{}{\relkey{course\_code}, dept\_name}: a department's name stored
once \emph{per course}, which is the redundancy you were removing. Every
relation is still in BCNF, so no normal-form test fails --- but the department
has no table of its own, and \fd{dept\_code}{dept\_name} now lives in
\emph{no} single table, so nothing can enforce it: two CS courses could carry
two different department names. \chref{decomposition} calls that losing
dependency preservation. Handle the deepest dependency first and each fact
lands with its own determinant.

\medskip
\textbf{Six relations, every one in BCNF --- and they are exactly the six
tables \chref{mapping} produced from the ER diagram.} Two routes, one from
design and one from dependencies, arrived at the same schema. That agreement
is the best evidence a design is right.
\end{worked}

\dsfig{%
\begin{tikzpicture}[
  wide/.style={rounded corners=3pt, draw=accent, line width=1.1pt,
               fill=accent!7, text=accent!70!black, font=\scriptsize,
               align=left, text width=4.4cm, inner sep=4pt},
  tb/.style={rounded corners=2pt, draw=#1, line width=0.9pt, fill=#1!7,
             text=#1!60!black, font=\scriptsize, align=left,
             text width=4.5cm, inner sep=3pt},
  arr/.style={-{Stealth[length=2.2mm]}, primary!50, line width=0.9pt},
  lbl/.style={font=\scriptsize\itshape, text=neutral, align=center}
]
  \node[wide] (r) at (-4.9,0) {\textbf{report} --- 14 attributes\\[2pt]
       roll\_no, name, programme, course\_code, title, credit\_hours,
       dept\_code, dept\_name, semester, section\_no, room, emp\_id,
       instructor, grade\\[3pt]
       \textbf{60{,}037 rows}; a department's name stored
       \textbf{7{,}936} times};

  \node[tb=primary]  (s) at (2.6, 2.25) {\textbf{student}(\relkey{roll\_no},
       name, programme)};
  \node[tb=primary]  (c) at (2.6, 1.35) {\textbf{course}(\relkey{course\_code},
       title, credit\_hours, dept\_code)};
  \node[tb=greenhl]  (d) at (2.6, 0.45) {\textbf{department}(\relkey{dept\_code},
       dept\_name)};
  \node[tb=purplehl] (x) at (2.6,-0.45) {\textbf{section}(\relkey{course\_code},
       \relkey{semester}, \relkey{section\_no}, room, emp\_id)};
  \node[tb=greenhl]  (i) at (2.6,-1.35) {\textbf{instructor}(\relkey{emp\_id},
       instructor)};
  \node[tb=accent]   (e) at (2.6,-2.25) {\textbf{enrolment}(\relkey{roll\_no},
       \relkey{course\_code}, \relkey{semester}, \relkey{section\_no},
       grade)};

  \foreach \t in {s,c,x,e}{\draw[arr] (r.east) -- (\t.west);}
  \draw[arr, greenhl!60] (c.east) .. controls (5.4,1.35) and (5.4,0.45)
       .. (d.east);
  \draw[arr, greenhl!60] (x.east) .. controls (5.4,-0.45) and (5.4,-1.35)
       .. (i.east);

  \node[lbl, text=primary!70!black] at (-1.2,2.35) {partial: 2NF};
  \node[lbl, text=greenhl!60!black, anchor=west] at (5.6,0.9)
       {transitive:\\3NF};
  \node[lbl, text=greenhl!60!black, anchor=west] at (5.6,-0.9)
       {transitive:\\3NF};
\end{tikzpicture}}%
{\textit{report} decomposed to BCNF. Splitting on each partial dependency
gives four relations; splitting on the two transitive ones gives two more. The
result is the six tables \chref{mapping} derived from the ER diagram by an
entirely different route.}{bcnf}

\begin{conceptbox}[3NF against BCNF: the one case where they differ]
They differ only when a relation has two \emph{overlapping} composite
candidate keys --- rare, but it exists. A classic case:
\relation{teaches}{student, course, instructor}, where each instructor teaches
one course (\fd{instructor}{course}) and each student has one instructor per
course (\fd{student, course}{instructor}).

\smallskip
\fd{instructor}{course}: \textit{instructor} is not a superkey, so BCNF is
violated. But \textit{course} is prime --- it is in the key \{student,
course\} --- so 3NF is satisfied. Decomposing to BCNF gives
\relation{}{\relkey{instructor}, course} and \relation{}{\relkey{student},
\relkey{instructor}}, and \textbf{loses} the dependency
\fd{student, course}{instructor}: no single table can now enforce it.

\smallskip
That trade --- BCNF, or keeping every dependency enforceable --- is
\chref{decomposition}'s subject. For the university schema it never arises.
\end{conceptbox}

\begin{pitfall}
\begin{tight}
  \item \textbf{Thinking normalisation is about saving bytes.} It saved a
        factor of 1.7. It removed 7{,}935 copies of one department name.
        The copies are the point.
  \item \textbf{Believing a wide table reads faster.} Measured: 162$\times$
        slower on a lookup without an index; a tie with one.
  \item \textbf{Comparing a \sqlkw{CREATE TABLE AS} copy with indexed tables.}
        The copy has no index. Measure like for like.
  \item \textbf{Testing 2NF against one key when there are several.} A partial
        dependency on \emph{any} candidate key violates it.
  \item \textbf{Stopping at 3NF without checking BCNF.} They usually coincide;
        check rather than assume.
  \item \textbf{Splitting a relation and losing the join.} Every decomposition
        must keep the determinant in both halves --- \chref{decomposition}
        explains why.
\end{tight}
\end{pitfall}

\begin{lab}[Break the wide table three ways, then decompose it]
\begin{lstlisting}[language=bash]
# --- 1. Seed, and rebuild report from Chapter 17 --------------------
python code/schema/seed.py university.db
sqlite3 university.db
\end{lstlisting}
\begin{lstlisting}[language=sql]
-- --- 2. Recreate report (Chapter 17's lab, step 2), then ---------
SELECT COUNT(*) FROM report;

-- --- 3. The update anomaly -------------------------------------
.timer on
UPDATE report SET dept_name = 'Computing' WHERE dept_code = 'CS';
-- Rows changed? (.changes on, or SELECT changes();)
UPDATE department SET name = 'Computing' WHERE dept_code = 'CS';

-- --- 4. Leave it half done --------------------------------------
UPDATE report SET dept_name = 'Information Tech.'
 WHERE dept_code = 'IT' AND semester = 'FALL2025';
SELECT dept_name, COUNT(*) FROM report
 WHERE dept_code = 'IT' GROUP BY dept_name;
-- Two answers. Write the query a programmer would need to find out
-- which name is current. Can it be written at all?

-- --- 5. The insertion anomaly -----------------------------------
-- Add a course nobody is enrolled in yet, to report.
-- Try it. What would you have to put in roll_no and grade?

-- --- 6. The deletion anomaly ------------------------------------
SELECT course_code, COUNT(*) FROM report
 GROUP BY course_code ORDER BY 2 LIMIT 1;
-- Delete every row for that course. Then ask report its title.

-- --- 7. Read cost, like for like --------------------------------
SELECT course_code, title FROM report
 WHERE roll_no = (SELECT MIN(roll_no) FROM student);
CREATE INDEX idx_report_roll ON report(roll_no);
SELECT course_code, title FROM report
 WHERE roll_no = (SELECT MIN(roll_no) FROM student);
-- Two timings. Which part of the difference was the join?
\end{lstlisting}
\begin{lstlisting}[language=bash]
# --- 8. On paper: decompose report to BCNF ---------------------------
# One violating dependency at a time. Name each new relation, its key,
# and which normal form the split repaired. Compare your six relations
# with Figure 9.1. Then run the measured version:
python code/m18_normalforms.py
\end{lstlisting}

\smallskip
\textbf{What to notice.} Step 4 is the one with no repair. Once a table holds
two versions of a fact and nothing records which came later, no query can tell
you which is true. That is not a performance problem, and no index fixes it.
It is the reason normal forms exist.
\end{lab}

\begin{checkpoint}
\begin{tightnum}
  \item For each of the three anomalies, give the specific operation on
        \textit{report} that exhibits it, the measured consequence, and the
        single dependency responsible.
  \item \relation{R}{\relkey{A}, \relkey{B}, C, D} has \fd{A}{C} and
        \fd{C}{D}. State its normal form, identify each violation by kind, and
        decompose it to BCNF.
  \item The wide table read a transcript 162$\times$ slower than the
        normalised schema and a department report 1.5$\times$ faster. Explain
        both results, and say which one would change if \textit{report} were
        given an index on \texttt{roll\_no}, and to what.
\end{tightnum}
\end{checkpoint}

\begin{chaptersummary}
\begin{tight}
  \item Three anomalies --- \textbf{update}, \textbf{insertion},
        \textbf{deletion} --- share one cause: a fact whose determinant is
        smaller than the key is stored once per key value.
  \item \textbf{Measured:} one department's name was held \textbf{7{,}936}
        times; renaming it changed \textbf{7{,}869} rows against 1; a
        half-finished update left \textbf{two names} for one department.
  \item \textbf{Measured:} a new course could not be inserted into the wide
        table at all, and deleting a course's enrolments deleted the course.
  \item \textbf{Measured:} like for like, the wide table was
        \textbf{1.7$\times$} the size --- 1.2$\times$ only when compared
        without its index.
  \item \textbf{Measured:} the wide table was \textbf{162$\times$ slower} on a
        lookup with no index and tied with one; it was \textbf{1.5$\times$
        faster} only on a full-scan report.
  \item \textbf{1NF}: atomic values. \textbf{2NF}: no partial dependency.
        \textbf{3NF}: no transitive dependency on a non-prime attribute.
        \textbf{BCNF}: every determinant is a superkey.
  \item Decompose on a violating \fd{X}{Y} into $X \cup Y$ and $R - Y$, and
        repeat.
  \item \textit{report} decomposes to exactly the six tables \chref{mapping}
        derived from the ER diagram --- two independent routes, one schema.
  \item 3NF and BCNF differ only with overlapping composite keys, where BCNF
        can cost an enforceable dependency.
\end{tight}
\end{chaptersummary}

\begin{reviewq}
  \item \textbf{(MCQ)} A non-prime attribute depending on part of a composite
        key violates: (a)~1NF; (b)~2NF; (c)~3NF but not 2NF; (d)~BCNF only.
  \item \textbf{(MCQ)} In BCNF, for every non-trivial \fd{X}{A}: (a)~$A$ is
        prime; (b)~$X$ is a superkey; (c)~$X$ is a candidate key;
        (d)~$A$ is in $X$.
  \item \textbf{(Short)} Why can no query repair a table holding two versions
        of the same fact?
  \item \textbf{(Short)} Give the dependency in \textit{report} that is
        transitive, and say which normal form it violates.
  \item \textbf{(Short)} Why was the wide table 1.2$\times$ the size as
        created but 1.7$\times$ like for like?
  \item \textbf{(Applied)} A spreadsheet records
        \relation{loan}{member\_id, member\_name, book\_id, title, author,
        date\_out, date\_back}, one row per loan. Write its dependencies,
        identify its key, classify every violation, and decompose it to BCNF,
        naming each relation's key.
  \item \textbf{(Applied)} A manager proposes keeping \textit{report} as the
        only table ``because reports are faster''. Using every figure in the
        measured box, write a reply that concedes what is true, refutes what is
        not, and proposes a design that keeps the 1.5$\times$ for reports
        without keeping the anomalies --- a design \chref{decomposition}
        measures.
\end{reviewq}

%=====================================================================
\chapter{Decomposition, and Denormalising on Purpose}\label{ch:decomposition}
%=====================================================================
\locator{5}{Part V --- Getting the Design Right}

\begin{outcomes}
By the end of this chapter you will be able to:
\begin{tight}
  \item \textbf{Define} a lossless-join decomposition and \textbf{apply} the
        one-line test for it.
  \item \textbf{Show} what a lossy decomposition does to a query, by counting
        its spurious tuples.
  \item \textbf{Define} dependency preservation and \textbf{identify} a
        decomposition that loses a dependency.
  \item \textbf{Decide} between BCNF and 3NF when they conflict.
  \item \textbf{Measure} a denormalisation --- what it saves per read, what it
        costs per write --- and \textbf{compute} where it breaks even.
\end{tight}
\end{outcomes}

\begin{prereq}
\chref{normalforms}: the normal forms, and its measurement of the wide table.
\chref{fds} for closure, which the lossless test uses. You need
\texttt{code/m19\_decomposition.py}.
\end{prereq}

\begin{keyterms}
decomposition $\bullet$ lossless join $\bullet$ lossy join $\bullet$ spurious
tuple $\bullet$ dependency preservation $\bullet$ 3NF synthesis $\bullet$
denormalisation $\bullet$ summary table $\bullet$ trigger $\bullet$
break-even $\bullet$ read--write ratio
\end{keyterms}

\section{Two Things a Split Must Not Lose}

\chref{normalforms} decomposed \textit{report} into six tables by splitting on
violating dependencies. Any split \emph{can} be made; not every split is
correct. A decomposition must keep two things, and they are independent.

\begin{definitionbox}[Lossless join]
A decomposition of $R$ into $R_1$ and $R_2$ is \term{lossless} if joining them
back gives exactly $R$ --- no tuple missing, \textbf{no tuple added}.

\smallskip
\textbf{The test.} It is lossless if and only if the attributes they share
determine all of one side:
\[
  (R_1 \cap R_2) \rightarrow R_1 \quad\text{or}\quad
  (R_1 \cap R_2) \rightarrow R_2 .
\]
In words: the shared attributes are a superkey of one of the halves. Compute
$(R_1 \cap R_2)^{+}$, as in \chref{fds}, and look.

\smallskip
A \term{lossy} decomposition never loses tuples --- it \emph{adds}
\term{spurious tuples}, rows that look like facts and were never true.
\end{definitionbox}

\begin{definitionbox}[Dependency preservation]
A decomposition \term{preserves dependencies} if every dependency of $R$ can be
checked within a single one of the $R_i$ --- or follows from dependencies that
can.

\smallskip
It matters because a database enforces rules one table at a time. A
dependency that spans two tables can only be checked by joining them, and
nothing in the schema does that on every insert.
\end{definitionbox}

\begin{measured}[A split that invents facts, a split that loses a rule, and a summary table priced]
\textsf{\footnotesize \texttt{code/m19\_decomposition.py}, seeded database
(4{,}000 students), SQLite 3.51.2, Intel i5-6600, Windows 11.}

\medskip
\textbf{1. Lossy against lossless}

\begin{tabular}{@{}L{9.6cm}R{2.8cm}@{}}
\multicolumn{2}{@{}l}{\textbf{lossy:} $R_1$(roll\_no, course\_code) and
  $R_2$(course\_code, grade)} \\
graded (roll\_no, course\_code, grade) facts & 39{,}152 \\
$R_1 \bowtie R_2$ on course\_code & 182{,}920 \\
\textbf{spurious tuples --- never true} & \textbf{143{,}768} \\
\multicolumn{2}{@{}l}{\textbf{lossless:} $S_1$(roll\_no, name, programme) and
  $S_2$(roll\_no, semester, cnic)} \\
student rows, and $S_1 \bowtie S_2$ & 4{,}000 and 4{,}000 \\
\end{tabular}

\smallskip
The lossy join returned \textbf{4.7$\times$} the real facts: every student now
appears to hold every grade that anybody got in the course. The test explains
it in one line --- $R_1 \cap R_2 = \{$course\_code$\}$, a key of neither half;
$S_1 \cap S_2 = \{$roll\_no$\}$, the key of both.

\medskip
\textbf{2. Lossless, but a dependency is lost}

Two tables, \textit{course\_dept}(course\_code, dept\_code) and
\textit{course\_dname}(course\_code, dept\_name), both keyed on course\_code ---
so the join is lossless. Two courses, both CS. A second name for CS:

\begin{tabular}{@{}L{9.6cm}L{2.8cm}@{}}
insert \texttt{('CS-150', 'Computing')} & \textbf{accepted} \\
the join: CS $\rightarrow$ & \texttt{'Computer Science'} \\
the join: CS $\rightarrow$ & \texttt{'Computing'} \\
\end{tabular}

\smallskip
\fd{dept\_code}{dept\_name} is violated, and nothing could refuse it: the rule
spans two tables, and a \sqlkw{CHECK} sees one row of one table. This is
exactly the decomposition \chref{normalforms} warned the wrong splitting order
produces.

\medskip
\textbf{3. Denormalisation, priced} --- enrolments and credit hours per
department

\begin{tabular}{@{}L{8.4cm}R{2.6cm}@{}}
read: computed by two joins & 69.79\,ms \\
read: from a summary table & 0.060\,ms \\
\quad \textbf{saved per read} & \textbf{69.73\,ms} \\
write: one enrolment, no trigger & 28.464\,ms \\
write: one enrolment, trigger maintaining the summary & 32.143\,ms \\
\quad \textbf{added per write} & \textbf{3.679\,ms} \\
summary agreed with the joins, before and after 400 writes & yes, yes \\
\end{tabular}

\smallskip
\textbf{Break-even: one read every 19 writes.} Above that ratio the summary
table is cheaper overall; below it, computing on demand is. A department report
is read far more often than once per nineteen enrolments, so for this report
the summary table wins --- by a margin that is not close.

\smallskip
Most of the 28\,ms per write is the \sqlkw{COMMIT} --- making one enrolment
durable --- and the trigger's two \sqlkw{UPDATE}s ride inside that same
transaction. That is why the trigger adds 13\%, not 200\%.
\end{measured}

\dsfig{%
\begin{tikzpicture}[
  cell/.style={draw=primary!35, line width=0.5pt, minimum width=1.25cm,
               minimum height=0.40cm, font=\scriptsize\ttfamily,
               text=primary!65!black, inner sep=0pt},
  hd/.style={cell, fill=primary, text=white, font=\scriptsize\bfseries},
  bad/.style={cell, fill=accent!14, text=accent},
  ttl/.style={font=\scriptsize\bfseries, text=#1, anchor=south},
  nt/.style={font=\scriptsize\itshape, text=neutral, align=center}
]
  % original
  \node[ttl=primary] at (-5.8,0.5) {the facts};
  \foreach \c [count=\i] in {roll, course, grade}{
    \node[hd] at (-7.05+\i*1.25-1.25+1.25,0.2) {\c};}
  \foreach \r/\a/\b/\g in {1/Ayesha/CS-103/A, 2/Bilal/CS-103/C}{
    \node[cell] at (-5.80,-0.4*\r+0.2) {\a};
    \node[cell] at (-4.55,-0.4*\r+0.2) {\b};
    \node[cell] at (-3.30,-0.4*\r+0.2) {\g};}
  \node[nt] at (-4.55,-1.05) {2 rows};

  % R1, R2
  \node[ttl=secondary!70!black] at (-0.6,1.65) {$R_1$};
  \node[hd] at (-1.25,1.35) {roll};  \node[hd] at (0.0,1.35) {course};
  \node[cell] at (-1.25,0.95) {Ayesha}; \node[cell] at (0.0,0.95) {CS-103};
  \node[cell] at (-1.25,0.55) {Bilal};  \node[cell] at (0.0,0.55) {CS-103};
  \node[ttl=secondary!70!black] at (-0.6,-0.35) {$R_2$};
  \node[hd] at (-1.25,-0.65) {course}; \node[hd] at (0.0,-0.65) {grade};
  \node[cell] at (-1.25,-1.05) {CS-103}; \node[cell] at (0.0,-1.05) {A};
  \node[cell] at (-1.25,-1.45) {CS-103}; \node[cell] at (0.0,-1.45) {C};

  \draw[-{Stealth[length=2mm]}, primary!50, line width=0.9pt]
       (-2.55,-0.2) -- (-2.0,-0.2);
  \draw[-{Stealth[length=2mm]}, primary!50, line width=0.9pt]
       (0.75,-0.2) -- (1.35,-0.2) node[midway, above,
       font=\scriptsize\bfseries, text=primary] {$\bowtie$};

  % rejoined
  \node[ttl=accent] at (3.35,0.5) {$R_1 \bowtie R_2$ on course};
  \foreach \c [count=\i] in {roll, course, grade}{
    \node[hd] at (1.6+\i*1.25-1.25+0.5,0.2) {\c};}
  \foreach \r/\a/\g/\s in {1/Ayesha/A/cell, 2/Ayesha/C/bad,
                           3/Bilal/A/bad, 4/Bilal/C/cell}{
    \node[\s] at (2.10,-0.4*\r+0.2) {\a};
    \node[\s] at (3.35,-0.4*\r+0.2) {CS-103};
    \node[\s] at (4.60,-0.4*\r+0.2) {\g};}
  \node[nt, text width=3.9cm, text=accent] at (3.35,-1.95)
       {4 rows: \textbf{2 never true}. Ayesha did not get a C; Bilal did not
        get an A};
\end{tikzpicture}}%
{A lossy split on the smallest possible scale. The two halves share only
\textit{course}, which is a key of neither, so the join pairs every student in
the course with every grade in it. On the seeded data this is 143{,}768
invented facts.}{lossy}

\begin{alertbox}[What this does and does not license]
\begin{tight}
  \item \textbf{Always test a decomposition for losslessness.} A lossy one
        produces 4.7 false facts for every true one here, and no query
        afterwards can tell them apart.
  \item \textbf{The summary table is a denormalisation, and it is correct only
        while the trigger runs.} Drop the trigger, insert through a path that
        bypasses it, or restore an old backup of one table, and the summary is
        silently wrong. That risk is what the 69\,ms buys.
  \item \textbf{Denormalise by adding, never by replacing.} The normalised
        tables stayed. The summary is derived from them and can be rebuilt
        from them at any moment --- which is the difference between this and
        \chref{normalforms}'s \textit{report}, where the wide table was the
        only copy.
  \item \textbf{Decide with a ratio, not a feeling.} One read per nineteen
        writes is a number; ``reports are slow'' is not.
\end{tight}
\end{alertbox}

\section{When BCNF and Dependency Preservation Conflict}

\begin{conceptbox}[The trade, and the rule for making it]
Every relation has a \textbf{lossless, dependency-preserving decomposition into
3NF}. The \term{3NF synthesis} algorithm finds one: take a minimal cover
(\chref{fds}), make one relation per dependency, add one relation holding a
candidate key if none already does.

\smallskip
A lossless BCNF decomposition always exists too --- but it is \emph{not}
always dependency-preserving. \chref{normalforms}'s
\relation{teaches}{student, course, instructor} is the case: BCNF forces
\fd{student, course}{instructor} out of every single table.

\smallskip
\textbf{The rule:} prefer BCNF; when it would lose a dependency you need
enforced, stop at 3NF and accept the small redundancy that remains. A
redundancy you can see is safer than a rule nothing enforces. For the
university schema the question never arises --- the six tables are in BCNF
\emph{and} preserve every dependency.
\end{conceptbox}

\begin{worked}[Testing two decompositions of \textit{section}]
\relation{section}{\relkey{course\_code}, \relkey{semester},
\relkey{section\_no}, room, emp\_id}, with
\fd{course\_code, semester, section\_no}{room, emp\_id}.

\medskip
\textbf{Decomposition A:} $A_1$(course\_code, semester, section\_no, room)
and $A_2$(course\_code, semester, section\_no, emp\_id).

\smallskip
$A_1 \cap A_2 = \{$course\_code, semester, section\_no$\}$, whose closure is
everything: a superkey of both. \textbf{Lossless.} The dependency splits into
\fd{\dots}{room} in $A_1$ and \fd{\dots}{emp\_id} in $A_2$ --- each checkable in
one table. \textbf{Preserving.} Correct, and pointless: it is two tables where
one would do, and every split is a join somebody pays for on every read.

\medskip
\textbf{Decomposition B:} $B_1$(course\_code, room) and $B_2$(room, semester,
section\_no, emp\_id).

\smallskip
$B_1 \cap B_2 = \{$room$\}$. Does \textit{room} determine all of either half?
No: one room hosts many courses and many sections. \textbf{Lossy.} Joined back
on room, every course would appear in every section that ever met in the same
room --- the pattern of part~1 of the measurement, in a different table.

\medskip
\textbf{The lesson.} The test took one closure each time. B ``looks
reasonable'' --- room is a natural thing to split on --- and is wrong in a way
the data would never report.
\end{worked}

\section{Triggers, and the Summary Table Written Out}

\begin{examplebox}[The summary table that was measured]
\begin{lstlisting}[language=sql]
CREATE TABLE dept_stats (
  dept_code   TEXT PRIMARY KEY NOT NULL REFERENCES department(dept_code),
  n_enrol     INTEGER NOT NULL,
  sum_credits INTEGER NOT NULL
) STRICT;

CREATE TRIGGER stats_ins AFTER INSERT ON enrolment BEGIN
  UPDATE dept_stats
     SET n_enrol     = n_enrol + 1,
         sum_credits = sum_credits +
           (SELECT credit_hours FROM course
             WHERE course_code = NEW.course_code)
   WHERE dept_code =
           (SELECT dept_code FROM course
             WHERE course_code = NEW.course_code);
END;
\end{lstlisting}

\smallskip
\texttt{NEW} is the row just inserted. A matching \texttt{stats\_del} trigger
on \sqlkw{DELETE} uses \texttt{OLD}. \textbf{An \sqlkw{UPDATE} trigger is
needed too}, and the harness does not have one: an enrolment moved from one
course to another would leave both departments' figures wrong. That gap is
\chref{views}' starting point, and the reason a view is often the better
answer.
\end{examplebox}

\begin{pitfall}
\begin{tight}
  \item \textbf{Splitting on an attribute that is not a key of either half.}
        Lossy: 143{,}768 spurious tuples here.
  \item \textbf{Expecting a lossy join to lose rows.} It adds them. The answer
        is too big, not too small, and looks plausible.
  \item \textbf{A rule that spans two tables.} No \sqlkw{CHECK} can enforce
        it; it needs a trigger, or a design that keeps it in one table.
  \item \textbf{Denormalising by replacing the normalised tables.} Add the
        summary; keep the source of truth.
  \item \textbf{Forgetting the \sqlkw{UPDATE} trigger.} Insert and delete are
        maintained; a changed row silently is not.
  \item \textbf{Deciding to denormalise without the ratio.} Measure reads and
        writes; compute the break-even.
\end{tight}
\end{pitfall}

\begin{lab}[Invent 143{,}768 facts, then price a summary table]
\begin{lstlisting}[language=bash]
# --- 1. The seeded database ----------------------------------------
python code/schema/seed.py university.db
sqlite3 university.db
\end{lstlisting}
\begin{lstlisting}[language=sql]
-- --- 2. A lossy split -------------------------------------------
CREATE TEMP TABLE r1 AS SELECT DISTINCT roll_no, course_code
                          FROM enrolment WHERE grade IS NOT NULL;
CREATE TEMP TABLE r2 AS SELECT DISTINCT course_code, grade
                          FROM enrolment WHERE grade IS NOT NULL;
SELECT COUNT(*) FROM (SELECT DISTINCT roll_no, course_code, grade
                        FROM enrolment WHERE grade IS NOT NULL);
SELECT COUNT(*) FROM r1 JOIN r2 USING (course_code);
-- How many facts were invented? Apply the one-line test to explain.

-- --- 3. Pick one student and see what the lossy join says ------
SELECT grade FROM r1 JOIN r2 USING (course_code)
 WHERE roll_no = (SELECT MIN(roll_no) FROM r1) LIMIT 10;
SELECT grade FROM enrolment
 WHERE roll_no = (SELECT MIN(roll_no) FROM r1) AND grade IS NOT NULL;

-- --- 4. The summary table and its trigger -----------------------
-- Type in dept_stats and stats_ins from this chapter, then fill it:
INSERT INTO dept_stats
  SELECT d.dept_code, COUNT(e.roll_no), COALESCE(SUM(c.credit_hours),0)
    FROM department d
    LEFT JOIN course c    ON c.dept_code   = d.dept_code
    LEFT JOIN enrolment e ON e.course_code = c.course_code
   GROUP BY d.dept_code;

-- --- 5. Time both reads -----------------------------------------
.timer on
SELECT * FROM dept_stats;
SELECT c.dept_code, COUNT(*), SUM(c.credit_hours)
  FROM enrolment e JOIN course c ON c.course_code = e.course_code
 GROUP BY c.dept_code;

-- --- 6. Break it --------------------------------------------------
UPDATE enrolment SET course_code =
       (SELECT course_code FROM course WHERE dept_code = 'MA' LIMIT 1)
 WHERE rowid = (SELECT MIN(rowid) FROM enrolment);
-- Re-run both reads from step 5. Do they still agree? Why not?
\end{lstlisting}
\begin{lstlisting}[language=bash]
# --- 7. The measured version ----------------------------------------
python code/m19_decomposition.py
\end{lstlisting}

\smallskip
\textbf{What to notice.} Step 6 is the price of the 69\,ms. The summary table is
correct exactly as long as every change to \textit{enrolment} passes through a
trigger that knows about it --- and the harness, like most first attempts,
forgot \sqlkw{UPDATE}. (Step 6 may also be refused if the move would duplicate
an enrolment's key; pick another row if so.)
\end{lab}

\begin{checkpoint}
\begin{tightnum}
  \item Apply the lossless test to $R_1$(roll\_no, course\_code) and
        $R_2$(course\_code, grade), and explain from the result why the join
        produced 4.7 times the real facts.
  \item Give a decomposition of the university schema that is lossless but not
        dependency-preserving, name the dependency it loses, and show an insert
        that violates the dependency without any table refusing it.
  \item The summary table saves 69.73\,ms per read and costs 3.679\,ms per
        write. Show the break-even calculation, and say what would change it
        if the database moved to a fast SSD where a commit took 1\,ms.
\end{tightnum}
\end{checkpoint}

\begin{chaptersummary}
\begin{tight}
  \item A decomposition must be \textbf{lossless} (the join gives back exactly
        $R$) and should be \textbf{dependency-preserving} (every rule checkable
        in one table). The two are independent.
  \item \textbf{Lossless test:} the shared attributes must be a superkey of one
        half. One closure decides it.
  \item \textbf{Measured:} a lossy split of \textit{enrolment} produced
        \textbf{143{,}768 spurious tuples} --- 4.7$\times$ the real facts. A
        lossy join adds rows; it does not lose them.
  \item \textbf{Measured:} a lossless split that lost
        \fd{dept\_code}{dept\_name} \textbf{accepted two names for CS}. A rule
        spanning two tables cannot be enforced by a \sqlkw{CHECK}.
  \item 3NF synthesis is always lossless and preserving; BCNF is always
        lossless but not always preserving. Prefer BCNF; stop at 3NF when BCNF
        would cost a rule you need.
  \item \textbf{Measured:} a trigger-maintained summary table saved
        \textbf{69.73\,ms per read} for \textbf{3.679\,ms per write} --- break
        even at \textbf{one read per 19 writes}.
  \item Most of a write's cost is the \sqlkw{COMMIT}; a trigger rides inside
        it, which is why it added 13\%.
  \item Denormalise by \emph{adding} derived data beside the normalised
        tables, never by replacing them --- and maintain it on insert, delete
        \emph{and} update.
\end{tight}
\end{chaptersummary}

\begin{reviewq}
  \item \textbf{(MCQ)} A lossy decomposition, when joined back, produces:
        (a)~fewer tuples; (b)~the same tuples; (c)~extra, spurious tuples;
        (d)~an error.
  \item \textbf{(MCQ)} $R_1 \cap R_2 = \{A\}$ is lossless if: (a)~$A$ is a
        superkey of $R_1$ or of $R_2$; (b)~$A$ is in both; (c)~$A$ is
        indexed; (d)~$R_1$ and $R_2$ have the same size.
  \item \textbf{(Short)} Why can a \sqlkw{CHECK} constraint not enforce a
        dependency that spans two tables?
  \item \textbf{(Short)} State the rule for choosing between BCNF and 3NF.
  \item \textbf{(Short)} What three triggers does a summary table need, and
        what goes wrong without the third?
  \item \textbf{(Applied)} Decompose
        \relation{booking}{room, date, slot, course\_code, emp\_id, capacity}
        (\chref{fds}'s review question) into BCNF, test each split for
        losslessness by closure, and say whether every dependency is
        preserved.
  \item \textbf{(Applied)} The university asks for a ``students per
        programme per semester'' figure on its home page, read 5{,}000 times a
        day; about 300 enrolments change a day. Design the summary table and
        its triggers, compute the break-even using this chapter's measured
        costs, and say whether you would build it.
\end{reviewq}

%=====================================================================
\chapter{Views, Triggers and Keeping Data Honest}\label{ch:views}
%=====================================================================
\locator{5}{Part V --- Getting the Design Right}

\begin{outcomes}
By the end of this chapter you will be able to:
\begin{tight}
  \item \textbf{Create} a view and \textbf{explain} what is stored when you
        do.
  \item \textbf{Show}, from a compiled program, that a view costs nothing at
        run time.
  \item \textbf{Use} a view to provide an external schema and to hide a
        schema change from the programs that read it.
  \item \textbf{Write} \sqlkw{BEFORE}, \sqlkw{AFTER} and \sqlkw{INSTEAD OF}
        triggers, and \textbf{say} when each is the right one.
  \item \textbf{Choose} between a view and a trigger-maintained table by what
        each guarantees.
\end{tight}
\end{outcomes}

\begin{prereq}
\chref{architecture}: the external level, which a view implements.
\chref{decomposition}: the summary table and its missing \sqlkw{UPDATE}
trigger. You need \texttt{code/m20\_views.py}.
\end{prereq}

\begin{keyterms}
view $\bullet$ virtual table $\bullet$ base table $\bullet$ view expansion
$\bullet$ flattening $\bullet$ co-routine $\bullet$ updatable view $\bullet$
trigger $\bullet$ \sqlkw{BEFORE} $\bullet$ \sqlkw{AFTER} $\bullet$
\sqlkw{INSTEAD OF} $\bullet$ \texttt{NEW} $\bullet$ \texttt{OLD} $\bullet$
\sqlkw{RAISE} $\bullet$ audit table $\bullet$ business rule $\bullet$
materialised view
\end{keyterms}

\section{A Query With a Name}

\begin{definitionbox}[View]
A \term{view} is a named query. \sqlkw{CREATE VIEW} $v$ \sqlkw{AS}
\sqlkw{SELECT} \dots\ stores the \emph{text} of the \sqlkw{SELECT} in the data
dictionary (\chref{whodoeswhat}) and nothing else --- no rows.

\smallskip
A query that names $v$ is rewritten, before it is compiled, with $v$'s
\sqlkw{SELECT} in its place --- \term{view expansion}. The view is a \term{virtual table}; the tables it
reads are \term{base tables}.

\smallskip
So a view is always \textbf{current}: it is the query, run now, against
whatever the base tables hold now. It cannot be stale, because there is
nothing stored to go stale.
\end{definitionbox}

\dsfig{%
\begin{tikzpicture}[
  q/.style={rounded corners=3pt, draw=#1, line width=1pt, fill=#1!7,
            text=#1!60!black, font=\scriptsize\ttfamily, align=left,
            text width=5.0cm, inner sep=4pt},
  dd/.style={rounded corners=3pt, draw=goldhl, line width=1pt, fill=goldhl!8,
             text=goldhl!60!black, font=\scriptsize\ttfamily, align=left,
             text width=4.4cm, inner sep=4pt},
  st/.style={rounded corners=3pt, draw=neutral, line width=0.9pt,
             fill=neutral!8, text=neutral!80!black, font=\scriptsize\bfseries,
             align=center, text width=2.6cm, minimum height=0.8cm},
  arr/.style={-{Stealth[length=2.2mm]}, primary!55, line width=0.9pt},
  tg/.style={font=\scriptsize\itshape, text=neutral}
]
  \node[q=secondary] (u) at (-4.4,1.4)
       {SELECT * FROM transcript\\ WHERE roll\_no = ?};
  \node[tg, anchor=south] at (-4.4,2.05) {what the program sends};
  \node[dd] (v) at (2.0,1.4)
       {transcript = SELECT e.roll\_no, c.course\_code, \dots\\
        FROM enrolment e JOIN course c ON \dots};
  \node[tg, anchor=south] at (2.0,2.2) {what the dictionary holds --- text};
  \node[q=primary] (r) at (-1.2,-0.4)
       {SELECT e.roll\_no, c.course\_code, \dots\\
        FROM enrolment e JOIN course c ON \dots\\
        WHERE e.roll\_no = ?};
  \node[tg, anchor=north, text width=5.4cm, align=center] at (-1.2,-1.25)
       {the query after substitution: the view is gone, and the filter sits on
        the base table};
  \node[st] (c) at (4.6,-0.4) {compiled\\program};
  \node[tg, anchor=north, text width=3.0cm, align=center] at (4.6,-0.95)
       {23 instructions, as if written by hand};
  \draw[arr] (u.south) -- ++(0,-0.35) -| ([xshift=-1.0cm]r.north);
  \draw[arr] (v.south) -- ++(0,-0.35) -| ([xshift=1.2cm]r.north);
  \draw[arr] (r.east) -- (c.west);
\end{tikzpicture}}%
{What a view is, in one diagram. The dictionary holds only the view's text;
SQLite substitutes it into the query that names it, then compiles. By the time
anything runs there is no view left to cost anything, or to be out of date.}%
{viewexpand}

\begin{examplebox}[The three jobs a view does in this course]
\textbf{1. An external schema.} \chref{architecture}'s \texttt{fees\_view} ---
roll number and name, nothing else --- is the external level, built:

\begin{lstlisting}[language=sql]
CREATE VIEW fees_view AS SELECT roll_no, name FROM student;
\end{lstlisting}

\textbf{2. Logical data independence.} When \chref{evolution}'s migration
split a column, a view with the old shape let old programs keep working:

\begin{lstlisting}[language=sql]
CREATE VIEW student_v1 AS
SELECT roll_no, name, programme, semester,
       addr_house || ', ' || addr_street || ', ' || addr_city AS address
  FROM student;
\end{lstlisting}

\textbf{3. A query written once.} The transcript is a join every report needs.
Write it once and every query that names \texttt{transcript} gets it right.

\smallskip
\textbf{What a view does \emph{not} do in SQLite} is restrict anyone to it.
There are no users (\chref{whodoeswhat}), so a person who can open the file can
read \textit{student} directly. The external level exists here; its
enforcement waits for \chref{project}.
\end{examplebox}

\section{Measured: What a View Costs, and What a Trigger Costs}

\begin{measured}[A view costs nothing; a trigger costs less than the same work by hand]
\textsf{\footnotesize \texttt{code/m20\_views.py}, seeded database (4{,}000
students), SQLite 3.51.2, Intel i5-6600, Windows 11, mechanical disk.}

\medskip
\textbf{1. A view against the query it wraps} --- one student's transcript

\begin{tabular}{@{}L{6.4cm}R{2.0cm}R{1.6cm}R{2.6cm}@{}}
& \textbf{ms} & \textbf{rows} & \textbf{instructions} \\
the query written out & 0.091 & 15 & 23 \\
the same through a view & 0.104 & 15 & 23 \\
\end{tabular}

\smallskip
\textbf{The same 23 instructions in the same order.} Seventeen of them differ
at all, and every one differs only in which \emph{cursor number} it uses: the
view's program numbers its tables one higher. That is the whole difference ---
the view was substituted and compiled away before it could cost anything.

\medskip
\textbf{2. Does a filter on the view reach inside it?}

\begin{tabular}{@{}L{7.0cm}L{7.0cm}@{}}
\texttt{transcript WHERE roll\_no = ?} &
  \texttt{SEARCH e USING INDEX \dots\ (roll\_no=?)} \\
grouped view, no index on course\_code &
  \texttt{CO-ROUTINE} + \texttt{SCAN enrolment} \\
grouped view, with that index &
  \texttt{CO-ROUTINE} + \texttt{SEARCH \dots\ (course\_code=?)} \\
\end{tabular}

\smallskip
A grouped view cannot be flattened, so it runs as a co-routine --- but the
filter is still pushed into it. Without an index the plan shows a
\texttt{SCAN}, which looks like no push-down and is not: there was simply
nothing to search with. With the index, only one course's rows were ever
grouped (0.277\,ms through the view, 0.268\,ms written out).

\medskip
\textbf{3. What a trigger costs} --- an audit trigger recording every grade
change. Commit per row: 200 changes, median of 3. One commit: 30{,}000
changes, median of 5. Per row.

% Five columns inside a measured box must total 14.3 cm with their gaps:
% 4.2 + 2.2 + 1.8 + 2.2 + 2.2 = 12.6, plus four gaps of 0.42 = 14.3.
\begin{tabular}{@{}L{4.2cm}R{2.2cm}R{1.8cm}R{2.2cm}R{2.2cm}@{}}
& \textbf{1 per row} & \textbf{(range)} & \textbf{1 for all} & \textbf{(range)} \\
\sqlkw{UPDATE} only & 27.58\,ms & 27.4--34.7 & 0.0494\,ms & 0.041--0.064 \\
+ audit trigger & 31.29\,ms & 28.2--49.7 & 0.0622\,ms & 0.052--0.168 \\
+ the same \sqlkw{INSERT} by hand & 30.15\,ms & 29.6--34.9 & 0.0634\,ms &
  0.051--0.089 \\
\end{tabular}

\smallskip
\textbf{The trigger costs what the work inside it costs.} Against the same
\sqlkw{INSERT} written by hand, the ranges overlap in both modes --- neither is
reliably cheaper --- so the mechanism itself adds nothing measurable. What the
trigger adds to an update is real but small: about a hundredth of a
millisecond per row in bulk. With a commit per row it is lost in the disk
sync, whose run-to-run range is wider than the trigger's whole cost.

\smallskip
\textbf{A note on how this was established, because it took three tries.} The
first version ran each case once and blamed the cost on ``the second table'';
a hand-written control cost \emph{more}, so that was wrong. The second
version, on 200 rows, concluded the trigger was \emph{cheaper} than by hand ---
and \texttt{code/run\_all.py} caught that conclusion failing on a re-run: at
200 rows a bulk run takes 30\,ms in total, and the trigger's ``cost'' came out
\emph{negative} on two runs in three. A second full run of this version also
moved the trigger's share of an update from 26\% to 49\%. The one result that
held on every run is the one stated above. Small effects need many rows,
several runs, ranges and a control --- and a harness that re-checks the
handout's claims, because the second error would otherwise have been
published.

\medskip
\textbf{4. The view and \chref{decomposition}'s summary table, after one
enrolment moves from a CS course to an MA course}

\begin{tabular}{@{}L{3.0cm}R{2.4cm}R{2.4cm}R{2.8cm}@{}}
& \textbf{truth} & \textbf{view} & \textbf{summary table} \\
CS & 7{,}868 & 7{,}868 & \textbf{7{,}869} \\
MA & 7{,}174 & 7{,}174 & \textbf{7{,}173} \\
read cost & & 53.69\,ms & 0.076\,ms \\
\end{tabular}

\smallskip
The view is right, because it is the query run now. The table is wrong in both
rows, because it has insert and delete triggers and no update trigger --- the
gap \chref{decomposition} left. It is 700$\times$ faster to read and exactly as
correct as its triggers are complete.
\end{measured}

\begin{alertbox}[Choosing between a view and a maintained table]
\begin{tight}
  \item \textbf{Default to the view.} It costs nothing beyond the query it
        names, and it cannot be wrong. Measured.
  \item \textbf{Replace it with a maintained table only when the measured read
        cost forces you to} --- and then with triggers for \sqlkw{INSERT},
        \sqlkw{DELETE} \emph{and} \sqlkw{UPDATE}, plus a query that checks the
        table against the view, run on a schedule.
  \item \textbf{Do not fear triggers for their cost.} An audit trigger cost
        what the same \sqlkw{INSERT} written by the program cost --- about a
        hundredth of a millisecond a row in bulk. What they really cost is the
        attention needed to keep them complete.
  \item \textbf{Other systems offer a \term{materialised view}} --- a view
        whose result is stored and refreshed by the database. SQLite does not.
        The summary table with triggers is what that feature does, written out.
\end{tight}
\end{alertbox}

\section{Triggers}

\begin{definitionbox}[Trigger]
A \term{trigger} is a block of SQL the database runs automatically when a
named event happens to a table: \sqlkw{INSERT}, \sqlkw{UPDATE} (optionally
\sqlkw{OF} named columns) or \sqlkw{DELETE}.

\smallskip
It runs \sqlkw{BEFORE} or \sqlkw{AFTER} the change, once for each row
affected. Inside it, \texttt{NEW} is the row as it will be and \texttt{OLD} the
row as it was --- \texttt{NEW} has no meaning in a \sqlkw{DELETE} trigger, nor
\texttt{OLD} in an \sqlkw{INSERT} trigger.

\smallskip
\sqlkw{WHEN} limits it to rows meeting a condition. \texttt{RAISE(ABORT,
'message')} inside a trigger cancels the statement with an error.
\end{definitionbox}

\begin{worked}[A business rule no \sqlkw{CHECK} can express]
\textbf{The rule.} A student may not enrol in more than six sections in one
semester.

\medskip
\textbf{Step 1 --- why a \sqlkw{CHECK} cannot do it.} A \sqlkw{CHECK} sees one
row of one table. ``More than six'' is a fact about \emph{other} rows ---
the student's other enrolments --- so it is out of reach, exactly as
\chref{decomposition}'s two-table dependency was.

\medskip
\textbf{Step 2 --- \sqlkw{BEFORE INSERT}, so the bad row never arrives.}
\begin{lstlisting}[language=sql]
CREATE TRIGGER max_six BEFORE INSERT ON enrolment
WHEN (SELECT COUNT(*) FROM enrolment
       WHERE roll_no = NEW.roll_no AND semester = NEW.semester) >= 6
BEGIN
  SELECT RAISE(ABORT, 'a student may take at most six sections a term');
END;
\end{lstlisting}

\medskip
\textbf{Step 3 --- the gap.} This guards \sqlkw{INSERT}. An \sqlkw{UPDATE} that
moves an enrolment from Fall into a Spring that is already full is not
guarded. \textbf{Every rule implemented by trigger needs a trigger for every
event that can break it} --- the lesson of the summary table, again, and the
single most common defect in trigger code.

\medskip
\textbf{Step 4 --- why \sqlkw{BEFORE} and not \sqlkw{AFTER}.} Either would
work: \texttt{RAISE(ABORT)} in an \sqlkw{AFTER} trigger also undoes the
statement. \sqlkw{BEFORE} states the intent --- \emph{check, then allow} ---
and does no work on a row that will be refused.
\end{worked}

\begin{examplebox}[\sqlkw{INSTEAD OF}: making a view accept writes]
A view that accepts writes is called \term{updatable}. In SQLite
\textbf{no} view is: none can be written to directly --- not even
\texttt{fees\_view}, which reads a single table:
\texttt{cannot modify fees\_view because it is a view}. Other systems accept
writes to simple one-table views; SQLite accepts none, and an \sqlkw{INSTEAD
OF} trigger is the only way to say what a write to a view should do:

\begin{lstlisting}[language=sql]
CREATE TRIGGER fees_view_update INSTEAD OF UPDATE OF name ON fees_view
BEGIN
  UPDATE student SET name = NEW.name WHERE roll_no = OLD.roll_no;
END;

UPDATE fees_view SET name = 'Ayesha K. Khan'
 WHERE roll_no = '2021-IT-001';      -- now works
\end{lstlisting}

\smallskip
This completes the external level: a group of users can read \emph{and}
correct the part of the database they see, through the view, without ever
naming the base table.
\end{examplebox}

\begin{pitfall}
\begin{tight}
  \item \textbf{Believing a view stores rows.} It stores a \sqlkw{SELECT}. It
        cannot be stale.
  \item \textbf{Avoiding views for speed.} Measured: compiled away; same
        instructions as the query written out.
  \item \textbf{Reading \texttt{SCAN} in a grouped view's plan as ``no
        push-down''.} It may mean there was no index to search with.
  \item \textbf{A trigger for one event when the rule can be broken by three.}
        The summary table had no \sqlkw{UPDATE} trigger and was wrong after one
        \sqlkw{UPDATE}.
  \item \textbf{Using \texttt{NEW} in a \sqlkw{DELETE} trigger.} There is no
        new row.
  \item \textbf{Treating a view as security in SQLite.} No users, so no
        restriction.
  \item \textbf{Concluding from one timing run.} The first version of this
        chapter's harness did, and its explanation was wrong.
\end{tight}
\end{pitfall}

\begin{lab}[Prove a view is free, catch a summary table lying, and enforce a rule]
\begin{lstlisting}[language=bash]
# --- 1. The seeded database ----------------------------------------
python code/schema/seed.py university.db
sqlite3 university.db
\end{lstlisting}
\begin{lstlisting}[language=sql]
-- --- 2. A view, and the query it wraps --------------------------
CREATE VIEW transcript AS
SELECT e.roll_no, c.course_code, c.title, e.semester, e.grade
  FROM enrolment e JOIN course c ON c.course_code = e.course_code;
EXPLAIN QUERY PLAN SELECT * FROM transcript
 WHERE roll_no = (SELECT MIN(roll_no) FROM student);
-- The filter reached the base table. Where can you see that?

-- --- 3. What was stored -----------------------------------------
SELECT type, name, sql FROM sqlite_master WHERE name = 'transcript';
-- Text. No rows. That is the whole of a view.

-- --- 4. An external schema --------------------------------------
CREATE VIEW fees_view AS SELECT roll_no, name FROM student;
SELECT * FROM fees_view LIMIT 3;
SELECT * FROM student LIMIT 3;
-- The second still works. What would a server DBMS do instead?

-- --- 5. A rule no CHECK can express -----------------------------
CREATE TRIGGER max_six BEFORE INSERT ON enrolment
WHEN (SELECT COUNT(*) FROM enrolment WHERE roll_no = NEW.roll_no
                                       AND semester = NEW.semester) >= 6
BEGIN SELECT RAISE(ABORT, 'at most six sections a term'); END;
SELECT roll_no, COUNT(*) FROM enrolment WHERE semester = 'FALL2026'
 GROUP BY roll_no HAVING COUNT(*) = 6 LIMIT 1;
-- Try to enrol that student in one more FALL2026 section. Record
-- the message. Then find a way around the trigger with an UPDATE.

-- --- 6. An audit trail ------------------------------------------
CREATE TABLE grade_audit (roll_no TEXT, course_code TEXT,
  old_grade TEXT, new_grade TEXT,
  changed_at TEXT DEFAULT (datetime('now')));
CREATE TRIGGER audit_grade AFTER UPDATE OF grade ON enrolment
WHEN OLD.grade IS NOT NEW.grade BEGIN
  INSERT INTO grade_audit (roll_no, course_code, old_grade, new_grade)
  VALUES (OLD.roll_no, OLD.course_code, OLD.grade, NEW.grade);
END;
UPDATE enrolment SET grade = 'A'
 WHERE rowid = (SELECT MIN(rowid) FROM enrolment WHERE grade = 'B');
SELECT * FROM grade_audit;
\end{lstlisting}
\begin{lstlisting}[language=bash]
# --- 7. The measured version ----------------------------------------
python code/m20_views.py
\end{lstlisting}

\smallskip
\textbf{What to notice.} Step 5's second half is the chapter's warning, found
by hand: a rule enforced on \sqlkw{INSERT} can be broken by \sqlkw{UPDATE}.
Writing the \sqlkw{BEFORE UPDATE} trigger that closes the gap is part of the
project in \chref{project}.
\end{lab}

\begin{checkpoint}
\begin{tightnum}
  \item A view's compiled program had the same 23 instructions as the query it
        wraps, differing only in cursor numbers. Explain what SQLite did with
        the view to produce that, and why it means the view cannot be stale.
  \item The summary table read in 0.076\,ms and the view in 53.69\,ms, and
        after one \sqlkw{UPDATE} only the view was right. State the condition
        under which you would choose the table, and list everything you would
        have to build alongside it.
  \item Write the trigger that closes the gap in the six-section rule for
        \sqlkw{UPDATE}, and say which of \texttt{NEW} and \texttt{OLD} it must
        use, and why.
\end{tightnum}
\end{checkpoint}

\begin{chaptersummary}
\begin{tight}
  \item A \textbf{view} stores a \sqlkw{SELECT}, not rows. It is substituted
        into each query that names it, so it is always current.
  \item \textbf{Measured:} a view compiled to \textbf{the same 23
        instructions} as the query it wraps, differing only in cursor
        numbers. It costs nothing.
  \item \textbf{Measured:} a filter on a view reaches the base table, even
        through a grouped view run as a co-routine --- visible once there is an
        index to search with.
  \item Views provide the \textbf{external level}, \textbf{logical data
        independence}, and a query written once. In SQLite they do not restrict
        access.
  \item A \textbf{trigger} runs per row, \sqlkw{BEFORE} or \sqlkw{AFTER} an
        event, with \texttt{NEW} and \texttt{OLD}; \sqlkw{RAISE(ABORT)}
        refuses the change.
  \item \textbf{Measured:} an audit trigger cost \textbf{what the same
        \sqlkw{INSERT} written by hand cost} --- the ranges overlapped on every
        run --- about a hundredth of a millisecond a row in bulk. With a commit
        per row its cost was inside the disk-sync noise.
  \item Triggers enforce rules a \sqlkw{CHECK} cannot reach --- but only for
        the events they are written for. Every event that can break the rule
        needs one.
  \item \textbf{Measured:} after one \sqlkw{UPDATE}, the view was right and
        the trigger-maintained table was wrong. Default to the view.
  \item \sqlkw{INSTEAD OF} triggers make a view writable.
\end{tight}
\end{chaptersummary}

\begin{reviewq}
  \item \textbf{(MCQ)} \sqlkw{CREATE VIEW} stores: (a)~a copy of the rows;
        (b)~the text of the query; (c)~an index; (d)~a trigger.
  \item \textbf{(MCQ)} In a \sqlkw{DELETE} trigger, which row reference is
        available? (a)~\texttt{NEW}; (b)~\texttt{OLD}; (c)~both;
        (d)~neither.
  \item \textbf{(Short)} Why can a view never be stale, while a summary table
        can?
  \item \textbf{(Short)} Give a rule that needs a trigger because no
        \sqlkw{CHECK} can express it, and say why.
  \item \textbf{(Short)} When is \sqlkw{INSTEAD OF} the only kind of trigger
        that will do?
  \item \textbf{(Applied)} Write the views that give the examinations office,
        the fees office and the student portal their external schemas from
        \chref{architecture}'s review question, an \sqlkw{INSTEAD OF} trigger
        letting the portal correct a student's phone number, and say what
        SQLite cannot enforce about any of it.
  \item \textbf{(Applied)} The audit trigger added about 0.013\,ms per row in
        bulk.
        Estimate what auditing every grade change costs per semester on the
        seeded data, decide whether that matters, and list the three trigger
        events an audit of \textit{enrolment} would need to be complete.
\end{reviewq}


\part{Making It Work, and What Comes Next}
%=====================================================================
% PART VI OPENER
%=====================================================================
\thispagestyle{plain}
\structuralfigures{P6.}

\begin{center}
\begin{tcolorbox}[enhanced, width=0.92\textwidth, colback=goldhl!5,
  colframe=goldhl, boxrule=0pt, leftrule=5pt, arc=3pt, top=8pt, bottom=8pt]
\large\itshape\color{goldhl!60!black}
Everything so far has been one person, one database, one question at a time. Part~VI is what changes when that stops being true: when the table is large enough that how you look matters, when two people write at once, and when somebody who should not be reading the data is. It closes with the two questions a database course owes its students --- what a real server does that this one does not, and when the relational model is the wrong answer.
\end{tcolorbox}
\end{center}

\vspace{6pt}

\dsfigh{%
\begin{tikzpicture}[
  cbox/.style={rounded corners=3pt, draw=goldhl, line width=1pt, fill=goldhl!7,
               text=goldhl!55!black, font=\small\bfseries, align=center,
               minimum width=4.0cm, minimum height=1.0cm},
  hbox/.style={rounded corners=3pt, draw=goldhl, line width=1.8pt, fill=goldhl!16,
               text=goldhl!55!black, font=\small\bfseries, align=center,
               minimum width=4.0cm, minimum height=1.0cm},
  arr/.style={-{Stealth[length=2.4mm]}, goldhl!60, line width=1pt},
  note/.style={font=\scriptsize\itshape, text=goldhl!55!black, align=center,
               fill=white, inner sep=2pt}
]
  \node[hbox] (c0) at (0.00,0)  {Ch.\ 21\\Indexes};
  \node[cbox] (c1) at (5.20,0)  {Ch.\ 22\\Transactions};
  \node[cbox] (c2) at (10.40,0) {Ch.\ 23\\Security, NoSQL, Project};
  \draw[arr] (c0) -- (c1);
  \draw[arr] (c1) -- (c2);
  \node[note] at (2.60,-1.0) {when the table is big};
  \node[note] at (7.80,-1.0) {when there are two of you};
  \node[font=\scriptsize\itshape, text=accent, align=center] at (0.00,1.15)
       {the rule everybody repeats, and the row count at which it stops being
        true};
\end{tikzpicture}}%
{Reading path through Part~VI. Chapter 21 is the measured chapter students
remember, because it contradicts something they were told in their first week
and shows them exactly where the contradiction begins.}{part6map}

\newpage

%=====================================================================
\chapter{Indexes, and Where They Stop Helping}\label{ch:indexes}
%=====================================================================
\locator{6}{Part VI --- Making It Work, and What Comes Next}

\begin{outcomes}
By the end of this chapter you will be able to:
\begin{tight}
  \item \textbf{Explain} what an index is, using the B-tree, without
        assuming any course in data structures.
  \item \textbf{Predict} whether a query can use a given index, including a
        composite one, from the order of its columns.
  \item \textbf{Find} the share of a table above which an index is slower than
        reading everything, by measurement.
  \item \textbf{Explain} why SQLite's optimiser chose a twenty-times slower
        plan, and \textbf{override} it.
  \item \textbf{State} what an index costs, and \textbf{decide} whether a
        proposed one is worth it.
\end{tight}
\end{outcomes}

\begin{prereq}
\chref{whodoeswhat}: the optimiser and the buffer manager. \chref{firstdatabase}:
the database is a file of 4{,}096-byte pages. \chref{joins},
\chref{subqueries} and \chref{grouping}, each of which ended with an index
making something faster --- this chapter is about why, and about where it
stops. \textbf{No data-structures course is assumed.} You need
\texttt{code/m21\_indexes.py}.
\end{prereq}

\begin{keyterms}
index $\bullet$ B-tree $\bullet$ node $\bullet$ fan-out $\bullet$ depth
$\bullet$ leaf $\bullet$ rowid $\bullet$ covering index $\bullet$ composite
index $\bullet$ leftmost prefix $\bullet$ selectivity $\bullet$ crossover
$\bullet$ sequential read $\bullet$ random read $\bullet$ \sqlkw{ANALYZE}
$\bullet$ \texttt{sqlite\_stat1} $\bullet$ STAT4 $\bullet$ \sqlkw{INDEXED BY}
$\bullet$ \sqlkw{NOT INDEXED}
\end{keyterms}

\section{What an Index Is}

A table in SQLite is stored in order of its rowid. To find the student with
roll number \texttt{2024-IT-00317} in a table with no index on roll number,
the only method is to look at every row --- a \sqlkw{SCAN}. Four chapters have
shown plans that said \sqlkw{SEARCH} instead, and every one of them was using
the same idea.

\begin{definitionbox}[Index]
An \term{index} on a column is a separate structure, kept in the same file,
that holds every value of that column \textbf{in sorted order}, each paired
with the rowid of the row it came from.

\smallskip
Because it is sorted, a value can be found without looking at most of them ---
exactly as a word is found in a dictionary without reading the dictionary.
Because each entry carries a rowid, the full row can then be fetched directly.

\smallskip
SQLite keeps every index as a \term{B-tree}. You do not need to know how one is
built to use it; you need to know one number about it, below.
\end{definitionbox}

\begin{conceptbox}[The B-tree, in one paragraph and one number]
A B-tree is a tree of pages --- its \term{nodes}. The top page holds a few hundred keys that divide
the whole range of values into a few hundred parts, each with a pointer to a
page below; that page divides its part into a few hundred more, and so on,
until the bottom pages --- the \term{leaves} --- hold the actual entries.

\smallskip
The number that matters is the \term{fan-out}: how many children each page
points to. A 4{,}096-byte page holding short integer keys has a fan-out of a
few hundred. So:

\begin{tabular}{@{}L{3.6cm}L{10.2cm}@{}}
1 level   & a few hundred entries \\
2 levels  & tens of thousands \\
3 levels  & millions \\
4 levels  & hundreds of millions \\
\end{tabular}

\smallskip
\textbf{Finding any value means reading one page per level} --- three or four
reads, whether the table has a thousand rows or a million. That is the whole
reason part~1 of the measurement below shows a lookup costing the same at every
size.
\end{conceptbox}

\dsfig{%
\begin{tikzpicture}[
  pg/.style={rounded corners=2pt, draw=#1, line width=0.9pt, fill=#1!8,
             text=#1!60!black, font=\scriptsize\ttfamily, minimum height=0.55cm,
             inner sep=3pt},
  lf/.style={pg=greenhl, minimum width=1.9cm},
  ln/.style={draw=neutral!60, line width=0.7pt},
  hot/.style={draw=accent, line width=1.4pt, -{Stealth[length=2.2mm]}},
  lv/.style={font=\scriptsize\itshape, text=neutral, anchor=east},
  nt/.style={font=\scriptsize\itshape, text=neutral, align=left}
]
  \node[pg=primary, minimum width=4.6cm] (r) at (0,2.4)
       {\textbar\ 250k \textbar\ 500k \textbar\ 750k \textbar};
  \node[pg=secondary, minimum width=3.0cm] (i1) at (-4.5,1.1)
       {\textbar\ 60k \textbar\ 120k \textbar\ \dots};
  \node[pg=secondary, minimum width=3.0cm] (i2) at (-0.9,1.1)
       {\textbar\ 310k \textbar\ 370k \textbar\ \dots};
  \node[font=\scriptsize, text=neutral] at (2.1,1.1) {\dots};
  \node[pg=secondary, minimum width=3.0cm] (i4) at (4.6,1.1)
       {\textbar\ 810k \textbar\ 870k \textbar\ \dots};
  \foreach \c in {i1,i2,i4}{\draw[ln] (r.south) -- (\c.north);}

  \node[lf] (l1) at (-2.6,-0.25) {317 512 $\rightarrow$ rowid};
  \node[lf] (l2) at (-0.4,-0.25) {317 513 $\rightarrow$ rowid};
  \node[font=\scriptsize, text=neutral] at (1.4,-0.25) {\dots};
  \draw[ln] (i2.south) -- (l1.north);
  \draw[ln] (i2.south) -- (l2.north);

  \node[pg=neutral, minimum width=3.6cm] (row) at (-1.5,-1.55)
       {the row itself, in the table};

  \draw[hot] ([xshift=2mm]r.south west) ++(1.3,0) -- (i2.north);
  \draw[hot] (i2.south) -- (l1.north);
  \draw[hot] (l1.south) -- (row.north);

  \node[lv] at (-6.1,2.4) {root};
  \node[lv] at (-6.1,1.1) {interior};
  \node[lv] at (-6.1,-0.25) {leaves};
  \node[lv] at (-6.1,-1.55) {table};
  \node[nt, text width=4.0cm, anchor=west] at (2.4,-0.95)
       {\textcolor{accent}{\textbf{finding k = 317\,512}}: one page at each
        level, then one jump to the row --- four reads among a million rows};
\end{tikzpicture}}%
{An index on \textit{k} over a million rows, three levels deep. Each page
divides its range among a few hundred children, so the path from the root to
any key is three pages long, and one more read fetches the row. The scan it
replaces reads every page of the table.}{btree}

\section{Measured: Where an Index Wins, and Where It Loses}

\begin{measured}[A 3{,}501$\times$ win, a 21$\times$ loss, and an optimiser that could not tell the difference]
\textsf{\footnotesize \texttt{code/m21\_indexes.py}, a table of up to
1{,}000{,}000 rows of about 80 bytes, column \textit{k} a random permutation
(so index order and table order are unrelated). SQLite 3.51.2, Intel i5-6600,
16\,GB, Windows 11, mechanical disk. Best of three.}

\medskip
\textbf{1. A point lookup} --- \texttt{SELECT v FROM t WHERE k = ?}

\begin{tabular}{@{}R{2.6cm}R{2.6cm}R{2.6cm}R{2.6cm}@{}}
\textbf{rows} & \textbf{scan} & \textbf{index} & \textbf{speed-up} \\
10{,}000    & 1.37\,ms   & 0.053\,ms & 26$\times$ \\
100{,}000   & 19.03\,ms  & 0.053\,ms & 360$\times$ \\
1{,}000{,}000 & 199.90\,ms & 0.057\,ms & \textbf{3{,}501$\times$} \\
\end{tabular}

\smallskip
The scan grows in proportion to the table. The index barely moves: a hundred
times the rows added one level, at most, to a three-level tree.

\medskip
\textbf{2. The crossover} --- \texttt{SUM(length(v)) WHERE k BETWEEN 0 AND
}$x$, each plan forced, over a growing share of a million rows

\begin{tabular}{@{}R{1.6cm}R{2.2cm}R{2.4cm}R{2.4cm}L{1.6cm}L{2.6cm}@{}}
\textbf{share} & \textbf{rows} & \textbf{via index} & \textbf{full scan}
& \textbf{faster} & \textbf{optimiser chose} \\
0.1\% & 1{,}000    & 7.1\,ms    & 178.1\,ms & index & index \\
1\%   & 10{,}000   & 90.6\,ms   & 207.0\,ms & index & index \\
2\%   & 20{,}000   & 182.0\,ms  & 210.0\,ms & index & index \\
5\%   & 50{,}000   & 436.4\,ms  & 190.3\,ms & scan  & \textbf{index} \\
10\%  & 100{,}000  & 784.3\,ms  & 202.5\,ms & scan  & \textbf{index} \\
20\%  & 200{,}000  & 1{,}549.9\,ms & 220.1\,ms & scan & \textbf{index} \\
50\%  & 500{,}000  & 3{,}955.3\,ms & 287.5\,ms & scan & \textbf{index} \\
100\% & 1{,}000{,}000 & 7{,}934.8\,ms & 381.0\,ms & scan & \textbf{index} \\
\end{tabular}

\smallskip
\textbf{The index stops winning at about 2\% of the table.} Below that it reads
only the rows it needs. Above it, every matching row is a separate jump to a
different page of the table, a \term{random read} --- and a scan, which
reads each page once and in order, a \term{sequential read}, wins by more
as more rows are wanted. At the whole table the index plan
was \textbf{21$\times$ slower}.

\medskip
\textbf{3. Why the optimiser chose the index every time}

\begin{tabular}{@{}L{7.8cm}L{6.0cm}@{}}
STAT4 histograms compiled into this SQLite & \textbf{no} \\
what \sqlkw{ANALYZE} recorded in \texttt{sqlite\_stat1} & \texttt{'1000000 1'} \\
the whole table, as written & 7{,}952.4\,ms (\texttt{SEARCH \dots\ ik}) \\
the whole table, with \texttt{+k} & 407.2\,ms (\texttt{SCAN t}) \\
\end{tabular}

\smallskip
\texttt{'1000000 1'} is everything the optimiser knew: a million rows, one row
per value of \textit{k}. It knew \textit{k} was unique. It did \emph{not} know
how the values were spread, so it could not tell \texttt{BETWEEN 0 AND 999}
from \texttt{BETWEEN 0 AND 999999}, and treated every range on an indexed
column as selective. The histograms that would have told it --- STAT4 --- are
not compiled into the SQLite that Python ships.

\medskip
\textbf{4. What indexes cost} --- 100{,}000 rows; insert 10{,}000 more; median
of 3

\begin{tabular}{@{}L{3.2cm}R{3.0cm}R{3.0cm}R{3.2cm}@{}}
& \textbf{file bytes} & \textbf{insert} & \textbf{(range)} \\
0 indexes & 8{,}232{,}960  & 71.5\,ms  & 47.6--79.0 \\
1 on \textit{k} & 9{,}383{,}936 & 64.8\,ms & 62.0--69.2 \\
+ 1 on \textit{v} (68-byte text) & 17{,}272{,}832 & 137.9\,ms & 123.8--150.2 \\
+ 1 on (\textit{k}, \textit{v}) & 25{,}485{,}312 & 141.5\,ms & 135.2--168.9 \\
\end{tabular}

\smallskip
\textbf{Read the ranges first.} Zero and one index overlap: the new rows arrived
in key order, so each new entry went on the right-hand end of the index --- the
cheapest case, as for roll numbers issued in sequence. The index on the wide
text column is the one that cost: it doubled the insert time and added more
bytes than the table had.
\end{measured}

\begin{alertbox}[What the crossover means in practice]
\begin{tight}
  \item \textbf{An index is for finding a few rows, and ``a few'' must be
        measured.} About 2\% for this million-row table. On the seeded
        \textit{enrolment} table --- 60{,}037 narrower rows --- an index on
        \textit{grade} selecting 21\% ran in 8.94\,ms against a scan's
        9.40\,ms: a tie, so the crossover there sits near a fifth. The share
        moves with table size, row width and cache; the shape --- a share past
        which the index loses --- does not.
  \item \textbf{The optimiser can be wrong, and here it was wrong every time
        past 2\%} --- not through a bug, but because it was never told how the
        data was spread. ``The database will choose the best plan'' assumes it
        has the information to.
  \item \textbf{Two ways to overrule it.} Write \texttt{+k} instead of
        \texttt{k} in the condition, which hides the column from the index for
        that one query; or name the table \texttt{t NOT INDEXED}. The opposite,
        \texttt{t INDEXED BY ik}, insists on a named index --- the harness
        used both to measure each plan at every share. Use them only on a
        query you have measured.
  \item \textbf{Low-selectivity columns rarely deserve an index.} An index on
        \textit{programme} (four values) or \textit{grade} (six) selects a
        quarter or a fifth of the table for every value --- at or past the
        crossover. Measured above, the grade index only tied a scan: it bought
        nothing, and it still costs bytes and a write on every change.
\end{tight}
\end{alertbox}

\section{Using an Index Well}

\begin{definitionbox}[Composite index, leftmost prefix, covering index]
A \term{composite index} on $(a, b)$ is sorted by $a$, and by $b$ within each
$a$ --- like a telephone directory sorted by surname, then first name.

\smallskip
It can serve a condition on $a$, or on $a$ and $b$ together, but \textbf{not
on $b$ alone}: the $b$ values are scattered through every $a$. This is the
\term{leftmost-prefix} rule, and \chref{subqueries}'s measurement turned on it
--- \textit{enrolment}'s primary key is (roll\_no, course\_code, \dots), so it
could find a student's enrolments and could not find a course's.

\smallskip
A \term{covering index} contains every column a query needs, so the table
itself is never read. \texttt{USING COVERING INDEX} in a plan means exactly
that, and it removes the per-row jump that loses past the crossover.
\end{definitionbox}

\begin{examplebox}[Ordering: what an index gives for free, and what it cannot]
From part~3 of the measurement, on 100{,}000 rows with an index on \textit{k}:

\medskip
\begin{tabular}{@{}L{6.4cm}R{2.0cm}L{5.4cm}@{}}
\texttt{ORDER BY k LIMIT 10} & 0.08\,ms & reads in index order, stops at 10 \\
\texttt{ORDER BY v LIMIT 10} & 23.87\,ms & \texttt{TEMP B-TREE}: sorts every row \\
\texttt{ORDER BY k \% 7, k LIMIT 10} & 24.74\,ms & an expression: sorts too \\
\end{tabular}

\smallskip
An index is a sorted copy of \emph{one} ordering. A query that wants that
order gets it for nothing --- and, with \sqlkw{LIMIT}, stops after ten rows
instead of reading a hundred thousand. Any other order, \textbf{including an
expression of the indexed column itself}, means sorting everything first. The
third row is the one people do not expect.
\end{examplebox}

\begin{pitfall}
\begin{tight}
  \item \textbf{Indexing every column.} Each is another B-tree kept in step
        on every write. Measured: the text index doubled insert time.
  \item \textbf{Indexing a column with few values.} Every lookup selects a
        large share --- past the crossover.
  \item \textbf{Expecting an index on $(a, b)$ to serve a condition on $b$.}
        Leftmost prefix only.
  \item \textbf{Wrapping the indexed column in an expression.}
        \texttt{WHERE lower(name) = 'ali'} cannot use an index on
        \textit{name}. (SQLite allows an index on the expression itself.)
  \item \textbf{Trusting the optimiser's choice without measuring.} It picked
        a 21$\times$ slower plan here, after \sqlkw{ANALYZE}.
  \item \textbf{Measuring an index on a table whose order matches the
        column.} Then ``random'' lookups are sequential and the crossover
        vanishes. The harness shuffles \textit{k} for exactly this reason.
\end{tight}
\end{pitfall}

\begin{lab}[Find the crossover on your own machine]
\begin{lstlisting}[language=bash]
# --- 1. The measured version first: it builds a million rows --------
python code/m21_indexes.py --quick
\end{lstlisting}
\begin{lstlisting}[language=sql]
-- --- 2. On the university database, the indexes you already have -
.open university.db
SELECT name, tbl_name, sql FROM sqlite_master WHERE type = 'index';
-- Some have no sql: they were made by PRIMARY KEY and UNIQUE.

-- --- 3. Leftmost prefix, seen in a plan --------------------------
EXPLAIN QUERY PLAN SELECT * FROM enrolment WHERE roll_no = 'x';
EXPLAIN QUERY PLAN SELECT * FROM enrolment WHERE course_code = 'x';
-- Same primary-key index. One SEARCH, one SCAN. Why?

-- --- 4. A covering index ----------------------------------------
CREATE INDEX idx_enr_cg ON enrolment(course_code, grade);
EXPLAIN QUERY PLAN
SELECT grade, COUNT(*) FROM enrolment WHERE course_code = 'CS-101'
 GROUP BY grade;
-- Find COVERING. Which table was never read?

-- --- 5. An index on a column with six values --------------------
CREATE INDEX idx_enr_grade ON enrolment(grade);
.timer on
SELECT COUNT(*), AVG(section_no) FROM enrolment WHERE grade = 'B';
SELECT COUNT(*), AVG(section_no) FROM enrolment WHERE +grade = 'B';
-- Which was faster? What share of the table did 'B' select?

-- --- 6. What ANALYZE knows --------------------------------------
ANALYZE;
SELECT * FROM sqlite_stat1;
-- Read each stat string: rows, then rows per key value.

-- --- 7. An order the index cannot give --------------------------
EXPLAIN QUERY PLAN SELECT * FROM student ORDER BY name LIMIT 10;
CREATE INDEX idx_student_name ON student(name);
EXPLAIN QUERY PLAN SELECT * FROM student ORDER BY name LIMIT 10;
EXPLAIN QUERY PLAN SELECT * FROM student ORDER BY lower(name) LIMIT 10;
\end{lstlisting}
\begin{lstlisting}[language=bash]
# --- 8. The full version, a million rows (about a minute) -----------
python code/m21_indexes.py
\end{lstlisting}

\smallskip
\textbf{What to notice.} Step 5 is the crossover in your own schema. Grade B
selects 21\% of enrolments, and on the seeded table the two forms come out
level --- 8.94 against 9.40\,ms on the measuring machine --- while the
optimiser chooses the index without hesitation. Compare that with the 2\% the
million-row table gave: the same index, the same optimiser, a crossover ten
times further out, because a small table is cheap to scan. \textbf{The rule is
not ``2\%''. The rule is that the share exists and has to be measured.}
\end{lab}

\begin{checkpoint}
\begin{tightnum}
  \item A lookup took 0.053\,ms at 10{,}000 rows and 0.057\,ms at a million,
        while a scan went from 1.37 to 199.90\,ms. Explain both using the
        fan-out and depth of a B-tree, without assuming any data-structures
        course.
  \item The optimiser chose the index for a query that selected the whole
        table, after \sqlkw{ANALYZE}, and it was 21$\times$ slower. Explain
        what information it lacked, where that information would have come
        from, and two ways to make it scan instead.
  \item An index on \textit{enrolment}(roll\_no, course\_code) exists. For each
        of \texttt{WHERE roll\_no = ?}, \texttt{WHERE course\_code = ?} and
        \texttt{WHERE roll\_no = ? AND course\_code = ?}, say whether it can be
        used and why.
\end{tightnum}
\end{checkpoint}

\begin{chaptersummary}
\begin{tight}
  \item An \textbf{index} is every value of a column, in sorted order, each
        with the rowid of its row --- kept by SQLite as a \textbf{B-tree}.
  \item A B-tree's \textbf{fan-out} is a few hundred, so a million entries fit
        in three levels and a lookup reads one page per level.
  \item \textbf{Measured:} a point lookup was \textbf{3{,}501$\times$} faster
        with an index at a million rows, and cost the same at every size.
  \item \textbf{Measured: the index stopped winning at about 2\%} of the table.
        Past that, a jump per row loses to reading every page once, in order;
        at the whole table the index was \textbf{21$\times$ slower}.
  \item \textbf{Measured: the optimiser chose the index at every share}, after
        \sqlkw{ANALYZE}, because \texttt{sqlite\_stat1} records only row counts
        and STAT4 histograms are not compiled into Python's SQLite.
  \item Override a bad plan with \texttt{+column} or \sqlkw{NOT INDEXED} ---
        on a query you have measured.
  \item A \textbf{composite} index serves only its \textbf{leftmost prefix}; a
        \textbf{covering} index means the table is never read.
  \item An index gives its own order free, including for \sqlkw{LIMIT}; any
        other order, even an expression of the indexed column, costs a sort.
  \item \textbf{Measured:} indexes cost bytes and writes; a wide text index
        doubled insert time, while an integer index on in-order keys cost
        nothing measurable.
\end{tight}
\end{chaptersummary}

\begin{reviewq}
  \item \textbf{(MCQ)} An index on $(a, b)$ can be used for a condition on:
        (a)~$b$ alone; (b)~$a$ alone; (c)~either; (d)~neither.
  \item \textbf{(MCQ)} \texttt{USING COVERING INDEX} in a plan means:
        (a)~the index is unique; (b)~the table was not read; (c)~the index
        covers the whole table; (d)~a temporary index was built.
  \item \textbf{(Short)} Why does a query that wants half a table run faster as
        a scan than through an index?
  \item \textbf{(Short)} What does \texttt{sqlite\_stat1}'s string
        \texttt{'1000000 1'} tell the optimiser, and what does it not?
  \item \textbf{(Short)} Why is an index on \textit{grade} unlikely to help?
  \item \textbf{(Applied)} For the university schema, propose the three indexes
        you would add, justify each with a query from Part~IV that it serves
        and the plan it would change, and name one index you would refuse to
        add, with the reason.
  \item \textbf{(Applied)} Re-run part~2 of the harness with \textit{k} in the
        same order as the rowid instead of shuffled (change one line), and
        report what happens to the crossover. Explain the result in terms of
        sequential and random reads, and say why the shuffled version is the
        honest measurement.
\end{reviewq}

%=====================================================================
\chapter{Transactions, Concurrency and Recovery}\label{ch:transactions}
%=====================================================================
\locator{6}{Part VI --- Making It Work, and What Comes Next}

\begin{outcomes}
By the end of this chapter you will be able to:
\begin{tight}
  \item \textbf{State} the four ACID properties and \textbf{give} a failure
        that each one prevents.
  \item \textbf{Show} a lost update happening, and \textbf{prevent} it with
        the right kind of transaction.
  \item \textbf{Explain} SQLite's locks and why it allows one writer at a time.
  \item \textbf{Distinguish} the deadlock SQLite can have from the one a server
        DBMS has, and \textbf{say} how each is resolved.
  \item \textbf{Back up} a live database safely, and \textbf{state} what the
        journal does after a crash.
\end{tight}
\end{outcomes}

\begin{prereq}
\chref{whynotfiles}: atomicity and concurrent access were two of the six
problems. \chref{whodoeswhat}: the transaction, lock and recovery managers.
\chref{evolution}: \sqlkw{BEGIN} and \sqlkw{COMMIT} in a migration. No
operating-systems course is assumed. You need
\texttt{code/m22\_transactions.py}.
\end{prereq}

\begin{keyterms}
transaction $\bullet$ \sqlkw{BEGIN} $\bullet$ \sqlkw{COMMIT} $\bullet$
\sqlkw{ROLLBACK} $\bullet$ \sqlkw{SAVEPOINT} $\bullet$ ACID $\bullet$
atomicity $\bullet$ consistency $\bullet$ isolation $\bullet$ durability
$\bullet$ autocommit $\bullet$ lost update $\bullet$ dirty read $\bullet$
serialisable $\bullet$ lock $\bullet$ \sqlkw{SHARED} $\bullet$
\sqlkw{RESERVED} $\bullet$ \sqlkw{EXCLUSIVE} $\bullet$ \sqlkw{DEFERRED}
$\bullet$ \sqlkw{IMMEDIATE} $\bullet$ deadlock $\bullet$ \texttt{SQLITE\_BUSY}
$\bullet$ \texttt{busy\_timeout} $\bullet$ rollback journal $\bullet$ WAL
$\bullet$ \texttt{synchronous}
\end{keyterms}

\section{Promises About Groups of Changes}

Every constraint in Part~II was a promise about a row. A transaction is a
promise about a \emph{group} of changes: that they happen together, or not at
all, and that nothing anybody else does at the same moment can see them half
done. \chref{whynotfiles} listed atomicity and concurrent access among the six
problems of files. This chapter is where the database keeps those promises ---
and measures what keeping them costs.

\begin{definitionbox}[Transaction, and ACID]
A \term{transaction} is a sequence of operations treated as one unit, begun
with \sqlkw{BEGIN} and ended with \sqlkw{COMMIT} (keep it) or \sqlkw{ROLLBACK}
(undo it all). Outside one, SQLite is in \term{autocommit}: every statement is
its own transaction.

\smallskip
\textbf{Atomicity}: all of it happens, or none. \textbf{Consistency}: it takes
the database from one state satisfying every constraint to another.
\textbf{Isolation}: concurrent transactions do not see each other's
unfinished work. \textbf{Durability}: once committed, it survives a crash.
\end{definitionbox}

\begin{measured}[A crash, a lost update, a deadlock, and the price of durability]
\textsf{\footnotesize \texttt{code/m22\_transactions.py}; two real connections
to one file, real failures. SQLite 3.51.2, Intel i5-6600, Windows 11,
mechanical disk.}

\medskip
\textbf{1. Atomicity} --- move a student between sections, a \sqlkw{DELETE}
then an \sqlkw{INSERT}, and fail between them

\begin{tabular}{@{}L{4.4cm}L{3.4cm}L{6.0cm}@{}}
& \textbf{enrolled in} & \\
autocommit & \textbf{nothing} & the \sqlkw{DELETE} was already permanent \\
\sqlkw{BEGIN} \dots\ \sqlkw{ROLLBACK} & section 1 & the two were one unit \\
\end{tabular}

\medskip
\textbf{2. The lost update} --- two clerks each add one to a seat count of 40:
A reads, B reads, A writes, B writes

\begin{tabular}{@{}L{3.4cm}L{8.6cm}R{1.6cm}@{}}
& \textbf{what happened} & \textbf{final} \\
no transaction & both read 40, both wrote 41 & \textbf{41} \\
\sqlkw{BEGIN DEFERRED} & B's write refused: \texttt{database is locked}; then
  A's \sqlkw{COMMIT} refused too; B rolled back; A committed & 41 \\
\sqlkw{BEGIN IMMEDIATE} & B could not begin until A had finished; then read 41
  and wrote 42 & \textbf{42} \\
\end{tabular}

\smallskip
With no transaction an enrolment simply vanished, and nobody was told. With
\sqlkw{DEFERRED} the two clerks \textbf{deadlocked}: B held a read lock and
wanted to write; A held the write lock and needed B's read lock released before
it could commit. Neither could move until B gave up --- so B's +1 was lost
again, \emph{but B was told}, and can retry. \sqlkw{IMMEDIATE} counted both.

\medskip
\textbf{3. \texttt{SQLITE\_BUSY}} --- A holds the write lock for 300\,ms; B
tries to write

\begin{tabular}{@{}L{4.4cm}L{6.0cm}R{3.0cm}@{}}
\textbf{B's busy\_timeout} & \textbf{outcome} & \textbf{B waited} \\
0\,ms (the default) & failed at once: \texttt{database is locked} & 0\,ms \\
1{,}000\,ms & succeeded after A committed & 455\,ms \\
\end{tabular}

\smallskip
B waited for A's 300\,ms, A's commit, and the interval at which SQLite's busy
handler retries --- and then went ahead.

\medskip
\textbf{4. Durability, priced} --- 200 inserts

\begin{tabular}{@{}L{7.4cm}R{2.8cm}R{2.8cm}@{}}
& \textbf{total} & \textbf{per row} \\
a commit each, \texttt{synchronous = FULL} & 4{,}590.9\,ms & 22.954\,ms \\
one commit for all, \texttt{synchronous = FULL} & 27.8\,ms & 0.139\,ms \\
a commit each, \texttt{synchronous = OFF} & 239.2\,ms & 1.196\,ms \\
\end{tabular}

\smallskip
\textbf{One commit for all 200 was 165$\times$ faster than one each.} A commit
waits until the disk confirms the data is really on it; on a mechanical disk
that wait is 23\,ms and is almost the whole cost. \texttt{synchronous = OFF}
skips the wait --- 19$\times$ faster, still 9$\times$ slower than batching ---
and gives up durability: a power cut can lose transactions that were reported
as committed, or corrupt the file. Batching keeps the guarantee \emph{and} is
faster.
\end{measured}

\begin{alertbox}[Four rules from the measurement]
\begin{tight}
  \item \textbf{Anything that is two statements and must be one fact goes in a
        transaction.} Measured: autocommit left a student enrolled in nothing.
  \item \textbf{A read-then-write goes in \sqlkw{BEGIN IMMEDIATE}.} It takes
        the write lock before the read, so nobody else can read the old value
        in between. \sqlkw{DEFERRED} (the default) deadlocked.
  \item \textbf{Set \texttt{busy\_timeout} on every connection.} The default
        is 0, so a second writer is refused instead of waiting. It is per
        connection, like \chref{constraints}'s \texttt{foreign\_keys}, and for
        the same reason has to be set every time.
  \item \textbf{Group writes into transactions; never turn durability off to
        go faster.} The guarantee and the speed are not in conflict here.
\end{tight}
\end{alertbox}

\section{How SQLite Isolates: One Lock for the Whole File}

\begin{definitionbox}[SQLite's locks]
SQLite locks the \emph{database file}, not rows or tables. A connection
progresses through:

\smallskip
\begin{tabular}{@{}L{2.8cm}L{11.0cm}@{}}
\sqlkw{SHARED} & may read. Any number of connections may hold it. \\
\sqlkw{RESERVED} & intends to write. Only one at a time; readers continue. \\
\sqlkw{EXCLUSIVE} & is committing. Needed to write to the file; requires
  every \sqlkw{SHARED} lock to be released first. \\
\end{tabular}

\smallskip
\sqlkw{BEGIN DEFERRED} takes no lock until the first read (then
\sqlkw{SHARED}) or write (then \sqlkw{RESERVED}). \sqlkw{BEGIN IMMEDIATE}
takes \sqlkw{RESERVED} at once. \sqlkw{BEGIN EXCLUSIVE} takes
\sqlkw{EXCLUSIVE} at once.

\smallskip
The consequence: \textbf{one writer at a time, for the whole database}, and
transactions that are \term{serialisable} --- the result is always as if they
had run one after another. There are no weaker isolation levels to choose
from, because there is nothing to relax --- and so no \term{dirty read}:
no connection can ever see another's uncommitted change.

\smallskip
A \term{savepoint} marks a point inside a transaction to roll back to
without abandoning the whole of it: \texttt{SAVEPOINT s1;} \dots\
\texttt{ROLLBACK TO s1;} undoes only the work since \texttt{s1}.
\end{definitionbox}

\dsfig{%
\begin{tikzpicture}[
  tx/.style={rounded corners=3pt, draw=#1, line width=1pt, fill=#1!8,
             text=#1!60!black, font=\scriptsize\bfseries, align=center,
             minimum width=2.2cm, minimum height=0.8cm},
  rs/.style={rounded corners=3pt, draw=neutral, line width=0.9pt,
             fill=neutral!8, text=neutral!80!black, font=\scriptsize,
             align=center, minimum width=2.4cm, minimum height=0.8cm},
  hold/.style={-{Stealth[length=2.2mm]}, greenhl!70!black, line width=1.1pt},
  want/.style={-{Stealth[length=2.2mm]}, accent, line width=1.1pt, dashed},
  ttl/.style={font=\scriptsize\bfseries, text=primary, anchor=south},
  lb/.style={font=\scriptsize\itshape, text=#1}
]
  % --- SQLite: one file, a lock upgrade
  \node[ttl] at (-3.9,1.85) {SQLite: one lock, upgraded};
  \node[tx=secondary] (a) at (-5.6,0.9) {A};
  \node[tx=purplehl]  (b) at (-2.2,0.9) {B};
  \node[rs] (f) at (-3.9,-0.6) {the database file};
  \draw[hold] (a.south) -- node[lb=greenhl!60!black, left, pos=0.45]
       {holds RESERVED} (f.north west);
  \draw[hold] (b.south) -- node[lb=greenhl!60!black, right, pos=0.45]
       {holds SHARED} (f.north east);
  \draw[want] (a.east) -- node[lb=accent, above] {needs B to let go} (b.west);
  \node[lb=neutral, text width=6.0cm, align=center] at (-3.9,-1.55)
       {A cannot commit while B reads; B cannot write while A holds
        RESERVED. \textbf{Measured} --- SQLite refuses, one rolls back.};

  % --- server: two rows, a cycle
  \node[ttl] at (3.9,1.85) {a server DBMS: two rows, a cycle};
  \node[tx=secondary] (c) at (2.2,0.9) {T1};
  \node[tx=purplehl]  (d) at (5.6,0.9) {T2};
  \node[rs] (r1) at (2.2,-0.6) {row: section 1};
  \node[rs] (r2) at (5.6,-0.6) {row: section 2};
  \draw[hold] (c) -- node[lb=greenhl!60!black, left] {locks} (r1);
  \draw[hold] (d) -- node[lb=greenhl!60!black, right] {locks} (r2);
  \draw[want] (c.south east) -- node[lb=accent, above, sloped, pos=0.35]
       {wants} (r2.north west);
  \draw[want] (d.south west) -- node[lb=accent, above, sloped, pos=0.35]
       {wants} (r1.north east);
  \node[lb=neutral, text width=6.0cm, align=center] at (3.9,-1.55)
       {Each waits for the other's row for ever. The server finds the cycle
        and kills one. \textbf{Cannot happen in SQLite}: there are no row
        locks.};
\end{tikzpicture}}%
{Two kinds of deadlock. The left one was produced by the harness with two
\sqlkw{DEFERRED} transactions. The right one needs row-level locks, which
SQLite does not have, so it is drawn here rather than run --- \chref{project}
is where a server's locking is compared properly.}{deadlock}

\begin{conceptbox}[What SQLite does not have, and what you would get instead]
A server DBMS --- PostgreSQL, MySQL, Oracle --- locks rows, and most use
\term{multi-version concurrency control}: a writer creates a new version of a
row while readers keep seeing the old one, so readers and writers never block
each other. It offers several \emph{isolation levels}, from \textsc{read
committed} to \textsc{serializable}, trading correctness for concurrency.

\smallskip
SQLite has none of that, and for one database file on one machine it does not
need it. What it costs is concurrency: many readers, \textbf{one} writer. For a
university's results system with a few clerks it is plenty; for a booking site
with a thousand simultaneous purchases it is the reason to move to a server,
and \chref{project} makes that comparison.
\end{conceptbox}

\section{Recovery, and Backup}

\begin{definitionbox}[The journal, and what happens after a crash]
Before changing a page of the database file, SQLite writes the page's
\emph{original} contents to a \term{rollback journal}, \texttt{university.db-journal}.
Only then does it change the file. At \sqlkw{COMMIT} it deletes the journal.

\smallskip
After a crash, the next connection finds a journal still present, knows a
transaction was interrupted, and copies the original pages back. The database
returns to its state before the transaction began --- atomicity and
durability, kept by one extra file.

\smallskip
\term{WAL} mode (\texttt{PRAGMA journal\_mode = WAL}) reverses the idea:
changes are appended to a write-ahead log and copied into the database later.
Readers then never block the writer, which is the nearest SQLite comes to a
server's concurrency.
\end{definitionbox}

\begin{examplebox}[Backing up a database that is in use]
\textbf{The wrong way:} \texttt{copy university.db backup.db} while a program
is writing. If the copy happens mid-transaction it captures a half-written
file without its journal --- a backup that may not open, discovered on the day
it is needed.

\smallskip
\textbf{Two right ways}, both of which take a consistent snapshot:
\begin{lstlisting}[language=sql]
.backup backup.db                     -- in the sqlite3 shell
VACUUM INTO 'backup.db';              -- as SQL, from any program
\end{lstlisting}

\smallskip
And \chref{evolution}'s rule applies: \textbf{open the backup} and run
\texttt{PRAGMA integrity\_check;} on it. A backup that has never been restored
is a hope.
\end{examplebox}

\begin{pitfall}
\begin{tight}
  \item \textbf{Two statements that must be one fact, in autocommit.} Measured:
        a student enrolled in nothing.
  \item \textbf{Read-then-write in \sqlkw{BEGIN DEFERRED}.} Measured: a
        deadlock. Use \sqlkw{IMMEDIATE}.
  \item \textbf{Leaving \texttt{busy\_timeout} at 0.} A second writer is
        refused instead of waiting.
  \item \textbf{A commit per row in a bulk load.} Measured: 165$\times$ slower.
  \item \textbf{\texttt{synchronous = OFF} for speed.} It gives up durability
        and is still slower than batching.
  \item \textbf{Copying a live database file as a backup.} Use
        \texttt{.backup} or \sqlkw{VACUUM INTO}, then open the result.
  \item \textbf{Python's default \texttt{isolation\_level}.} The
        \texttt{sqlite3} module silently issues \sqlkw{BEGIN} for you before
        writes. The harness sets \texttt{isolation\_level=None} so every
        \sqlkw{BEGIN} is one you wrote.
\end{tight}
\end{pitfall}

\begin{lab}[Lose an enrolment, deadlock two clerks, and back up safely]
\begin{lstlisting}[language=bash]
# --- 1. Two shells on ONE file: open two terminals ------------------
sqlite3 txn.db          # terminal A
sqlite3 txn.db          # terminal B
\end{lstlisting}
\begin{lstlisting}[language=sql]
-- --- 2. In A only: a table to fight over ------------------------
CREATE TABLE section (id TEXT PRIMARY KEY, seats_taken INTEGER) STRICT;
INSERT INTO section VALUES ('CS-103/1', 40);

-- --- 3. The lost update, by hand: type in this order -------------
-- A:  SELECT seats_taken FROM section;         -- 40
-- B:  SELECT seats_taken FROM section;         -- 40
-- A:  UPDATE section SET seats_taken = 41;
-- B:  UPDATE section SET seats_taken = 41;
-- A:  SELECT seats_taken FROM section;         -- 41. Two clerks, +1.

-- --- 4. The deadlock: reset to 40, then type in this order -------
-- A:  BEGIN DEFERRED;  SELECT seats_taken FROM section;
-- B:  BEGIN DEFERRED;  SELECT seats_taken FROM section;
-- A:  UPDATE section SET seats_taken = 41;     -- works
-- B:  UPDATE section SET seats_taken = 41;     -- database is locked
-- A:  COMMIT;                                  -- also locked. Why?
-- B:  ROLLBACK;
-- A:  COMMIT;                                  -- now it works

-- --- 5. The fix: reset to 40, then ------------------------------
-- A:  BEGIN IMMEDIATE;  SELECT seats_taken FROM section;
-- B:  BEGIN IMMEDIATE;                         -- refused: wait
-- A:  UPDATE section SET seats_taken = 41;  COMMIT;
-- B:  BEGIN IMMEDIATE;  SELECT seats_taken FROM section;   -- 41
-- B:  UPDATE section SET seats_taken = 42;  COMMIT;

-- --- 6. Waiting instead of failing ------------------------------
-- B:  .timeout 5000
-- Repeat step 5. What does B do now instead of refusing?

-- --- 7. A safe backup -------------------------------------------
.backup txn_backup.db
\end{lstlisting}
\begin{lstlisting}[language=bash]
# --- 8. Open the backup and check it --------------------------------
sqlite3 txn_backup.db "PRAGMA integrity_check; SELECT * FROM section;"

# --- 9. The measured version ----------------------------------------
python code/m22_transactions.py
\end{lstlisting}

\smallskip
\textbf{What to notice.} Step 4's second \sqlkw{COMMIT} is the one to explain.
A already holds the write lock; to commit it needs the \sqlkw{EXCLUSIVE} lock,
which requires every reader to leave --- and B is a reader, waiting to become a
writer. Neither can move. SQLite cannot wait its way out of that, so it
refuses, and one of the two must roll back.
\end{lab}

\begin{checkpoint}
\begin{tightnum}
  \item For each ACID property, give the measured result in this chapter that
        shows it being kept or broken, and say what would have happened
        without it.
  \item Two clerks deadlocked under \sqlkw{BEGIN DEFERRED} and both increments
        counted under \sqlkw{BEGIN IMMEDIATE}. Explain both outcomes in terms of
        which lock each connection held and wanted at each step.
  \item A bulk load of 200 rows took 4.6 seconds with a commit per row and
        28\,ms in one transaction. Explain the difference, say why
        \texttt{synchronous = OFF} is not the fix, and say what it would cost
        the university if it were used.
\end{tightnum}
\end{checkpoint}

\begin{chaptersummary}
\begin{tight}
  \item A \textbf{transaction} groups changes into one unit with
        \sqlkw{BEGIN}, \sqlkw{COMMIT} and \sqlkw{ROLLBACK}. Outside one, every
        statement is its own (\textbf{autocommit}).
  \item \textbf{ACID}: atomicity, consistency, isolation, durability.
  \item \textbf{Measured:} a failure between a \sqlkw{DELETE} and an
        \sqlkw{INSERT} left a student \textbf{enrolled in nothing} in
        autocommit; inside a transaction, nothing changed.
  \item \textbf{Measured:} two read-then-write clerks \textbf{lost an update}
        with no transaction, \textbf{deadlocked} under \sqlkw{DEFERRED}, and
        were both counted under \sqlkw{IMMEDIATE}.
  \item SQLite locks the whole file --- \sqlkw{SHARED}, \sqlkw{RESERVED},
        \sqlkw{EXCLUSIVE} --- so it has one writer at a time and is always
        serialisable.
  \item \textbf{Measured:} with \texttt{busy\_timeout} 0 a second writer was
        refused; with 1{,}000\,ms it waited 455\,ms and succeeded. Set it on
        every connection.
  \item A server's two-row deadlock cycle cannot occur in SQLite, which has no
        row locks; SQLite's deadlock is two connections on the one file lock.
  \item \textbf{Measured:} one commit for 200 inserts was \textbf{165$\times$}
        faster than one each. \texttt{synchronous = OFF} gave up durability and
        was still 9$\times$ slower than batching.
  \item The \textbf{rollback journal} holds original pages until commit, and is
        replayed after a crash. Back up with \texttt{.backup} or
        \sqlkw{VACUUM INTO}, never by copying a live file --- then open the
        backup.
\end{tight}
\end{chaptersummary}

\begin{reviewq}
  \item \textbf{(MCQ)} In autocommit mode, each statement is: (a)~never
        committed; (b)~its own transaction; (c)~part of one transaction per
        connection; (d)~rolled back on exit.
  \item \textbf{(MCQ)} \sqlkw{BEGIN IMMEDIATE} acquires: (a)~no lock;
        (b)~\sqlkw{SHARED}; (c)~\sqlkw{RESERVED}; (d)~a row lock.
  \item \textbf{(Short)} Why does a read-then-write need its write lock before
        the read?
  \item \textbf{(Short)} Why can SQLite never have a two-row deadlock cycle?
  \item \textbf{(Short)} What does the rollback journal contain, and what is
        done with it after a crash?
  \item \textbf{(Applied)} Write the transaction that moves a student from one
        section to another and adjusts both sections' seat counts. Say which
        kind of \sqlkw{BEGIN} you would use and why, what happens if it fails
        after the second statement, and how a second clerk doing the same move
        at the same moment would be affected.
  \item \textbf{(Applied)} The university's results are entered by six clerks
        at once during the last week of term. Using the measured lock
        behaviour, explain what each clerk would experience with
        \texttt{busy\_timeout} 0 and with 5{,}000\,ms, estimate how long a
        clerk might wait if each commit took 23\,ms, and say at what number of
        clerks you would recommend moving to a server DBMS --- and why the
        answer is not a fixed number.
\end{reviewq}

%=====================================================================
\chapter{Security, NoSQL, and the Semester Project}\label{ch:project}
%=====================================================================
\locator{6}{Part VI --- Making It Work, and What Comes Next}

\begin{outcomes}
By the end of this chapter you will be able to:
\begin{tight}
  \item \textbf{Explain} SQL injection, \textbf{demonstrate} it against your
        own database, and \textbf{prevent} it with parameters.
  \item \textbf{Write} the \sqlkw{GRANT} and \sqlkw{REVOKE} statements a server
        would need for the university's external schemas, and \textbf{say} what
        SQLite offers instead.
  \item \textbf{Describe} the four families of NoSQL database and
        \textbf{state} what each gives up to get what it offers.
  \item \textbf{Store and query} a document inside SQLite using its JSON
        functions, and \textbf{compare} it with the relational form.
  \item \textbf{Decide} whether a given system needs a server DBMS, using the
        limits measured in this course.
\end{tight}
\end{outcomes}

\begin{prereq}
Appendix~A, section A.6 --- \textbf{the scope statement, which governs this
chapter's laboratory}. \chref{whodoeswhat} for the authorisation manager SQLite
lacks; \chref{views} for external schemas; \chref{transactions} for the
one-writer limit. The finished \texttt{university.db}.
\end{prereq}

\begin{keyterms}
SQL injection $\bullet$ parameterised query $\bullet$ placeholder $\bullet$
least privilege $\bullet$ \sqlkw{GRANT} $\bullet$ \sqlkw{REVOKE} $\bullet$
role $\bullet$ encryption at rest $\bullet$ NoSQL $\bullet$ key--value store
$\bullet$ document store $\bullet$ wide-column store $\bullet$ graph database
$\bullet$ schema-on-read $\bullet$ horizontal scaling $\bullet$ eventual
consistency $\bullet$ CAP $\bullet$ JSON
\end{keyterms}

\section{Security: the Component SQLite Does Not Have}

\chref{whodoeswhat} drew six components and said SQLite was missing one: the
authorisation manager. HEC's 2017 and 2023 outlines do not mention security at
all; HEC 2025's fifth learning outcome requires it, which is why it is here.
It has two halves --- keeping out people who should not be in, and making sure
that people who \emph{are} in cannot make the database do something it should
not --- and SQLite can help with only the second.

\begin{alertbox}[Scope --- this governs the laboratory below]
Everything in this chapter runs against \textbf{a database file you created,
on a machine you own}: your \texttt{university.db}, with invented data, which
you can delete. \textbf{Running any of it against a system you do not own and
have not been authorised in writing to test is a criminal offence in Pakistan}
under the \textbf{Prevention of Electronic Crimes Act 2016}, sections~3 and~4.
The purpose of the demonstration is that you never write the vulnerable code,
not that you can find it in somebody else's.
\end{alertbox}

\begin{definitionbox}[SQL injection]
\term{SQL injection} happens when a program builds a SQL statement by pasting
in text that a user supplied. If the text contains SQL, the database runs it
--- as part of the statement the programmer wrote, with the programmer's
permissions.

\smallskip
A \term{parameterised query} sends the statement and the values separately.
The statement has \term{placeholders} (\texttt{?}); the values are bound to
them afterwards and are \emph{never} parsed as SQL, whatever they contain.
\end{definitionbox}

\begin{worked}[One login check, built two ways]
A student portal checks a roll number before showing marks.

\medskip
\textbf{The vulnerable version} --- text pasted into SQL:
\begin{lstlisting}[language=python]
roll = input("Roll number: ")
sql = f"SELECT name FROM student WHERE roll_no = '{roll}'"
rows = cx.execute(sql).fetchall()
\end{lstlisting}

\textbf{Step 1 --- an ordinary input.} The roll number \texttt{2023-IT-00003}
produces the statement
\begin{lstlisting}[language=sql]
SELECT name FROM student WHERE roll_no = '2023-IT-00003'
\end{lstlisting}
and one row comes back. Correct.

\medskip
\textbf{Step 2 --- an input that is SQL.} \texttt{x' OR '1'='1} gives
\begin{lstlisting}[language=sql]
SELECT name FROM student WHERE roll_no = 'x' OR '1'='1'
\end{lstlisting}
The quote in the input closed the string the programmer opened; \texttt{OR
'1'='1'} is true for every row. On the seeded database it returns \textbf{all
4{,}000 students}.

\medskip
\textbf{Step 3 --- the parameterised version.}
\begin{lstlisting}[language=python]
rows = cx.execute("SELECT name FROM student WHERE roll_no = ?",
                  (roll,)).fetchall()
\end{lstlisting}
The same input now looks for a student whose roll number is, literally,
\texttt{x' OR '1'='1}. There is none: \textbf{0 rows}. The quote is a character
in a value, not the end of a string, because the value was never parsed.

\medskip
\textbf{The rule.} Never build SQL from input with \texttt{+}, \texttt{\%} or
an f-string. Every value goes through a placeholder. There is no escaping
function that is as safe, and no input validation that replaces it.
\end{worked}

\begin{examplebox}[The external schemas of Chapter 2, as a server would enforce them]
SQLite cannot run any of this. It is written out because it is what the
external level of \chref{architecture} \emph{means} once a database has users:

\begin{lstlisting}[language=sql]
CREATE ROLE fees_office;
CREATE ROLE exams_office;
CREATE ROLE portal;

GRANT SELECT ON fees_view TO fees_office;          -- names and rolls only
GRANT SELECT, UPDATE (grade) ON enrolment TO exams_office;
GRANT SELECT ON transcript TO portal;              -- a view, never a table

REVOKE ALL ON student FROM portal;                 -- least privilege
\end{lstlisting}

\smallskip
\term{Least privilege}: each user gets exactly what its job needs and no more.
Then even a successful SQL injection through the portal can only read what
\texttt{portal} may read --- the damage is bounded by the grant, not by the
programmer's care.

\smallskip
\textbf{What SQLite offers instead} is the operating system: whoever can read
the file can read all of it. File permissions, an encrypted disk, or an
encrypted SQLite build (SQLCipher) --- \term{encryption at rest} --- protect the file; nothing protects one
table from another.
\end{examplebox}

\section{NoSQL: What It Gives Up to Get What It Offers}

\begin{definitionbox}[Four families]
\term{NoSQL} names databases that do not use the relational model. The four
common families:

\smallskip
\begin{tabular}{@{}L{3.0cm}L{4.6cm}L{6.2cm}@{}}
\textbf{Family} & \textbf{Stores} & \textbf{Good at} \\
key--value & a value per key, opaque & caching, sessions: get and put by key \\
document & a nested document (JSON) per key & records read whole, with
  varying shape \\
wide-column & rows with many sparse columns & huge write volumes across many
  machines \\
graph & nodes and edges & paths: friends of friends, prerequisites of
  prerequisites \\
\end{tabular}

\smallskip
What they usually give up, to different degrees: \textbf{joins},
\textbf{a declared schema} (the data's shape is checked when read, not when
written --- \term{schema-on-read}), and \textbf{multi-record ACID
transactions}. What they gain is the ability to spread one database across
many machines (\term{horizontal scaling}). Spreading it costs
something too: the \term{CAP} result says that when the network between
machines fails, a system must choose between staying consistent and staying
available, and many choose \term{eventual consistency} --- copies may
disagree for a while, and agree once updates have spread.
\end{definitionbox}

\begin{worked}[One student's record as a document, inside SQLite]
SQLite has JSON functions, so a document store can be shown without
installing one.

\medskip
\textbf{Step 1 --- build the document.} One student, with their enrolments
nested inside:
\begin{lstlisting}[language=sql]
CREATE TABLE student_doc (roll_no TEXT PRIMARY KEY, doc TEXT NOT NULL);
INSERT INTO student_doc
SELECT s.roll_no,
       json_object('name', s.name, 'programme', s.programme,
         'enrolments', (SELECT json_group_array(
             json_object('course', e.course_code, 'term', e.semester,
                         'grade', e.grade))
            FROM enrolment e WHERE e.roll_no = s.roll_no))
  FROM student s;
\end{lstlisting}

\textbf{Step 2 --- what it is good at.} One student's whole record is one row,
read in one lookup, with no join:
\begin{lstlisting}[language=sql]
SELECT doc FROM student_doc WHERE roll_no = '2023-IT-00003';
\end{lstlisting}

\textbf{Step 3 --- what it is bad at.} ``How many A grades were given in
CS-101?'' was one \sqlkw{GROUP BY} in \chref{grouping}. Now every document must
be opened and unpacked:
\begin{lstlisting}[language=sql]
SELECT COUNT(*) FROM student_doc, json_each(doc, '$.enrolments') j
 WHERE json_extract(j.value, '$.course') = 'CS-101'
   AND json_extract(j.value, '$.grade')  = 'A';
\end{lstlisting}
On the seeded database both forms find \textbf{161}. The relational query took
9.8\,ms; the document form took \textbf{78.4\,ms}, eight times as long, because
it parsed every one of 4{,}000 documents to look inside each.

\textbf{Step 4 --- what it has given up.} Nothing stops a document with
\texttt{"grade": "Z"}, or \texttt{"grde"}, or a course that does not exist ---
\chref{constraints}'s three model-based constraints are gone, and a course's
title, if stored, is copied into every document that mentions it:
\chref{normalforms}'s update anomaly, by design.

\medskip
\textbf{The lesson.} A document store is the right model when records are read
and written \emph{whole} and rarely questioned \emph{across}. A university's
results are the opposite --- nearly every question in Part~IV crosses students
--- which is why this course is relational. That is a decision about the
\emph{questions}, not about which technology is newer.
\end{worked}

\dsfig{%
\begin{tikzpicture}[
  tb/.style={rounded corners=2pt, draw=#1, line width=0.9pt, fill=#1!7,
             text=#1!60!black, font=\scriptsize\ttfamily, align=left,
             inner sep=3pt},
  ttl/.style={font=\scriptsize\bfseries, text=#1, anchor=south},
  lk/.style={-{Stealth[length=2mm]}, neutral!60, line width=0.7pt},
  nt/.style={font=\scriptsize\itshape, text=#1, align=center}
]
  \node[ttl=primary] at (-4.0,2.2) {relational: the fact lives once};
  \node[tb=primary] (s) at (-5.3,1.3) {student\\2023-IT-00003\\Ayesha Q.};
  \node[tb=accent]  (e) at (-2.7,1.3) {enrolment\\00003 CS-101 A\\00003 MA-150 B\\\dots};
  \node[tb=greenhl] (c) at (-4.0,-0.35) {course\\CS-101 \ PH Course 1\\MA-150 \ \dots};
  \draw[lk] (e.west) -- (s.east);
  \draw[lk] (e.south) -- (c.north east);
  \node[nt=primary, text width=5.2cm] at (-4.0,-1.45)
       {``A grades in CS-101'': one index search on one table. 9.8\,ms};

  \node[ttl=purplehl] at (3.6,2.2) {document: the student lives once};
  \node[tb=purplehl, text width=5.4cm] (d) at (3.6,0.75)
       {\textbraceleft{} "name": "Ayesha Q.",\\
        \ \ "enrolments": [\\
        \ \ \ \ \textbraceleft"course":"CS-101","grade":"A"\textbraceright,\\
        \ \ \ \ \textbraceleft"course":"MA-150","grade":"B"\textbraceright,
        \dots ] \textbraceright};
  \node[nt=purplehl, text width=5.6cm] at (3.6,-0.85)
       {one lookup returns the whole record ---\\
        but ``A grades in CS-101'' opens all 4{,}000. 78.4\,ms};
  \node[nt=accent, text width=5.6cm] at (3.6,-1.6)
       {and nothing stops \texttt{"grde": "Z"}};
\end{tikzpicture}}%
{One student, held two ways. The relational form stores each fact once and
answers questions across students with a search; the document form stores each
student once and answers questions about one student in a single read. Which
is right depends on which kind of question the system is asked.}{docrel}

\begin{conceptbox}[When to leave SQLite --- using what this course measured]
\begin{tabular}{@{}L{4.4cm}L{4.4cm}L{5.0cm}@{}}
\textbf{Limit} & \textbf{Measured in} & \textbf{Move to a server when} \\
one writer at a time & \chref{transactions} & many users write at the same
  moment, all day \\
no users or privileges & \chref{whodoeswhat}, this chapter & different people
  must see different parts \\
on one machine & \chref{firstdatabase} & the data must be reached over a
  network by many programs \\
optimiser without STAT4 & \chref{indexes} & large, skewed tables where plan
  choice matters \\
\end{tabular}

\smallskip
And stay with it when none of those apply: SQLite is the most widely deployed
database in the world --- in every phone, every browser --- precisely because
for one program and one file it has nothing to administer. \chref{firstdatabase}
counted that as zero installation steps, zero ports and zero accounts; this
table is the other side of the same trade.
\end{conceptbox}

\section{The Semester Project}

\begin{definitionbox}[What you submit]
The project is the \texttt{university.db} you have built since
\chref{firstdatabase}, finished. Submit one folder containing:

\begin{tightnum}
  \item \texttt{design.pdf} --- your ER diagram with every cardinality and
        participation constraint (\chref{cardinality}), and the dependencies
        of every table (\chref{fds}).
  \item \texttt{v1\_create.sql}, \texttt{v2\_constraints.sql} and
        \texttt{v3\_final.sql} --- the schema and each migration
        (\chref{evolution}), each runnable from an empty file.
  \item \texttt{queries.sql} --- twelve queries answering the questions on
        the next page, each preceded by a comment stating the question.
  \item \texttt{triggers.sql} --- the six-section rule for \sqlkw{INSERT}
        \emph{and} \sqlkw{UPDATE} (\chref{views}), and an audit trigger.
  \item \texttt{report.pdf} --- six pages: three decisions you made and the
        evidence for each, including at least one measurement of your own.
  \item \texttt{university.db} itself, after \texttt{PRAGMA
        integrity\_check} and \texttt{PRAGMA foreign\_key\_check} both pass.
\end{tightnum}
\end{definitionbox}

\begin{examplebox}[The twelve questions]
\begin{tightnum}
  \item Every student's transcript, with course titles and credit hours.
  \item Each student's grade-point average over graded courses only.
  \item The top three students in each programme, ties included.
  \item Sections with no instructor --- the 73 of \chref{joins}.
  \item Instructors teaching nothing, written with \sqlkw{NOT EXISTS}.
  \item Students who have never received an A.
  \item Every prerequisite of a course, to any depth.
  \item Students enrolled in a course without having passed its direct
        prerequisites.
  \item Enrolments and average credit hours per department.
  \item Each course's pass rate, ungraded enrolments excluded and counted.
  \item A running total of credit hours per student, in correct term order.
  \item One question of your own that needs a window function.
\end{tightnum}

\smallskip
For each one, \texttt{queries.sql} must also contain its
\sqlkw{EXPLAIN QUERY PLAN}, and your report must name any query whose plan you
changed with an index, with both timings.
\end{examplebox}

\begin{conceptbox}[How it is marked]
\begin{tabular}{@{}L{6.0cm}R{1.6cm}L{6.2cm}@{}}
\textbf{Component} & \textbf{\%} & \textbf{What earns the marks} \\
design: ER diagram and dependencies & 20 & every constraint argued, not just
  drawn \\
schema and migrations & 20 & runs from empty; every constraint enforced,
  checked \\
the twelve queries & 30 & correct with nulls; \sqlkw{NOT IN} on a nullable
  column scores 0 \\
triggers & 10 & every event that can break the rule is covered \\
report and measurement & 20 & claims backed by numbers you produced \\
\end{tabular}

\smallskip
A query that returns the right answer on the seeded data and the wrong one when
a null appears is wrong. Test every one of them with a null.
\end{conceptbox}

\begin{pitfall}
\begin{tight}
  \item \textbf{Building SQL from input.} Parameters, always.
  \item \textbf{Escaping instead of parameterising.} Escaping is a list of
        characters somebody remembered; a parameter is never parsed.
  \item \textbf{Treating a view as access control in SQLite.} No users, so no
        restriction.
  \item \textbf{Choosing NoSQL because it is newer.} Choose by the questions:
        across records, relational; whole records, document.
  \item \textbf{Submitting a schema that only runs on your machine.} Run every
        script from an empty file before handing in.
  \item \textbf{Testing queries only on data without nulls.} The FALL2026
        grades are null for a reason.
\end{tight}
\end{pitfall}

\begin{lab}[Inject your own database, then make it impossible]
Read Appendix~A, section A.6, first. Every step uses \texttt{university.db}
on your own machine.

\begin{lstlisting}[language=bash]
# --- 1. The seeded database, which you can delete -------------------
python code/schema/seed.py university.db
\end{lstlisting}
\begin{lstlisting}[language=python]
# --- 2. The vulnerable query, in a file called inject.py ------------
import sqlite3
cx = sqlite3.connect("university.db")
for roll in ["2023-IT-00003", "x' OR '1'='1"]:
    sql = f"SELECT COUNT(*) FROM student WHERE roll_no = '{roll}'"
    print(repr(roll), "->", cx.execute(sql).fetchone()[0], "rows")

# --- 3. The same two inputs, parameterised --------------------------
for roll in ["2023-IT-00003", "x' OR '1'='1"]:
    n = cx.execute("SELECT COUNT(*) FROM student WHERE roll_no = ?",
                   (roll,)).fetchone()[0]
    print(repr(roll), "->", n, "rows (parameterised)")
\end{lstlisting}
\begin{lstlisting}[language=sql]
-- --- 4. A document, and a question it is bad at -----------------
.open university.db
CREATE TABLE student_doc (roll_no TEXT PRIMARY KEY, doc TEXT NOT NULL);
-- Type in Step 1 of the worked example, then:
SELECT json_extract(doc, '$.name'),
       json_array_length(doc, '$.enrolments')
  FROM student_doc LIMIT 3;

-- --- 5. Compare one cross-record question both ways --------------
.timer on
SELECT COUNT(*) FROM enrolment
 WHERE course_code = 'CS-101' AND grade = 'A';
SELECT COUNT(*) FROM student_doc, json_each(doc, '$.enrolments') j
 WHERE json_extract(j.value, '$.course') = 'CS-101'
   AND json_extract(j.value, '$.grade') = 'A';
-- Same count? Record both timings.

-- --- 6. Put in a document no relational table would accept --------
INSERT INTO student_doc VALUES ('X-1',
  '{"name":"Nobody","enrolments":[{"course":"NONE","grde":"Z"}]}');
-- Accepted. List the constraints from Chapter 8 it breaks.
\end{lstlisting}
\begin{lstlisting}[language=bash]
# --- 7. Run your project's checks before submitting ------------------
sqlite3 university.db "PRAGMA integrity_check; PRAGMA foreign_key_check;"
\end{lstlisting}

\smallskip
\textbf{What to notice.} Step 2's second line and step 3's second line are the
same input, giving 4{,}000 rows and 0. Nothing about the input changed; only
whether it was ever parsed as SQL. And step 6 is the other half of the
chapter: a model that accepts anything checks nothing.
\end{lab}

\begin{checkpoint}
\begin{tightnum}
  \item Explain exactly why \texttt{x' OR '1'='1} returned every student when
        pasted into the query and none when bound as a parameter, naming the
        character that made the difference and what happened to it in each
        case.
  \item Write the \sqlkw{GRANT} statements giving the examinations office
        permission to read every enrolment and change only grades, and say
        what SQLite would let that office do with the same file.
  \item For the question ``all of one student's grades'' and the question
        ``the number of A grades in CS-101'', say which model --- relational or
        document --- answers each more naturally, and why.
\end{tightnum}
\end{checkpoint}

\begin{chaptersummary}
\begin{tight}
  \item \textbf{SQL injection}: text pasted into SQL is parsed as SQL.
        \texttt{x' OR '1'='1} returned \textbf{4{,}000 students} pasted and
        \textbf{0} as a parameter.
  \item A \textbf{parameterised query} sends values separately; they are never
        parsed. There is no escaping function as safe.
  \item A server enforces external schemas with roles, \sqlkw{GRANT} and
        \sqlkw{REVOKE}, and \textbf{least privilege} bounds the damage of any
        flaw. SQLite has no users: the file's permissions are the only
        control.
  \item \textbf{NoSQL} families --- key--value, document, wide-column, graph
        --- usually give up joins, a declared schema and multi-record ACID to
        gain horizontal scaling.
  \item A document is read whole in one lookup and questioned across records
        only by unpacking every one; it accepts any shape, so none of
        \chref{constraints}'s constraints hold.
  \item Choose the model by the questions. A university's questions cross
        records, which is why this course is relational.
  \item Leave SQLite when the measured limits bite: many simultaneous writers,
        per-user access, network access, large skewed tables. Otherwise stay.
  \item The project is the \texttt{university.db} built all term, with its
        design, migrations, twelve queries, triggers and a report backed by
        your own measurements.
\end{tight}
\end{chaptersummary}

\begin{reviewq}
  \item \textbf{(MCQ)} A parameterised query prevents SQL injection because
        the value is: (a)~escaped; (b)~validated; (c)~never parsed as SQL;
        (d)~encrypted.
  \item \textbf{(MCQ)} Which NoSQL family suits ``friends of friends''?
        (a)~key--value; (b)~document; (c)~wide-column; (d)~graph.
  \item \textbf{(Short)} State the principle of least privilege and say how it
        limits the damage of an injection that does succeed.
  \item \textbf{(Short)} What does ``schema-on-read'' mean, and what does a
        document store therefore not check on write?
  \item \textbf{(Short)} Give two limits measured in this course that would
        justify moving the university's database to a server.
  \item \textbf{(Applied)} Rewrite a program of your own from the term that
        builds any SQL with string formatting so that every value is a
        parameter. Show one input that changed the old version's behaviour
        and does not change the new one's.
  \item \textbf{(Applied)} A startup proposes storing the university's
        records as one JSON document per student ``because it is simpler''.
        Using this chapter's worked example and the twelve project questions,
        say which questions become harder and why, what integrity is lost, and
        under what changed requirements the startup would be right.
\end{reviewq}


%---------------------------------------------------------------------
\appendix
% \part* under titlesec already emits its own contents line, so do not add
% one explicitly here or the entry appears twice.
\part*{Appendices}
%=====================================================================
\chapter{The Toolkit: Setting It Up and Keeping It}\label{app:toolkit}
%=====================================================================
\normalfigures

This appendix is the whole of the course's setup. It is short, and that is the
point: the longest-running complaint about database courses is that the first
fortnight goes on installing something. Here there is nothing to install.

\section{What You Need}

\rowcolors{2}{white}{lightbg}
\begin{longtable}{@{}L{3.4cm}L{4.2cm}L{6.8cm}@{}}
\toprule
\rowcolor{primary!15}
\textbf{Thing} & \textbf{Version used here} & \textbf{How to get it} \\
\midrule
\endhead
\textbf{SQLite} & 3.51.2 & \textbf{You already have it.} It ships inside
Python, inside most operating systems, and inside a great deal of software you
already run \\
A terminal & any & Already present \\
Python & 3.13 & Only for the measurement harnesses in \texttt{code/}. No
chapter's \emph{content} needs it \\
A text editor & any & For writing \texttt{.sql} files \\
\bottomrule
\end{longtable}

\begin{conceptbox}[Check, in one line, before you read any further]
\begin{lstlisting}[language=bash]
python -c "import sqlite3; print(sqlite3.sqlite_version)"
\end{lstlisting}

If that prints a version number --- any version from 3.37 onwards --- you have
everything this course needs. If you also have the \texttt{sqlite3} command on
your path, so much the better; if not, every lab works through Python's
built-in shell instead, and \chref{firstdatabase} shows both.
\end{conceptbox}

\section{Why SQLite, and What It Costs You}

\begin{alertbox}[This is a curriculum decision, not a shortcut]
\textbf{HEC names no database product in any edition of this curriculum.} The
2017 and 2023 outlines say a student should use SQL ``in any DBMS''; the 2025
outcomes say ``popular database management systems''. There is no edition in
which choosing SQLite departs from the published syllabus.

\smallskip
What it buys is the first fortnight. There is no server to start, no port, no
password, no user account, no configuration file and nothing to go wrong
before the subject begins. A database is one file. You can copy it, email it,
put it in version control, delete it and start again --- and in \chref{normalforms}
you will do exactly that.

\smallskip
\textbf{What it costs is stated here and again in \chref{project}}, because a
tool's limits are part of what you are learning:

\begin{tight}
  \item \textbf{No users, no \sqlkw{GRANT}, no \sqlkw{REVOKE}.} SQLite has no
        access control of any kind: whoever can open the file can do anything
        to it. \chref{project} writes the SQL a server DBMS would use and runs
        it nowhere.
  \item \textbf{No stored procedures.} Triggers (\chref{views}) carry the
        argument about putting logic in the database.
  \item \textbf{No row-level locking and no isolation levels.} SQLite lets one
        writer at a time and makes the others wait. \chref{transactions} can
        therefore show you a lost update prevented, a busy timeout and a failed
        upgrade from a deferred transaction --- but a genuine two-row deadlock
        is drawn rather than run.
\end{tight}

\smallskip
Every one of those three is a thing a \emph{server} adds, and knowing which
three they are is more useful than having had them switched on all semester
without noticing.
\end{alertbox}

\section{The Ten Commands You Will Actually Use}

SQLite's shell has two kinds of input: SQL, which ends with a semicolon, and
\emph{dot-commands}, which do not and which are the shell's own.

\rowcolors{2}{white}{lightbg}
\begin{longtable}{@{}L{4.6cm}L{10.0cm}@{}}
\toprule
\rowcolor{primary!15}
\textbf{Command} & \textbf{What it does} \\
\midrule
\endhead
\texttt{sqlite3 university.db} & Open (or create) a database and enter the
shell \\
\texttt{.tables} & List the tables \\
\texttt{.schema student} & Show the \texttt{CREATE} statement for one table ---
the fastest way to find out what a table actually promises \\
\texttt{.headers on} & Print column names above results. Turn this on first,
every time \\
\texttt{.mode box} & Draw results as a table instead of pipe-separated text \\
\texttt{.read setup.sql} & Run a file of SQL. How every lab in this book starts \\
\texttt{.output results.txt} & Send what follows to a file \\
\texttt{.import data.csv student} & Load a CSV into a table \\
\texttt{.backup backup.db} & Take a consistent copy \emph{while the database is
open} --- \chref{transactions} explains why that is not the same as copying the
file \\
\texttt{.quit} & Leave \\
\bottomrule
\end{longtable}

\begin{examplebox}[The two lines to type at the start of every session]
\begin{lstlisting}[language=bash]
sqlite3 university.db
\end{lstlisting}
then, inside the shell:
\begin{lstlisting}[language=sql]
.headers on
.mode box
PRAGMA foreign_keys = ON;
\end{lstlisting}

\smallskip
The first two make results readable. \textbf{The third is not cosmetic.}
SQLite does not enforce foreign keys unless you ask it to, \emph{every time you
connect}, and \chref{constraints} is largely about what happens when you
forget. Put all three in a file called \texttt{init.sql} and start every
session with \texttt{.read init.sql} until it is a habit.
\end{examplebox}

\section{Resetting to a Known State}

A database is one file, which makes this the simplest reset in the series:

\begin{lstlisting}[language=bash]
# --- 1. Keep a copy before anything risky -------------------------
cp university.db university.db.bak

# --- 2. Start again from nothing ----------------------------------
rm university.db
sqlite3 university.db ".read code/schema/v1_create.sql"

# --- 3. And refill it ----------------------------------------------
python code/schema/seed.py 1000
\end{lstlisting}

\begin{conceptbox}[Do this deliberately once, in Week 2]
Delete your database on purpose, rebuild it from the script, and confirm it is
the same. A student who has done that once is a student who will keep their
schema in a file rather than in the shell's history --- which is the single
habit that decides whether the Week~16 project can be handed to somebody else
and run.
\end{conceptbox}

\section{The Machine Every Timing Was Taken On}

Eleven chapters quote a measured figure. Every one was produced by running the
named file in \texttt{code/} on this machine:

\rowcolors{2}{white}{lightbg}
\begin{longtable}{@{}L{3.6cm}L{11.0cm}@{}}
\toprule
\rowcolor{primary!15}
\textbf{Component} & \textbf{Specification} \\
\midrule
\endhead
Processor & Intel Core i5-6600, 4 cores, 3.30\,GHz \\
Memory & 16\,GB \\
Operating system & Windows 11, 64-bit \\
SQLite & 3.51.2 \\
Python & 3.13.13 \\
Storage & Mechanical disk, not a solid-state drive --- which matters for
\chref{indexes} and is the reason its figures are larger than yours may be \\
\bottomrule
\end{longtable}

\begin{alertbox}[Your numbers will differ, and the ratios should not]
A measurement is a fact about the machine it ran on. If your laptop has a
solid-state disk it will beat these figures, perhaps by a factor of ten on the
chapters that touch storage.

\smallskip
\textbf{What should survive the move is the shape:} the row count at which an
index starts to pay (\chref{indexes}), the direction normalisation moves query
time (\chref{normalforms}), the fact that one writer blocks another
(\chref{transactions}). If a \emph{ratio} in this book does not reproduce on
your machine, that is worth investigating and worth telling your lecturer
about. If an absolute number does not, that is expected.
\end{alertbox}

\section{Scope: What You May Point Any of This At}

\begin{alertbox}[Read this before Chapter 23's laboratory]
Everything in this course runs against \textbf{a database file you created, on
a machine you own}. \chref{project} discusses SQL injection and shows what a
badly built query permits --- against your own university database, with
invented data, in a file you can delete.

\smallskip
\textbf{Running the same thing against any system you do not own and have not
been authorised in writing to test is a criminal offence in Pakistan} under the
\textbf{Prevention of Electronic Crimes Act 2016}, sections~3 and~4, which
cover unauthorised access to an information system and to data.

\smallskip
There is no educational exemption, and a university assignment is not
authorisation. ``I was learning'', ``I did no harm'' and ``the site was
insecure'' are not defences.

\smallskip
In practice this is an easy rule to follow, because the only file named
anywhere in this book is \texttt{university.db} and it is yours.
\end{alertbox}

%=====================================================================
\chapter{Just Enough Sets, Logic and Logarithms}\label{app:logic}
%=====================================================================
\normalfigures

This appendix is the mathematics the course uses, and nothing more. It is short
because the course needs little: sets for the relational model and the
algebra, logic for \sqlkw{WHERE} and for nulls, one idea about logarithms for
indexes, and the definition of a function for Part~V. Each section names the
chapter that depends on it. Questions 6--9, 11 and 12 of the Prerequisite
Self-Check are answered here.

\section{Sets}

\begin{definitionbox}[Set, element, subset]
A \term{set} is a collection of distinct things, its \term{elements}, with no
order. $\{1,2,3\}$ and $\{3,1,2\}$ are the same set; $\{1,1,2\}$ is just
$\{1,2\}$.

\smallskip
$x \in A$ means $x$ is an element of $A$. $A \subseteq B$ ($A$ is a
\term{subset} of $B$) means every element of $A$ is in $B$. The empty set is
$\varnothing$. $|A|$ is the number of elements.

\smallskip
\textbf{Why it matters:} a relation is a \emph{set} of tuples
(\chref{relational}). ``No duplicates'' and ``no order'' are not rules somebody
chose for databases; they are what the word ``set'' means.
\end{definitionbox}

\begin{examplebox}[The three operations --- Self-Check question 6]
$A = \{1,2,3\}$, $B = \{2,3,4\}$.

\medskip
\begin{tabular}{@{}L{3.6cm}L{3.0cm}L{4.8cm}L{2.4cm}@{}}
\textbf{Operation} & \textbf{Result} & \textbf{Contains} & \textbf{SQL} \\
union $A \cup B$ & $\{1,2,3,4\}$ & in $A$ or in $B$ (or both) & \sqlkw{UNION} \\
intersection $A \cap B$ & $\{2,3\}$ & in both & \sqlkw{INTERSECT} \\
difference $A - B$ & $\{1\}$ & in $A$ and not in $B$ & \sqlkw{EXCEPT} \\
\end{tabular}

\smallskip
Difference is not symmetric: $B - A = \{4\}$. And
$A \cap B = A - (A - B)$: take away from $A$ everything not in $B$, and what is
left is in both. \chref{algebra} uses exactly this to show intersection is not
a fundamental operator.
\end{examplebox}

\begin{definitionbox}[Ordered pair, Cartesian product --- Self-Check question 7]
An \term{ordered pair} $(a, b)$ has a first element and a second; unlike a
set, $(a,b) \neq (b,a)$ when $a \neq b$.

\smallskip
The \term{Cartesian product} $A \times B$ is the set of every pair $(a, b)$
with $a \in A$ and $b \in B$. It has $|A| \times |B|$ elements.

\smallskip
$\{1,2\} \times \{x, y, z\} = \{(1,x), (1,y), (1,z), (2,x), (2,y), (2,z)\}$ ---
six pairs. In \chref{algebra} the product of 4{,}000 students and 60{,}037
enrolments is 240{,}148{,}000 pairs, which is why no database ever computes a
join that way.
\end{definitionbox}

\section{Relations and Functions}

\begin{definitionbox}[Relation, in the mathematical sense]
A \term{relation} between sets $A$ and $B$ is any subset of $A \times B$ ---
a chosen collection of pairs. ``Is enrolled in'', between students and
sections, is the set of (student, section) pairs that are true.

\smallskip
With more sets the pairs become tuples: a relation over $A$, $B$ and $C$ is a
subset of $A \times B \times C$. That is precisely Codd's definition, and why a
database table is called a relation. The \emph{domains} of \chref{relational}
are the sets $A$, $B$, $C$.
\end{definitionbox}

\begin{definitionbox}[Function --- Self-Check question 12]
A \term{function} $f$ from $A$ to $B$ is a relation in which \textbf{every
element of $A$ is paired with exactly one element of $B$}. One input, one
output. $f(x) = x^2$ is a function: each $x$ has one square. ``The square
roots of $x$'' is not: 4 has two.

\smallskip
Many inputs may share an output --- $(-2)^2 = 2^2$ --- that is allowed. What
is not allowed is one input with two outputs.

\smallskip
\textbf{This is the whole foundation of \chref{fds}.} The functional
dependency \fd{roll\_no}{name} says ``name is a function of roll\_no'': a roll
number has exactly one name. Two students sharing a name is fine (many inputs,
one output). One roll number with two names is the violation.
\end{definitionbox}

\section{Logic}

\begin{definitionbox}[Propositions, and the three connectives]
A \term{proposition} is a statement that is true or false: ``$x > 5$''.

\medskip
\begin{tabular}{@{}cc|ccc@{}}
$P$ & $Q$ & $P$ \sqlkw{AND} $Q$ & $P$ \sqlkw{OR} $Q$ & \sqlkw{NOT} $P$ \\
\midrule
T & T & T & T & F \\
T & F & F & T & F \\
F & T & F & T & T \\
F & F & F & F & T \\
\end{tabular}

\smallskip
\sqlkw{AND} is true only when both are; \sqlkw{OR} when at least one is.
\sqlkw{OR} is \emph{inclusive}: true when both are.
\end{definitionbox}

\begin{examplebox}[Two Self-Check questions, answered]
\textbf{Question 8 --- is ``$x > 5$ and $x < 3$'' ever true?} No. No number is
both greater than 5 and less than 3, so the statement is a
\term{contradiction}: false for every $x$. A \sqlkw{WHERE} clause that is a
contradiction returns no rows, and no error --- a common bug when a condition
was meant to be \sqlkw{OR}.

\medskip
\textbf{Question 9 --- ``not ($P$ and $Q$)'' with $P$ true, $Q$ false.} $P$
\sqlkw{AND} $Q$ is false, so its negation is \textbf{true}. \term{De Morgan's
laws} make this mechanical:
\[
  \neg(P \wedge Q) \equiv \neg P \vee \neg Q \qquad
  \neg(P \vee Q) \equiv \neg P \wedge \neg Q
\]
``Not both'' means ``at least one is not''. \chref{select} uses the second law
to explain \sqlkw{NOT IN}: \texttt{x NOT IN (a, b)} is \texttt{NOT (x = a OR x
= b)}, which is \texttt{x <> a AND x <> b}.
\end{examplebox}

\begin{conceptbox}[A third truth value --- what SQL adds]
Two-valued logic has no room for ``I do not know''. SQL needs one, because a
comparison with a null cannot be true or false. It adds \unk{}, and the tables
extend (\chref{select} has them in full). The two facts that matter:

\begin{tight}
  \item \texttt{F AND} anything is F, and \texttt{T OR} anything is T --- a
        missing value cannot change those.
  \item \texttt{NOT} \unk{} is \unk{}. So De Morgan's laws still hold, but a
        condition and its negation no longer cover every row --- 19{,}992
        enrolments fall through the gap on the seeded data.
\end{tight}
\end{conceptbox}

\begin{definitionbox}[Quantifiers, and their SQL]
$\exists x\, P(x)$, ``there exists'', is true if $P$ holds for \emph{at least
one} $x$. $\forall x\, P(x)$, ``for all'', is true if $P$ holds for
\emph{every} $x$.

\smallskip
They are linked by negation: $\neg \forall x\, P(x) \equiv \exists x\, \neg
P(x)$ --- ``not every student passed'' means ``some student did not''.

\smallskip
SQL has $\exists$ directly, as \sqlkw{EXISTS} (\chref{subqueries}). It has no
\sqlkw{FOR ALL}; ``every'' is written as ``there is no counter-example'':
\sqlkw{NOT EXISTS} ( \dots\ \sqlkw{NOT} \dots\ ). ``Students who passed every
course they took'' becomes ``students for whom there is no course they took
and failed''.
\end{definitionbox}

\section{Logarithms: How Many Halvings}

\begin{definitionbox}[The logarithm as a count --- Self-Check question 11]
How many times must 1{,}000{,}000 be halved to reach 1? Halving 20 times gives
$1{,}000{,}000 / 2^{20} = 1{,}000{,}000 / 1{,}048{,}576 < 1$, and 19 times is
not quite enough. So about \textbf{20}.

\smallskip
That count is the \term{logarithm}: $\log_2 1{,}000{,}000 \approx 20$. In
general $\log_b n$ is how many times you divide $n$ by $b$ to reach 1 --- or,
the other way round, the power of $b$ that gives $n$.

\smallskip
The property that matters: \textbf{multiplying $n$ adds only a little to its
logarithm.} $\log_2$ of a thousand is about 10; of a million, 20; of a
billion, 30. A thousand times more data costs ten more halvings.
\end{definitionbox}

\begin{examplebox}[Why it is the whole story of Chapter 21]
A B-tree page divides its range into a few hundred parts --- divides by, say,
$b = 300$ rather than by 2. Its depth is $\log_{300} n$:

\medskip
\begin{tabular}{@{}R{3.6cm}R{3.4cm}R{3.4cm}@{}}
\textbf{rows $n$} & $\log_2 n$ & $\log_{300} n$ \\
1{,}000 & 10 & 1.2 \\
1{,}000{,}000 & 20 & 2.4 \\
1{,}000{,}000{,}000 & 30 & 3.6 \\
\end{tabular}

\smallskip
Rounding up, a million rows need a tree three levels deep, and a billion only
four. That is why \chref{indexes} measured a lookup at 0.053\,ms for 10{,}000
rows and 0.057\,ms for a million: a hundredfold more data added at most one
level. A scan, which reads everything, has no logarithm in it --- a hundred
times the rows is a hundred times the work.
\end{examplebox}

% Not a checkpoint environment: every checkpoint in this handout is answered
% in Appendix C, which is per chapter. These are answered here instead.
\begin{conceptbox}[Three to try, with answers]
\begin{tightnum}
  \item With $A = \{a, b, c, d\}$ and $B = \{c, d, e\}$, write $A \cup B$,
        $A \cap B$, $A - B$ and $B - A$, and give $|A \times B|$.
  \item Is ``each student has exactly one date of birth'' a function from
        students to dates? Is ``each date has exactly one student born on
        it''? Say which, if either, could be a functional dependency.
  \item Write ``every section has at least one student'' using only
        \sqlkw{NOT EXISTS}, by first rewriting it as ``there is no section
        such that \dots''.
\end{tightnum}

\smallskip
\textbf{Answers.} (1)~$\{a,b,c,d,e\}$; $\{c,d\}$; $\{a,b\}$; $\{e\}$;
$4 \times 3 = 12$. (2)~The first is a function, and is
\fd{roll\_no}{date\_of\_birth}. The second is not --- one date has many
students --- so \fd{date\_of\_birth}{roll\_no} does not hold. (3)~``There is no
section with no enrolment'':
\begin{lstlisting}[language=sql]
NOT EXISTS (
  SELECT 1 FROM section s
   WHERE NOT EXISTS (
     SELECT 1 FROM enrolment e
      WHERE e.course_code = s.course_code
        AND e.semester    = s.semester
        AND e.section_no  = s.section_no))
\end{lstlisting}

\end{conceptbox}

%=====================================================================
\chapter{Answers to the Checkpoints}\label{app:answers}
%=====================================================================
\normalfigures

Three answers per chapter, in order. Attempt each checkpoint from memory before
reading these. Several are judgement calls rather than facts, and say so; where
an answer quotes a number, the number comes from the chapter's harness or from
the seeded database, and you can reproduce it.

% The slide pipeline parses this appendix: \ansch{<number> --- <title>} must be
% followed immediately by an ans environment of exactly three items.
\newcommand{\ansch}[1]{\subsection*{\color{primary}Chapter #1}}
\newenvironment{ans}
  {\begin{enumerate}[leftmargin=*, itemsep=4pt, topsep=3pt,
                     label=\textbf{\arabic*.}]}
  {\end{enumerate}}

%---------------------------------------------------------------------
\ansch{1 --- Why Not Just Files?}
\begin{ans}
  \item \textbf{It escapes most of them, for now.} With one clerk per file and
        no shared data there is no redundancy, no inconsistency between copies
        and no concurrent access. It still has \emph{isolation} (``which
        borrowers are also in arrears?'' needs both files), \emph{integrity}
        (rules live in each clerk's head) and \emph{atomicity} (a half-saved
        file). The answer changes the moment any fact --- a student's name ---
        is recorded in both files, or a second clerk edits either one.
  \item A \textbf{database} is the data together with its description; a
        \textbf{DBMS} is the software that manages it. ``We bought a database''
        might mean a product (SQLite, Oracle) with no data in it, or a
        populated set of tables, or a whole system --- three purchases with
        different costs, owners and risks. You cannot install, back up or
        migrate it until you know which.
  \item Disk is cheap; \textbf{disagreement} is not. Redundancy's real cost is
        \emph{inconsistency}: two copies of one fact that can diverge. Example:
        a student's name changes and is corrected in \texttt{students.csv} but
        not \texttt{fees.csv}. Each file is internally correct; the system now
        holds two names for one person, and the fee receipt goes to someone
        who no longer exists.
\end{ans}

\ansch{2 --- Three Levels, and Why They Are Separate}
\begin{ans}
  \item Adding an index: \textbf{internal}, physical independence. Adding a
        column: \textbf{conceptual}, logical independence (programs that never
        named it carry on). Renaming a column: \textbf{conceptual}, and logical
        independence \emph{fails} unless a view keeps the old name. A faster
        disk: \textbf{internal}, physical independence. A view with no home
        addresses: \textbf{external}; no independence is involved --- it is
        the external level being used.
  \item A schema is the \emph{form}; an instance is what is \emph{written on
        it}. A blank admission form is the schema: it says there will be a
        name, a date of birth, a programme. This year's 4{,}000 filled-in
        forms are the instance. Reprinting the form with a new box is a schema
        change and is rare; filling one in is an instance change and happens
        every minute.
  \item The \textbf{external level} is affected: views can define it, but
        nothing can confine a user to one, because SQLite has no users. You can
        still write the views and have programs use only them. You cannot stop
        anyone who can open the file from reading the base tables. In an
        application the program itself becomes the gatekeeper --- it holds the
        file, authenticates people, and only ever queries the views --- which
        is enforcement by convention, not by the database.
\end{ans}

\ansch{3 --- What Is Inside a DBMS, and Who Does What}
\begin{ans}
  \item The \textbf{optimiser} (a different plan --- perhaps statistics were
        refreshed, or an index dropped) and the \textbf{buffer manager} (Monday
        the pages were in memory; Friday they were not). Tell them apart with
        \sqlkw{EXPLAIN QUERY PLAN}: if the plan changed, it is the optimiser;
        if the plan is identical, run the query twice --- a fast second run
        means the first was paying for disk reads.
  \item \textbf{Advantage:} any tool can discover a database it has never
        seen with an ordinary \sqlkw{SELECT}, and the schema is backed up with
        the data. \textbf{Problem:} the dictionary can be damaged by the same
        failures as the data, and in SQLite anyone who can write the file can
        in principle corrupt \texttt{sqlite\_master} and with it every table.
  \item DBA: when to back up. Designer: what the keys are. Programmer: how a
        query is parameterised. End user: what question is asked. In this
        course the designer, programmer and DBA decisions are all yours; the
        DBA decisions that need users and privileges simply cannot be made,
        and backup becomes \texttt{.backup} run by you.
\end{ans}

\ansch{4 --- Your First Database, in One Command}
\begin{ans}
  \item Page 1 (4{,}096 bytes) holds the 100-byte database header and the
        schema table \texttt{sqlite\_master} --- the data dictionary. Page 2
        (4{,}096 bytes) is the root page of the new table's B-tree, empty.
        Fewer is impossible: storage is allocated in whole pages, the
        dictionary must exist for the file to describe itself, and the table
        must have a root.
  \item From 1{,}000 to 10{,}000 rows the file grew $262{,}144 - 32{,}768 =
        229{,}376$ bytes, about 25.5 bytes a row, so 2{,}000 rows predicts
        $32{,}768 + 25{,}500 \approx 58{,}300$, rounded up to a whole page:
        61{,}440 (15 pages). \textbf{Measured: 57{,}344 bytes, 14 pages.} The
        naive guess --- double 32{,}768 --- gives 65{,}536 and is further off,
        because part of the file is fixed overhead (the header page and the
        tree's interior page) that does not grow with the rows.
  \item
\begin{lstlisting}[language=sql]
CREATE TABLE instructor (
  emp_id      TEXT PRIMARY KEY NOT NULL,
  name        TEXT NOT NULL,
  department  TEXT NOT NULL,
  designation TEXT
) STRICT;
\end{lstlisting}
        \texttt{emp\_id} is \textbf{text} because \texttt{E-0114} has
        a letter and a leading zero and is never used in arithmetic; it is
        \textbf{not null} because an identifier that is missing identifies
        nothing. \texttt{designation} may be null: a new appointment may not
        have one yet.
\end{ans}

\ansch{5 --- The Entity-Relationship Model}
\begin{ans}
  \item ``Almost every'' means some have two, and one column cannot hold two
        values that can be searched. The first student with two numbers forces
        either a lost number or a comma-separated list, which no query can
        search correctly. Instead, a \texttt{student\_phone(roll\_no, phone)}
        table with both columns in the key; it costs one join whenever a
        number is shown.
  \item Candidate keys: $\{$emp\_id$\}$ (issued uniquely by the university),
        $\{$cnic$\}$ (national identity, unique by law), $\{$email$\}$ (issued
        one per employee). \textbf{Primary key: emp\_id} --- unique, always
        present (CNIC may be missing for foreign staff), never changing (email
        changes with a name or a reorganisation), short. \texttt{office\_no},
        \texttt{name} and \texttt{department} are not keys: offices are
        shared, names repeat, departments hold many people.
  \item A grade is a fact about \emph{a student in a particular section} ---
        neither alone determines it. On \ent{student} there is room for one
        grade per student for their whole degree. On \ent{section} there is
        room for one grade per section, the same for every student in it. Only
        the pairing of the two has a grade, and the pairing is the
        relationship.
\end{ans}

\ansch{6 --- Cardinality, Participation and Weak Entities}
\begin{ans}
  \item \textbf{It will refuse a section with no instructor.} A section is
        created in June for the autumn timetable; the lecturer is appointed in
        August. For two months the database will not accept the section ---
        and somebody will invent an instructor called ``TBA''.
  \item \textbf{M:N}, with \textbf{partial} participation of \ent{student}
        (a student may be enrolled in nothing) and \textbf{partial}
        participation of \ent{section} (a section may have no students). With
        $(1,N)$ at the section end, participation of \ent{section} becomes
        \textbf{total}: a section with no students could not exist, so a new
        section could not be created until its first enrolment arrived.
  \item Section numbers repeat across courses --- every course has a section
        1 --- so \texttt{section\_no} identifies a section only within its
        owner. Partial key: \texttt{section\_no}. Full key: $\{$course\_code,
        semester, section\_no$\}$. With a surrogate \texttt{section\_id} you
        must still declare \sqlkw{UNIQUE} (course\_code, semester,
        section\_no); forget it and the same section can be created twice,
        with students split between the copies.
\end{ans}

\ansch{7 --- The Relational Model}
\begin{ans}
  \item \textbf{Degree 5} counts attributes and belongs to the schema;
        \textbf{cardinality 200} counts tuples and belongs to the instance.
        Cardinality changes without a schema change, with every
        \sqlkw{INSERT} or \sqlkw{DELETE}. Degree changes only with a schema
        change: \sqlkw{ALTER TABLE} \dots\ \sqlkw{ADD} or \sqlkw{DROP COLUMN}.
  \item \texttt{SELECT COUNT(*) FROM enrolment WHERE roll\_no = ?} --- if the
        same enrolment can be inserted twice, a student's course count is
        inflated. Declare the composite \sqlkw{PRIMARY KEY} (roll\_no,
        course\_code, semester, section\_no), which makes the duplicate
        impossible.
  \item One null marker carries both meanings, and SQLite has no second kind.
        \textbf{Design 1:} a \texttt{grade\_status} column (\texttt{'not sat'},
        \texttt{'lost'}, \texttt{'graded'}) beside a nullable grade --- one
        column more, and a rule that the two agree. \textbf{Design 2:} a
        separate \texttt{exam\_incident} table recording lost papers --- no
        nulls overloaded, but a join to tell the two cases apart.
\end{ans}

\ansch{8 --- Keys and Integrity Constraints}
\begin{ans}
  \item \texttt{PRAGMA foreign\_keys;} tells you whether \emph{this
        connection} enforces --- not whether anyone else's does, or ever did.
        \texttt{PRAGMA foreign\_key\_list(t);} shows the declarations exist.
        \texttt{PRAGMA foreign\_key\_check;} lists every orphan already
        present: rows means enforcement was off at some point; none means no
        violation survives, \emph{not} that enforcement was on. Nothing in the
        file records whether it ever was.
  \item \ent{section}$\rightarrow$\ent{course}: \sqlkw{CASCADE} --- a
        section is a weak entity and means nothing without its course.
        \ent{section}$\rightarrow$\ent{instructor}: \sqlkw{SET NULL} --- the
        section survives a resignation; participation is partial.
        \ent{enrolment}$\rightarrow$\ent{student}: \sqlkw{RESTRICT} --- a
        student with a grade history must never be silently deleted.
  \item It will accept a \textbf{null} primary key and store it, violating
        \textbf{entity integrity}. Fix it with \sqlkw{NOT NULL} on the column,
        or \sqlkw{STRICT} on the table --- measured in \chref{constraints},
        each one alone refuses the null.
\end{ans}

\ansch{9 --- Mapping an ER Diagram to Tables}
\begin{ans}
  \item On \ent{department}, as \texttt{head\_emp\_id}: every department must
        have a head (total), but most instructors head nothing (partial), so
        the totally-participating side can declare the column \sqlkw{NOT NULL}.
        On \ent{instructor} instead it would be \texttt{heads\_dept\_code},
        null for all but eight of a hundred instructors.
  \item \texttt{roll\_no}: who. \texttt{course\_code}: which course.
        \texttt{semester}: when. \texttt{section\_no}: which group within that
        offering. Without \texttt{semester} a student could not appear in the
        same course and section number in two terms --- so repeating a failed
        course becomes impossible, or overwrites the old grade.
  \item The table \textit{prerequisite}, whose primary key is the pair and
        both of whose columns reference \textit{course}; the unstated rule is
        the \sqlkw{CHECK}:
\begin{lstlisting}[language=sql]
CREATE TABLE prerequisite (
  course_code TEXT NOT NULL REFERENCES course(course_code),
  requires    TEXT NOT NULL REFERENCES course(course_code),
  PRIMARY KEY (course_code, requires),
  CHECK (course_code <> requires)    -- nothing requires itself
) STRICT;
\end{lstlisting}
 (A two-step cycle needs a trigger;
        \chref{subqueries} shows why it also needs \sqlkw{UNION}.)
\end{ans}

\ansch{10 --- SQL: Creating the Schema}
\begin{ans}
  \item The marks were 95, 72, \texttt{'absent'}, 88 and 100. Affinity left
        \texttt{'absent'} as text; \sqlkw{AVG} read it as 0, so the sum was
        $95 + 72 + 0 + 88 + 100 = 355$ and, because it is not null, the count
        was 5: $355 / 5 = 71$. Proof: \texttt{SELECT typeof(score), COUNT(*)
        FROM marks GROUP BY 1;} shows the text row, and \texttt{SELECT
        SUM(score), COUNT(score) FROM marks;} returns 355 and 5.
  \item A column constraint follows one column and can name only it; a table
        constraint follows all the columns and can name several. In
        \texttt{v1\_create.sql}: the composite \sqlkw{PRIMARY KEY} of
        \textit{section} (three columns), and the three-column
        \sqlkw{FOREIGN KEY} of \textit{enrolment} --- neither belongs to one
        column, so neither can be written beside one.
  \item \textbf{Compatibility}: other tools and older SQLite versions (before
        3.37) cannot read a \sqlkw{STRICT} table, and it forbids type names
        like \texttt{VARCHAR} or \texttt{DATETIME} that ported SQL uses. That
        is a real cost when the file is shared with old software; for this
        course's database it does not apply, and the measured cost was zero.
\end{ans}

\ansch{11 --- Changing a Schema That Is Already Full}
\begin{ans}
  \item Nullable \texttt{email}: \textbf{cheap} --- no existing row changes.
        \texttt{email} \sqlkw{NOT NULL}: \textbf{expensive} --- every row needs
        a value, so backfill and rebuild. \texttt{semester} text to integer:
        \textbf{expensive} --- a rebuild, and step~0 must find any value that
        will not convert. Splitting a whole address: \textbf{impossible} to do
        exactly --- the parts were never recorded separately, and parsing free
        text guesses.
  \item After the migration, backfilled values are indistinguishable from
        values entered since; a rollback that drops the column cannot tell
        which were invented, and one that keeps them keeps invented data. To
        make it reversible, record at migration time which rows were
        backfilled --- a \texttt{migration\_backfill(roll\_no)} table, or a
        flag column --- so the reverse can remove exactly those.
  \item If the migration fails after \sqlkw{COMMIT} but the version is set
        afterwards, that is merely late. The real danger is the reverse: set
        after a migration that failed and was rolled back, and the file claims
        to be version 2 while holding the version 1 schema. Every later
        migration then assumes columns that are not there.
\end{ans}

\ansch{12 --- Relational Algebra}
\begin{ans}
  \item A relation is a \emph{set}, so $\pi$ removes duplicates: four
        programmes, four tuples. A SQL result is a \emph{bag} and keeps every
        row. \sqlkw{DISTINCT} reconciles them:
        \texttt{SELECT DISTINCT programme FROM student} returns 4.
  \item $\pi_{roll}(\sigma_{semester=2}(\mathit{student})) -
        \pi_{roll}(\sigma_{semester='FALL2026'}(\mathit{enrolment}))$, that is:
\begin{lstlisting}[language=sql]
SELECT roll_no FROM student   WHERE semester = 2
EXCEPT
SELECT roll_no FROM enrolment WHERE semester = 'FALL2026';
\end{lstlisting}
        ``Nothing'' is the difference. The tempting wrong version, a join with
        \texttt{e.semester <> 'FALL2026'}, returns second-year students enrolled
        in \emph{some other} term --- including those also enrolled in
        FALL2026. Note the two meanings of \texttt{semester}: \chref{joins}
        returns to them.
  \item $4{,}000 \times 60{,}037 = 240{,}148{,}000$ intermediate tuples, of
        which 60{,}037 survive the selection. \textbf{Closure} --- every
        intermediate result is a relation --- is what makes algebraic laws such
        as ``select before you multiply'' valid, so the optimiser may compute
        the join by index lookups and still return exactly the same relation.
\end{ans}

\ansch{13 --- Joins: the Algebra and the SQL}
\begin{ans}
  \item The two tables share \textbf{two} column names: \texttt{roll\_no} and
        \texttt{semester}. \sqlkw{NATURAL JOIN} required both to be equal, and
        a year of study (\texttt{8}) never equals a term label
        (\texttt{'FALL2025'}). In the query: \texttt{JOIN \dots\ USING
        (roll\_no)} or \sqlkw{ON}. In the schema: rename
        \texttt{student.semester} to \texttt{year\_of\_study} --- it is the one
        whose name is misleading.
  \item They are the 73 sections whose \texttt{emp\_id} is null.
        \texttt{SELECT s.* FROM section s LEFT JOIN instructor i ON
        s.emp\_id = i.emp\_id WHERE i.emp\_id IS NULL;} The decision was making
        \ent{section}'s participation in \ent{teaches} \textbf{partial} in
        \chref{cardinality} --- deliberately, so a section can exist before its
        lecturer is appointed.
  \item Without an index, every pair of 4{,}000 students was compared ---
        about 8 million --- to find 28{,}517 matches; with one, only the
        matches were visited: 108$\times$ faster. Growth ratios mislead because
        the number of \emph{matches} itself grew four-fold per doubling (names
        come from a fixed pool), so even the indexed join looked quadratic down
        its column. Only a comparison at equal size isolates search cost.
\end{ans}

\ansch{14 --- Selecting Rows, and the Trouble with NULL}
\begin{ans}
  \item $60{,}037 - 7{,}268 - 32{,}777 = 19{,}992$ rows: the ungraded FALL2026
        enrolments, whose grade is null. For them both \texttt{= 'A'} and
        \texttt{<> 'A'} are \unk{}, and \sqlkw{WHERE} keeps only \emph{true}.
        \texttt{WHERE grade <> 'A' OR grade IS NULL} returns all 52{,}769 that
        are not an A (in SQLite \texttt{WHERE grade IS NOT 'A'} says the same).
  \item \texttt{x NOT IN (a, NULL)} means \texttt{x <> a AND x <> NULL}. The
        second comparison is \unk{}. If the first is false the whole is false;
        if the first is true, \texttt{T AND U = U}. So for every $x$ the
        condition is false or \unk{} --- never true --- and \sqlkw{WHERE}
        discards every row.
  \item It decides that \textbf{a missing mark is a zero} --- an absent
        student failed. Scores 80, 90, null: \sqlkw{AVG} gives
        $(80+90)/2 = 85$; \sqlkw{AVG}(\sqlkw{COALESCE}) gives $(80+90+0)/3 =
        56.7$. Write \sqlkw{AVG}(score), and report the count of missing marks
        beside it, so the decision about them is made by the examinations
        office rather than by a function.
\end{ans}

% Plain text: the slide parser reads this heading with [^}]*, so a
% brace inside the title would end it early and lose all three answers.
\ansch{15 --- Sub-queries, EXISTS and Set Operations}
\begin{ans}
  \item The \sqlkw{IN} sub-query of step~2 is \textbf{uncorrelated}. Its
        \sqlkw{EXISTS} is \textbf{correlated}, on \texttt{s.roll\_no}. In
        step~3, \sqlkw{NOT IN} and the two \sqlkw{EXCEPT} branches are
        uncorrelated; \sqlkw{NOT EXISTS} is correlated, on \texttt{i.emp\_id}.
        Steps~4 and~5 are correlated, on \texttt{c.course\_code}. The recursive
        CTE's step refers to itself, not to an outer query.
  \item \sqlkw{NOT IN} returned 0. Its list contained the 73 null
        \texttt{emp\_id}s, so each instructor's test included \texttt{x <>
        NULL}, which is \unk{} --- and no row's condition could be true.
        \sqlkw{EXCEPT} is a set operation, not a comparison: it removes a
        left-hand row only if an equal row appears on the right, and a null on
        the right could only remove a null on the left. \texttt{instructor.emp\_id}
        has none.
  \item Without an index, each of 48 courses scanned all 60{,}037
        enrolments --- 2.9 million row visits --- while \sqlkw{GROUP BY} read
        them once. With an index, each course became one index search and the
        correlated form read only what it counted, while \sqlkw{GROUP BY} still
        read everything. The advice is really ``avoid \emph{unindexed}
        correlated sub-queries''; correlation itself costs nothing.
\end{ans}

\ansch{16 --- Grouping, Aggregation and Window Functions}
\begin{ans}
  \item 2.383 divides the points by the \textbf{40{,}045} graded enrolments;
        1.590 divides the same points by all \textbf{60{,}037}, counting each
        ungraded enrolment as a zero. Publish 2.383, labelled as the average
        over graded courses, with the 19{,}992 in progress reported separately:
        an unfinished course is not an F.
  \item \texttt{course\_code LIKE 'CS-\%'} mentions only a grouping column, so
        it gives the same answer whether rows are removed before grouping or
        whole groups after --- and SQLite moved it into \sqlkw{WHERE} itself.
        \texttt{COUNT(*) > 1000} needs the count, which does not exist until
        grouping is finished, so it can only be tested after. Write
        \sqlkw{WHERE} because it states what you mean; the rewrite is an
        optimisation, not a promise.
  \item ``Each student's rank within their programme.''
        \texttt{SELECT roll\_no, programme, gpa, RANK() OVER (PARTITION BY
        programme ORDER BY gpa DESC) FROM per\_student;} \sqlkw{GROUP BY}
        programme would collapse each programme to one row and lose every
        student; the window keeps all of them and adds the rank.
\end{ans}

\ansch{17 --- Functional Dependencies}
\begin{ans}
  \item $\{$emp\_id$\}^{+}$: start \{emp\_id\}; \fd{emp\_id}{instructor} adds
        instructor; nothing else applies: \textbf{\{emp\_id, instructor\}}.
        $\{$course\_code, semester, section\_no$\}^{+}$: pass 1 adds room and
        emp\_id (the section rule) and title, credit\_hours, dept\_code (the
        course rule); pass 2 adds instructor (via emp\_id) and dept\_name (via
        dept\_code); pass 3 adds nothing. Ten attributes; not a superkey ---
        roll\_no, name, programme and grade are missing.
  \item $\{$roll\_no, course\_code$\}^{+}$: pass 1 adds name, programme (from
        roll\_no) and title, credit\_hours, dept\_code (from course\_code);
        pass 2 adds dept\_name. \textbf{dept\_name is in the closure, so the
        dependency follows.}
  \item Step 1: already single right sides. Step 2: in \fd{AC}{D}, is $C$
        extraneous? $\{A\}^{+} = \{A, B, C\}$ by \fd{A}{B} and \fd{A}{C}, then
        \fd{AC}{D} gives $D$ --- so \fd{A}{D} holds and $C$ is extraneous.
        Step 3: \fd{A}{C} is redundant (\fd{A}{B}, \fd{B}{C}). Result
        $\{$\fd{A}{B}, \fd{B}{C}, \fd{A}{D}$\}$; $\{A\}^{+} = \{A,B,C,D\}$ in
        both.
\end{ans}

\ansch{18 --- Normal Forms: 1NF to BCNF}
\begin{ans}
  \item \textbf{Update}: renaming department CS changed 7{,}869 rows (one in
        \textit{department}), and a half-done rename left two names --- caused
        by \fd{dept\_code}{dept\_name}. \textbf{Insertion}: a new course with no
        enrolments had no row to live in --- \fd{course\_code}{title, \dots}
        with a key that needs a student. \textbf{Deletion}: removing PH-350's
        1{,}122 enrolments removed the course --- the same dependency.
  \item Key $\{A, B\}$. \fd{A}{C} is \textbf{partial} (violates 2NF);
        \fd{C}{D} is \textbf{transitive} (violates 3NF). So R is in 1NF only.
        Deepest first: on \fd{C}{D}, $R_1(\underline{C}, D)$, remainder
        $(A, B, C)$; on \fd{A}{C}, $R_2(\underline{A}, C)$, remainder
        $R_3(\underline{A}, \underline{B})$. All three are in BCNF.
  \item The transcript was a \emph{lookup}: the normalised \textit{enrolment}
        could search its primary key by roll\_no, while \textit{report} had no
        index at all and was read end to end. The department report was a
        \emph{full scan}, cheaper over one table than over three. With an index
        on roll\_no the transcript result changes to \textbf{a tie} (0.08 against
        0.09\,ms); the full-scan result does not change.
\end{ans}

\ansch{19 --- Decomposition, and Denormalising on Purpose}
\begin{ans}
  \item $R_1 \cap R_2 = \{$course\_code$\}$. Its closure (under the rules
        these two tables can express) does not contain roll\_no or grade, so
        it is a superkey of neither: \textbf{lossy}. Joining on course\_code
        pairs every student who took a course with every grade given in it ---
        182{,}920 rows from 39{,}152 true facts.
  \item Split \textit{course} into \texttt{course\_dept(course\_code,
        dept\_code)} and \texttt{course\_dname(course\_code, dept\_name)}:
        lossless (course\_code is the key of both), but \fd{dept\_code}{dept\_name}
        lies in neither. Inserting \texttt{('CS-150', 'Computing')} when
        another CS course says \texttt{'Computer Science'} is accepted ---
        measured in \chref{decomposition}.
  \item Pays off when reads per write exceed $3.679 / 69.73 = 0.053$: one
        read in every 19 writes. On a 1\,ms-commit SSD, writes get cheaper in
        general, and the part of the 3.679\,ms that came from making a second
        table's page durable at each commit shrinks; \chref{views} measured the
        trigger's own work at about a hundredth of a millisecond. The added
        cost per write falls, so the
        break-even falls too --- the summary table pays off at even fewer
        reads.
\end{ans}

\ansch{20 --- Views, Triggers and Keeping Data Honest}
\begin{ans}
  \item SQLite stored only the view's \sqlkw{SELECT} text, substituted it into
        the query that named it, and compiled the result. What runs is the
        hand-written query with renumbered cursors. Nothing was stored, so
        nothing can be out of date: each use is the query run now.
  \item Choose the table when measured read cost matters --- here 700$\times$
        --- and reads greatly outnumber writes. Alongside it: triggers on
        \sqlkw{INSERT}, \sqlkw{DELETE} \emph{and} \sqlkw{UPDATE} of every
        table it summarises; a query comparing it with the equivalent view,
        run on a schedule; and a script that rebuilds it from the base tables
        when they disagree.
  \item
\begin{lstlisting}[language=sql]
CREATE TRIGGER max_six_upd
BEFORE UPDATE OF semester, roll_no ON enrolment
WHEN (SELECT COUNT(*) FROM enrolment
       WHERE roll_no = NEW.roll_no
         AND semester = NEW.semester) >= 6
BEGIN SELECT RAISE(ABORT, 'at most six sections a term'); END;
\end{lstlisting}
        It uses \texttt{NEW},
        because the limit concerns the term the enrolment is moving
        \emph{into}; \texttt{OLD} is the term it is leaving.
\end{ans}

\ansch{21 --- Indexes, and Where They Stop Helping}
\begin{ans}
  \item Each page of a B-tree divides its range among a few hundred children,
        so three levels reach millions of entries and a lookup reads one page
        per level --- three or four reads at 10{,}000 rows or a million.
        Hence 0.053 against 0.057\,ms. A scan reads every page, and a hundred
        times the rows is roughly a hundred times the pages: 1.37 to
        199.90\,ms.
  \item It lacked the \textbf{distribution} of \textit{k}: \texttt{sqlite\_stat1}
        recorded only the row count and one row per key, so every range looked
        selective. That would come from \textbf{STAT4} histograms, which the
        SQLite shipped with Python does not include. To force the scan: write
        \texttt{+k} in the condition, or \texttt{FROM t NOT INDEXED}.
  \item \texttt{roll\_no = ?}: \textbf{yes} --- the leftmost column.
        \texttt{course\_code = ?}: \textbf{no} --- course\_code values are
        scattered through every roll\_no. Both: \textbf{yes} --- the full
        prefix, the most selective use.
\end{ans}

\ansch{22 --- Transactions, Concurrency and Recovery}
\begin{ans}
  \item \textbf{Atomicity}: the failed move left a student in nothing under
        autocommit, unchanged inside a transaction. \textbf{Consistency}:
        constraints from \chref{constraints} refuse any end state that breaks
        them. \textbf{Isolation}: \sqlkw{IMMEDIATE} counted both clerks where
        no transaction lost one. \textbf{Durability}: \texttt{synchronous =
        FULL} waited 23\,ms per commit for the disk; without it a power cut can
        lose a reported commit.
  \item \sqlkw{DEFERRED}: both took \sqlkw{SHARED} to read; A's write took
        \sqlkw{RESERVED}; B's write wanted \sqlkw{RESERVED} and was refused;
        A's \sqlkw{COMMIT} wanted \sqlkw{EXCLUSIVE}, which needs B's
        \sqlkw{SHARED} gone --- deadlock until B rolled back.
        \sqlkw{IMMEDIATE}: A took \sqlkw{RESERVED} at \sqlkw{BEGIN}, so B
        could not begin; after A committed, B began, read 41 and wrote 42.
  \item Each commit waited about 23\,ms for the disk to confirm; 200 waits is
        4.6\,s, and one transaction waits once. \texttt{synchronous = OFF}
        removes the wait by abandoning the guarantee --- and was still
        9$\times$ slower than batching. Used for results entry, a power cut in
        the last week of term could lose grades the clerks saw saved, or
        corrupt the file.
\end{ans}

\ansch{23 --- Security, NoSQL, and the Semester Project}
\begin{ans}
  \item The \textbf{single quote}. Pasted, it closed the string the programmer
        had opened, so \texttt{OR '1'='1'} became SQL and was true for all
        4{,}000 rows. Bound as a parameter, the value was never parsed: the
        quote was just a character in a roll number that no student has, so 0
        rows.
  \item
\begin{lstlisting}[language=sql]
GRANT SELECT         ON enrolment TO exams_office;
GRANT UPDATE (grade) ON enrolment TO exams_office;
\end{lstlisting}
        In SQLite
        the same office, if it can open the file, can read and change every
        table, including fees and addresses --- there are no users to grant to.
  \item ``All of one student's grades'': the \textbf{document} --- one lookup
        returns the whole record. ``A grades in CS-101'': \textbf{relational}
        --- one indexed condition on one table, against unpacking all 4{,}000
        documents (9.8 against 78.4\,ms on the seeded data). The model should
        follow the questions the system is asked most.
\end{ans}

%=====================================================================
\chapter{Glossary}\label{app:glossary}
%=====================================================================
\normalfigures

Every key term in the handout, in one line, with the chapter that defines it.
The definition here is a reminder; the chapter's is the one to learn.

\rowcolors{2}{white}{lightbg}
\begin{longtable}{@{}L{3.9cm}L{8.9cm}R{1.2cm}@{}}
\toprule
\rowcolor{primary!15}
\textbf{Term} & \textbf{Meaning} & \textbf{Ch.} \\
\midrule
\endhead
ACID & atomicity, consistency, isolation, durability: what a transaction
  promises & 22 \\
aggregate function & computes one value from many rows: \sqlkw{COUNT},
  \sqlkw{SUM}, \sqlkw{AVG}, \sqlkw{MIN}, \sqlkw{MAX} & 16 \\
anomaly & a change that a redundant table makes wrong: update, insertion or
  deletion & 18 \\
ANSI/SPARC architecture & one database described three times: external,
  conceptual, internal & 2 \\
anti-join & the rows of $R$ with \emph{no} partner in $S$ & 15 \\
Armstrong's axioms & reflexivity, augmentation, transitivity: derive every
  dependency that follows & 17 \\
atomic value & one indivisible value per attribute per tuple & 7 \\
atomicity & a transaction happens wholly or not at all & 1, 22 \\
attribute & a named property of an entity type, or a column of a relation & 5,
  7 \\
attribute closure ($X^{+}$) & every attribute that $X$ determines & 17 \\
automatic index & an index SQLite builds for one query and then discards & 13 \\
autocommit & every statement is its own transaction & 22 \\
B-tree & the sorted tree of pages SQLite stores every table and index as & 21 \\
backfill & filling a new column for rows that already exist & 11 \\
bare column & a column in a grouped query that is neither grouped nor
  aggregated & 16 \\
BCNF & every determinant is a superkey & 18 \\
break-even & the read--write ratio at which a denormalisation pays for itself
  & 19 \\
buffer manager & decides which pages are held in memory & 3 \\
busy\_timeout & how long a connection waits for a lock before failing & 22 \\
candidate key & a minimal superkey & 5, 7 \\
cardinality (of a relation) & the number of tuples & 7 \\
cardinality ratio & 1:1, 1:N or M:N --- how many on each side & 6 \\
Cartesian product & every pairing of a tuple of $R$ with a tuple of $S$ & 12 \\
cascade & a referential action: delete the children with the parent & 8 \\
CHECK & a constraint every value must satisfy, within one row & 8, 10 \\
closure (algebra) & every operator's result is a relation, so expressions nest
  & 12 \\
column / table constraint & a constraint on one column / on several & 10 \\
common table expression & a named sub-query, written with \sqlkw{WITH} & 15 \\
composite attribute & an attribute with meaningful parts & 5 \\
composite key & a key of more than one attribute & 5 \\
conceptual schema & the whole database described once, without storage & 2 \\
concurrent access & more than one user changing data at the same time & 1, 22 \\
correlated sub-query & one that refers to a column of the outer query & 15 \\
covering index & an index holding every column a query needs & 21 \\
crossover & the share of a table past which an index is slower than a scan &
  21 \\
data dictionary & the tables in which the DBMS stores the schema & 3 \\
data independence & changing one level without changing those above it & 2 \\
database & data together with its description & 1 \\
DBA & database administrator: decides how the database runs & 3 \\
DBMS & the software that creates, maintains and controls a database & 1 \\
deadlock & two transactions each waiting for a lock the other holds & 22 \\
decomposition & splitting a relation into smaller ones & 18, 19 \\
degree (of a relation) & the number of attributes & 7 \\
degree (of a relationship) & how many entity types take part & 5 \\
denormalisation & storing derived or repeated data on purpose, for speed & 19 \\
dependency preservation & every dependency checkable within one table & 19 \\
derived attribute & one computed from others, normally not stored & 5 \\
determinant & the left side of a functional dependency & 17 \\
difference ($-$) & tuples in $R$ and not in $S$ & 12 \\
dirty read & reading another transaction's uncommitted change & 22 \\
domain & the set of values an attribute may take & 7 \\
domain constraint & every value comes from its attribute's domain & 8 \\
durability & a committed transaction survives a crash & 22 \\
entity / type / set & a thing; a kind of thing; all of them at one moment & 5 \\
entity integrity & no part of a primary key is null & 8 \\
equijoin & a theta join using only equality & 13 \\
EXPLAIN QUERY PLAN & shows the plan the optimiser chose & 3, 13 \\
external schema & what one group of users sees: a view & 2, 20 \\
fan-out & how many children a B-tree page points to & 21 \\
first normal form & every value atomic & 18 \\
foreign key & attributes whose values must match a primary key elsewhere, or be
  null & 7, 8 \\
functional dependency & \fd{X}{Y}: tuples agreeing on $X$ agree on $Y$ & 17 \\
generated column & a column computed from others; \sqlkw{VIRTUAL} or
  \sqlkw{STORED} & 10 \\
identifying relationship & joins a weak entity to its owner & 6 \\
index & a column's values in sorted order, each with its row's rowid & 21 \\
inner / outer join & drops / keeps rows with no partner & 13 \\
instance & the data in a database at one moment & 2 \\
internal schema & how the data is actually stored & 2 \\
isolation & concurrent transactions do not see each other's unfinished work &
  22 \\
leftmost prefix & a composite index serves only conditions on its first
  columns & 21 \\
lossless join & a decomposition that joins back to exactly the original & 19 \\
lost update & two read-then-writes, one overwriting the other & 1, 22 \\
mapping & translating between adjacent levels, or an ER diagram to tables & 2,
  9 \\
migration & a numbered script moving a schema from one version to the next &
  11 \\
minimal cover & the smallest equivalent set of dependencies & 17 \\
multi-valued attribute & one with several values per entity; becomes a table &
  5, 9 \\
natural join & an equijoin on every shared attribute name & 13 \\
NoSQL & databases outside the relational model & 23 \\
null & the absence of a value; neither zero nor empty & 7, 14 \\
optimiser & chooses the cheapest plan for a statement & 3 \\
page & the fixed-size unit of storage: 4{,}096 bytes here & 4 \\
parameterised query & values bound separately and never parsed as SQL & 23 \\
partial / total participation & an entity need not / must take part & 6 \\
partial dependency & on part of a candidate key: violates 2NF & 18 \\
partial key & distinguishes weak entities of one owner & 6 \\
primary key & the candidate key chosen to identify rows & 5, 7 \\
projection ($\pi$) & chooses columns and removes duplicates & 12 \\
recursive relationship & an entity type related to itself & 5 \\
referential action & what happens to children when a parent changes & 8 \\
referential integrity & a reference points at something that exists & 8 \\
relation & a set of tuples over the same attributes & 7 \\
relation schema & a relation's name and attributes, key underlined & 4, 7 \\
relational algebra & operators that take and return relations & 12 \\
role name & names an end of a relationship, essential when recursive & 5 \\
rollback journal & original pages saved until commit, replayed after a crash &
  22 \\
rowid & the hidden 64-bit key of every ordinary SQLite table & 10 \\
savepoint & a point inside a transaction to roll back to & 22 \\
schema & the description of a database; changes rarely & 2 \\
second normal form & 1NF with no partial dependency & 18 \\
selection ($\sigma$) & chooses rows & 12 \\
self-join & a table joined to itself under two aliases & 13 \\
semi-join & the rows of $R$ with at least one partner in $S$ & 15 \\
serialisable & the result is as if transactions ran one after another & 22 \\
SQL injection & input pasted into SQL and run as SQL & 23 \\
STAT4 & histograms that let the optimiser estimate a range; not in Python's
  SQLite & 21 \\
storage class & what a value actually is: null, integer, real, text, blob &
  10 \\
STRICT & a table whose declared types are enforced & 10 \\
sub-query & a query inside another & 15 \\
superkey & attributes no two tuples agree on & 5, 7 \\
surrogate key & an invented identifier with no meaning in the world & 5 \\
third normal form & no transitive dependency on a non-prime attribute & 18 \\
three-valued logic & true, false and \unk{} & 14 \\
transaction & a group of changes treated as one & 22 \\
transitive dependency & $K \rightarrow X \rightarrow Y$, with $X$ not a key:
  violates 3NF & 18 \\
trigger & SQL run automatically on an insert, update or delete & 19, 20 \\
tuple & one row: a value from each attribute's domain & 7 \\
type affinity & a declared type is a preference; unconvertible values are kept
  & 10 \\
union-compatible & same number of attributes, compatible domains, same order &
  12 \\
\unk{} & the truth value of any comparison with a null & 14 \\
view & a named query, stored as text and expanded when used & 20 \\
WAL & write-ahead log: readers never block the writer & 22 \\
weak entity type & one with no key of its own & 6 \\
window function & a value per row from its partition, without collapsing rows
  & 16 \\
\bottomrule
\end{longtable}

%=====================================================================
\chapter{Resources}\label{app:resources}
%=====================================================================
\normalfigures

\section{The Four Books HEC Names, and Which Editions to Use}

HEC's 2017 and 2023 outlines list the same four reference books, in the same
words, with an edition for each and \textbf{no publication year for any of
them}. An edition number without a year cannot be checked against a library
catalogue or a bookshop without first finding out when it appeared --- and it
does not say whether a newer one has replaced it. The 2023 booklet adds
``(or use any other standard and latest books)'', which licenses exactly what
this appendix does: name each book at the edition to use, with its year.

\medskip
\begin{longtable}{@{}L{6.2cm}L{2.4cm}L{5.4cm}@{}}
\toprule
\rowcolor{primary!15}
\textbf{Book, as HEC lists it} & \textbf{HEC names} & \textbf{Use} \\
\midrule
\endhead
Connolly and Begg, \emph{Database Systems: A Practical Approach to Design,
Implementation, and Management} & 6th & 6th edition, Pearson, 2015 --- still
the latest \\
Garcia-Molina, Ullman and Widom, \emph{Database Systems: The Complete Book} &
2nd & 2nd edition, Pearson Prentice Hall, 2009 --- still the latest \\
Silberschatz, Korth and Sudarshan, \emph{Database System Concepts} & 6th &
\textbf{7th edition}, McGraw-Hill, 2019 --- HEC names a superseded edition \\
Ramakrishnan and Gehrke, \emph{Database Management Systems} & 3rd & 3rd
edition, McGraw-Hill, 2003 --- still the latest, and the oldest of the four \\
\bottomrule
\end{longtable}

\begin{conceptbox}[Which to read for what]
\begin{tight}
  \item \textbf{Silberschatz} is the one to own. Its chapters map most closely
        onto this course's order --- relational model, SQL, ER design,
        normalisation, indexing, transactions --- and the 7th edition covers
        the NoSQL material of \chref{project}.
  \item \textbf{Connolly and Begg} is the most patient with design: read it
        beside Part~II if \chref{ermodel} and \chref{cardinality} felt fast.
  \item \textbf{Ramakrishnan and Gehrke} is the strongest on what is inside a
        DBMS --- buffer management, B-trees, query evaluation --- and is the
        book for \chref{indexes} and \chref{transactions} if you want the
        mechanism behind the measurement.
  \item \textbf{Garcia-Molina, Ullman and Widom} is the most mathematical;
        its treatment of functional dependencies and decomposition is the one
        to read for \chref{fds} to \chref{decomposition}.
\end{tight}

\smallskip
None of the four is about SQLite. Where this handout describes SQLite's
behaviour --- affinity, \sqlkw{PRAGMA} \texttt{foreign\_keys}, file locking ---
the textbooks describe the standard, and SQLite's own documentation (below) is
the authority.
\end{conceptbox}

\section{The Tool's Own Documentation}

Every SQLite behaviour this handout measures is documented by its authors,
and the documentation is the place to settle any disagreement with a
textbook. All are at \texttt{sqlite.org}; the page names are given so they can
be found from the site's search.

\medskip
\begin{longtable}{@{}L{6.0cm}L{8.0cm}@{}}
\toprule
\rowcolor{primary!15}
\textbf{Page} & \textbf{Settles} \\
\midrule
\endhead
\emph{Datatypes In SQLite} & type affinity and storage classes ---
  \chref{ddl}'s \texttt{'hello'} in an \texttt{INTEGER} column \\
\emph{STRICT Tables} & what \sqlkw{STRICT} enforces, and since which version \\
\emph{SQLite Foreign Key Support} & why \texttt{foreign\_keys} is off by
  default, and that it is per connection --- \chref{constraints} \\
\emph{The SQLite Query Optimizer Overview} & automatic indexes, push-down,
  and the \texttt{+column} trick --- \chref{joins} to \chref{indexes} \\
\emph{EXPLAIN QUERY PLAN} & how to read \sqlkw{SCAN}, \sqlkw{SEARCH} and
  \sqlkw{TEMP B-TREE} \\
\emph{ANALYZE} & \texttt{sqlite\_stat1}, and STAT4 when it is compiled in \\
\emph{File Locking And Concurrency In SQLite Version 3} & \sqlkw{SHARED},
  \sqlkw{RESERVED}, \sqlkw{EXCLUSIVE} --- \chref{transactions} \\
\emph{Atomic Commit In SQLite} & the rollback journal and recovery after a
  crash \\
\emph{Write-Ahead Logging} & WAL mode and what it changes for readers \\
\emph{ALTER TABLE} & the four forms and the twelve-step procedure ---
  \chref{evolution} \\
\emph{JSON Functions And Operators} & \texttt{json\_object},
  \texttt{json\_each} --- \chref{project} \\
\bottomrule
\end{longtable}

\medskip
For the Python side, the \texttt{sqlite3} module's documentation at
\texttt{docs.python.org} explains \texttt{isolation\_level} --- the setting
\chref{transactions}' harness changes so that every \sqlkw{BEGIN} is one you
wrote.

\section{The Original Papers}

Three papers founded the subject. All three are short, and the first is still
the clearest statement of why a database should be relational.

\begin{tight}
  \item E.~F.~Codd, ``A Relational Model of Data for Large Shared Data
        Banks'', \emph{Communications of the ACM} 13(6), 1970 ---
        \chref{relational} and \chref{algebra}.
  \item P.~P.~Chen, ``The Entity-Relationship Model --- Toward a Unified View
        of Data'', \emph{ACM Transactions on Database Systems} 1(1), 1976 ---
        \chref{ermodel}.
  \item ANSI/X3/SPARC Study Group on Database Management Systems,
        \emph{Interim Report}, 1975 --- the three-level architecture of
        \chref{architecture}.
\end{tight}

\section{Law}

\begin{tight}
  \item \emph{Prevention of Electronic Crimes Act 2016}, Pakistan, sections 3
        and 4 --- unauthorised access to an information system and to data.
        The scope statement of Appendix~A, section A.6, and of \chref{project}
        rests on it.
\end{tight}

\section{This Handout's Own Code}

Every measured figure was produced by a file in \texttt{code/}, on the machine
described in Appendix~A, section A.5:

\medskip
\begin{longtable}{@{}L{5.2cm}L{8.8cm}@{}}
\toprule
\rowcolor{primary!15}
\textbf{File} & \textbf{Produces} \\
\midrule
\endhead
\texttt{schema/v1\_create.sql} & the university schema, from \chref{mapping} \\
\texttt{schema/v2\_constraints.sql} & the migration of \chref{evolution} \\
\texttt{schema/seed.py} & reproducible data at any scale: 4{,}000 students,
  60{,}037 enrolments by default \\
\texttt{m04\_firstdatabase.py} & \chref{firstdatabase}: time to a first row,
  bytes per page \\
\texttt{m08\_constraints.py} & \chref{constraints}: foreign keys off by default \\
\texttt{m10\_ddl.py} & \chref{ddl}: affinity against \sqlkw{STRICT} \\
\texttt{m13\_joins.py} & \chref{joins}: \sqlkw{NATURAL JOIN}, missing sections,
  automatic indexes \\
\texttt{m15\_subqueries.py} & \chref{subqueries}: anti-joins and correlation \\
\texttt{m16\_grouping.py} & \chref{grouping}: nulls, sorts, windows \\
\texttt{m18\_normalforms.py} & \chref{normalforms}: the wide table against BCNF \\
\texttt{m19\_decomposition.py} & \chref{decomposition}: lossy joins and a
  priced summary table \\
\texttt{m20\_views.py} & \chref{views}: what views and triggers cost \\
\texttt{m21\_indexes.py} & \chref{indexes}: the crossover \\
\texttt{m22\_transactions.py} & \chref{transactions}: a lost update, a
  deadlock, the price of a commit \\
\texttt{run\_all.py} & runs every harness and checks each result's
  \emph{shape} still holds \\
\texttt{sync\_labs.py} & checks that every file and command a laboratory names
  exists \\
\bottomrule
\end{longtable}

%=====================================================================
\chapter{The Boundary With Other Courses}\label{app:boundary}
%=====================================================================
\normalfigures

This course sits early in the degree --- second semester --- and three later
courses depend on it. It therefore assumes little and promises a great deal,
and both are recorded here: what is assumed, what is delivered, and to whom,
so that a lecturer can see what may be taken for granted on each side and a
student can see why a familiar topic reappears.

\section{The Edition History}

\rowcolors{2}{white}{lightbg}
\begin{longtable}{@{}L{1.5cm}L{3.1cm}L{3.8cm}L{1.9cm}L{3.0cm}@{}}
\toprule
\rowcolor{primary!15}
\textbf{Edition} & \textbf{Status} & \textbf{Degrees} & \textbf{Semester}
& \textbf{Content published?} \\
\midrule
\endhead
\textbf{2017} & Computing Core, 3+1 & BSCS, BSSE, BSIT & IV, IV, \textbf{V}
  & yes \\
\textbf{2023} & Computing Core, 4 (3-3) & all ten computing BS, and ADCS &
  \textbf{2} & yes --- the 2017 text \\
2025 & Major Compulsory, 3+1 & BSCS and ADCS only --- \textbf{no BSIT} & III
  & \textbf{no} \\
\bottomrule
\end{longtable}

\smallskip
The department teaches CS-103 at \textbf{BSIT, second semester, 4 credit hours
(3+1)}, which matches HEC 2023. HEC 2017 placed it at BSIT semester V. The 2023
course content is the 2017 text word for word; 2025 removed it. The Course
Specification at the front records all six defects found in these outlines.

\begin{alertbox}[One booklet, two prerequisites]
HEC 2017's course description page gives \textbf{``Prerequisites: None''}.
The same booklet's scheme-of-studies tables, on every page that lists the
course, give \textbf{``Data Structures \& Algorithms''}. HEC 2023 gives
``None''.

\smallskip
At second semester Data Structures has not been taken, so this handout
\textbf{assumes no data-structures course}. The two chapters that would most
naturally lean on one are written to stand alone: \chref{indexes} explains a
B-tree by its fan-out and depth, with Appendix~B's logarithm, and never by an
algorithm; \chref{transactions} explains locks without operating-systems
vocabulary. Recorded here so that a reviewer comparing the two documents finds
the answer already written.
\end{alertbox}

\section{Programming Fundamentals --- the Real Prerequisite}

\begin{conceptbox}[The rule]
\textbf{Programming Fundamentals owns the program. This course owns what the
program keeps.}

\smallskip
A student arriving here should be able to write a loop, read a file, and call
a function. Nothing more is used: Python appears in the harnesses, which
students run and read rather than write, and in two laboratories ---
\chref{whynotfiles}'s file-based system and \chref{project}'s injection
demonstration --- where the programs are a dozen lines.
\end{conceptbox}

\section{Data Structures --- Not Assumed}

\rowcolors{2}{white}{lightbg}
\begin{longtable}{@{}L{5.4cm}L{3.4cm}L{5.6cm}@{}}
\toprule
\rowcolor{primary!15}
\textbf{Topic} & \textbf{Owned by} & \textbf{Here} \\
\midrule
\endhead
B-trees: insertion, splitting, balance & Data Structures & \emph{not taught};
  \chref{indexes} needs only fan-out and depth \\
hashing & Data Structures & mentioned, not used \\
sorting algorithms & Data Structures & \chref{grouping} measures what a sort
  \emph{costs} --- \sqlkw{TEMP B-TREE} --- not how it works \\
logarithms & Appendix~B & taught here, as ``how many halvings'' \\
\bottomrule
\end{longtable}

\smallskip
\textbf{The useful repetition:} a student who takes Data Structures later will
meet the B-tree's insertion algorithm, and will already know why it matters ---
\chref{indexes} measured a lookup costing 0.057\,ms in a million rows, and a
scan 199.90\,ms.

\section{E-Commerce Application Development --- Downstream}

\begin{conceptbox}[The rule, as that course states it]
\textbf{``Database Systems owns the model. This course owns the consequences of
somebody else's model.''}

\smallskip
E-Commerce Application Development's own Appendix~F says that, and lists what
it assumes: \emph{``SQL selects and joins, primary and foreign keys,
normalisation, what an index is.''} This course must deliver all four, and
does.
\end{conceptbox}

\rowcolors{2}{white}{lightbg}
\begin{longtable}{@{}L{4.6cm}L{4.6cm}L{5.2cm}@{}}
\toprule
\rowcolor{primary!15}
\textbf{E-Commerce assumes} & \textbf{Delivered in} & \textbf{And measured} \\
\midrule
\endhead
SQL selects and joins & \chref{select}, \chref{joins}, \chref{subqueries},
  \chref{grouping} & \sqlkw{NOT IN} returning 0 where \sqlkw{NOT EXISTS}
  returned 3 \\
primary and foreign keys & \chref{relational}, \chref{constraints} & foreign
  keys off by default; a \texttt{TEXT PRIMARY KEY} storing null \\
normalisation & \chref{fds} to \chref{decomposition} & a wide table
  162$\times$ slower on a lookup; 143{,}768 spurious tuples \\
what an index is & \chref{indexes} & the crossover --- which E-Commerce
  then finds again, in a product catalogue \\
\bottomrule
\end{longtable}

\smallskip
\textbf{The useful repetition}, from the downstream side: E-Commerce tells its
students that ``every Database Systems student has been told that indexes make
queries faster'', and makes that a measurement with a crossover in it. In this
course they have already measured it --- so the E-Commerce chapter becomes a
second instance of a known result, on a schema they did not design, rather than
a surprise.

\section{Information Security --- Later}

\begin{conceptbox}[The rule]
\textbf{Information Security owns the principles. This course owns one
database, and what can be done to it from inside.}

\smallskip
\chref{project} shows SQL injection against the student's own file, under the
scope statement of Appendix~A, and stops there. Threat modelling,
authentication, cryptography and the law are Information Security's; this
course shows the one defence that belongs to the person writing the query ---
the parameter --- and the access control a server adds that SQLite lacks.
\end{conceptbox}

\section{Software Engineering --- Later}

\begin{tight}
  \item \textbf{Assumed there, delivered here:} that a schema is designed
        before it is built (\chref{ermodel} to \chref{mapping}), and changed by
        numbered, reviewed migrations (\chref{evolution}).
  \item \textbf{Taught properly there, used informally here:} version
        control. \chref{evolution}'s first rule --- never edit a migration that
        has run anywhere --- is a version-control rule as much as a database
        one, and Software Engineering is where the tooling behind it is taught.
\end{tight}

%=====================================================================
\chapter{The Sixteen-Week Schedule}\label{app:schedule}
%=====================================================================
\normalfigures

This appendix is the \textbf{source of record} for the teaching schedule. The
one-page \texttt{DBS\_Fall2026\_Lecture\_Schedule.pdf} distributed to students
is generated from it, and the slide pipeline parses the table below. Change
this first; regenerate the one-page version afterwards.

Twenty-three chapters are taught across sixteen weeks, so seven weeks carry
two chapters. The pairings are not arbitrary: in every case the two are either
one topic the printed book splits for length --- the two halves of the ER
model, the two halves of decomposition --- or a pair in which the second is
unusable without the first.

Two weeks are deliberately left alone. \textbf{Chapter~18, Normal Forms, has a
week to itself}, although it would pair naturally with Chapter~17. Its
laboratory breaks the running database and rebuilds it, and that does not fit
in half a week alongside anything else. \textbf{Chapter~21, Indexes}, likewise:
the measurement at a million rows takes most of a laboratory session to run
and the rest of it to argue about.

\section{The Schedule}

% MACHINE-READ. The slide builder takes the FIRST longtable in this file and
% requires: five columns; a bare integer week; chapters as comma-separated
% digits or N--M ranges; literal --- for an empty assessment cell; and no
% chapter appearing in two weeks.
\rowcolors{2}{white}{lightbg}
% column 3 holds the bold head "Chapters", which needs 1.8 cm at 10.95 pt
\begin{longtable}{@{}L{1.1cm}L{3.2cm}L{1.8cm}L{4.4cm}L{3.9cm}@{}}
\toprule
\rowcolor{primary!15}
\textbf{Wk} & \textbf{Lecture topic} & \textbf{Chapters} & \textbf{Laboratory} &
\textbf{Assessment} \\
\midrule
\endhead
1 & Why not files; the three-level architecture & 1, 2 & Break a file-based
system on purpose; find the three levels in SQLite & --- \\
2 & Who does what; your first database & 3, 4 & Install nothing; create a
database and measure what it cost & --- \\
3 & The entity-relationship model & 5 & Entities and attributes for the
university, on paper & --- \\
4 & Cardinality, participation, weak entities & 6 & The university ER diagram,
finished and argued over & --- \\
5 & The relational model; keys and constraints & 7, 8 & Orphan rows in a
correctly declared schema & \textbf{Quiz 1 (Ch 1--6)} \\
6 & Mapping a diagram to tables & 9 & Your diagram becomes eight tables, by
rule & \textbf{Assignment 1 set} \\
7 & Creating and changing a schema & 10, 11 & \texttt{CREATE TABLE}; then
change one that is already full & --- \\
8 & Relational algebra & 12 & Midterm, then algebra by hand on small relations
& \textbf{Midterm (Ch 1--11)} \\
9 & Joins; selecting rows and NULL & 13, 14 & Lose 73 sections and 19{,}992
rows to joins and NULL, then find them all & \textbf{Assignment 1 due} \\
10 & Sub-queries; grouping and aggregation & 15, 16 & One question seven ways;
an average moved by a third & \textbf{Project proposal} \\
11 & Functional dependencies & 17 & The wide report table: its dependencies,
its closures and its key & --- \\
12 & Normal forms: 1NF to BCNF & 18 & Break the running database, then
normalise it and measure both & \textbf{Quiz 2 (Ch 12--17)} \\
13 & Decomposition; views and triggers & 19, 20 & A lossy split that invents
facts; a summary table caught lying & \textbf{Assignment 2 set} \\
14 & Indexes, and where they stop helping & 21 & One million rows, with and
without an index & --- \\
15 & Transactions, concurrency and recovery & 22 & Two writers, one database,
and a lost update & \textbf{Assignment 2 due} \\
16 & Security, NoSQL and the project & 23 & Project presentations &
\textbf{Project due; presentations} \\
\bottomrule
\end{longtable}

\section{How the Assessment Adds Up}

HEC publishes the assessment instruments for this course --- sessional
examination, home assignments, quizzes, project, presentations, final
examination --- without weights. The weights below are the department's, and
are stated here so that a student can see the whole shape of the semester on
one page.

\rowcolors{2}{white}{lightbg}
\begin{longtable}{@{}L{4.6cm}L{2.0cm}L{8.4cm}@{}}
\toprule
\rowcolor{primary!15}
\textbf{Instrument} & \textbf{Weight} & \textbf{What it examines} \\
\midrule
\endhead
Two quizzes & 10\% & Weeks 5 and 12. Short, closed-book, on the chapters named
in the schedule \\
Two assignments & 10\% & Set in weeks 6 and 13. Both are done on your own
database and submitted as a change to it plus a short write-up \\
\textbf{Laboratory} & \textbf{15\%} & Assessed weekly from Week~2. This course
carries a \textbf{credited laboratory} --- one of its four credit hours --- and
the marks follow the credit \\
Midterm examination & 20\% & Week 8, Chapters 1--11 \\
Semester project & 20\% & The database itself, built across the whole semester
(\chref{project}), submitted in week 16 with its presentation \\
Final examination & 25\% & All chapters, with emphasis after the midterm \\
\bottomrule
\end{longtable}

\begin{alertbox}[The laboratory is a quarter of the course, not a formality]
This is the only course in the series with a \emph{credited} laboratory: HEC
specifies 4 credit hours as 3+1 in 2017 and 2025, and as 4 (3--3) in 2023. One
of those four hours is the lab, and 15\% of the marks follow it.

\smallskip
That is also why the laboratories in this handout are longer than those in its
sibling courses. They are not homework to be fitted around the lectures; they
are a timetabled hour with a demonstrator in the room, and every one of them
adds something to the database you will submit in Week~16.

\smallskip
\chref{project} states what is required and how it is marked. Read it in
Week~3, not in Week~15.
\end{alertbox}


\end{document}
